
%%FILE 3-majority/blueprint/src/content.tex
% Content of the blueprint for the 3-majority opinion-dynamics formalization.
% Mirrors latex/three_majority.tex in spirit, but every statement here also
% carries \lean{...} / \leanok / \uses{...} tags matched to the *actual*
% Lean declarations in ThreeMajority/*.lean (namespace `ThreeMajority`), so
% that leanblueprint can build the dependency graph and check coverage.
% Where the Lean proof takes a different route than the paper draft in
% latex/three_majority.tex, this file follows the Lean proof.

\begin{abstract}
\noindent
We give a complete, elementary proof that the \emph{3-majority}
opinion-dynamics process on $n$ fully-mixing agents, each holding one of two
opinions, reaches consensus within $O(\log n)$ rounds with probability
$1-O(1/n)$, provided the initial imbalance between the two opinions exceeds
a fixed constant fraction of $n$ (at least $60\%$ vs.\ $40\%$). The
probability space is finite and uniform throughout; the only analytic
ingredient beyond finite sums is a single self-contained exponential-moment
(Chernoff-type) concentration lemma for sums of independent
$\{0,1\}$-valued random variables, proved from the elementary bound
$1+x \le e^x$ --- no measure theory, no PMF/kernels, no appeal to Mathlib's
probability library. This is a genuinely different regime from the
companion \texttt{rumor\_spread} (push protocol) development: the
opinion-$1$ count is \emph{not} monotone in the round index (an unlucky
round can shrink it), so a counting argument over ``good rounds'' does not
suffice --- honest concentration around the process's mean trajectory is
needed in \emph{both} the growth and saturation phases. This argument has
been \emph{fully formalized} (no \texttt{sorry}) in Lean~4 + Mathlib; the
main theorem is \texttt{ThreeMajority.majority3\_consensus\_whp}.

Section~\ref{sec:bcekmn-any-start} formalizes, for any number of colours, the comparison of
3-Majority with the Voter model of Berenbrink et al.\ (PODC 2017) and their bound
$O(n^{3/4}\log^{7/8}n)$ on the consensus time from any configuration. Its second phase is
Theorem~3.1 of Becchetti et al.\ (SODA 2016), from $k\le n^{1/3-\varepsilon}$ colours, which
Section~\ref{sec:bcnpt16-many-opinions} formalizes; together they give the bound without
hypothesis.
\end{abstract}

\section{The model}\label{sec:model}

Fix $n \ge 1$ and let $V = \{0,1,\dots,n-1\}$ be the set of agents, each
holding one of two opinions, $1$ (``red'') or $0$ (``blue''). Write
$I \subseteq V$ for the set of agents currently holding opinion $1$.

\begin{definition}[Round randomness]\label{def:round}
  \lean{ThreeMajority.Tgt3}
  \leanok
  A \emph{round configuration} is a function $r \colon V \to V \times V
  \times V$ assigning to every agent $v$ an ordered triple of \emph{samples},
  each drawn independently and uniformly from $V$, \emph{with} replacement
  (an agent may sample itself, or the same agent twice --- harmless, since a
  majority of three Boolean values is always well defined). Let
  $R_n = V^{3V}$ denote the (finite) set of round configurations, with the
  uniform probability measure.
\end{definition}

\begin{definition}[Majority step]\label{def:step}
  \lean{ThreeMajority.step, ThreeMajority.sampleCount, ThreeMajority.sampleCountOf, ThreeMajority.mem_step}
  \leanok
  \uses{def:round}
  For an opinion-$1$ set $I \subseteq V$ and $r \in R_n$,
  \[
    \mathrm{step}(I,r) \;=\; \{\, v \in V : \text{at least two of } r(v)\text{'s three
    coordinates lie in } I \,\}.
  \]
  Each agent independently samples $3$ agents and adopts the majority of
  their current opinions.
\end{definition}

\begin{definition}[Process]\label{def:process}
  \lean{ThreeMajority.run}
  \leanok
  \uses{def:step}
  Fix a horizon $T \in \N$ and an initial set $I_0 \subseteq V$. The sample
  space is $\Omega = R_n^{\,T}$ with the uniform (i.i.d.) measure. For
  $r = (r_0,\dots,r_{T-1}) \in \Omega$ define $I_{t+1} =
  \mathrm{step}(I_t, r_t)$. We write $X_t = |I_t|$ and $U_t = n - X_t$. In
  the formalization, a trajectory is a \texttt{List} of round configurations
  and $\mathrm{run}\ I\ l$ folds \texttt{step} over the list $l$.
\end{definition}

\begin{lemma}[Monotone coupling]\label{lem:mono}
  \lean{ThreeMajority.step_mono, ThreeMajority.run_mono, ThreeMajority.sampleCount_mono}
  \leanok
  \uses{def:step, def:process}
  For a \emph{fixed} round $r$, the map $I \mapsto \mathrm{step}(I,r)$ is
  monotone: $I \subseteq I'$ implies $\mathrm{step}(I,r) \subseteq
  \mathrm{step}(I',r)$; iterating over a trajectory, $\mathrm{run}$ is
  monotone in the same sense. This is purely combinatorial (no probability
  involved).
\end{lemma}

\begin{remark}
  Unlike the push protocol, where the analogous monotonicity
  (Lemma~\ref{lem:mono} there) is load-bearing throughout, here the
  per-round probabilistic estimates below (Lemmas~\ref{lem:growthround},
  \ref{lem:satclosed}, etc.) are instead stated and proved directly for
  \emph{any} set $I$ satisfying a hypothesis on $|I|$, using monotonicity of
  the real polynomial $p$ of Section~\ref{sec:map} on $[0,1]$ rather than a
  reference-embedding argument via Lemma~\ref{lem:mono}. Lemma~\ref{lem:mono}
  is recorded because it is the exact analogue of the push protocol's
  structural fact, but it is not otherwise used by the chain of lemmas
  leading to the main theorem.
\end{remark}

\section{Main theorem}\label{sec:main-statement}

Throughout, $\log$ denotes the natural logarithm.

\begin{theorem}[3-majority reaches consensus in $O(\log n)$ rounds w.h.p.]
  \label{thm:main}
  \lean{ThreeMajority.majority3_consensus_whp}
  \leanok
  \uses{lem:mainfailclean}
  Let $n$ be such that $\log n \ge 30$ (in particular $n \ge 2$), let
  $I_0 \subseteq V$ with $|I_0| \ge \tfrac35 n$, and set
  \[
    T \;=\; 10 + (T_{2a}(n) + 2), \qquad T_{2a}(n) = \ceil{6\log n}.
  \]
  Then $\Pr[\, I_T = V \,] \;\ge\; 1 - \dfrac{500}{n}$.
\end{theorem}

\begin{remark}
  As in the push protocol's write-up, the threshold $\log n \ge 30$ and
  the constant $500$ are deliberately crude: the proof is organized to keep
  every inequality an easy direct computation, not to be sharp. Reaching
  this modest, two-digit threshold on $\log n$ (rather than an
  astronomically large one) relies on two ingredients in
  Section~\ref{sec:ineq}: a degree-$11$ instance of the generic
  exponential-vs-power bound (Lemma~\ref{lem:exppow}), and a log-log bound
  tangent at a reference point near the threshold itself rather than at the
  fixed point $e$ (Lemma~\ref{lem:loglog}). It is known
  (Doerr--Goldberg--Minder--Sauerwald--Scheideler 2011;
  Becchetti--Clementi--Natale--Pasquale--Silvestri, SODA 2015) that an
  initial imbalance of only $\Theta(\sqrt{n\log n})$ already suffices for
  $O(\log n)$-round consensus w.h.p.; reaching that sharp threshold
  requires tracking the process on the scale of its own standard deviation
  throughout, a substantially more delicate multi-scale argument. Here a
  constant-fraction initial imbalance keeps every step to a single,
  reusable Chernoff bound (Section~\ref{sec:chernoff}) applied at a handful
  of explicit parameters.
\end{remark}

The proof has two phases, as in the push protocol, but for a structurally
different reason: \textbf{growth} (Section~\ref{sec:growth}), where the
opinion-$1$ fraction self-amplifies from $60\%$ to $75\%$; and
\textbf{saturation} (Section~\ref{sec:saturation}), where the last
dissenting agents are eliminated in three stages. Because the process is
not monotone in the round index, \emph{both} phases need the Chernoff bound
of Section~\ref{sec:chernoff} --- unlike the push protocol, where only
saturation needed a (much weaker) expectation argument.

\section{Finite uniform probability and independence}\label{sec:prelim}

As in the push protocol, all randomness is uniform over finite types, so
instead of measure theory or \texttt{PMF} we work with plain finite sums.

\begin{definition}[Average]\label{def:avg}
  \lean{ThreeMajority.avg}
  \leanok
  For a finite type $\alpha$ and $f \colon \alpha \to \R$,
  $\mathrm{avg}(f) := \bigl(\sum_a f(a)\bigr)/|\alpha|$ is the expectation of
  $f$ under the uniform distribution on $\alpha$.
\end{definition}

\begin{definition}[Trajectory expectation]\label{def:expList}
  \lean{ThreeMajority.expList}
  \leanok
  \uses{def:avg}
  Exactly as in the push protocol: $\mathrm{expList}\ \alpha\ T\ F$ is the
  expectation of a functional $F$ of a trajectory of $T$ i.i.d.\ uniform
  draws from $\alpha$, defined by recursion on $T$ (the head draw is
  averaged out first), so conditioning on the first round is a definitional
  unfolding.
\end{definition}

\begin{lemma}[Basic properties]\label{lem:avg-basic}
  \lean{ThreeMajority.avg_nonneg, ThreeMajority.avg_le_avg, ThreeMajority.avg_add,
    ThreeMajority.avg_sub, ThreeMajority.avg_const_mul, ThreeMajority.avg_sum,
    ThreeMajority.avg_indicator, ThreeMajority.avg_const,
    ThreeMajority.expList_nonneg, ThreeMajority.expList_le_expList,
    ThreeMajority.expList_add, ThreeMajority.expList_const_mul,
    ThreeMajority.expList_const, ThreeMajority.expList_append}
  \leanok
  \uses{def:avg, def:expList}
  $\mathrm{avg}$ and $\mathrm{expList}$ are monotone, additive, homogeneous,
  and constants pass through unchanged; $\mathrm{expList}$ over a
  concatenated horizon splits as a nested $\mathrm{expList}$
  (\texttt{expList\_append}), used to glue phases in
  Section~\ref{sec:proof-main}.
\end{lemma}

\begin{lemma}[Independence over a product space]\label{lem:indep}
  \lean{ThreeMajority.avg_mul_prod, ThreeMajority.avg_fst_mul, ThreeMajority.avg_snd_mul,
    ThreeMajority.avg_equiv, ThreeMajority.avg_prod_pi, ThreeMajority.avg_eval}
  \leanok
  \uses{def:avg}
  \textbf{New relative to the push protocol}, since that development only
  ever averaged over a single uniform draw at a time. For finite types
  $\beta,\delta$ and $g \colon \beta\to\R$, $h\colon\delta\to\R$:
  $\mathrm{avg}(p \mapsto g(p_1)h(p_2)) = \mathrm{avg}(g)\,\mathrm{avg}(h)$
  over the product $\beta\times\delta$ (\texttt{avg\_mul\_prod}), with
  marginal consequences \texttt{avg\_fst\_mul}/\texttt{avg\_snd\_mul}.
  Iterating over a $\mathrm{Fin}\ n$-indexed product
  (\texttt{avg\_prod\_pi}): for $f \colon \mathrm{Fin}\ n \to \gamma \to \R$,
  \[
    \mathrm{avg}\Bigl(x \mapsto \textstyle\prod_i f_i(x_i)\Bigr)
      \;=\; \prod_i \mathrm{avg}(f_i)
  \]
  over $x \colon \mathrm{Fin}\ n \to \gamma$ --- the discrete,
  measure-theory-free form of ``coordinate projections on a product space
  are independent,'' with a one-coordinate marginal corollary
  (\texttt{avg\_eval}). This is exactly the independence fact that lets the
  $n$ agents' per-round updates be combined into a single concentration
  bound (Section~\ref{sec:chernoff}).
\end{lemma}

\section{The majority map}\label{sec:map}

Write $x = X_t/n \in [0,1]$ for the opinion-$1$ fraction, and
$\mathrm{ind}(I)$ for the real-valued membership indicator of $I$, so
$\mathrm{avg}(\mathrm{ind}(I)) = |I|/n$.

\begin{lemma}[Majority-of-three as a polynomial]\label{lem:majpoly}
  \lean{ThreeMajority.ind, ThreeMajority.avg_ind, ThreeMajority.avg_ind_eq,
    ThreeMajority.majorityIndicator_eq}
  \leanok
  \uses{def:avg, def:step}
  For $a,b,c \in \{0,1\}$: $\mathbf 1[a+b+c\ge2] = ab+bc+ac-2abc$. Applied
  pointwise to the three samples of a fixed agent, this expresses the
  majority indicator as a polynomial in the three membership indicators.
\end{lemma}

\begin{definition}[Per-agent contribution]\label{def:Ymaj}
  \lean{ThreeMajority.Y_maj, ThreeMajority.Y_maj_zero_one, ThreeMajority.card_step_eq_sum}
  \leanok
  \uses{lem:majpoly}
  $Y_{\mathrm{maj}}(I,v,s) \in \{0,1\}$ is $1$ iff a majority of the sample
  triple $s$ lies in $I$ (independent of $v$, kept as a function of $v$ to
  match the shape the Chernoff bound of Section~\ref{sec:chernoff} expects).
  Then $|\mathrm{step}(I,r)| = \sum_v Y_{\mathrm{maj}}(I,v,r(v))$
  (\texttt{card\_step\_eq\_sum}): the per-agent, $\{0,1\}$-valued
  decomposition that Section~\ref{sec:chernoff}'s tail bounds need.
\end{definition}

\begin{lemma}[Exact one-round drift]\label{lem:drift}
  \lean{ThreeMajority.avg_ind_mul_fst_snd1, ThreeMajority.avg_ind_mul_snd1_snd2,
    ThreeMajority.avg_ind_mul_fst_snd2, ThreeMajority.avg_ind_mul_triple,
    ThreeMajority.avg_Y_maj_eq, ThreeMajority.sum_avg_Y_maj, ThreeMajority.avg_card_step}
  \leanok
  \uses{def:Ymaj, lem:indep}
  Conditionally on a fixed $I$ with $|I|=m$, the next count
  $X_{t+1} = |\mathrm{step}(I,r)|$ satisfies
  \[
    \E_r[X_{t+1}] \;=\; n\,p(x), \qquad p(x) := 3x^2-2x^3, \qquad x = m/n,
  \]
  proved exactly (no approximation) via Lemma~\ref{lem:majpoly},
  Definition~\ref{def:Ymaj} and the independence toolkit of
  Lemma~\ref{lem:indep} (each pairwise/triple product of membership
  indicators over distinct coordinates of the product space
  $\mathrm{Tgt3}\ n$ factors by \texttt{avg\_mul\_prod}).
\end{lemma}

\begin{definition}[Dissent contribution]\label{def:Ydis}
  \lean{ThreeMajority.Y_dis, ThreeMajority.Y_dis_zero_one, ThreeMajority.card_dissent_eq_sum,
    ThreeMajority.sum_avg_Y_dis}
  \leanok
  \uses{def:Ymaj, lem:drift}
  $Y_{\mathrm{dis}} := 1 - Y_{\mathrm{maj}}$, the complementary per-agent
  dissent contribution, so $n - |\mathrm{step}(I,r)| = \sum_v
  Y_{\mathrm{dis}}(I,v,r(v))$ with mean $n - n\,p(x)$
  (\texttt{sum\_avg\_Y\_dis}); by the algebraic identity
  $1-p(x) = (1-x)^2(1+2x)$ this is the quantity controlled in
  Section~\ref{sec:saturation}.
\end{definition}

\begin{remark}
  The algebraic identities for $p$ used throughout below --- $p(x)-x =
  x(1-x)(2x-1)$, $1-p(x) = (1-x)^2(1+2x)$, and the growth/saturation
  regional bounds $p(\tfrac12+\beta)-\tfrac12 \ge \tfrac54\beta$ for
  $0\le\beta\le\tfrac14$ and $u(3-2u)\le\tfrac58$ for $0\le u\le\tfrac14$ ---
  are not isolated as separate Lean lemmas: each use site
  (\texttt{growth\_round}, \texttt{saturation\_round\_closed},
  \texttt{saturation\_mean\_quad}) re-derives the specific polynomial
  inequality it needs inline via \texttt{ring}/\texttt{nlinarith}, in the
  same spirit as the push protocol's reverse-Markov bound, which is also
  not isolated as its own named lemma.
\end{remark}

\section{Elementary inequalities}\label{sec:ineq}

\begin{lemma}[Quadratic-tight logarithm bound]\label{lem:logquad}
  \lean{ThreeMajority.log_ge_quadratic}
  \leanok
  For $\tfrac12 \le x \le 1$: $(x-1)-(x-1)^2 \le \log x$. (Mathlib's crude
  linear bound $\log x \ge x-1$ is false for $x<1$ and, worse, plugging the
  correct one-sided bound $1-\tfrac1x\le\log x$ into the Chernoff exponent
  $k-\mu-k\log(k/\mu)$ gives exactly $0$ when $k/\mu\to1$ --- useless. This
  quadratic correction, derived from Mathlib's degree-$3$ Taylor lower bound
  for $\exp$ by elementary algebra, is what makes the growth-phase Chernoff
  bound (Lemma~\ref{lem:growthround}) bite.)
\end{lemma}

\begin{lemma}[Exponential beats any fixed power]\label{lem:exppow}
  \lean{Real.pow_div_factorial_le_exp, ThreeMajority.exp_ge_cube}
  \leanok
  For $x\ge0$ and any $k\in\N$: $x^k/k! \le \exp(x)$ --- the single
  degree-$k$ term of the Taylor truncation, extracted by dropping every
  other (nonnegative, since $x\ge0$) term of the sum. The cubic case
  $k=3$ (\texttt{exp\_ge\_cube}) is the special case used by
  Lemma~\ref{lem:logquad}. Choosing $k$ large enough relative to a given
  lower bound on $\log n$ lets $n$ be shown to dominate any fixed
  polynomial in $\log n$ using only a \emph{modest} such lower bound
  (Stage 2b, Lemma~\ref{lem:stage2b}, and the final numeric bound
  Lemma~\ref{lem:mainfailclean}) --- a single low-degree truncation, such
  as the cubic case alone, would instead force $\log n$ to be
  astronomically large.
\end{lemma}

\begin{lemma}[Tangent-line log bound]\label{lem:logtangent}
  \lean{ThreeMajority.log_le_tangent_div, ThreeMajority.exp_three_ge_twenty}
  \leanok
  For $L>0$ and any reference point $c>0$: $\log L \le \log c - 1 + L/c$,
  from the crude bound $\log t \le t - 1$ applied to $t = L/c$; tight at
  $L=c$ (equality there), unlike a bound tangent at a single globally-fixed
  point, so choosing $c$ near the $L$ actually in play gives a far sharper
  bound. Also records $\exp(3)\ge20$ (cubing the standard $10$-digit lower
  bound on $e$), used to instantiate $c=\exp(3)$ below.
\end{lemma}

\begin{lemma}[Log-log bound]\label{lem:loglog}
  \lean{ThreeMajority.log_log_le_of_pos}
  \leanok
  \uses{lem:logtangent}
  For $\log x>0$: $\log\log x \le 2 + (\log x)/20$, from
  Lemma~\ref{lem:logtangent} tangent at the reference point
  $c=\exp(3)\approx20.09$ (so $\log c = 3$ exactly), combined with
  $\exp(3)\ge20$ to replace $1/\exp(3)$ by the cruder but rational $1/20$.
  Far tighter, once $\log x$ is more than a few units above $e$, than a
  bound tangent at the fixed point $c=e$ (the crude bound $\log t\le t-1$
  applied directly at $t=(\log x)/e$) --- exactly the regime
  Lemma~\ref{lem:stage2b} needs. Used only inside Lemma~\ref{lem:stage2b}'s
  numeric estimate.
\end{lemma}

\section{An elementary Chernoff bound}\label{sec:chernoff}

The only concentration tool used in the whole development, and the only
place beyond elementary combinatorics where real exponentials appear. These
bounds are proved once in the shared \texttt{dynamics} library
(\texttt{Dynamics/Chernoff.lean}, namespace \texttt{Dynamics}), which the
plurality and median developments use as well.

\begin{lemma}[MGF bound]\label{lem:mgf}
  \lean{Dynamics.avg_chernoff_mgf}
  \leanok
  \uses{lem:indep}
  Let $Y_1,\dots,Y_n \colon \gamma \to \{0,1\}$ be the per-agent
  contributions on the product space $\mathrm{Fin}\ n \to \gamma$ (one
  independent draw of $\gamma$ per agent), with $\mu=\sum_i\mathrm{avg}(Y_i)$
  and $X(x) := \sum_i Y_i(x_i)$. Then for every real $t$ (no sign
  restriction):
  \[
    \E[\exp(tX)] \;\le\; \exp\bigl(\mu(e^t-1)\bigr).
  \]
  Proved via the independence toolkit (Lemma~\ref{lem:indep}) applied to
  $x \mapsto \prod_i \exp(tY_i(x_i))$, and the elementary bound
  $1+y\le e^y$ (Mathlib's \texttt{Real.add\_one\_le\_exp}) at each factor.
\end{lemma}

\begin{lemma}[Chernoff tail bounds]\label{lem:chernoff}
  \lean{Dynamics.avg_chernoff_upper_of_mgf, Dynamics.avg_chernoff_lower_of_mgf}
  \leanok
  \uses{lem:mgf}
  In the setting of Lemma~\ref{lem:mgf}: for $t\ge0$ and any $k$,
  $\Pr[X\ge k] \le \exp(\mu(e^t-1)-tk)$ (\texttt{avg\_chernoff\_upper\_of\_mgf}, Markov
  applied to $\exp(tX)$); symmetrically for $t\le0$,
  $\Pr[X\le k] \le \exp(\mu(e^t-1)-tk)$ (\texttt{avg\_chernoff\_lower\_of\_mgf}).
\end{lemma}

\begin{lemma}[Closed-form tail bounds at the optimal $t$]\label{lem:chernofflog}
  \lean{Dynamics.avg_chernoff_upper_log, Dynamics.avg_chernoff_lower_log}
  \leanok
  \uses{lem:chernoff}
  Plugging $t=\log(k/\mu)$ into Lemma~\ref{lem:chernoff}: for $k\ge\mu>0$,
  \[
    \Pr[X\ge k] \;\le\; \exp\bigl(k-\mu-k\log(k/\mu)\bigr)
    \qquad(\mu=\E X),
  \]
  and symmetrically for $0<k\le\mu$,
  $\Pr[X\le k] \le \exp(k-\mu-k\log(k/\mu))$.
  Both closed forms are \emph{mean-scaled}: the exponent depends on the true
  $\mu$ directly, not on a crude worst-case range/variance proxy, which
  matters late in the saturation phase once few at-risk agents remain
  (Lemma~\ref{lem:stage2b}). The Lean statements are the generalized
  variants, which take only an \emph{upper} bound $\mu_{\mathrm{ub}}\ge\mu$
  (\texttt{avg\_chernoff\_upper\_log}, used by the saturation phase) or a
  \emph{lower} bound $\mu_{\mathrm{lb}}\le\mu$ (\texttt{avg\_chernoff\_lower\_log}, used by
  the growth phase) on the true mean (the exact forms are the case of equality) --- monotonicity of the exponent in $\mu$ at the optimal $t$
  makes both directions sound, and this is the form each phase actually
  needs, since the exact mean depends on $x=|I|/n$ while the phases only
  assume an inequality on $|I|$.
\end{lemma}

\section{Growth phase}\label{sec:growth}

Set $\beta_j = \tfrac1{10}(1.1)^j$ for $j=0,\dots,10$; $\beta_0=0.1$ (the
hypothesis $X_0\ge0.6n$) and $\beta_{10}>0.25$.

\begin{definition}[Bias schedule]\label{def:growthbias}
  \lean{ThreeMajority.growthBias, ThreeMajority.growthBias_zero, ThreeMajority.growthBias_nonneg,
    ThreeMajority.growthBias_mono, ThreeMajority.growthBias_le_of_le,
    ThreeMajority.growthBias_le_quarter, ThreeMajority.growthBias_ten}
  \leanok
  $\mathrm{growthBias}(j) := \tfrac1{10}(\tfrac{11}{10})^j$, the deterministic
  bias target after $j$ growth rounds; monotone, $\le\tfrac14$ for $j\le9$,
  and $>\tfrac14$ at $j=10$.
\end{definition}

\begin{lemma}[Single growth round]\label{lem:growthround}
  \lean{ThreeMajority.growth_round}
  \leanok
  \uses{lem:drift, lem:chernofflog, def:growthbias}
  Starting from opinion-$1$ fraction $\ge\tfrac12+\beta$
  ($\tfrac1{10}\le\beta\le\tfrac14$), the probability of \emph{failing} to
  reach the next bias target (guaranteed growth factor $\tfrac{11}{10}$,
  comfortably below the true drift factor $\tfrac54$ from
  Lemma~\ref{lem:drift}) is at most $\exp(-n/10^7)$. Proved by computing a
  lower bound $\mu_{\mathrm{lb}} = n(\tfrac12+\tfrac54\beta)$ on the true
  mean via monotonicity of $p$ on $[\tfrac12+\beta,1]$ (an inline
  \texttt{nlinarith} computation, not a separate lemma), then applying the
  lower-tail bound \texttt{avg\_tail\_le\_log\_ge} at
  $k=n(\tfrac12+\tfrac{11}{10}\beta)$, and bounding the resulting exponent
  using Lemma~\ref{lem:logquad}.
\end{lemma}

\begin{lemma}[Union bound over growth rounds]\label{lem:growthfail}
  \lean{ThreeMajority.growth_fail_le}
  \leanok
  \uses{lem:growthround}
  Starting from bias $\mathrm{growthBias}(i)$, after $j$ further rounds
  (with $i+j\le10$) the probability of failing to reach
  $\mathrm{growthBias}(i+j)$ is at most $j\exp(-n/10^7)$, by induction on
  $j$ via the $\mathrm{expList}$-conditioning idiom (mirroring the push
  protocol's \texttt{Main.lean}).
\end{lemma}

\begin{lemma}[Phase 1]\label{lem:phase1}
  \lean{ThreeMajority.growth_phase1}
  \leanok
  \uses{lem:growthfail}
  Starting from opinion-$1$ fraction $\ge\tfrac35$, after $10$ rounds the
  opinion-$1$ fraction exceeds $\tfrac34$, except with probability
  $\le10\exp(-n/10^7)$.
\end{lemma}

\section{Saturation phase}\label{sec:saturation}

Write $U_t = n-X_t$. By Lemma~\ref{lem:phase1}, $U_{10}\le n/4$ except with
probability $\le10\exp(-n/10^7)$. Getting $U_t$ all the way to
\emph{exactly} $0$ needs three stages, since a fixed-ratio Chernoff bound
needs the mean itself to be $\Omega(\log n)$ for a polynomially-small
per-round failure probability.

\begin{lemma}[General per-round upper-tail bound]\label{lem:satclosed}
  \lean{ThreeMajority.saturation_round_closed, ThreeMajority.saturation_round_generic}
  \leanok
  \uses{def:Ydis, lem:chernofflog}
  Given a reference dissent bound $U_t\le M\le n/4$ and any target
  $k\ge\tfrac58M$ (an upper bound $\mu_{\mathrm{ub}}=\tfrac58M$ on the true
  one-round mean, via $1-p(x)=(1-x)^2(1+2x)\le\tfrac58(1-x)$ for
  $x\ge\tfrac34$), the probability the next dissent count exceeds $k$ is at
  most $\exp(k-\tfrac58M-k\log(k/(\tfrac58M)))$
  (\texttt{saturation\_round\_closed}); the same conclusion holds for any
  directly-supplied upper bound $\mu_{\mathrm{ub}}$ on the true mean
  (\texttt{saturation\_round\_generic}, used by Stage 2b where the
  quadratic estimate below is the tight one).
\end{lemma}

\begin{lemma}[Quadratic mean bound]\label{lem:satquad}
  \lean{ThreeMajority.saturation_mean_quad}
  \leanok
  \uses{def:Ydis}
  For any dissent bound $U_t\le M\le n$, the true one-round mean is at most
  $3M^2/n$ (the crude bound $1-p(x)\le3(1-x)^2$, dropping the $\tfrac58$
  refinement). Used both by Stage 2b and by Stage 2c (at $M=10$).
\end{lemma}

\begin{lemma}[Numeric per-round bound]\label{lem:satnumeric}
  \lean{ThreeMajority.saturation_round_numeric}
  \leanok
  \uses{lem:satclosed, lem:logquad}
  Lemma~\ref{lem:satclosed}'s exponent, simplified via Mathlib's
  quadratic-tight Pad\'e bound for $\log$
  (\texttt{Real.le\_log\_one\_add\_of\_nonneg}) to the closed form
  $-\lambda^2/(k+\tfrac58M)$, $\lambda=k-\tfrac58M$.
\end{lemma}

\begin{lemma}[Stage 2a per-round contraction]\label{lem:satcontract}
  \lean{ThreeMajority.saturation_round_contract}
  \leanok
  \uses{lem:satnumeric}
  Guaranteed contraction factor $0.7$ (vs.\ the true $\tfrac58$): failure
  probability $\le\exp(-M/250)$.
\end{lemma}

\begin{definition}[Deterministic dissent target]\label{def:satafter}
  \lean{ThreeMajority.satAfter, ThreeMajority.satAfter_ge_floor, ThreeMajority.satAfter_zero,
    ThreeMajority.satAfter_le_of_le, ThreeMajority.satAfter_contract_le}
  \leanok
  $\mathrm{satAfter}(n,M,i) := \max(500\log n,\ (0.7)^i M)$, the target
  dissent bound after $i$ Stage 2a rounds, clamped below at the floor
  $500\log n$; written in closed form (not recursively) so shifts by one
  round agree definitionally on the unclamped branch.
\end{definition}

\begin{lemma}[Stage 2a, union bound]\label{lem:satfail}
  \lean{ThreeMajority.saturation_fail_le}
  \leanok
  \uses{lem:satcontract, def:satafter}
  Starting from a dissent bound $M$ ($500\log n\le M\le n/4$), offset $i$
  rounds in, the probability the dissent count exceeds
  $\mathrm{satAfter}(n,M,i+j)$ after $j$ further rounds is at most
  $j\,n^{-2}$, by induction on $j$ via the same $\mathrm{expList}$-
  conditioning idiom as Lemma~\ref{lem:growthfail}, with the tail direction
  flipped (dissent shrinking rather than opinion-$1$ count growing).
\end{lemma}

\begin{lemma}[Stage 2a, wrapped]\label{lem:stage2a}
  \lean{ThreeMajority.T2a, ThreeMajority.satAfter_quarter_le_floor, ThreeMajority.saturation_stage2a}
  \leanok
  \uses{lem:satfail}
  With $T_{2a}(n) := \ceil{6\log n}$: starting from any dissent bound
  $\le n/4$, after $T_{2a}(n)$ rounds the dissent count exceeds the floor
  $500\log n$ with probability at most $T_{2a}(n)\cdot n^{-2}$.
\end{lemma}

\begin{lemma}[Stage 2b]\label{lem:stage2b}
  \lean{ThreeMajority.saturation_stage2b}
  \leanok
  \uses{lem:satquad, lem:satclosed, lem:exppow, lem:loglog}
  Provided $\log n\ge30$: whenever $U_t\le500\log n$,
  $\Pr[U_{t+1}>10] \le \exp(-\log n) = 1/n$. One round, jumping from the
  $\Theta(\log n)$ floor directly to the fixed constant $10$: since the
  mean is already tiny ($\Theta((\log n)^2/n)$), this single round has an
  astronomically small failure probability once $n$ is large enough to beat
  the relevant polynomial in $\log n$, which a degree-$11$ instance of
  Lemma~\ref{lem:exppow}, combined with the tighter tangent-line bound of
  Lemma~\ref{lem:loglog}, already secures at this modest threshold.
\end{lemma}

\begin{lemma}[Stage 2c]\label{lem:stage2c}
  \lean{ThreeMajority.saturation_stage2c}
  \leanok
  \uses{lem:satquad}
  Whenever $U_t\le10$: $\Pr[U_{t+1}\ge1] \le 300/n$, via ordinary Markov
  applied to the nonnegative-integer $U_{t+1}$ ($\E[U_{t+1}]\le300/n$ by
  Lemma~\ref{lem:satquad} at $M=10$) --- no concentration needed for the
  very last step.
\end{lemma}

\section{Proof of the main theorem}\label{sec:proof-main}

\begin{lemma}[Stages 2b+2c combined]\label{lem:sat2b2c}
  \lean{ThreeMajority.ne_univ_iff_pos_dissent, ThreeMajority.saturation_2b2c}
  \leanok
  \uses{lem:stage2b, lem:stage2c}
  Two rounds, from $U_t\le500\log n$ to \emph{exactly} $U=0$, with total
  failure probability $\le\exp(-\log n)+300/n$, by chaining
  Lemma~\ref{lem:stage2b} then Lemma~\ref{lem:stage2c} via
  $\mathrm{expList}$-conditioning.
\end{lemma}

\begin{lemma}[All three saturation stages combined]\label{lem:sat2a2b2c}
  \lean{ThreeMajority.saturation_2a2b2c}
  \leanok
  \uses{lem:stage2a, lem:sat2b2c}
  $T_{2a}(n)+2$ rounds, from $U_t\le n/4$ to exactly $U=0$, with total
  failure probability $\le T_{2a}(n)\cdot n^{-2} + \exp(-\log n) + 300/n$.
\end{lemma}

\begin{theorem}[Main failure bound, raw form]\label{lem:mainfail}
  \lean{ThreeMajority.majority3_consensus_fail_le}
  \leanok
  \uses{lem:phase1, lem:sat2a2b2c}
  For $\log n\ge30$ and $|I_0|\ge\tfrac35n$: over $10+(T_{2a}(n)+2)$
  rounds,
  \[
    \Pr[I_T\ne V] \;\le\; 10e^{-n/10^7}
      + \bigl(T_{2a}(n)\,n^{-2} + e^{-\log n} + 300/n\bigr),
  \]
  by gluing Lemma~\ref{lem:phase1} (growth) to
  Lemma~\ref{lem:sat2a2b2c} (saturation), exactly as the push protocol's
  \texttt{prNotAllInformed\_le} glues its own two phases.
\end{theorem}

\begin{theorem}[Clean numeric bound]\label{lem:mainfailclean}
  \lean{ThreeMajority.hfloor4_of_hbig, ThreeMajority.majority3_consensus_fail_le_clean}
  \leanok
  \uses{lem:mainfail, lem:exppow}
  Under the same hypotheses, the four error terms of
  Theorem~\ref{lem:mainfail} are each bounded by $O(1/n)$ using
  Lemma~\ref{lem:exppow} (to show $n$ dominates the relevant polynomials in
  $\log n$ once $\log n\ge30$), giving the single clean bound
  $\Pr[I_T\ne V] \le 500/n$.
\end{theorem}

With $T = 10+(T_{2a}(n)+2)$, Theorem~\ref{thm:main}
(\texttt{majority3\_consensus\_whp}) follows from
Theorem~\ref{lem:mainfailclean} by
$\Pr[I_T=V] = 1-\Pr[I_T\ne V] \ge 1-500/n$.

\section{\texorpdfstring{3-Majority from any configuration in $O(n^{3/4}\log^{7/8} n)$ rounds}{3-Majority from any configuration}}\label{sec:bcekmn-any-start}

This section formalizes the upper bound for 3-Majority with any number of colours of
Berenbrink, Clementi, Elsässer, Kling, Mallmann-Trenn and Natale, \emph{Ignore or comply? On
breaking symmetry in consensus} (PODC 2017) \cite{BCEKMN17}: Theorem~4 of arXiv:1702.04921 (v1),
with the comparison framework of its Section~2 and its Lemmas~2 to~4
(\href{https://github.com/formal-dynamics/leanamics/blob/main/3-majority/FORMALIZATION_DIFFERENCES.md#3-majority-from-any-configuration-anystart-roadmap-maj-6-b}{differences
from the paper}).

\textbf{The second phase.} The proof of Theorem~4 applies Theorem~3.1 of Becchetti, Clementi,
Natale, Pasquale and Trevisan \cite{BCNPT16}, which is not part of \cite{BCEKMN17}.
Definition~\ref{def:as-phase2} states that result, Theorem~\ref{thm:as-main} takes it as a
hypothesis, exactly as the paper's proof uses it, and Section~\ref{sec:bcnpt16-many-opinions}
(roadmap MAJ-12~(a)) proves it (Theorem~\ref{thm:mo-phase2}), so that
Theorem~\ref{thm:as-main-unconditional}, Theorem~4 of \cite{BCEKMN17}, has no hypothesis.

\subsection{Model and majorization}

Agents $V=\{0,\dots,n-1\}$ hold colours in a finite set $\Sigma$; a configuration is a map
$c\colon V\to\Sigma$, with $c_a$ agents of colour $a$ and $\kappa(c)=|\{a : c_a>0\}|$ remaining
colours. Consensus is $\kappa(c)\le1$.

\begin{definition}[3-Majority and Voter with any number of colours]\label{def:as-model}
\lean{ThreeMajority.majColour,ThreeMajority.stepCol,ThreeMajority.runCol,ThreeMajority.voterStep,ThreeMajority.voterRun,ThreeMajority.colourCount,ThreeMajority.numColours}\leanok
\uses{def:round}
In a round of 3-Majority every agent draws an ordered sample triple $(s_1,s_2,s_3)$ uniformly
with replacement (Definition~\ref{def:round}); if $s_2$ and $s_3$ have the same colour it adopts
that colour, otherwise the colour of $s_1$. In a round of Voter every agent adopts the colour of
one uniformly random agent (itself included).
\end{definition}

The paper breaks a three-way tie by a uniformly random sample; since the samples are i.i.d.,
taking the first one gives the same law (Lemma~\ref{lem:as-alpha}).

\begin{lemma}[Two colours]\label{lem:as-bool}
\lean{ThreeMajority.stepCol_bool}\leanok\uses{def:as-model,def:step}
With colours $\{0,1\}$, a round of 3-Majority is the step of Definition~\ref{def:step}.
\end{lemma}
\begin{proof}\leanok
Case analysis on the colours of the three samples.
\end{proof}

\begin{definition}[Majorization, Schur-convex observables]\label{def:as-majorization}
\lean{ThreeMajority.Majorizes,ThreeMajority.countVec,ThreeMajority.SchurConvex}\leanok
\uses{def:as-model}
For $x,y\in\R^\Sigma$, $x\succeq y$ if $\sum x=\sum y$ and every sum of $y$ over a set
$S\subseteq\Sigma$ is at most a sum of $x$ over a set of the same size (equivalently, for every $j$
the $j$ largest entries of $y$ sum to at most the $j$ largest entries of $x$). An observable $\varphi$
of configurations is Schur-convex if $(c_a)_a\succeq(c'_a)_a$ implies $\varphi(c')\le\varphi(c)$.
\end{definition}

\begin{lemma}[Majorization and colours]\label{lem:as-colours}
\lean{ThreeMajority.numColours_le_of_majorizes,ThreeMajority.schurConvex_numColours_le}\leanok
\uses{def:as-majorization}
If $(c_a)_a\succeq(c'_a)_a$ then $\kappa(c)\le\kappa(c')$. Hence the indicator of
$\kappa(c)\le\kappa$ is Schur-convex.
\end{lemma}
\begin{proof}\leanok
The colours of $c'$ carry all $n$ agents, so a set of $\kappa(c')$ colours carries all agents of $c$.
\end{proof}

\subsection{Comparison of anonymous consensus processes}

\begin{definition}[AC-processes and dominance, Definitions~1 and~2]\label{def:as-ac}
\lean{ThreeMajority.acKernel,ThreeMajority.Dominates}\leanok\uses{def:as-majorization}
An anonymous consensus process with process function $\alpha$ moves from $c$ by letting every
agent draw its new colour independently from the probability vector $\alpha(c)$. The process
$\alpha$ dominates $\alpha'$ if $c\succeq c'$ implies $\alpha(c)\succeq\alpha'(c')$.
\end{definition}

\begin{proposition}[Proposition~1, Rinott]\label{prop:as-rinott}
\lean{ThreeMajority.multinomial_schurConvex,ThreeMajority.expect_le_of_transfer,ThreeMajority.expect_le_of_partial_sums,ThreeMajority.expect_le_of_majorizes}\leanok
\uses{def:as-majorization}
If $p\succeq q$ are probability vectors on $\Sigma$ and $\varphi$ is Schur-convex, the
expectation of $\varphi$ when the $n$ agents draw their colours i.i.d.\ from $q$ is at most the
expectation when they draw from $p$.
\end{proposition}
\begin{proof}\leanok
The paper cites this result without proof. For one Robin Hood transfer between colours $a$ and
$b$, the expectation, as a function of the mass $\theta$ of $a$ (with $a$ and $b$ sharing a fixed
total $s$), is nondecreasing on $[s/2,s]$: its derivative is a sum over agents, and swapping $a$
and $b$ for all agents and then for one agent pairs its terms into products of two factors with
the same sign. A chain of transfers (the Hardy--Littlewood--P\'olya construction) leads from $p$
to $q$ when $q$ is sorted and the partial sums of $p$ dominate those of $q$; sorting and
relabelling the colours, under which Schur-convex observables are invariant, reduce the general
case to this one.
\end{proof}

\begin{theorem}[Theorem~2, distributional form]\label{thm:as-comparison}
\lean{ThreeMajority.ac_comparison,ThreeMajority.ac_numColours}\leanok
\uses{def:as-ac,prop:as-rinott,lem:as-colours}
If $\alpha$ dominates $\alpha'$ and $c\succeq c'$, then for every Schur-convex $\varphi$ and every
$T$, $\E_{c'}[\varphi(P'_T)]\le\E_c[\varphi(P_T)]$. In particular, from the same configuration,
the dominating process has at most $\kappa$ colours at time $T$ with at least the probability of
the dominated one.
\end{theorem}
\begin{proof}\leanok
Induction on $T$. With $u_T$, $v_T$ the expectations at time $T$ for the dominated and the
dominating process, $w(y)=\max\{u_T(z):z\preceq y\}$ is Schur-convex and $u_T\le w\le v_T$ on
comparable configurations, so Proposition~\ref{prop:as-rinott} applied to $w$ gives the step
from $T$ to $T+1$. The second statement is the case of the indicator of
Lemma~\ref{lem:as-colours}. The paper's coupling (Lemma~1, from Strassen's theorem) is not
needed.
\end{proof}

\subsection{3-Majority is at least as fast as Voter}

\begin{lemma}[Equations~(1) and~(2)]\label{lem:as-alpha}
\lean{ThreeMajority.alpha3M,ThreeMajority.alphaVoter,ThreeMajority.alphaVoter_weight,ThreeMajority.alpha3M_weight,ThreeMajority.apply_ofStep_stepCol,ThreeMajority.apply_ofStep_voterStep}\leanok
\uses{def:as-model,def:as-ac}
Both processes are AC-processes. With $x=c/n$, Voter has $\alpha^V_a(c)=x_a$ and 3-Majority has
$\alpha^{3M}_a(c)=x_a(1+x_a-\|x\|_2^2)$.
\end{lemma}
\begin{proof}\leanok
The agents' samples are independent. For 3-Majority,
$[\mathrm{maj}=a]=[x_2=a][x_3=a]+[x_1=a]-[x_2=x_3][x_1=a]$, and the number of pairs of agents of
equal colour is $\sum_b c_b^2$.
\end{proof}

\begin{lemma}[Lemma~2]\label{lem:as-vs-voter}
\lean{ThreeMajority.dominates_alpha3M_alphaVoter,ThreeMajority.voter_le_threeMaj}\leanok
\uses{lem:as-alpha,thm:as-comparison}
3-Majority dominates Voter. Hence, from the same configuration, after $T$ rounds 3-Majority has
at most $\kappa$ colours with at least the probability that Voter has.
\end{lemma}
\begin{proof}\leanok
For a set $S$ of the largest entries of $x$, $\|x\|_2^2\sum_S x\le\sum_S x^2$, since
$\sum_{i\in S}\sum_j x_ix_j(x_i-x_j)\ge0$; this gives $\alpha^{3M}(c)\succeq\alpha^V(c)$, and with
$c\succeq c'$ also $\alpha^{3M}(c)\succeq\alpha^V(c')$. Conclude with
Theorem~\ref{thm:as-comparison}.
\end{proof}

\subsection{The Voter bound}

\begin{lemma}[Lemma~4 and Equation~(6), duality]\label{lem:as-dual}
\lean{ThreeMajority.walkStep,ThreeMajority.numColours_voterRun_le,ThreeMajority.voter_dual}\leanok
\uses{def:as-model}
Let coalescing random walks start on all $n$ nodes of the complete graph with self-loops. After
$T$ rounds, Voter has more than $k$ colours with at most the probability that more than $k$
walks remain.
\end{lemma}
\begin{proof}\leanok
Pathwise, the colours of Voter after rounds $y_1,\dots,y_T$ are colours of the nodes occupied
by the walks moved by $y_T,\dots,y_1$; the rounds are i.i.d., so reversing their order does not
change the law. The paper proves an equality under a coupling and uses only this inequality.
\end{proof}

\begin{lemma}[Equation~(7) and the expected number of walks]\label{lem:as-walks}
\lean{ThreeMajority.walk_drift,ThreeMajority.walk_drift_paper,ThreeMajority.walk_expect_le}\leanok
\uses{lem:as-dual}
If $X_t$ walks remain, $\E[X_{t+1}\mid X_t=x]\le x-x(x-1)/(3n)$ for every $x$, hence
$\le x-x^2/(10n)$ for $x\ge2$ (Equation~(7)). From any set of walks, $\E X_t\le1+3n/t$ for
$t\ge1$.
\end{lemma}
\begin{proof}\leanok
$\E[X_{t+1}\mid X_t=x]=n(1-(1-1/n)^x)$, the expected number of occupied cells when $x$ balls go
into $n$ cells; bound $(1-1/n)^x$ by the Bonferroni inequality of degree~3. The map
$x\mapsto x-x(x-1)/(3n)$ is concave and nondecreasing on $[0,n+1]$, so its tangent-line bound and
induction on $t$ give $\E X_t\le1+3n/t$. This replaces the paper's variable drift argument
for $\E[T^k_C]\le20n/k$.
\end{proof}

\begin{lemma}[Lemma~3]\label{lem:as-voter-whp}
\lean{ThreeMajority.voter_reduce_whp}\leanok\uses{lem:as-walks}
For $n\ge2$, $k\ge1$ and $T\ge24(n/k)\log n$, from any configuration Voter has more than $k$
colours after $T$ rounds with probability at most $1/n$.
\end{lemma}
\begin{proof}\leanok
Blocks of $\lceil6n/k\rceil$ rounds: by Markov's inequality on $X-1$ (the walks never vanish),
each block leaves more than $k$ walks with probability at most $1/2$, from any start. Restart
over $\lfloor T/\lceil6n/k\rceil\rfloor\ge\log_2 n$ blocks (for $k<n$; for $k\ge n$ there is
nothing to prove) and use Lemma~\ref{lem:as-dual}.
\end{proof}

\subsection{The main theorem}

\begin{lemma}[Phase~1]\label{lem:as-phase1}
\lean{ThreeMajority.threeMaj_reduce_whp}\leanok\uses{lem:as-vs-voter,lem:as-voter-whp}
For $n\ge2$, $k\ge1$ and $T\ge24(n/k)\log n$, from any configuration 3-Majority has more than
$k$ colours after $T$ rounds with probability at most $1/n$.
\end{lemma}
\begin{proof}\leanok
Lemma~\ref{lem:as-vs-voter} and Lemma~\ref{lem:as-voter-whp}.
\end{proof}

\begin{definition}[Phase~2: the cited result]\label{def:as-phase2}
\lean{ThreeMajority.Bcnpt16Phase2}\leanok\uses{def:as-model}
For $\varepsilon>0$, $\mathrm{Phase2}(\varepsilon)$ is the proposition (Theorem~3.1 of
\cite{BCNPT16}, Theorem~8 of \cite{BCEKMN17}): there are $C,N$ such that for $n\ge N$, from any
configuration with colours in $\{0,\dots,n-1\}$ and $k\le n^{1/3-\varepsilon}$ colours,
3-Majority has not reached consensus after any
$T\ge C(k^2\sqrt{\log n}+k\log n)(k+\log n)$ rounds with probability at most $1/n$.
This proposition is a hypothesis of Theorem~\ref{thm:as-main}; Theorem~\ref{thm:mo-phase2}
proves it for every $\varepsilon>0$ (Section~\ref{sec:bcnpt16-many-opinions}). The cited paper's
``w.h.p.''\ is $1-O(n^{-\lambda})$ for some $\lambda>0$, which gives the form $1/n$ by repeating
the phase a constant number of times; the proof of Theorem~\ref{thm:mo-phase2} gives the form
$1/n$ directly.
\end{definition}

\begin{lemma}[The asymptotics of Theorem~4]\label{lem:as-asymptotics}
\lean{ThreeMajority.phase1Colours,ThreeMajority.anyStart_asymptotics}\leanok
For every $A\ge0$ there are $C,N$ such that for $n\ge N$ and
$k=\lfloor n^{1/4}\log^{1/8}n\rfloor$: $n\ge2$, $1\le k\le n^{1/3-1/24}$, and
$24(n/k)\log n+1+A(k^2\sqrt{\log n}+k\log n)(k+\log n)\le Cn^{3/4}\log^{7/8}n$.
\end{lemma}
\begin{proof}\leanok
From $\log y\le y^{1/4}$ and $\log^{1/8}y\le y^{1/24}$ for $y$ large.
\end{proof}

\begin{theorem}[Theorem~4 from Phase~2]\label{thm:as-main}
\lean{ThreeMajority.threeMaj_anyStart_consensus}\leanok
\uses{lem:as-phase1,def:as-phase2,lem:as-asymptotics}
Assume $\mathrm{Phase2}(1/24)$. There are $C,N$ such that for $n\ge N$, from any configuration
with colours in $\{0,\dots,n-1\}$, 3-Majority has not reached consensus after any
$T\ge Cn^{3/4}\log^{7/8}n$ rounds with probability at most $2/n$.
\end{theorem}
\begin{proof}\leanok
Let $k=\lfloor n^{1/4}\log^{1/8}n\rfloor$ and $T_1=\lceil24(n/k)\log n\rceil$. After $T_1$ rounds,
more than $k$ colours remain with probability at most $1/n$ (Lemma~\ref{lem:as-phase1});
otherwise the current configuration has at most $k\le n^{1/3-1/24}$ colours, and
$\mathrm{Phase2}(1/24)$ bounds the failure of the remaining $T-T_1$ rounds by $1/n$.
Lemma~\ref{lem:as-asymptotics} fits both phases in $Cn^{3/4}\log^{7/8}n$ rounds. The proof in the
paper needs a minor correction: it applies Lemma~3 with $k=n^{1/4}$, which gives
$O(n^{3/4}\log n)$ rounds for Phase~1; the choice $k=n^{1/4}\log^{1/8}n$ balances the two phases
and gives the bound as stated.
\end{proof}

\begin{theorem}[Theorem~4]\label{thm:as-main-unconditional}
\lean{ThreeMajority.threeMaj_anyStart_consensus_unconditional}\leanok
\uses{thm:as-main,thm:mo-phase2}
There are $C,N$ such that for $n\ge N$, from any configuration with colours in
$\{0,\dots,n-1\}$, 3-Majority has not reached consensus after any $T\ge Cn^{3/4}\log^{7/8}n$
rounds with probability at most $2/n$.
\end{theorem}
\begin{proof}\leanok
Theorem~\ref{thm:as-main}, whose hypothesis $\mathrm{Phase2}(1/24)$ is
Theorem~\ref{thm:mo-phase2} for $\varepsilon=1/24$.
\end{proof}

\section{3-Majority with many opinions}\label{sec:bcnpt16-many-opinions}

This section formalizes the adversary-free case of Becchetti, Clementi, Natale, Pasquale and
Trevisan, \emph{Stabilizing consensus with many opinions} (SODA 2016) \cite{BCNPT16}: its
Theorem~3.1, with its Lemmas~2.1 and~3.2 to~3.5 and Lemma~B.1, in the numbering of
arXiv:1508.06782 (v1). The
\href{https://github.com/formal-dynamics/leanamics/blob/main/3-majority/FORMALIZATION_DIFFERENCES.md#3-majority-with-many-opinions-manyopinions-roadmap-maj-12-a}{differences
from the paper} are listed in the repository. Theorem~3.1 is the proposition $\mathrm{Phase2}(\varepsilon)$ of
Definition~\ref{def:as-phase2} (Theorem~\ref{thm:mo-phase2}), so it makes Theorem~4 of
\cite{BCEKMN17} unconditional (Theorem~\ref{thm:as-main-unconditional}).

The process is the 3-Majority of Definition~\ref{def:as-model}, whose rule is, pointwise, the rule
of the paper (three samples uniform with repetition; on three different colours, the first one),
with colours in $\{0,\dots,n-1\}$. We write $c_a$ for the support of colour $a$, $C'_a$ for its
support after one round, $j$ or $k$ for the number of colours, and $m=16\sqrt{jn\log n}$. Where
the paper assumes $j\le n^{1/3-\varepsilon}$, the statements below assume $Cj^3\log n\le n$ for a
constant $C$, which is what the proofs use; it holds for $j\le n^{1/3-\varepsilon}$ and $n$
large. The proofs follow the paper's, with the deviations and the minor corrections noted in each
proof.

\subsection{One round}

\begin{lemma}[Lemma~2.1]\label{lem:mo-expect}
\lean{ThreeMajority.numColours_eq_card_image,ThreeMajority.colourCount_eq_card_filter,ThreeMajority.stepCol_apply,ThreeMajority.adoptInd,ThreeMajority.colourCount_stepCol_eq_sum,ThreeMajority.avg_colourCount_stepCol,ThreeMajority.avg_colourCount_stepCol_le}\leanok
\uses{def:as-model}
For any finite set of colours, any configuration $c$ and any colour $a$,
\[
  \E[C'_a]=c_a\Bigl(1+\frac{c_a}{n}-\sum_b\frac{c_b^2}{n^2}\Bigr)
  \le c_a\Bigl(1+\frac{c_a}{n}-\frac{1}{\kappa(c)}\Bigr).
\]
\end{lemma}
\begin{proof}\leanok
$C'_a$ is the sum over the agents of independent indicators of adopting $a$, each with mean
$x_a(1+x_a-\|x\|_2^2)$ for $x=c/n$ (as in Lemma~\ref{lem:as-alpha}). The inequality is the
Cauchy--Schwarz bound $\sum_b c_b^2\ge n^2/\kappa(c)$.
\end{proof}

\begin{lemma}[One-round tails and anti-concentration]\label{lem:mo-tails}
\lean{ThreeMajority.avg_colourCount_stepCol_ge,ThreeMajority.avg_colourCount_stepCol_le_tail,ThreeMajority.avg_pz_zero_one,ThreeMajority.avg_abs_colourCount_stepCol_sub_ge,ThreeMajority.sum_variance_adoptInd_ge}\leanok
\uses{lem:mo-expect,lem:chernoff,lem:indep}
Let $\sigma_a^2$ be the sum of the variances of the agents' indicators of adopting $a$.
\begin{enumerate}
  \item If $\E[C'_a]\le\mu<k$, then $\Pr(C'_a\ge k)\le\exp(-(k-\mu)^2/(k+\mu))$; if
    $0<k<\mu\le\E[C'_a]$, then $\Pr(C'_a\le k)\le\exp(-(\mu-k)^2/(2\mu))$.
  \item If $\sigma_a^2\ge1$, then $\E|C'_a-\theta|\ge(9/256)\,\sigma_a$ for every $\theta\in\R$.
  \item If $c$ has $j\ge2$ colours, all with more than $n/j-nu$ agents, where $u\ge0$ and
    $ju\le1/250$, then $\sigma_a^2\ge n/(4j)$ for every colour $a$ of $c$.
\end{enumerate}
\end{lemma}
\begin{proof}\leanok
(1) The Chernoff bounds of Section~\ref{sec:chernoff} for a sum of independent indicators.
(2) If $|\E[C'_a]-\theta|\ge\sigma_a/4$, this is Jensen's inequality. Otherwise, by the
Paley--Zygmund inequality (from the second and fourth moments of a sum of independent centred
$\{0,1\}$-coordinates) the centred sum $S$ satisfies $\sigma_a^2\le4S^2$ with probability at least
$9/64$, and there $|C'_a-\theta|\ge\sigma_a/2-\sigma_a/4$. (3) The adoption probability $p$ of
Lemma~\ref{lem:mo-expect} lies between $(1-1/250)^2/j$ and $(1/2+1/250)(1+1/250)$, and
$\sigma_a^2=np(1-p)$.
\end{proof}

\subsection{Hitting times under a drift}

\begin{lemma}[Lemma~3.2, finite horizon]\label{lem:mo-drift}
\lean{ThreeMajority.hitProb_additive_drift}\leanok
Let $K$ be a Markov kernel on a finite set, $I$ a set with $K(x,I)=1$ for $x\in I$, $B$ a target
set with hitting time $\tau$, and $f$ a real function. Let $\beta,\eta\ge0$, and assume that
$f(x)\le M$ for $x\in I$ and that, for $x\in I\setminus B$, $f(x)\le\beta$,
$\E[f(X_1)\mid X_0=x]\ge f(x)+\lambda$ and $\Pr(f(X_1)>\beta\mid X_0=x)\le\eta$. Then for every
$a\in I$ with $f(a)\ge0$ and every $t$, $\lambda t\Pr_a(\tau>t)\le\beta+M\eta t$.
\end{lemma}
\begin{proof}\leanok
Let $X'$ be the chain stopped at $B$. Summing the drift over the steps before $\tau$,
$\E_a[f(X'_t)]\ge f(a)+\lambda\E_a[\min(t,\tau)]\ge\lambda t\Pr_a(\tau>t)$. By a union bound over
the steps, $f(X'_t)>\beta$ with probability at most $\eta t$, and $f\le M$, so
$\E_a[f(X'_t)]\le\beta+M\eta t$. The paper bounds $\E[\tau]$ by optional stopping, under a
hypothesis on the overshoot at time $\tau$ (Lemma~B.1); the finite horizon and the one-step
overshoot bound avoid both.
\end{proof}

\begin{lemma}[Hit, then absorb]\label{lem:mo-strong-markov}
\lean{ThreeMajority.hitProb_mul_le_event}\leanok
Let $P$ be an absorbing set for $K$. If from every state of $B$ the chain is in $P$ after $s$
steps with probability at least $q$, then from every state $a$ it is in $P$ at time $t+s$ with
probability at least $q\Pr_a(\tau\le t)$.
\end{lemma}
\begin{proof}\leanok
The strong Markov property at $\tau$, proved by induction on $t$ without stopping times: from a
state of $B$, the hypothesis and absorption; otherwise, condition on the first step.
\end{proof}

\subsection{Symmetry breaking}

\begin{definition}[Broken symmetry]\label{def:mo-broken}
\lean{ThreeMajority.BrokenSymmetry,ThreeMajority.brokenSymmetry_iff}\leanok
\uses{def:as-model}
A configuration $x$ has broken symmetry for $j$ colours if $\kappa(x)<j$, or $\kappa(x)\le j$ and
some colour has between $1$ and $n/j-16\sqrt{jn\log n}$ agents.
\end{definition}

The paper's threshold is $n/j-\sqrt{jn\log n}$; the factor $16$ is needed in
Lemma~\ref{lem:mo-stage1}. The death of a colour also breaks the symmetry (the paper's minimum
is over the initial colours, so a dead colour has support $0$).

\begin{lemma}[Equation~(3)]\label{lem:mo-min-drift}
\lean{ThreeMajority.min_count_drift}\leanok
\uses{lem:mo-expect,lem:mo-tails}
There are $\delta>0$ and $C$ such that, for every configuration $x$ with $j=\kappa(x)\ge2$
colours, $Cj^3\log n\le n$ and every colour of $x$ above $n/j-m$ agents, the smallest support
among the colours of $x$ (a colour that dies counting $0$) decreases in expectation by at least
$\delta\sqrt n/(j\sqrt j)$ in one round.
\end{lemma}
\begin{proof}\leanok
Take $\delta=1/500$, $C=1024$ and $w=\sqrt{n/j}$. If the smallest support $c_{\min}$ is at most
$n/j-w/250$, Lemma~\ref{lem:mo-expect} for a smallest colour, with $c_{\min}\ge n/(2j)$
(Equation~(5)), gives the drift. Otherwise all colours are above $n/j-w/250$. The supports sum to
$n$, so the smallest new support is at most $n/j-\frac1{2j}\sum_a|C'_a-n/j|$, and
$\E|C'_a-n/j|\ge(9/256)\sigma_a\ge(9/512)w$ by Lemma~\ref{lem:mo-tails}; hence the expected
smallest new support is at most $n/j-9w/1024<c_{\min}-w/1000$. The paper argues in this case with
the Central Limit Theorem and negative association instead.
\end{proof}

\begin{lemma}[Equation~(6) and Lemma~B.1, one step]\label{lem:mo-overshoot}
\lean{ThreeMajority.min_count_overshoot}\leanok
\uses{lem:mo-expect,lem:mo-tails}
There is $C$ such that, under the hypotheses of Lemma~\ref{lem:mo-min-drift}, one round puts some
colour of $x$ below $n/j-64\sqrt{jn\log n}$ agents with probability at most $1/n^2$.
\end{lemma}
\begin{proof}\leanok
Take $C=2^{14}$. Every support is at most $n/j+(j-1)m$, so by Lemma~\ref{lem:mo-expect} every
colour of $x$ has $\E[C'_a]\ge n/j-2m$. The Chernoff lower tail at $n/j-4m$
(Lemma~\ref{lem:mo-tails}) and a union bound over the at most $n$ colours give $1/n^2$. The paper
bounds the overshoot at the hitting time instead (Lemma~B.1); its proof needs a minor correction,
since it treats the smallest support as a Markov chain, while it is a function of the
configuration. The one-step form avoids this argument.
\end{proof}

\begin{lemma}[Lemma~3.3, symmetry breaking]\label{lem:mo-breaking}
\lean{ThreeMajority.symmetry_breaking,ThreeMajority.breakPotential}\leanok
\uses{def:mo-broken,lem:mo-drift,lem:mo-min-drift,lem:mo-overshoot}
There is $C$ such that, from every configuration $c$ with $j\ge2$ colours and $Cj^3\log n\le n$,
3-Majority reaches broken symmetry for $j$ colours within any $t\ge Cj^2\sqrt{\log n}$ rounds with
probability at least $1/2$.
\end{lemma}
\begin{proof}\leanok
Lemma~\ref{lem:mo-drift} with $I=\{\kappa\le j\}$, $B$ the configurations with broken symmetry,
$f=n/j$ minus the smallest support among the $j$ colours (the paper's $Y_t$),
$\lambda=\delta\sqrt n/(j\sqrt j)$ (Lemma~\ref{lem:mo-min-drift}), $\beta=64\sqrt{jn\log n}$ and
$\eta=1/n^2$ (Lemma~\ref{lem:mo-overshoot}), and $M=n/j$. For $t\ge Cj^2\sqrt{\log n}$ (the proof
takes $C=145\,409$), $\lambda t\ge4\beta$ and $M\eta t\le\lambda t/4$, so $\Pr(\tau>t)\le1/2$.

The paper states the conclusion at the fixed time $t$, while its proof bounds the hitting time,
which is what is stated here; its composition in Theorem~3.1 needs a minor correction, the strong
Markov property at the hitting time (Lemma~\ref{lem:mo-strong-markov}). The statement also needs
a minor correction: it needs the hypothesis $j^3\log n=O(n)$, as do the paper's Lemmas~3.4
and~3.5 (otherwise
the threshold $n/j-\sqrt{jn\log n}$ is negative and the event has probability $0$). Finally, its
bound $\E[\tau]<j^2\sqrt{\log n}$ should read $\E[\tau]\le(2\alpha/\varepsilon)j^2\sqrt{\log n}$.
\end{proof}

\subsection{Dropping}

\begin{lemma}[Lemma~3.4, dropping stage~1]\label{lem:mo-stage1}
\lean{ThreeMajority.dropping_stage1,ThreeMajority.stage1Level,ThreeMajority.stage1_step}\leanok
\uses{lem:mo-expect,lem:mo-tails}
There is $C$ such that, for $j\ge2$ with $Cj^3\log n\le n$, from every configuration with at most
$j$ colours and every colour $a$ with at most $n/j-16\sqrt{jn\log n}$ agents, at every time $t$
with $Cj\log n\le t\le n$ colour $a$ has more than $Cj^2\log n$ agents with probability at most
$1/n$.
\end{lemma}
\begin{proof}\leanok
Take $C=240$. As in Appendix~B of the paper, the gap $n/j-c_a$ grows by a factor $1+1/(4j)$ per
round until $c_a\le n/(2j)$, and then $c_a$ shrinks by a factor $1-1/(4j)$ per round down to a
floor. Here both stages are one moving level,
$\max\{n/j-m(1+1/(4j))^s,\ (n/(2j))(1-1/(4j))^{s-s_A},\ K\}$ at round $s$, with
$K=\min\{240j^2\log n,n/(2j)\}$ and $s_A=\lceil5j\log n\rceil$. By Lemma~\ref{lem:mo-expect} and
the Chernoff upper tail (Lemma~\ref{lem:mo-tails}), $c_a$ exceeds the level of the next round
with probability at most $1/n^2$; after $s_A+\lceil4j\log n\rceil\le240j\log n$ rounds the level
is $K$, and a union bound over the $t\le n$ rounds gives $1/n$.

The proof in the paper needs a minor correction. At its threshold $n/j-\sqrt{jn\log n}$, and
again near its floor $j^2\log n$, the one-round bound of its Equation~(7) is only about
$n^{-1/12}$, and the union bound over the rounds exceeds $1$ for large $j$ (for instance
$j\ge20n^{1/10}$, which $j\le n^{1/3-\varepsilon}$ allows). The threshold $16\sqrt{jn\log n}$,
which Lemma~\ref{lem:mo-breaking} reaches within the same $O(j^2\sqrt{\log n})$ rounds, and the
floor $240j^2\log n$ make each round fail with probability at most $1/n^2$.
\end{proof}

\begin{lemma}[Lemma~3.5, dropping stage~2]\label{lem:mo-stage2}
\lean{ThreeMajority.dropping_stage2,ThreeMajority.avg_min_colourCount_stepCol_le,ThreeMajority.expList_min_colourCount_le}\leanok
\uses{lem:mo-expect,lem:mo-tails}
There is $C$ such that, for $j\ge2$ with $Cj^3\log n\le n$, from every configuration with at most
$j$ colours and every colour $a$ with at most $n/(2j)$ agents, colour $a$ has died after any
$t\ge Cj\log n$ rounds with probability at least $1/2$.
\end{lemma}
\begin{proof}\leanok
Take $C=1024$. As in the paper's sketch, with the potential $\psi=\min\{c_a,n/(2j)\}$:
Lemma~\ref{lem:mo-expect} gives $\E[\psi']\le(1-1/(2j))\psi+(n/(4j^2))[c_a>n/(2j)]$, and by the
Chernoff upper tail $c_a$ escapes above $n/(2j)$ with probability at most $1/n^4$ per round.
Hence $\E[\psi_s]\le(1-1/(2j))^s\,n/(2j)+(n/(4j^2))\,n^{-4}s^2$, which is at most $1/2$ at
$s=\lceil2j\log n\rceil$, and Markov's inequality gives $\Pr(c_a\ge1)\le1/2$. Later times follow
since a dead colour stays dead.
\end{proof}

\subsection{Consensus}

\begin{lemma}[One phase]\label{lem:mo-phase}
\lean{ThreeMajority.phase_colour_dies}\leanok
\uses{lem:mo-breaking,lem:mo-stage1,lem:mo-stage2,lem:mo-strong-markov}
There is $C$ such that, from every configuration with $j\ge2$ colours and $Cj^3\log n\le n$,
after any $t\ge C(j^2\sqrt{\log n}+j\log n)$ rounds fewer than $j$ colours remain with probability
at least $1/5$.
\end{lemma}
\begin{proof}\leanok
From a configuration with broken symmetry, either fewer than $j$ colours remain, or a colour has
at most $n/j-16\sqrt{jn\log n}$ agents; then Lemma~\ref{lem:mo-stage1} brings it to at most
$240j^2\log n\le n/(2j)$ agents except with probability $1/n$, and Lemma~\ref{lem:mo-stage2} kills
it with probability at least $1/2$: in total at least $\frac12(1-\frac1n)\ge\frac25$. Fewer than
$j$ colours is absorbing, so Lemma~\ref{lem:mo-strong-markov} at the hitting time of broken
symmetry and Lemma~\ref{lem:mo-breaking} give $\frac12\cdot\frac25=\frac15$ (the proof takes
$C=293\,349$). The paper's proof of Theorem~3.1 needs a minor correction here: its probability
$1/4$ ignores the failure probability of Lemma~3.4; any positive constant suffices.
\end{proof}

\begin{lemma}[Amplification over phases]\label{lem:mo-phases}
\lean{ThreeMajority.consensus_of_phases}\leanok
\uses{def:as-model}
If from every configuration with between $2$ and $k$ colours a phase of $L$ rounds removes a
colour with probability at least $p\le1$, then from every configuration with at most $k$ colours
consensus fails after $M$ phases with probability at most $(1-p/2)^M2^k$.
\end{lemma}
\begin{proof}\leanok
The potential equal to $2^{\kappa}$ without consensus and to $0$ at consensus: colours never
increase, and with probability at least $p$ one of them disappears, so over one phase its
expectation is multiplied by at most $(1-p)+p/2=1-p/2$. This makes precise the paper's
``standard probabilistic arguments''.
\end{proof}

\begin{theorem}[Theorem~3.1, for $Ck^3\log n\le n$]\label{thm:mo-main}
\lean{ThreeMajority.consensus_many_opinions,ThreeMajority.expList_not_consensus_le}\leanok
\uses{lem:mo-phase,lem:mo-phases}
There is $C$ such that, from every configuration with $k$ colours and $Ck^3\log n\le n$,
3-Majority has not reached consensus after any $T\ge C(k^2\sqrt{\log n}+k\log n)(k+\log n)$
rounds with probability at most $1/n$.
\end{theorem}
\begin{proof}\leanok
Phases of $L=\lceil C_1(k^2\sqrt{\log n}+k\log n)\rceil$ rounds, with $C_1$ the constant of
Lemma~\ref{lem:mo-phase}, which applies to every configuration with $2\le j\le k$ colours. By
Lemma~\ref{lem:mo-phases} with $p=1/5$ and $M=10(k+\lceil\log n\rceil)$ phases, consensus fails
after $ML$ rounds with probability at most $(9/10)^M2^k\le1/n$, and consensus is absorbing. The
proof takes $C=15(C_1+1)=4\,400\,250$.
\end{proof}

\begin{theorem}[Theorem~3.1]\label{thm:mo-phase2}
\lean{ThreeMajority.bcnpt16Phase2_iff,ThreeMajority.bcnpt16Phase2_of_pos}\leanok
\uses{thm:mo-main,def:as-phase2}
For every $\varepsilon>0$, $\mathrm{Phase2}(\varepsilon)$ holds: there are $C,N$ such that for
$n\ge N$, from every configuration with colours in $\{0,\dots,n-1\}$ and $k\le n^{1/3-\varepsilon}$
colours, 3-Majority has not reached consensus after any
$T\ge C(k^2\sqrt{\log n}+k\log n)(k+\log n)$ rounds with probability at most $1/n$.
\end{theorem}
\begin{proof}\leanok
Theorem~\ref{thm:mo-main} with the same $C$: $k^3\le n^{1-3\varepsilon}$ and
$|C|\log n\le n^{3\varepsilon}$ for $n\ge N$, so $Ck^3\log n\le n$.
\end{proof}

With $\varepsilon=1/24$, Theorem~\ref{thm:mo-phase2} is the hypothesis of
Theorem~\ref{thm:as-main}, which gives Theorem~\ref{thm:as-main-unconditional}.

\section{The Lean formalization}\label{sec:lean}

The argument above is fully formalized (no \texttt{sorry}) in the
accompanying repository (library \texttt{ThreeMajority}, namespace
\texttt{ThreeMajority}). File structure (dependency order):
\texttt{Prob.lean} (Section~\ref{sec:prelim}: aliases of the finite-uniform
\texttt{avg}/\texttt{expList} machinery of the shared \texttt{dynamics}
library, including the independence toolkit \texttt{avg\_prod\_pi} and
friends), \texttt{Model.lean}
(Section~\ref{sec:model}), \texttt{Bounds.lean} (Section~\ref{sec:ineq}),
the shared \texttt{Dynamics/Chernoff.lean} (Section~\ref{sec:chernoff}),
\texttt{OneRound.lean}
(Section~\ref{sec:map}), \texttt{Growth.lean} (Section~\ref{sec:growth}),
\texttt{Saturation.lean} (Section~\ref{sec:saturation}), \texttt{Main.lean}
(Section~\ref{sec:proof-main} and the main theorem). The files
\texttt{AnyStartModel}, \texttt{AnyStartMajorization}, \texttt{AnyStartRinott},
\texttt{AnyStartComparison}, \texttt{AnyStartVsVoter}, \texttt{AnyStartVoter},
\texttt{AnyStartAsymptotics} and \texttt{AnyStartMain} formalize
Section~\ref{sec:bcekmn-any-start}, on the shared finite \texttt{Distribution} and
\texttt{Kernel} layer of \texttt{dynamics}; the files \texttt{ManyOpinionsExpect},
\texttt{ManyOpinionsMoments}, \texttt{ManyOpinionsDrift}, \texttt{ManyOpinionsTails},
\texttt{ManyOpinionsBreaking}, \texttt{ManyOpinionsDropping} and \texttt{ManyOpinionsMain}
formalize Section~\ref{sec:bcnpt16-many-opinions}. Design constraint
preserved from \texttt{rumor\_spread}: no measure theory, no
\texttt{PMF}/\texttt{ENNReal}, no Mathlib \texttt{ProbabilityTheory}
kernels --- everything is finite sums, plus \texttt{Real.exp}/\texttt{Real.log}
where the Chernoff bound genuinely needs them (structurally central here,
unlike the push protocol, where \texttt{Real.exp} appears only in a numeric
side lemma).

\begin{thebibliography}{9}
\bibitem{DGMSS11} B.~Doerr, L.A.~Goldberg, L.~Minder, T.~Sauerwald,
C.~Scheideler. \emph{Stabilizing consensus with the power of two choices}.
SPAA 2011.
\bibitem{BCNPS15} L.~Becchetti, A.~Clementi, E.~Natale, F.~Pasquale,
R.~Silvestri. \emph{Plurality consensus in the gossip model}. SODA 2015.
\bibitem{BCEKMN17} P.~Berenbrink, A.~Clementi, R.~Elsässer, P.~Kling,
F.~Mallmann-Trenn, E.~Natale. \emph{Ignore or comply? On breaking symmetry in consensus}.
PODC 2017. arXiv:1702.04921.
\bibitem{BCNPT16} L.~Becchetti, A.~Clementi, E.~Natale, F.~Pasquale, L.~Trevisan.
\emph{Stabilizing consensus with many opinions}. SODA 2016. arXiv:1508.06782.
\end{thebibliography}


%%FILE averaging/blueprint/src/content.tex
\begin{abstract}
We formalize the averaging dynamics on a finite graph, in which every node replaces its value by
the average of its neighbours' values: the degree-weighted sum is conserved, the values stay in
the initial range, and on a connected graph with an odd closed walk all values converge to the
degree-weighted average of the initial values; on a connected bipartite graph they need not
converge. The distance to the limit decays like $\lambda^t$, with $\lambda$ the second largest
absolute eigenvalue of the random walk (Lovász); in the random sequential version, where one random
edge averages its endpoints, the expected state evolves by the matrix $I-L/2m$. On a
well-clustered regular graph started from random signs, the sign of the last change recovers the two clusters exactly (Becchetti, Clementi, Natale, Pasquale, Trevisan, SODA 2017).
\end{abstract}

\section{The dynamics}
\begin{definition}[Averaging]\label{def:avg}
\lean{Averaging.avgStep,Averaging.avgIter,Averaging.degAvg}\leanok
$x_{t+1}(v)=\frac1{\deg v}\sum_{u\sim v}x_t(u)$, and $\bar x=\sum_v\deg(v)\,x_0(v)/\sum_v\deg(v)$.
\end{definition}
\begin{lemma}[Conservation and range]\label{lem:basic}
\lean{Averaging.degree_weighted_sum_step,Averaging.avgStep_le,Averaging.le_avgStep,Averaging.degree_weighted_sum_iter}\leanok\uses{def:avg}
$\sum_v\deg(v)\,x_{t}(v)$ does not depend on $t$; without isolated nodes, a round never leaves an
interval containing all current values.
\end{lemma}
\begin{proof}\leanok
Double counting over ordered adjacent pairs; an average of values in an interval stays in it.
\end{proof}

\section{Convergence}
\begin{lemma}[Walk weights]\label{lem:walks}
\lean{Averaging.transW,Averaging.transW_nonneg,Averaging.transW_sum,Averaging.avgIter_eq_sum_transW,Averaging.transW_pos_of_walk,Averaging.exists_uniform_walk_length}\leanok\uses{def:avg}
$x_t(v)=\sum_u W_t(v,u)\,x_0(u)$ with $W_t$ stochastic and $W_t(v,u)>0$ when a walk of length $t$
joins $v$ and $u$. On a connected graph with an odd closed walk, some length $L$ joins every pair.
\end{lemma}
\begin{proof}\leanok
Induction on $t$; for $L$, fix the parity with the odd closed walk and pad with back-and-forth steps.
\end{proof}
\begin{theorem}[Convergence to the degree-weighted average]\label{thm:conv}
\lean{Averaging.fMax_antitone,Averaging.fMin_monotone,Averaging.range_avgIter_le,Averaging.gap_iter,Averaging.tendsto_degAvg}\leanok\uses{lem:basic,lem:walks}
On a connected graph with an odd closed walk, $x_t(v)\to\bar x$ for every $v$.
\end{theorem}
\begin{proof}\leanok
All entries of $W_L$ are at least some $\delta>0$, so $\max-\min$ contracts by $1-|V|\delta$ every
$L$ rounds (Doeblin); the common limit is $\bar x$ by conservation.
\end{proof}
\begin{theorem}[Bipartite graphs]\label{thm:bip}
\lean{Averaging.not_tendsto_of_colorable}\leanok\uses{def:avg}
On a connected bipartite graph with at least two nodes, the values $\pm1$ of a proper
2-colouring change sign at every round, so they converge nowhere.
\end{theorem}
\begin{proof}\leanok
All neighbours of a node have the other colour.
\end{proof}

\section{Rate of convergence}\label{sec:lovasz}
Theorem~5.1 of Lovász, \emph{Random walks on graphs: a survey} (1993), in the form of
Theorem~33 of the survey of Becchetti, Clementi and Natale, \emph{Consensus dynamics: an
overview} (SIGACT News 2020). Two printed details need a minor correction: Lovász's text defines
$\lambda$ as $\min\{|\lambda_2|,|\lambda_n|\}$, a slip for $\max$ (with $\min$ the bound would
fail, e.g.\ on a connected bipartite graph), and Theorem~33 of the survey prints
$\sqrt{d(v)/d(v)}$ for $\sqrt{d(v)/d(u)}$; the corrected form is formalized
(\href{https://github.com/formal-dynamics/leanamics/blob/main/averaging/FORMALIZATION_DIFFERENCES.md#rate-bound-and-sequential-averaging-rate-sequential-avg-1}{details}).
\begin{definition}[Random walk and its spectrum]\label{def:walk}
\lean{Averaging.walkMatrix,Averaging.normAdjMatrix,Averaging.normAdjMatrix_isHermitian,Averaging.walkLambda,Averaging.walkStationary}\leanok\uses{def:avg}
$P=D^{-1}A$ is the transition matrix of the random walk, $N=D^{-1/2}AD^{-1/2}$ its symmetric
version with eigenvalues $\lambda_1\ge\lambda_2\ge\dots\ge\lambda_n$,
$\lambda=\max\{|\lambda_2|,|\lambda_n|\}$ ($0$ if $n<2$), and $\pi(v)=\deg(v)/2m$ with $m$ the
number of edges.
\end{definition}
\begin{lemma}[Matrix form of the dynamics]\label{lem:walk-bridges}
\lean{Averaging.walkMatrix_pow_apply,Averaging.avgIter_eq_walkMatrix_pow_mulVec,Averaging.charpoly_walkMatrix}\leanok\uses{def:walk,lem:walks}
$P^t(u,v)=W_t(u,v)$ is the probability that the walk from $u$ is at $v$ after $t$ steps, and
$x_t=P^tx_0$. If every degree is positive, $P$ and $N$ have the same characteristic polynomial,
so $\lambda$ is computed from the spectrum of $P$.
\end{lemma}
\begin{proof}\leanok
Induction on $t$; $P=SNS^{-1}$ with $S=D^{-1/2}$.
\end{proof}
\begin{theorem}[Lovász's bound, Survey Theorem 33]\label{thm:lovasz}
\lean{Averaging.abs_walkMatrix_pow_sub_walkStationary_le}\leanok\uses{def:walk,lem:walk-bridges}
If every degree is positive, $|P^t(u,v)-\pi(v)|\le\sqrt{\deg(v)/\deg(u)}\,\lambda^t$ for all
nodes $u,v$ and all $t$.
\end{theorem}
\begin{proof}\leanok
$P^t(u,v)=\sqrt{\deg(v)/\deg(u)}\,N^t(u,v)$ and $\pi(v)=\sqrt{\deg(v)/\deg(u)}\,s(u)s(v)$ with
$s=\sqrt\pi$ a unit eigenvector of $N$ for the eigenvalue $1$. At most one eigenvalue of $N$
exceeds $\lambda$ in absolute value, and then it is $1$ with eigenvector $s$; so $N$ contracts
$s^\perp$ by $\lambda$, and $N^t(u,v)-s(u)s(v)=\langle Qe_u,N^tQe_v\rangle$ with $Q=I-ss^\top$
is at most $\lambda^t$ by Cauchy--Schwarz.
\end{proof}
\begin{lemma}[Non-bipartite graphs]\label{lem:lambda-lt-one}
\lean{Averaging.walkLambda_lt_one}\leanok\uses{def:walk,lem:walk-bridges,thm:conv}
On a connected graph with an odd closed walk, $\lambda<1$.
\end{lemma}
\begin{proof}\leanok
An eigenvalue $\mu$ of $N$ with $|\mu|=1$ gives a vector $y$ with $x_t=\mu^ty$ for the
dynamics started at $y$; by Theorem~\ref{thm:conv} $\mu=-1$ is impossible and $\mu=1$ forces $y$
constant, so $1$ is a simple eigenvalue and all others have absolute value below $1$.
\end{proof}
\begin{corollary}[Rate of the averaging dynamics]\label{cor:avg-rate}
\lean{Averaging.abs_avgIter_sub_walkStationary_le}\leanok\uses{thm:lovasz}
If every degree is positive,
$|x_t(u)-\sum_v\pi(v)x_0(v)|\le\lambda^t\sum_v\sqrt{\deg(v)/\deg(u)}\,|x_0(v)|$.
\end{corollary}
\begin{proof}\leanok
$x_t(u)-\sum_v\pi(v)x_0(v)=\sum_v(P^t(u,v)-\pi(v))x_0(v)$ and Theorem~\ref{thm:lovasz}.
\end{proof}

\section{Sequential averaging}\label{sec:bgps-gossip}
The random sequential (gossip) averaging of Boyd, Ghosh, Prabhakar and Shah, \emph{Randomized
gossip algorithms} (IEEE Trans.\ Inf.\ Theory 52, 2006), the source of equations (4) and (5) of
the survey of Becchetti, Clementi and Natale.
\begin{definition}[Edge step]\label{def:edge-matrix}
\lean{Averaging.Sequential.edgeMatrix,Averaging.Sequential.kernelMatrix,Averaging.Sequential.gossipMeanMatrix,Averaging.Sequential.seqRun,Averaging.Sequential.meanMatrix}\leanok
A step on the edge $(i,j)$ multiplies the state by $W_{ij}=I-\frac12(e_i-e_j)(e_i-e_j)^\top$;
a run applies the steps of a list of oriented edges in order. For a stochastic matrix $P$
(given as a kernel), $\bar W_P=I-\bar D/2n+(P+P^\top)/2n$ with
$\bar D_{ii}=\sum_j(P_{ij}+P_{ji})$, and $\bar W=I-L/2m$ with $L$ the Laplacian.
\end{definition}
\begin{lemma}[Properties of an edge step]\label{lem:edge-matrix}
\lean{Averaging.Sequential.edgeMatrix_mulVec,Averaging.Sequential.edgeMatrix_mem_doublyStochastic,Averaging.Sequential.transpose_edgeMatrix_mul_self}\leanok\uses{def:edge-matrix}
$W_{ij}x$ replaces $x_i$ and $x_j$ by their average and keeps the other values; $W_{ij}$ is
doubly stochastic and a symmetric projection, $W_{ij}^\top W_{ij}=W_{ij}$.
\end{lemma}
\begin{proof}\leanok
Direct computation; $(e_i-e_j)^\top(e_i-e_j)=2$ for $i\ne j$.
\end{proof}
\begin{theorem}[Expected step, Survey equations (4) and (5)]\label{thm:mean-step}
\lean{Averaging.Sequential.avg_expect_edgeMatrix,Averaging.Sequential.gossipMeanMatrix_of_isRegular,Averaging.Sequential.avg_edgeMatrix,Averaging.Sequential.meanMatrix_of_isRegular}\leanok\uses{lem:edge-matrix,def:walk}
If a uniformly random node $i$ contacts $j$ with probability $P_{ij}$, then
$\mathbb E[W_{ij}]=\bar W_P$; if a uniformly random oriented edge is selected, then
$\mathbb E[W_{ij}]=\bar W$. On a $d$-regular graph with $d>0$ and $P=D^{-1}A$ both equal
$(1-1/n)I+P/n$.
\end{theorem}
\begin{proof}\leanok
Entrywise expansion of $W_{ij}(a,b)=\delta_{ab}-\frac12(\delta_{ia}-\delta_{ja})(\delta_{ib}-\delta_{jb})$;
there are $\deg(a)$ darts out of $a$ and $2m=nd$ darts in all.
\end{proof}
\begin{theorem}[First and second moments]\label{thm:moments}
\lean{Averaging.Sequential.avg_transpose_edgeMatrix_mul_self,Averaging.Sequential.avg_sum_sq_edgeMatrix_mulVec,Averaging.Sequential.expList_seqRun}\leanok\uses{thm:mean-step}
For a uniformly random oriented edge, $\mathbb E[W^\top W]=\bar W$ and
$\mathbb E\|Wx\|^2=x^\top\bar Wx$; after $t$ independent uniformly random oriented edges,
$\mathbb E[x_t]=\bar W^tx_0$.
\end{theorem}
\begin{proof}\leanok
$W^\top W=W$; for the first moment, induction on $t$ and linearity of expectation.
\end{proof}

\section{Strong reconstruction}\label{sec:bcnpt-reconstruction}
Becchetti, Clementi, Natale, Pasquale and Trevisan, \emph{Find your place: simple distributed
algorithms for community detection} (SODA 2017; SIAM J. Comput.\ 2020), Theorem~3.2.
\begin{definition}[Clustered regular graph]\label{def:clustered}
\lean{Averaging.IsBalancedPartition,Averaging.IsClusteredRegular,Averaging.clusterIndicator,Averaging.transitionMatrix,Averaging.maxAbsOtherEigenvalue,Averaging.color,Averaging.IsStrongReconstruction,Averaging.reconstructionTime}\leanok\uses{def:avg}
$G$ is $d$-regular on two clusters $V_1,V_2$ of $n$ nodes, each node having $b$ neighbours in the
other cluster; $\chi=\mathbb 1_{V_1}-\mathbb 1_{V_2}$, $P=A/d$ with eigenvalues
$\lambda_1\ge\lambda_2\ge\dots$, $\lambda=\max_{i\ge3}|\lambda_i|$. Node $v$ is blue at round $t$
if $x_t(v)\ge x_{t-1}(v)$; a coloring is a strong reconstruction if no color appears in both
clusters. $T(n,\delta)=\lceil\log(4n^3)/\log(1+\delta)\rceil+1$.
\end{definition}
\begin{lemma}[Lemma C.1]\label{lem:c1}
\lean{Averaging.avgIter_eq_transitionMatrix_pow_mulVec,Averaging.transitionMatrix_mulVec_clusterIndicator,Averaging.abs_avgIter_sub_le}\leanok\uses{def:clustered}
If $\lambda<1-2b/d$ and $x\in\{\pm1\}^{2n}$, then
$|x_t(v)-\alpha_1-\alpha_2(1-2b/d)^t\chi(v)|\le\lambda^t\sqrt{2n}$ with
$\alpha_1=\langle x,\mathbb 1\rangle/2n$, $\alpha_2=\langle x,\chi\rangle/2n$.
\end{lemma}
\begin{proof}\leanok
$\mathbb 1$ and $\chi$ are eigenvectors with eigenvalues $1$ and $1-2b/d$, both above $\lambda$;
at most two eigenvalues exceed $\lambda$ in absolute value, so by equality in Bessel's inequality
their eigenvectors lie in $\operatorname{span}\{\mathbb 1,\chi\}$, and $P$ contracts the
orthogonal complement of this plane by $\lambda$.
\end{proof}
\begin{theorem}[Deterministic core]\label{thm:core}
\lean{Averaging.sign_avgIter_sub_eq,Averaging.exists_sign_avgIter_sub_clusters,Averaging.isStrongReconstruction_color,Averaging.reconstructionTime_le}\leanok\uses{lem:c1}
On a connected clustered graph with $1-2b/d>(1+\delta)\lambda$, if $x\in\{\pm1\}^{2n}$ and
$\langle x,\chi\rangle\ne0$, then for $t\ge T(n,\delta)=O(\log n/\delta)$,
$\operatorname{sgn}(x_{t-1}(u)-x_t(u))=\operatorname{sgn}(\langle x,\chi\rangle\chi(u))$.
\end{theorem}
\begin{proof}\leanok
$|\langle x,\chi\rangle|\ge2$ by parity, $b\ge1$ by connectivity, $d<2n$, and
$(1+\delta)^{t-1}\ge4n^3$, so the main term of Lemma~\ref{lem:c1} dominates the errors.
\end{proof}
\begin{lemma}[Lemma B.1, exact form]\label{lem:b1}
\lean{Averaging.prob_dotProduct_clusterIndicator_eq_zero,Averaging.choose_div_four_pow_le}\leanok\uses{def:clustered}
For uniform $x\in\{\pm1\}^{2n}$, $P(\langle x,\chi\rangle=0)=\binom{2n}{n}/4^n\le1/\sqrt{\pi n}$.
\end{lemma}
\begin{proof}\leanok
$x\mapsto\{v: x(v)=\chi(v)\}$ maps the zeros bijectively to the $n$-subsets; Wallis' product.
\end{proof}
\begin{theorem}[Strong reconstruction, Theorem 3.2]\label{thm:strong}
\lean{Averaging.strong_reconstruction}\leanok\uses{thm:core,lem:b1}
With probability at least $1-1/\sqrt{\pi n}$, the coloring is a strong reconstruction at every
round $t\ge T(n,\delta)$.
\end{theorem}
\begin{proof}\leanok
Theorem~\ref{thm:core} and Lemma~\ref{lem:b1}.
\end{proof}

\section{Averaging whenever you meet}\label{sec:bcmnprt-opportunistic}
Becchetti, Clementi, Manurangsi, Natale, Pasquale, Raghavendra, Trevisan, \emph{Average whenever
you meet: opportunistic protocols for community detection} (ESA 2018, arXiv:1703.05045).
The proof of Theorem~4.1 needs a minor correction: its last step uses
$\mathbb E\|z^{(0)}\|^2\le 1/n$, while for uniform signs $\mathbb E\|z^{(0)}\|^2=n-2$. The
conclusion still holds, and the formal proof uses its own bookkeeping
(\href{https://github.com/formal-dynamics/leanamics/blob/main/averaging/FORMALIZATION_DIFFERENCES.md#averaging-whenever-you-meet-opportunistic-avg-3}{details}).
\begin{definition}[Opportunistic averaging, Algorithm 1]\label{def:opp}
\lean{Averaging.Opportunistic.signVec,Averaging.Opportunistic.edgeAvg,Averaging.Opportunistic.avgRun,Averaging.Opportunistic.oppStep,Averaging.Opportunistic.oppRun}\leanok
One uniformly random edge is activated per round; a node activated for the first time draws a
uniform sign, and both endpoints take the average of their two values. The rounds are i.i.d.
uniform darts of $G$.
\end{definition}
\begin{lemma}[Deferred decisions]\label{lem:opp-model}
\lean{Averaging.Opportunistic.avg_dart_edge,Averaging.Opportunistic.oppRun_isSome_iff,Averaging.Opportunistic.oppRun_getD,Averaging.Opportunistic.sum_avgRun}\leanok\uses{def:opp}
A uniform dart is a uniform edge; a node holds a value exactly when it has been activated, and
the protocol's values are those of the averaging process started from the coins $\sigma$; the sum
of the values is conserved.
\end{lemma}
\begin{proof}\leanok
Induction on the list of rounds; every edge carries two darts.
\end{proof}
\begin{theorem}[Expected step, equations (1), (2), (15)]\label{thm:opp-first}
\lean{Averaging.Opportunistic.stepMatrix,Averaging.Opportunistic.meanStepMatrix,Averaging.Opportunistic.stepMatrix_mulVec,Averaging.Opportunistic.stepMatrix_isSymm,Averaging.Opportunistic.stepMatrix_mul_self,Averaging.Opportunistic.avg_stepMatrix,Averaging.Opportunistic.meanStepMatrix_isSymm,Averaging.Opportunistic.meanStepMatrix_mem_doublyStochastic,Averaging.Opportunistic.avg_edgeAvg,Averaging.Opportunistic.expList_avgRun,Averaging.Opportunistic.avg_sum_sq_edgeAvg}\leanok\uses{def:opp}
$W_t$ is a symmetric projection, $\mathbb E[W]=\bar W=I-L/2m$ is symmetric and doubly stochastic,
$\mathbb E[x^{(t)}]=\bar W^tx^{(0)}$ and $\mathbb E\|Wx\|^2=x^\top\bar Wx$.
\end{theorem}
\begin{proof}\leanok
Entrywise computation over the darts; induction on $t$.
\end{proof}
\begin{definition}[Clustered regular graphs, Definition 2.2]\label{def:opp-clustered}
\lean{Averaging.Opportunistic.IsClusteredRegular,Averaging.Opportunistic.cutVec,Averaging.Opportunistic.projOne,Averaging.Opportunistic.projCut,Averaging.Opportunistic.projRest,Averaging.Opportunistic.ThirdEigenvalueLB}\leanok
$G$ is $d$-regular with communities $V_1$ and $V_2=V_1^c$ of size $n/2$, each node having $b$
neighbours in the other community, $0<b$, $2b<d<n$. The state splits as $x=x_\parallel+y+z$ along
$\mathbb 1$, $\chi=\mathbb 1_{V_1}-\mathbb 1_{V_2}$ and their orthogonal complement;
$\lambda_2=2b/d$, and $\lambda_3$ enters as a lower bound on the Rayleigh quotient of $L/d$ on
the complement.
\end{definition}
\begin{theorem}[Theorem 3.1 on clustered regular graphs]\label{thm:opp-31}
\lean{Averaging.Opportunistic.meanStepMatrix_mulVec_cutVec,Averaging.Opportunistic.expList_projCut,Averaging.Opportunistic.expList_avgRun_decomp,Averaging.Opportunistic.sum_sq_meanStepMatrix_pow_le,Averaging.Opportunistic.monotone_criterion,Averaging.Opportunistic.sign_criterion}\leanok\uses{thm:opp-first,def:opp-clustered}
$\bar W\chi=(1-\lambda_2/n)\chi$, the expected state decomposes exactly as in Lemma B.1, and the
monotonicity and sign criteria of Lemmas B.3 and B.4 hold at every node.
\end{theorem}
\begin{proof}\leanok
$\bar W$ contracts the complement of $\operatorname{span}\{\mathbb 1,\chi\}$ by
$1-\lambda_3/n$.
\end{proof}
\begin{lemma}[One round, Lemmas C.2 and C.3]\label{lem:opp-onestep}
\lean{Averaging.Opportunistic.IsClusteredRegular.avg_restSq_edgeAvg_le,Averaging.Opportunistic.IsClusteredRegular.avg_cutDev_edgeAvg_le}\leanok\uses{thm:opp-first,def:opp-clustered}
With $\beta=\langle x,\chi\rangle/n$: $\mathbb E\|z'\|^2\le(1-\lambda_3/n)\|z\|^2+\lambda_2\beta^2$ and,
for every $\gamma$, $\mathbb E\,n(\beta'-\gamma)^2\le(1-\lambda_2/n)\,n(\beta-\gamma)^2+\lambda_2\gamma^2
+(4\lambda_2/n^2)(n\beta^2+\|z\|^2)$.
\end{lemma}
\begin{proof}\leanok
$\mathbb E\|Wx\|^2=x^\top\bar Wx$ and Jensen for $\beta'$; the variance of $\beta'$ comes from the
$nb$ darts across the cut.
\end{proof}
\begin{theorem}[Second moment analysis, Lemma C.4 and Theorem 4.1]\label{thm:opp-41}
\lean{Averaging.Opportunistic.IsClusteredRegular.expList_secondMoment_le,Averaging.Opportunistic.IsClusteredRegular.secondMoment_bound_aux,Averaging.Opportunistic.secondMoment_bound}\leanok\uses{lem:opp-onestep}
If $100\lambda_2\log n\le\lambda_3$ and $3(n/\lambda_3)\log n\le t\le n/(4\lambda_2)$, then
$\mathbb E\|y^{(t)}+z^{(t)}-y^{(0)}\|^2\le3\lambda_2t/n$.
\end{theorem}
\begin{proof}\leanok
Unroll the two recursions from any state; average over the signs ($\mathbb E\|y^{(0)}\|^2=1$,
$\|z^{(0)}\|^2\le n$).
\end{proof}
\begin{theorem}[Non-ephemeral good nodes, Lemma 4.2]\label{thm:opp-42}
\lean{Averaging.Opportunistic.IsGood,Averaging.Opportunistic.card_good_ge,Averaging.Opportunistic.IsClusteredRegular.nonEphemeral_aux,Averaging.Opportunistic.IsClusteredRegular.prob_good_window,Averaging.Opportunistic.nonEphemeral_good}\leanok\uses{thm:opp-41}
If $\lambda_2/\lambda_3\le\lambda_3\varepsilon^4/(c\log^2n)$, then with probability at least
$1-\varepsilon$ at least $(1-3\varepsilon)n$ nodes are $\varepsilon$-good at every round of the
phase $6(n/\lambda_3)\log n\le t\le12(n/\lambda_3)\log n$.
\end{theorem}
\begin{proof}\leanok
Good nodes averaged only with good nodes of their community stay good; the bad rounds of the
phase are counted in expectation from the second-moment bounds restarted at the start of the
phase, and small initial cut components are handled by anti-concentration of $\langle\sigma,\chi\rangle$.
\end{proof}
\begin{theorem}[Sign recovery and weak reconstruction]\label{thm:opp-main}
\lean{Averaging.Opportunistic.IsWeakReconstruction,Averaging.Opportunistic.IsClusteredRegular.avg_abs_blockSum_le_mul,Averaging.Opportunistic.IsClusteredRegular.avg_blockSum_mul_neg,Averaging.Opportunistic.IsClusteredRegular.prob_sign_window,Averaging.Opportunistic.IsClusteredRegular.prob_weakReconstruction_window,Averaging.Opportunistic.sign_phase,Averaging.Opportunistic.weakReconstruction_phase}\leanok\uses{thm:opp-42}
Under the hypothesis of Lemma 4.2, with probability at least $1-4\varepsilon$ at least
$(1-4\varepsilon)n$ nodes have the sign of their community's initial average throughout the
phase, and with probability at least $1/2-6\varepsilon$ the signs of the values are a
$6\varepsilon$-weak reconstruction at every round of the phase.
\end{theorem}
\begin{proof}\leanok
A good node has the sign of $2S_h/n$ unless $|S_1|\le\varepsilon|S_2|$ or $|S_2|\le\varepsilon|S_1|$
(each of probability at most $\varepsilon+1/\sqrt{n/2+1}$); $S_1S_2<0$ with probability at least
$1/2-1/\sqrt{n/2+1}$.
\end{proof}


%%FILE crn/blueprint/src/content.tex
\begin{abstract}
We formalize the correspondence between chemical reaction networks and population protocols.
For a count-conserving bimolecular network ($A + B \to C + D$) with a common rate constant, the
jump chain of stochastic mass-action kinetics equals the population protocol that draws a
uniformly random ordered pair of distinct agents, conditioned on the pair reacting, as kernels on
count vectors. The worked instance is the approximate-majority network. We then formalize the easy
direction of the characterization of stably computable predicates (Angluin, Aspnes, Diamadi,
Fischer and Peralta, 2006): threshold and remainder predicates on input counts are stably
computable by population protocols, the class is closed under Boolean operations, and every
protocol is a count-conserving bimolecular network, so every Boolean combination of threshold and
remainder predicates is stably decided by such a network. Finally, on arbitrary connected
interaction graphs (Mertzios, Nikoletseas, Raptopoulos and Spirakis, 2014), a 4-state protocol
stably computes majority, no protocol with at most 3 states does, and the 3-state
approximate-majority protocol converges to the majority with probability at least $1/2$ from a
uniformly random placement.
\end{abstract}

\section{Chemical reaction networks}
\begin{definition}[Reactions, networks, counts]\label{def:crn}
\lean{Crn.Reaction,Crn.Network,Crn.Counts,Crn.Reaction.consumed,Crn.Reaction.produced,Crn.Counts.react}\leanok
A reaction $A + B \to C + D$ consumes the unordered pair $\{A, B\}$ of species and produces the
unordered pair $\{C, D\}$; a network is a nonempty finite set of reactions whose products differ
from their reactants. A count vector of $n$ molecules gives the number $x_a$ of molecules of each
species $a$, with $\sum_a x_a = n$; firing an applicable reaction $r$ replaces $x_a$ by
$x_a - \nu_a + \nu'_a$, with $\nu$, $\nu'$ the reactant and product multiplicities.
\end{definition}

\begin{definition}[Mass-action propensity]\label{def:propensity}
\lean{Crn.Reaction.propensity,Crn.Network.totalPropensity}\leanok\uses{def:crn}
With rate constant $k$, the propensity of $r$ at $x$ is $a_r(x) = k \prod_a \binom{x_a}{\nu_a}$
(Gillespie's combinatorial form), and $a_0(x) = \sum_r a_r(x)$.
\end{definition}

\begin{lemma}[Propensities]\label{lem:propensity}
\lean{Crn.Reaction.propensity_of_ne,Crn.Reaction.propensity_of_eq}\leanok\uses{def:propensity}
For $A + B \to \cdots$ with $A \ne B$ the propensity is $k\, x_A x_B$; for $A + A \to \cdots$ it
is $k \binom{x_A}{2}$.
\end{lemma}
\begin{proof}\leanok
Only the reactant species have nonzero multiplicity in the product.
\end{proof}

\begin{definition}[Jump chain]\label{def:jump}
\lean{Crn.Network.nextReaction,Crn.Network.jumpKernel}\leanok\uses{def:propensity}
From $x$ with $a_0(x) > 0$ the jump chain fires $r$ with probability $a_r(x)/a_0(x)$; from a
terminal state ($a_0(x) = 0$) it stays put.
\end{definition}

\begin{lemma}[The jump chain jumps]\label{lem:jump-self}
\lean{Crn.Network.jumpKernel_weight_self}\leanok\uses{def:jump}
If $a_0(x) \ne 0$, the jump chain from $x$ gives weight $0$ to $x$.
\end{lemma}
\begin{proof}\leanok
A reaction with nonzero propensity is applicable, and firing an applicable reaction whose
products differ from its reactants changes the counts.
\end{proof}

\section{Population protocols}
\begin{definition}[The population protocol of a network]\label{def:pp}
\lean{Crn.AgentPair,Crn.pairDist,Crn.counts,Crn.sampleDist,Crn.Reacts,Crn.outcome,Crn.ppStep,Crn.condStep,Crn.ppKernel}\leanok\uses{def:crn}
Each of $n \ge 2$ agents holds a species. One interaction draws a uniformly random ordered pair
$(u, v)$ of distinct agents and, independently, a uniformly random reaction $r$ of the network;
if the species of $u$ and $v$ are the reactants of $r$ the pair reacts and $r$ fires, otherwise
nothing changes. The protocol is observed on count vectors; conditioning on the pair reacting
gives a kernel on count vectors.
\end{definition}

\begin{lemma}[Key identity]\label{lem:pair}
\lean{Crn.pair_prob_of_ne,Crn.pair_prob_of_eq,Crn.pair_prob_eq_propensity}\leanok\uses{def:pp,lem:propensity}
The drawn pair has species $\{A, B\}$ with probability $x_A x_B / \binom{n}{2}$ if $A \ne B$, and
$\binom{x_A}{2} / \binom{n}{2}$ if $A = B$. Equivalently, it has the reactants of $r$ with
probability $a_r(x) / (k \binom{n}{2})$.
\end{lemma}
\begin{proof}\leanok
There are $n(n-1)$ ordered pairs of distinct agents, $2 x_A x_B$ of them with species $\{A, B\}$
for $A \ne B$, and $x_A(x_A - 1)$ with species $\{A, A\}$.
\end{proof}

\begin{theorem}[One interaction]\label{thm:step}
\lean{Crn.reactProb_eq,Crn.ppStep_expect}\leanok\uses{lem:pair,def:jump}
The drawn pair reacts with probability $p = a_0(x) / (k\, |N| \binom{n}{2})$, and one interaction
is a jump of the mass-action chain with probability $p$ and a no-op otherwise.
\end{theorem}
\begin{proof}\leanok
Average over the uniform reaction and apply the key identity to each reaction.
\end{proof}

\begin{theorem}[CRNs are population protocols]\label{thm:crn1}
\lean{Crn.condStep_eq_jumpKernel,Crn.jumpKernel_eq_ppKernel}\leanok\uses{thm:step,lem:jump-self}
For every rate constant $k > 0$ and $n \ge 2$, the population protocol conditioned on the drawn
pair reacting is the jump chain of mass-action kinetics, from every configuration and hence as
kernels on count vectors.
\end{theorem}
\begin{proof}\leanok
Conditioned on reacting, the interaction fires $r$ with probability proportional to $a_r(x)$ by
the key identity, which is the jump chain; from a configuration where no pair can react both sides
stay put.
\end{proof}

\section{Approximate majority}
\begin{definition}[The approximate-majority network]\label{def:am}
\lean{Crn.ApproxMajority.Species,Crn.ApproxMajority.network}\leanok\uses{def:crn}
Species $X$, $Y$ (the two opinions) and $B$ (blank), and the reactions $X + Y \to X + B$,
$X + Y \to Y + B$, $B + X \to X + X$, $B + Y \to Y + Y$ (Angluin, Aspnes and Eisenstat).
\end{definition}

\begin{theorem}[Approximate majority as a population protocol]\label{thm:am}
\lean{Crn.ApproxMajority.network_jump_expect,Crn.ApproxMajority.network_reactProb_eq,Crn.ApproxMajority.network_jumpKernel_eq_ppKernel}\leanok\uses{def:am,thm:crn1}
With $D = 2 x_X x_Y + x_B (x_X + x_Y) > 0$, the jump chain fires each $X + Y$ reaction with
probability $x_X x_Y / D$, $B + X \to X + X$ with probability $x_B x_X / D$ and $B + Y \to Y + Y$
with probability $x_B x_Y / D$; a drawn pair reacts with probability $D / (4 \binom{n}{2})$; and the
jump chain is the population protocol conditioned on reacting.
\end{theorem}
\begin{proof}\leanok
Instances of the general results, with $a_0(x) = k D$.
\end{proof}

\section{Stable computation}\label{sec:aadfp-stable}
This section and the two following ones formalize the easy direction of the characterization of
stably computable predicates of Angluin, Aspnes, Diamadi, Fischer and Peralta,
\emph{Computation in networks of passively mobile finite-state sensors} (Distributed Computing
18, 2006). Their Lemma~5 needs a minor correction for a population of a single agent: no
encounter happens, and if a lone agent keeps its state, its output stays $0$ even when the
predicate holds. The protocols
formalized here start each agent with its single-agent output and are correct for every $n\ge1$
(\href{https://github.com/formal-dynamics/leanamics/blob/main/crn/FORMALIZATION_DIFFERENCES.md#semilinear-predicates-are-stably-computable-stable-crn-3-easy-direction}{details}).
\begin{definition}[Population protocols with input and output]\label{def:protocol}
\lean{Crn.Protocol,Crn.Protocol.interact,Crn.Protocol.Step,Crn.Protocol.Reaches,Crn.Protocol.OutputStable}\leanok\uses{def:pp}
A population protocol has an input function $I : X \to Q$, an output function $O : Q \to
\{0, 1\}$ and a transition function $\delta : Q \times Q \to Q \times Q$. An encounter of an
ordered pair $(u, v)$ of distinct agents replaces their states $p, q$ by $\delta(p, q)$;
reachability is the reflexive-transitive closure of encounters. A configuration is
output-stable with output $b$ if in every configuration reachable from it every agent outputs $b$.
\end{definition}

\begin{definition}[Stable computation]\label{def:stable}
\lean{Crn.Protocol.StablyComputes,Crn.StablyComputable}\leanok\uses{def:protocol}
A protocol stably computes a predicate $\varphi$ on input count vectors if, for every nonempty
population and every input assignment $\iota$, every configuration reachable from $I \circ \iota$
can reach an output-stable configuration with output $\varphi$ of the input counts. A predicate is
stably computable if a protocol with finitely many states stably computes it.
\end{definition}

\begin{lemma}[Uniqueness]\label{lem:stable-unique}
\lean{Crn.Protocol.StablyComputes.unique}\leanok\uses{def:stable}
A protocol stably computes at most one predicate on nonzero count vectors.
\end{lemma}
\begin{proof}\leanok
From the initial configuration of a population with the given counts, reach an output-stable
configuration for the first predicate and then one for the second: they agree.
\end{proof}

\begin{theorem}[Boolean closure]\label{thm:boolean}
\lean{Crn.StablyComputable.map₂,Crn.StablyComputable.not,Crn.StablyComputable.and,Crn.StablyComputable.or}\leanok\uses{def:stable}
For every Boolean function $\xi$ of two arguments, if $\varphi$ and $\psi$ are stably computable
then so is $\xi(\varphi, \psi)$; in particular the class is closed under negation, conjunction and
disjunction.
\end{theorem}
\begin{proof}\leanok
Parallel composition: states are pairs, both components make their own transition in every
encounter, and the output of $(p_1, p_2)$ is $\xi(O_1 p_1, O_2 p_2)$.
\end{proof}

\section{Semilinear predicates}
\begin{theorem}[Threshold predicates]\label{thm:threshold}
\lean{Crn.threshold,Crn.stablyComputable_threshold,Crn.stablyComputable_le_sum}\leanok\uses{def:stable}
For integers $a_i$ and $c$, the predicates $\sum_i a_i x_i < c$ and $\sum_i a_i x_i \ge c$ are stably
computable.
\end{theorem}
\begin{proof}\leanok
Each agent holds a leader bit, an output bit and a count in $[-s, s]$, starting from $a_i$; an
encounter involving a leader merges the leaders into the initiator, which takes as much of the sum
of the two counts as fits in $[-s, s]$, and both agents set their output to whether the
initiator's count is below $c$. Unlike the original protocol, an input starts with output bit
$[a_i < c]$, which is needed for a population of one agent. The second form is the negation.
\end{proof}

\begin{theorem}[Remainder predicates]\label{thm:remainder}
\lean{Crn.remainder,Crn.stablyComputable_remainder}\leanok\uses{def:stable}
For integers $a_i$, $c$ and $m > 0$, the predicate $\sum_i a_i x_i \equiv c \pmod m$ is stably
computable.
\end{theorem}
\begin{proof}\leanok
The same leader protocol with counts in $\mathbb{Z}/m$: the initiator takes the sum of the two
counts and the responder gets $0$.
\end{proof}

\begin{definition}[Semilinear predicates]\label{def:semilinear}
\lean{Crn.IsSemilinearPred}\leanok\uses{thm:threshold,thm:remainder}
The Boolean combinations of threshold and remainder predicates.
\end{definition}

\begin{theorem}[Semilinear predicates are stably computable]\label{thm:semilinear}
\lean{Crn.IsSemilinearPred.isSemilinearSet,Crn.IsSemilinearPred.stablyComputable}\leanok\uses{def:semilinear,thm:boolean}
The truth set of every such predicate is a semilinear set of count vectors, and every such
predicate is stably computable.
\end{theorem}
\begin{proof}\leanok
Induction on the Boolean combination, with the two preceding theorems and Boolean closure; for
the first part, semilinear sets are closed under complement, intersection and union.
\end{proof}

\section{From protocols to networks}
\begin{definition}[Stable decision by a network]\label{def:crn-stable}
\lean{Crn.Counts.map,Crn.Network.Step,Crn.Network.Reaches,Crn.Network.OutputStable,Crn.Network.StablyComputes}\leanok\uses{def:crn}
A network with an injective input map $I : X \to S$ and an output map $O : S \to \{0, 1\}$
stably decides $\varphi$ if, from the initial count vector with $x_i$ molecules of $I(i)$ and nothing
else, every reachable count vector can reach one in which, now and in every count vector reachable
from it, every present species votes $\varphi(x)$.
\end{definition}

\begin{theorem}[Protocols are networks]\label{thm:transfer}
\lean{Crn.Protocol.exists_network,Crn.StablyComputable.exists_network,Crn.IsSemilinearPred.exists_network}\leanok\uses{def:crn-stable,thm:semilinear}
The transitions $p + q \to \delta_1(p, q) + \delta_2(p, q)$ that change $\{p, q\}$ form a
count-conserving bimolecular network whose reachable count vectors are the counts of the
configurations reachable in the protocol. Hence every stably computable predicate, and in
particular every Boolean combination of threshold and remainder predicates, is stably decided by
such a network.
\end{theorem}
\begin{proof}\leanok
Encounters change counts exactly by the corresponding reaction. Tagging the input states keeps
the input map injective and provides a transition that changes the multiset of states.
\end{proof}

\section{Majority on graphs with very small local memory (CRN-4)}\label{sec:mnrs-graph-majority}
This section formalizes Theorems~1, 2 and~4 of Mertzios, Nikoletseas, Raptopoulos and Spirakis,
\emph{Determining majority in networks with local interactions and very small local memory}
(ICALP 2014, arXiv:1404.7671, whose numbering we use; Distributed Computing 30, 2017): population
protocols run on a connected interaction graph, and two agents interact only along an edge. The
3-state protocol is that of Angluin, Aspnes and Eisenstat (Distributed Computing 21, 2008).
Graphs are undirected. The proofs of Theorems~1, 2 and~4 need minor corrections: in Theorem~1 the
type of two agents and one set index are corrected; in Theorem~2 ambassadors of the same color can
block each other, which a closest pair of opposite ambassadors avoids; in Theorem~4 the
monotonicity step is stated without proof (it follows from a monotone coupling), and the bound
$1/2$ also needs absorption with probability $1$
(\href{https://github.com/formal-dynamics/leanamics/blob/main/crn/FORMALIZATION_DIFFERENCES.md#majority-on-graphs-with-very-small-local-memory-graphmajority-crn-4}{details}).

\begin{definition}[Protocols on an interaction graph]\label{def:graph-stable}
\lean{Crn.dartPair,Crn.Protocol.GraphStep,Crn.Protocol.GraphReaches,Crn.Protocol.GraphOutputStable,Crn.Protocol.StablyComputesOnGraph,Crn.Protocol.GraphStablyComputes}\leanok\uses{def:stable}
On a graph $G$ on the agents, an encounter is an oriented edge $(u, v)$ of $G$ (initiator $u$,
responder $v$). Reachability and output stability on $G$ are defined as before with encounters
along edges of $G$. A protocol stably computes $\varphi$ on $G$ for inputs in a domain $D$ if, for
every input whose count vector lies in $D$, every configuration reachable on $G$ can reach an
output-stable configuration with output $\varphi$; it stably computes $\varphi$ on graphs if it
does so on every connected graph.
\end{definition}

\begin{lemma}[The complete graph]\label{lem:graph-top}
\lean{Crn.Protocol.graphStep_top,Crn.Protocol.graphReaches_top,Crn.Protocol.graphOutputStable_top,Crn.Protocol.stablyComputesOnGraph_top_iff,Crn.Protocol.GraphStablyComputes.stablyComputes}\leanok\uses{def:graph-stable}
On the complete graph these notions are those of Definition~\ref{def:stable}: stable
computation on the complete graphs of all nonempty populations, with the trivial domain, is
stable computation.
\end{lemma}
\begin{proof}\leanok
Every ordered pair of distinct agents is an oriented edge of the complete graph.
\end{proof}

\begin{lemma}[Embedding a population]\label{lem:graph-extend}
\lean{Crn.Protocol.GraphReaches.extend}\leanok\uses{def:graph-stable}
Reachability transfers along an injective map of agents that preserves edges, the other agents
being frozen.
\end{lemma}
\begin{proof}\leanok
Induction on reachability: an encounter along $(u, v)$ is simulated along $(f(u), f(v))$.
\end{proof}

\begin{definition}[Majority]\label{def:majority}
\lean{Crn.majority,Crn.HasMajority}\leanok\uses{def:graph-stable}
With two input types, red and green, majority is true if red has more agents than green, and the
domain excludes ties.
\end{definition}

\begin{definition}[The ambassador protocol]\label{def:ambassador}
\lean{Crn.AmbState,Crn.ambassadorδ,Crn.ambassador,Crn.card_ambState,Crn.ambassador_table,Crn.ambassador_table_rest,Crn.ambassador_symm}\leanok\uses{def:protocol}
A state is a color and a bit saying whether the vertex holds an ambassador; every agent starts
with an ambassador of its input color, and the output is the color. Two ambassadors of different
colors annihilate; an ambassador meeting a vertex without ambassador moves there and paints it;
nothing else changes. This is the table of Figure~1, with $4$ states, and the protocol is
symmetric.
\end{definition}

\begin{lemma}[Invariants]\label{lem:amb-invariant}
\lean{Crn.ambCount,Crn.ambCount_input,Crn.ambCount_interact_le,Crn.ambCount_sub_interact}\leanok\uses{def:ambassador}
The number of ambassadors of each color starts at the input count and never increases, and the
difference between red and green ambassadors never changes.
\end{lemma}
\begin{proof}\leanok
Count the two agents of an encounter, over the $16$ entries of the table.
\end{proof}

\begin{lemma}[Annihilation and spreading]\label{lem:amb-steps}
\lean{Crn.exists_reaches_annihilate,Crn.exists_reaches_allColor,Crn.graphOutputStable_of_allColor}\leanok\uses{lem:amb-invariant,def:graph-stable}
On a connected graph, while both colors have ambassadors, a configuration with one ambassador
of each color fewer is reachable. If all ambassadors have color $b$ and there is one, the all-$b$
coloring is reachable, and it is output-stable.
\end{lemma}
\begin{proof}\leanok
Take a red ambassador and a closest green one. The red ambassador walks along a shortest path,
moving onto vertices without ambassador and handing over to red ambassadors it meets, until it
meets the green one (strong induction on the distance). Spreading is similar, by induction on the
distance to a vertex of the other color and then on the number of such vertices.
\end{proof}

\begin{theorem}[MNRS14, Theorem 2]\label{thm:ambassador}
\lean{Crn.ambassador_graphStablyComputes,Crn.ambassador_stablyComputesOnGraph_top,Crn.ambassador_tie_stuck}\leanok\uses{lem:amb-steps,def:majority,lem:graph-top}
The ambassador protocol stably computes majority on every connected graph, for every input with
a majority. With one red and one green agent, a configuration is reachable from which no
output-stable configuration is reachable, so ties must be excluded.
\end{theorem}
\begin{proof}\leanok
From every reachable configuration, annihilate pairs until only ambassadors of the majority
color remain (there is one by the invariant), then spread that color. Configurations without
ambassadors are frozen.
\end{proof}

\begin{theorem}[MNRS14, Theorem 1]\label{thm:three-states}
\lean{Crn.not_stablyComputesOnGraph_top_majority,Crn.not_graphStablyComputes_majority,Crn.rank_majority_eq_four}\leanok\uses{lem:graph-extend,def:majority,lem:graph-top,thm:ambassador}
No protocol with at most $3$ states stably computes majority, ties excluded, on the complete
graphs of all nonempty populations, hence on all connected graphs. With
Theorem~\ref{thm:ambassador}, $4$ is the minimal number of states.
\end{theorem}
\begin{proof}\leanok
A single agent never interacts, so both outputs occur, and one output $b$ has a single state
$q_1$. On $3$ agents, the inputs $(b, b, \lnot b)$ and $(b, b, b)$ both reach the configuration in
which every agent is in $q_1$. Adding two agents of input $\lnot b$ and running the same encounters
on the first three agents, inputs with majority $\lnot b$ and $b$ reach a common configuration, a
contradiction.
\end{proof}

\begin{definition}[Approximate majority on a graph]\label{def:am-graph}
\lean{Crn.ApproxMajority.rule,Crn.ApproxMajority.protocol,Crn.ApproxMajority.rule_mem_network,Crn.ApproxMajority.exists_rule_of_mem_network,Crn.ApproxMajority.graphStep,Crn.ApproxMajority.graphKernel}\leanok\uses{def:am,def:protocol}
The one-way rule: the initiator never changes; an initiator $X$ or $Y$ turns an opposite responder
blank and recruits a blank responder. Its effective transitions are exactly the reactions of the
approximate-majority network. Each step activates a uniformly random oriented edge of $G$.
\end{definition}

\begin{definition}[Winning probabilities]\label{def:winprob}
\lean{Crn.ApproxMajority.Unanimous,Crn.ApproxMajority.placement,Crn.ApproxMajority.winProb,Crn.ApproxMajority.absorbProb,Crn.ApproxMajority.winProb_monotone,Crn.ApproxMajority.tendsto_winProb}\leanok\uses{def:am-graph}
With $m$ agents of type $Y$ placed on a uniformly random $m$-set and $X$ elsewhere,
$\mathrm{win}_t(s)$ is the probability that every agent is $s$ at time $t$. Unanimity is absorbing,
so $\mathrm{win}_t(s)$ is nondecreasing in $t$, and its supremum, the convergence probability, is
its limit.
\end{definition}

\begin{lemma}[Coupling and symmetry]\label{lem:am-coupling}
\lean{Crn.ApproxMajority.rank,Crn.ApproxMajority.rule_mono,Crn.ApproxMajority.graphStep_mono,Crn.ApproxMajority.event_allX_mono,Crn.ApproxMajority.swap,Crn.ApproxMajority.graphStep_swap,Crn.ApproxMajority.event_swap}\leanok\uses{def:winprob}
For the order $Y < B < X$ the rule is monotone in both states, so the probability that every agent
is $X$ at time $t$ is monotone in the initial configuration. Exchanging $X$ and $Y$ commutes with
the dynamics.
\end{lemma}
\begin{proof}\leanok
Run two ordered configurations with the same oriented edges; they stay ordered.
\end{proof}

\begin{lemma}[MNRS14, Corollary 1]\label{lem:hall}
\lean{Crn.ApproxMajority.exists_injective_of_regular,Crn.ApproxMajority.exists_injective_subset}\leanok
A family of $r$-element subsets of a finite universe, with at most as many sets as elements, in
which every element lies in the same number $d \ge 1$ of sets, has a system of distinct
representatives. In particular, if $n \le 2m$, every $m$-set $S$ of agents contains an
$(n - m)$-subset $A_S$, with $S \mapsto A_S$ injective on the $m$-sets.
\end{lemma}
\begin{proof}\leanok
Hall's theorem, with Hall's condition by double counting.
\end{proof}

\begin{theorem}[MNRS14, Theorem 4, at every finite time]\label{thm:am-graph}
\lean{Crn.ApproxMajority.winProb_minority_le_majority,Crn.ApproxMajority.absorbProb_minority_le_majority}\leanok\uses{lem:am-coupling,lem:hall}
On every graph with an edge, if $n \le 2m$, at every time $t$ the probability that the minority
$X$ has won is at most the probability that the majority $Y$ has won; hence the same holds for the
convergence probabilities.
\end{theorem}
\begin{proof}\leanok
For each $m$-set $S$, the configuration with $X$ exactly on $A_S$ lies below the one with $X$
exactly on $S$, which is the swap of the placement with $Y$ on $S$. By monotonicity and symmetry,
the minority wins from the first with probability at most the probability that the majority wins
from the placement on $S$. The first configuration is the placement on the $m$-set $A_S^c$, and
$S \mapsto A_S^c$ is a bijection of the $m$-sets, so averaging over $S$ gives the inequality
(eqs.~(18) to~(23)).
\end{proof}

\begin{theorem}[MNRS14, Theorem 4]\label{thm:am-half}
\lean{Crn.ApproxMajority.tendsto_event_unanimous,Crn.ApproxMajority.half_le_absorbProb_majority}\leanok\uses{thm:am-graph}
On a connected graph, from every configuration with a non-blank agent, the protocol is absorbed in
all-$X$ or all-$Y$ with probability tending to $1$. Hence, if $n \le 2m$ and $m \le n$, it converges
to the majority with probability at least $1/2$.
\end{theorem}
\begin{proof}\leanok
An initiator $s \ne B$ activated twice on the same responder makes it $s$, so an opinion present
spreads to all agents along the edges of a connected graph, and all-blank is never reached.
Absorption in a finite chain follows, and the two convergence probabilities sum to $1$.
\end{proof}

\section{Deterministic exact majority and plurality (CRN-5)}\label{sec:ghmss-exact-majority}
This section formalizes Theorems~3, 6 and~7 of Gąsieniec, Hamilton, Martin, Spirakis and
Stachowiak, \emph{Deterministic population protocols for exact majority and plurality} (OPODIS
2016, LIPIcs 70, paper 14). Colours are $k$-bit labels, so $C \le 2^k$ colours need $O(k)$ bits of
memory. Fig.~4 of the paper needs a minor correction: one entry of its table breaks Invariant~2,
and the rules of the text are used. The final argument of Section~3 needs a minor correction: it
uses, in the tie case, that every group keeps a strong state, an invariant the paper does not
state, which is added and proved. Section~5.2, which should inform all agents whether the most
frequent colour is unique, needs a major correction and is not formalized: on the colours
$1, 1, -1$ its bit never stabilizes, in any execution (a repaired rule passes exhaustive checks on
small populations). Theorem~7 is formalized as the protocol of Section~5.1 computes it: every
agent learns whether its own colour is the plurality colour, not the winner's label
(\href{https://github.com/formal-dynamics/leanamics/blob/main/crn/FORMALIZATION_DIFFERENCES.md#deterministic-exact-majority-and-plurality-exactmajority-crn-5}{details}).

\begin{definition}[Stable computation with general outputs]\label{def:em-marks}
\lean{Crn.Protocol.OutputStableAt,Crn.Protocol.StablyMarks,Crn.Protocol.StablyComputesWith,Crn.Protocol.stablyComputesWith_output_iff}\leanok\uses{def:stable}
For an output map $O : Q \to Y$ into any type and a target $g(x, i) \in Y$ depending on the input
counts $x$ and on an agent's input $i$, a protocol stably marks the agents with $g$ if, for every
nonempty population, every configuration reachable from the initial one can reach a configuration
from which every agent $v$ forever outputs $g(x, \iota(v))$. Stable computation of $f$ is the case
$g(x, i) = f(x)$; with the protocol's Boolean output it is Definition~\ref{def:stable}.
\end{definition}

\begin{definition}[The static protocol]\label{def:em-static}
\lean{Crn.StaticState,Crn.StaticState.δ,Crn.staticMajority,Crn.card_staticState}\leanok\uses{def:protocol}
Colours are in $\{-1, 0, 1\}$. The six states are the strong states $[x]$ and the weak states
$\langle x \rangle$, $x \in \{-1, 0, 1\}$, of weight $x$ and $0$. The states $[1]$ and $[-1]$
annihilate into $[0], [0]$; a strong $[x]$ turns a weak state into $\langle x \rangle$, and, for
$x \ne 0$, turns $[0]$ into $\langle x \rangle$. The output of $[x]$ and $\langle x \rangle$ is
$x$.
\end{definition}

\begin{lemma}[Lemmas 1 and 2]\label{lem:em-static}
\lean{Crn.StaticState.weight_sum_eq,Crn.StaticState.absWeight_sum_step_le,Crn.StaticState.exists_reaches_absWeight_sum_eq}\leanok\uses{def:em-static}
Along every execution the weight sum equals the colour sum $S$, the sum $R$ of the absolute
weights does not increase, and a configuration with $R = |S|$ is reachable.
\end{lemma}
\begin{proof}\leanok
Check the encounters (by computation over the $36$ pairs); while weights of both signs are
present, bring two of them together, which decreases $R$.
\end{proof}

\begin{theorem}[GHMSS16, Theorem 3]\label{thm:em-static}
\lean{Crn.staticMajority_stablyComputes}\leanok\uses{lem:em-static,def:em-marks}
The static protocol stably computes $\operatorname{sign}(\#1 - \#(-1))$: all agents eventually
and forever output $1$, $-1$, or $0$ for a tie.
\end{theorem}
\begin{proof}\leanok
Once $R = |S|$, every weight is $0$ or $s = \operatorname{sign} S$ and a strong agent $[s]$
exists (for $s = 0$, a strong state is never lost). It recruits every agent into the states of
output $s$, a set closed under the transitions.
\end{proof}

\begin{definition}[The dynamic protocol]\label{def:em-dynamic}
\lean{Crn.DynState,Crn.DynState.δ,Crn.DynState.recolour,Crn.dynamicMajority,Crn.card_dynState,Crn.Protocol.GStep,Crn.Protocol.GReaches}\leanok\uses{def:protocol}
The eight states are the strong states $[-2], \dots, [2]$ and the weak states
$\langle -1 \rangle, \langle 0 \rangle, \langle 1 \rangle$. Opposite strong states cancel as far
as possible, and strong states recruit weak ones and $[0]$ as in the static protocol; the
transitions do not read the colours. When an external force changes a colour $c$ to $c'$, the
state becomes $[w + c' - c]$, $w$ its weight. For Section~5 the protocol runs separately in groups
of agents, with encounters inside a group.
\end{definition}

\begin{definition}[Invariants of Section~3]\label{def:em-dyninv}
\lean{Crn.DynInv,Crn.DynInv.input}\leanok\uses{def:em-dynamic}
In every group the colour sum equals the weight sum, every agent's weight is within $1$ of its
colour, and every group contains a strong state. The initial configuration $[c_a]$ satisfies
them.
\end{definition}

\begin{lemma}[Lemmas 4 and 5]\label{lem:em-dynamic}
\lean{Crn.DynInv.gstep,Crn.DynInv.recolour,Crn.DynState.absWeight_sum_gstep_le,Crn.DynState.exists_greaches_noOpposite}\leanok\uses{def:em-dyninv}
The invariants are preserved by encounters inside groups and by colour changes. The sum of the
absolute weights does not increase, and a configuration without opposite weights in a group is
reachable.
\end{lemma}
\begin{proof}\leanok
Encounter by encounter (by computation over the finite state space); a colour change moves the
weight by the change of colour and makes the state strong.
\end{proof}

\begin{theorem}[GHMSS16, Section 3]\label{thm:em-dynamic}
\lean{Crn.DynInv.exists_stable,Crn.DynExtStep,Crn.dynamicMajority_stabilizes,Crn.dynamicMajority_stablyComputes}\leanok\uses{lem:em-dynamic,def:em-marks}
From every configuration satisfying the invariants, a configuration is reachable from which every
agent forever outputs the sign of the colour sum of its group. Hence after any finite interleaving
of encounters and colour changes the dynamic protocol stabilizes on the sign of the final colour
sum, and without the force it stably computes $\operatorname{sign}(\#1 - \#(-1))$.
\end{theorem}
\begin{proof}\leanok
After Lemma~5, every nonzero weight in a group has the sign $s$ of the group's colour sum, and a
strong agent of weight sign $s$ exists in the group (the third invariant gives it in the tie
case). It recruits the group into the states of output $s$.
\end{proof}

\begin{definition}[Driving the dynamic protocol]\label{def:em-drive}
\lean{Crn.Protocol.drive}\leanok\uses{def:em-dynamic}
Given a protocol $D$, a colour map from its states to $\{-1, 0, 1\}$ and a group map that $D$
never changes, the driven protocol runs $D$ and a copy of the dynamic protocol: in an encounter
$D$ makes its transition, the dynamic protocol makes its transition if the two agents are in the
same group, and each agent whose colour changed applies the state change of the force.
\end{definition}

\begin{theorem}[Composition]\label{thm:em-drive}
\lean{Crn.Protocol.drive_dynInv,Crn.Protocol.drive_stablyMarks}\leanok\uses{def:em-drive,thm:em-dynamic}
Along every execution of the driven protocol the invariants of Section~3 hold for the derived
colours. If $D$ stably marks the agents with their eventual colours, the driven protocol stably
marks every agent with the sign of the eventual colour sum of its group.
\end{theorem}
\begin{proof}\leanok
An encounter is an encounter of the dynamic protocol inside a group followed by two colour
changes. Once $D$ has stabilized, the colours are fixed and Theorem~\ref{thm:em-dynamic} applies
in every group.
\end{proof}

\begin{definition}[Absolute-Majority]\label{def:em-absolute}
\lean{Crn.Label,Crn.bitSign,Crn.bitsProtocol,Crn.majBit,Crn.absColour,Crn.absoluteMajority,Crn.absOutput,Crn.IsAbsMajority,Crn.absMajority,Crn.absTarget}\leanok\uses{def:em-static,def:em-drive}
Every agent keeps its label $l$ and runs $k$ copies of the static protocol on the bits $l[i]$,
read as $\pm 1$. They drive a copy of the dynamic protocol, with colour $1$ if every copy currently
reports the agent's own bit and $-1$ otherwise. An agent outputs the reported majority bits if its
dynamic protocol outputs $1$, and none otherwise. The specification is the colour held by more
than half of the agents, if any.
\end{definition}

\begin{theorem}[GHMSS16, Theorem 6]\label{thm:em-absolute}
\lean{Crn.bitsProtocol_stablyMarks,Crn.absMajority_eq_some_iff,Crn.absTarget_sum_sign_eq_one_iff,Crn.majBit_eq_of_isAbsMajority,Crn.absoluteMajority_stablyComputes,Crn.card_absState}\leanok\uses{def:em-absolute,thm:em-static,thm:em-drive}
Absolute-Majority stably computes the absolute majority colour, or reports that there is none,
with $2^k \cdot 6^k \cdot 8 \le 2^{4k+3}$ states.
\end{theorem}
\begin{proof}\leanok
The $k$ copies stabilize on the majority bits one at a time. Then the colour of an agent is $1$
exactly when its label consists of the majority bits, so the eventual colour sum is positive if
and only if that label is held by more than half of the agents, and the composition theorem
applies with a single group.
\end{proof}

\begin{definition}[Relative-Majority]\label{def:em-relative}
\lean{Crn.RelState,Crn.RelState.label,Crn.RelState.won,Crn.stageColour,Crn.relativeMajorityLevel,Crn.relativeMajority,Crn.IsGroupWinner,Crn.IsPlurality,Crn.duelColour,Crn.wins,Crn.bitAt,Crn.prefixMask}\leanok\uses{def:em-drive}
Stages $i = k - 1, \dots, 0$: the agents whose labels share the prefix $l[0..i-1]$ form a group,
and a copy of the dynamic protocol runs in each group, with colour $l[i]$ (as $\pm 1$) for the
agents that won all the stages above $i$ and $0$ for the others. An agent wins stage $i$ if its
colour is $-1$ and the stage outputs $-1$, or its colour is $1$ and the stage outputs $0$ or $1$.
The protocol is the $k$-fold iteration of the driven protocol, the colours being derived from the
states. The plurality colour is the most frequent one, ties broken in favour of the
lexicographically largest label.
\end{definition}

\begin{lemma}[The duel]\label{lem:em-duel}
\lean{Crn.duel_wins_iff,Crn.isGroupWinner_zero_iff}\leanok\uses{def:em-relative}
Once the stages above $i$ have stabilized, the duel of stage $i$ between the winners of the two
groups extending a prefix $p$ selects the winner of the group of $p$. The winner of the whole
population is the plurality colour.
\end{lemma}
\begin{proof}\leanok
The colour sum of the group of $p$ is the difference of the counts of the two group winners, and
ties go to the larger label.
\end{proof}

\begin{theorem}[GHMSS16, Theorem 7]\label{thm:em-relative}
\lean{Crn.relativeMajority_level_stablyMarks,Crn.relativeMajority_stablyMarks,Crn.card_relState_level,Crn.card_relState}\leanok\uses{lem:em-duel,thm:em-drive}
Relative-Majority stably marks the agents of the plurality colour: eventually and forever, every
agent outputs whether its own colour is the plurality colour. It has $2^k \cdot 8^k = 2^{4k}$
states.
\end{theorem}
\begin{proof}\leanok
Induction on the levels: level $m$ stably marks every agent with whether its label wins its group
of prefix length $k - m$, by the composition theorem and the duel lemma.
\end{proof}


%%FILE dynamics/blueprint/src/content.tex
\begin{abstract}
A shared library for finite stochastic dynamics using normalized real weights,
finite sums, and limits of finite-time probabilities. These are generic library
results, rather than statements attributed to a particular opinion-dynamics paper.
\end{abstract}

\section{Uniform expectations}
\begin{definition}[Finite uniform average]\label{def:avg}
\lean{Dynamics.avg}\leanok
For a finite type $A$, set $\operatorname{avg}(f)=\sum_{a\in A}f(a)/|A|$.
When $A$ is empty this is zero, following Lean's division convention.
\end{definition}
\begin{lemma}[Product independence]\label{lem:uniform-product}
\lean{Dynamics.avg_prod_pi,Dynamics.avg_equiv}\leanok\uses{def:avg}
Uniform averages factor over independent coordinates and are invariant under
bijections of the sample space.
\end{lemma}
\begin{lemma}[Finite trajectories]\label{lem:uniform-trajectory}
\lean{Dynamics.expList,Dynamics.expList_eq_avg_ofFn}\leanok\uses{def:avg}
The recursively defined expectation over $T$ independent draws equals the
uniform average over functions $\{0,\ldots,T-1\}\to A$.
\end{lemma}
\begin{proof}\leanok
Separate the first coordinate using the finite-function/product equivalence,
then induct on $T$. The empty sample space is handled separately.
\end{proof}
\begin{lemma}[Averages commute]\label{lem:avg-commute}
\lean{Dynamics.avg_avg_swap,Dynamics.expList_avg_comm,Dynamics.expList_zero_fun,Dynamics.expList_finset_sum}\leanok\uses{lem:uniform-trajectory}
Iterated uniform averages over two finite sets commute; in particular an outer
average commutes with the trajectory expectation $\mathrm{expList}$, which also
commutes with finite sums and vanishes on the zero functional.
\end{lemma}
\begin{proof}\leanok
Exchange the order of finite summation; for $\mathrm{expList}$, induct on $T$.
\end{proof}

\section{Weighted distributions and kernels}
\begin{definition}[Normalized real weights]\label{def:distribution}
\lean{Dynamics.Distribution,Dynamics.Distribution.expect,Dynamics.Distribution.prob}\leanok
A distribution consists of nonnegative real weights $p_a$ with $\sum_a p_a=1$.
Its expectation is $\sum_a p_a f(a)$, and an event probability is the expectation
of its indicator. Such a distribution cannot exist on an empty type.
\end{definition}
\begin{lemma}[Uniform agreement]\label{lem:uniform-agreement}
\lean{Dynamics.Distribution.uniform_expect}\leanok\uses{def:distribution,def:avg}
On a nonempty finite type the normalized uniform distribution agrees with the
existing finite average.
\end{lemma}
\begin{lemma}[Independent products and pushforward]\label{lem:weighted-product}
\lean{Dynamics.Distribution.independent_expect_prod,Dynamics.Distribution.independent_expect_eval,Dynamics.Distribution.map_expect}\leanok\uses{def:distribution}
Independent coordinates have product weights. Expectations of product
observables factor, and transporting weights along a function commutes with expectation.
\end{lemma}
\begin{proof}\leanok
Distribute products over finite sums; for pushforward interchange the two sums
and evaluate each fiber indicator.
\end{proof}
\begin{lemma}[Bridges between the expectation APIs]\label{lem:bridge}
\lean{Dynamics.avg_eq_expect,Dynamics.expList_eq_expect,Dynamics.Distribution.independent_uniform_expect,Dynamics.expList_eq_independent_expect}\leanok\uses{lem:uniform-trajectory,lem:uniform-agreement,lem:weighted-product}
The finite average is Mathlib's finite expectation, $\operatorname{avg}(f)=\E_{a\in A}f(a)$
(\texttt{Finset.expect}); both vanish on an empty type. On a nonempty type the independent
product of uniform distributions on $A$, indexed by a finite set $I$, is the uniform
distribution on $A^I$. Hence $\mathrm{expList}$ over $T$ rounds is the expectation of
$F$ over $T$ i.i.d.\ uniform draws, both as Mathlib's $\E_{\omega\colon\{0,\ldots,T-1\}\to A}$
and under the product distribution.
\end{lemma}
\begin{proof}\leanok
Mathlib's finite expectation is a sum divided by a cardinality. Every point of $A^I$ has
product weight $|A|^{-|I|}=|A^I|^{-1}$. Combine both with Lemma~\ref{lem:uniform-trajectory}.
\end{proof}
\begin{definition}[Transition kernel]\label{def:kernel}
\lean{Dynamics.Kernel,Dynamics.Kernel.iterate,Dynamics.Kernel.trajectory}\leanok\uses{def:distribution}
A finite kernel assigns a distribution to every state. Its operator iterates
compute expectations at finite times; trajectory functionals additionally
record the history.
\end{definition}
\begin{lemma}[Endpoint agreement]\label{lem:endpoint}
\lean{Dynamics.Kernel.trajectory_endpoint}\leanok\uses{def:kernel}
Ignoring history in a trajectory expectation recovers the iterated transition operator.
\end{lemma}

\section{Stationarity and absorption}
\begin{theorem}[Existence of stationary weights]\label{thm:stationary}
\lean{Dynamics.Kernel.exists_stationary}\leanok\uses{def:kernel}
Every stochastic kernel on a nonempty finite type has a distribution $p$ with $pK=p$.
\end{theorem}
\begin{proof}\leanok
Average the first $n+1$ evolved distributions. The stationarity defect of this
Ces\`aro average is $(p_{n+1}-p_0)/(n+1)$, which tends to zero coordinatewise.
Compactness of the finite probability simplex gives a convergent subsequence;
continuity of finite sums implies that its limit is stationary.
\end{proof}
\begin{lemma}[Uniform absorption block]\label{lem:block}
\lean{Dynamics.Kernel.exists_uniform_block}\leanok\uses{def:kernel}
Let $f$ be a zero-one survival indicator with $Kf\le f$. Suppose for each
state $a$ there is $n_a$ with $(K^{n_a}f)(a)<1$. There are $m>0$ and
$0\le q<1$ such that $K^m f\le qf$.
\end{lemma}
\begin{proof}\leanok
Take one plus the maximum of the finitely many access times. Antitonicity of
survival probabilities puts every state below one at that time. Their finite
maximum supplies $q$. At states with $f=0$, absorption keeps the expectation zero.
\end{proof}
\begin{theorem}[Geometric decay and convergence]\label{thm:absorption}
\lean{Dynamics.Kernel.geometric_blocks,Dynamics.Kernel.finite_absorption}\leanok\uses{lem:block}
Under these hypotheses, $K^{nm}f\le q^n f$, and $K^t f(a)\to0$ for every state $a$.
\end{theorem}
\begin{proof}\leanok
Induct over blocks using positivity and linearity. The geometric bound tends
to zero. For arbitrary $t$, compare with $m\lfloor t/m\rfloor$ by antitonicity.
\end{proof}
\begin{theorem}[Finite-horizon optional stopping]\label{thm:optional-stopping}
\lean{Dynamics.Kernel.event_eq_iterate,Dynamics.Kernel.eq_add_target_add_survival,Dynamics.Kernel.iterate_eq_event_add,Dynamics.Kernel.event_error_of_invariant,Dynamics.Kernel.tendsto_event_of_invariant,Dynamics.Kernel.iSup_event_of_invariant}\leanok\uses{def:kernel}
Let $\varphi$ equal $\varphi_A$ on an event $A$ and $\varphi_B$ on an event $B$, and let
$S$ be the event of lying in neither. If $\E_{x_0}\varphi(X_T)=\varphi(x_0)$ and
$m\le\varphi-\varphi_B\le M$ on $S$, then
$\varphi(x_0)-\varphi_B-(\varphi_A-\varphi_B)\Pr_{x_0}(X_T\in A)$ lies between
$m\Pr_{x_0}(X_T\in S)$ and $M\Pr_{x_0}(X_T\in S)$. Hence, if $\varphi_A\ne\varphi_B$,
the conservation law holds for every $T$ and $\Pr_{x_0}(X_T\in S)\to0$, then
$\Pr_{x_0}(X_T\in A)\to(\varphi(x_0)-\varphi_B)/(\varphi_A-\varphi_B)$; if moreover $A$
is absorbing, this is the supremum of the finite-time probabilities.
\end{theorem}
\begin{proof}\leanok
Pointwise, $\varphi-\varphi_B=(\varphi_A-\varphi_B)1_A+(\varphi-\varphi_B)1_S$ (on
$A\cap B$ both sides vanish). Take expectations at time $T$ and bound the last term by
monotonicity. For the limit, take $M=-m=\sum_a|\varphi(a)-\varphi_B|$ and squeeze. If $A$
is absorbing, the finite-time probabilities increase, so their limit is their supremum.
\end{proof}

\section{Rounds, concentration and phases}
\begin{definition}[Kernel of a round-based process]\label{def:ofstep}
\lean{Dynamics.Kernel.ofStep,Dynamics.Kernel.iterate_ofStep,Dynamics.Kernel.event_ofStep}\leanok\uses{def:kernel}
A deterministic update $\mathrm{step}\colon S\times R\to S$ driven by a uniform
round $r\in R$ defines the kernel $K(s)=\mathrm{law}(\mathrm{step}(s,r))$. Its
$T$-step expectations are expectations over $T$ i.i.d.\ uniform rounds.
\end{definition}
\begin{proof}\leanok
Induction on $T$, unfolding one round of $\mathrm{expList}$.
\end{proof}
\begin{theorem}[Time reversal of i.i.d.\ rounds]\label{thm:reverse}
\lean{Dynamics.reverse_ofFn,Dynamics.expList_comp_reverse,Dynamics.Kernel.iterate_ofStep_foldr}\leanok\uses{lem:uniform-trajectory,def:ofstep}
For every functional $F$ of $T$ i.i.d.\ uniform rounds,
$\E F(r_T,\dots,r_1)=\E F(r_1,\dots,r_T)$. In particular the $T$-step expectations of a
round-based process are also obtained by applying the rounds last to first.
\end{theorem}
\begin{proof}\leanok
Write both sides as uniform averages over $\mathrm{Fin}\,T\to R$ and reindex by the
involution $i\mapsto T-1-i$.
\end{proof}
\begin{theorem}[Hoeffding and Bernstein inequalities]\label{thm:concentration}
\lean{Dynamics.avg_hoeffding,Dynamics.avg_bernstein}\leanok\uses{def:avg}
Let $X=\sum_{i<n}Y_i(\omega_i)$ with independent uniform coordinates. If each
$Y_i$ is $\{0,1\}$-valued, $\Pr(X\ge\E X+\lambda)\le e^{-2\lambda^2/n}$. If
$Y_i-\E Y_i\le b$ and $\sum_i\mathrm{Var}\,Y_i\le\sigma^2$, then
$\Pr(X\ge\E X+\lambda)\le\exp\bigl(-\lambda^2/(2\sigma^2(1+b\lambda/(3\sigma^2)))\bigr)$.
\end{theorem}
\begin{proof}\leanok
Markov's inequality applied to $e^{tX}$, which factors over the coordinates.
Hoeffding's lemma bounds each Bernoulli factor; for Bernstein,
$e^y\le1+y+y^2/(2(1-u/3))$ for $y\le u<3$ bounds each factor, and $t$ is
chosen as $\lambda/(\sigma^2+b\lambda/3)$.
\end{proof}
\begin{theorem}[Chernoff bounds for independent Bernoulli trials]\label{thm:chernoff}
\lean{Dynamics.Distribution.bernoulli,Dynamics.Distribution.chernoff_upper_ratio,Dynamics.Distribution.chernoff_upper,Dynamics.Distribution.chernoff_lower_ratio,Dynamics.Distribution.chernoff_lower,Dynamics.Distribution.bernoulli_chernoff_upper_ratio,Dynamics.Distribution.bernoulli_chernoff_upper,Dynamics.Distribution.bernoulli_chernoff_lower_ratio,Dynamics.Distribution.bernoulli_chernoff_lower,Dynamics.avg_chernoff_upper_ratio,Dynamics.avg_chernoff_upper,Dynamics.avg_chernoff_lower_ratio,Dynamics.avg_chernoff_lower}\leanok\uses{def:distribution,def:avg}
Let $X=\sum_iX_i$ be a sum of independent $\{0,1\}$-valued trials with
$\Pr(X_i=1)=p_i$, and let $\mu_L\le\mu=\sum_ip_i\le\mu_H$. For $\delta>0$,
\[\Pr(X\ge(1+\delta)\mu_H)\le\Bigl(\frac{e^{\delta}}{(1+\delta)^{1+\delta}}\Bigr)^{\mu_H}
\le e^{-\delta^2\mu_H/(2+\delta)},\]
and for $0<\delta<1$,
\[\Pr(X\le(1-\delta)\mu_L)\le\Bigl(\frac{e^{-\delta}}{(1-\delta)^{1-\delta}}\Bigr)^{\mu_L}
\le e^{-\delta^2\mu_L/2}.\]
The trials are $\{0,1\}$-valued observables of the coordinates of an independent
product of distributions, independent biased coins (counting the heads in any
subset), or the coordinates of one uniform round (Mitzenmacher--Upfal,
\emph{Probability and Computing}, Theorems 4.4 and 4.5 with Exercise 4.7).
\end{theorem}
\begin{proof}\leanok
Markov's inequality applied to $e^{tX}$, which factors over the coordinates; each
factor is $1+p_i(e^t-1)\le e^{p_i(e^t-1)}$, so $\E e^{tX}\le e^{\mu_H(e^t-1)}$ for
$t\ge0$ and $\E e^{tX}\le e^{\mu_L(e^t-1)}$ for $t\le0$. Take $t=\log(1+\delta)$,
respectively $t=\log(1-\delta)$. The closed forms follow from
$\log(1+\delta)\ge2\delta/(2+\delta)$ (the first term of the series of
$\log\frac{1+x}{1-x}$ at $x=\delta/(2+\delta)$) and
$(1-\delta)\log(1-\delta)\ge-\delta+\delta^2/2$.
\end{proof}
\begin{theorem}[Lower tails, MGF and parametric Chernoff bounds]\label{thm:chernoffmgf}
\lean{Dynamics.avg_lower_tail_le_of_mgf,Dynamics.avg_hoeffding_lower,Dynamics.variance_le_avg_of_zero_one,Dynamics.avg_chernoff_mgf,Dynamics.avg_chernoff_upper_of_mgf,Dynamics.avg_chernoff_lower_of_mgf,Dynamics.avg_chernoff_upper_log,Dynamics.avg_chernoff_lower_log,Dynamics.avg_chernoff_lower_mul}\leanok\uses{thm:concentration,thm:chernoff}
Let $X=\sum_{i<n}Y_i(\omega_i)$ with independent uniform $\{0,1\}$-valued
coordinates and $\mu=\E X$. Then $\Pr(X\le\E X-\lambda)\le e^{-2\lambda^2/n}$, and a
$\{0,1\}$ coordinate has variance at most its mean. For every real $t$,
$\E e^{tX}\le\exp(\mu(e^t-1))$; hence $\Pr(X\ge k)\le\exp(\mu(e^t-1)-tk)$ for
$t\ge0$ and $\Pr(X\le k)\le\exp(\mu(e^t-1)-tk)$ for $t\le0$. If $\E X\le\mu\le k$
with $\mu>0$, then $\Pr(X\ge k)\le\exp(k-\mu-k\log(k/\mu))$; if $0<k\le\mu\le\E X$,
then $\Pr(X\le k)\le\exp(k-\mu-k\log(k/\mu))$. Finally
$\Pr(X\le(1-\delta)\E X)\le\exp(-\delta^2\E X/2)$ for $0\le\delta<1$.
\end{theorem}
\begin{proof}\leanok
The lower Hoeffding tail is the upper one for $1-Y_i$. For the MGF bound,
$e^{tY}=1+(e^t-1)Y$ on $\{0,1\}$ and $1+y\le e^y$, factor by factor; Markov's
inequality applied to $e^{tX}$ gives the tails, and $t=\log(k/\mu)$ the closed forms
(the exponent is monotone in $\mu$ in the right direction). The last bound is the
lower bound of Theorem~\ref{thm:chernoff} with $\mu_L=\E X$ (trivial for $\delta=0$).
\end{proof}
\begin{lemma}[Elementary tail bounds and monotonicity]\label{lem:tail}
\lean{Dynamics.Distribution.prob_mono,Dynamics.Distribution.expect_lt_one,Dynamics.avg_lt_one,Dynamics.avg_markov,Dynamics.avg_markov_one,Dynamics.Kernel.one_sub_avg_le_prob_ofStep,Dynamics.Kernel.iterate_monotone,Dynamics.Kernel.event_monotone,Dynamics.exp_neg_two_log,Dynamics.one_le_of_log_pos}\leanok\uses{def:distribution,def:ofstep}
Probability is monotone in the event. An observable $f\le1$ with $f(b)<1$ at a
point of positive weight has expectation $<1$. Markov's inequality
$\Pr(X\ge c)\le\E X/c$ for $X\ge0$, $c>0$; in particular
$\Pr(\sum_iY_i(\omega_i)\ge1)\le\sum_i\E Y_i$ for nonnegative coordinates. One
round of $\mathrm{step}$ lands in an event with probability at least
$1-\E[\mathrm{bad}]$ when $\mathrm{bad}\ge0$ is at least $1$ on every miss. If
$f\le Kf$, then $n\mapsto K^nf(a)$ is nondecreasing; in particular the occupation
probability of an absorbing event is nondecreasing in time. Finally
$e^{-2\log n}=1/n^2$ for $n\ge1$, and $\log n>0$ forces $n\ge1$.
\end{lemma}
\begin{proof}\leanok
Pointwise comparisons of indicators, integrated by monotonicity of the expectation;
for the kernel statements, induction on time with $K^{n+1}f=K^n(Kf)$.
\end{proof}
\begin{theorem}[Nested phases]\label{thm:phases}
\lean{Dynamics.Kernel.nested_phases}\leanok\uses{def:kernel}
Let $A_1\supseteq\cdots\supseteq A_T$ be sets of states such that from $A_i$
the chain stays in $A_i$ with probability at least $1-\varepsilon$ and, for
$i<T$, enters $A_{i+1}$ with probability at least $1-\nu$. From $A_1$, after
$\ell T$ steps the chain is in $A_T$ with probability at least
$1-T(\ell\varepsilon+\nu^\ell)$.
\end{theorem}
\begin{proof}\leanok
Track the indicator of lying outside a set. In $\ell$ steps the chain leaves
$A_i$ with probability at most $\ell\varepsilon$, and misses $A_{i+1}$ in all
$\ell$ steps with probability at most $\nu^\ell$. Sum over the $T$ phases.
\end{proof}
\section{Drift and graph rounds}
\begin{definition}[Time-dependent chain]\label{def:iterateseq}
\lean{Dynamics.Kernel.iterateSeq,Dynamics.Kernel.iterateSeq_const}\leanok\uses{def:kernel}
A family of kernels $(K_t)_{t\ge0}$ moves the chain from time $t$ to time
$t+1$ with $K_t$; $\mathrm{iterateSeq}(K,n,f)(a)$ is the expectation of $f$ at
time $n$ from $a$. A constant family gives the iterates of a single kernel.
\end{definition}
\begin{proof}\leanok
Induction on $n$.
\end{proof}
\begin{theorem}[Drift lemma; Berenbrink et al., ICALP 2016, Lemma 2.2]\label{thm:drift}
\lean{Dynamics.Kernel.drift_absorption_seq,Dynamics.Kernel.drift_absorption}\leanok
\uses{def:iterateseq}
Let $\Psi\ge0$, and for $t<T$ let $c_t\ge0$,
$\E[\Psi(X_{t+1})\mid X_t=x]\le\Psi(x)-c_t/\Psi(x)$ when $\Psi(x)>0$, and
$\E[\Psi(X_{t+1})\mid X_t=x]=0$ when $\Psi(x)=0$. If
$\sum_{t<T}c_t\ge4\Psi(x_0)^2$, then $\Pr_{x_0}(\Psi(X_T)=0)\ge1/2$.
\end{theorem}
\begin{proof}\leanok
Let $q=\Pr(\Psi(X_T)>0)$; survival is nonincreasing, so it is at least $q$
before $T$. With $\lambda=q/\Psi(x_0)$, the bound $1/y\ge2\lambda-\lambda^2y$
linearizes the drift: $\E\Psi(X_{t+1})\le(1+c_t\lambda^2)\E\Psi(X_t)
-2c_t\lambda\Pr(\Psi(X_t)>0)\le\E\Psi(X_t)-c_t\lambda q$. (The paper uses
Jensen's inequality here; this is its optimal-$\lambda$ form.) Summing,
$0<m q\le\E\Psi(X_T)\le\Psi(x_0)(1-4q^2)$ if $q>0$, where $m>0$ is the least
positive value of $\Psi$; hence $q\le1/2$.
\end{proof}
\begin{theorem}[Multiplicative drift; the argument of Lemma 2.4 there]\label{thm:multdrift}
\lean{Dynamics.Kernel.multiplicative_drift_seq,Dynamics.Kernel.multiplicative_drift}\leanok
\uses{def:iterateseq}
Let $\Psi\ge0$ with every positive value at least $\Psi_{\min}>0$, and
$\E[\Psi(X_{t+1})\mid X_t=x]\le(1-\delta_t)\Psi(x)$ for all $x$ and $t<T$.
Then $\Pr_{x_0}(\Psi(X_T)>0)\le\prod_{t<T}(1-\delta_t)\,\Psi(x_0)/\Psi_{\min}$.
\end{theorem}
\begin{proof}\leanok
By induction $\E\Psi(X_T)\le\prod_{t<T}(1-\delta_t)\Psi(x_0)$ (the factors are
nonnegative unless $\Psi\equiv0$); then Markov's inequality
$\mathbf 1_{\Psi>0}\le\Psi/\Psi_{\min}$.
\end{proof}
\begin{lemma}[Synchronous neighbour rounds]\label{lem:neighborround}
\lean{Dynamics.NeighborRound,Dynamics.avg_neighborSet,Dynamics.neighborRound_nonempty,Dynamics.card_neighborRound,Dynamics.avg_neighborRound_prod,Dynamics.avg_neighborRound_eval,Dynamics.avg_neighborRound_mul,Dynamics.avg_neighborRound_sum,Dynamics.avg_vertex_neighborRound}\leanok
\uses{def:avg}
In a round $r\in\prod_vN(v)$ every vertex picks a neighbour. Under the uniform
distribution the picks are independent and uniform:
$\E\prod_vf_v(r_v)=\prod_v\E_{u\in N(v)}f_v(u)$, and, without isolated
vertices, $\E g(r_v)=\E_{u\in N(v)}g(u)$. A uniform vertex $v$ with a uniform
round gives the sequential step ``uniform vertex, uniform neighbour $r_v$''.
\end{lemma}
\begin{proof}\leanok
$\sum_r\prod_vf_v(r_v)=\prod_v\sum_{u\in N(v)}f_v(u)$ and
$|\prod_vN(v)|=\prod_v\deg v$; the marginals take $f_w=1$ for $w\ne v$.
\end{proof}
\begin{lemma}[Uniform oriented edge]\label{lem:edgeround}
\lean{Dynamics.EdgeRound,Dynamics.edgeRound_nonempty,Dynamics.avg_edgeRound,Dynamics.avg_edgeRound_symm,Dynamics.avg_edgeRound_edge,Dynamics.avg_edgeRound_fst}\leanok
\uses{def:avg}
For a uniform dart $d$ of $G$,
$\E F(d_1,d_2)=\sum_v\sum_{u\in N(v)}F(v,u)/(2|E|)$; reversing $d$ preserves
the law, its underlying edge is a uniform edge, and its tail is $v$ with
probability $\deg v/(2|E|)$.
\end{lemma}
\begin{proof}\leanok
Group the darts by tail (the darts with tail $v$ are the neighbours of $v$)
or by edge (each edge carries two darts).
\end{proof}
\section{Hitting times}\label{sec:dgmss-hitting}
This section formalizes Claim~2.9 of Doerr, Goldberg, Minder, Sauerwald and Scheideler,
\emph{Stabilizing consensus with the power of two choices} (SPAA 2011), as
Theorem~\ref{thm:drifthitting}. Compared with the claim, the hypothesis $c_4\log q\le q$ is added
(the target is a possible value, which a non-asymptotic statement needs), and $c_5$ also depends
on $c_1,c_2,c_3$
(\href{https://github.com/formal-dynamics/leanamics/blob/main/dynamics/FORMALIZATION_DIFFERENCES.md#hitting-time-lemma-drifthitting-maj-8-doerr-et-al-spaa-2011-claim-29}{details}).
\begin{definition}[Hitting probability]\label{def:hitprob}
\lean{Dynamics.Kernel.hitProb}\leanok\uses{def:kernel}
$\mathrm{hitProb}(K,B,n)(a)=\Pr_a(\exists s\le n,\ X_s\in B)=\Pr_a(T_B\le n)$,
the expectation over the first $n$ transitions of the indicator that one of
the states at the times $0,\dots,n$ lies in $B$.
\end{definition}
\begin{lemma}[Geometric drift bound for hitting times]\label{lem:hitdrift}
\lean{Dynamics.Kernel.stopped,Dynamics.Kernel.one_sub_hitProb_eq_event_stopped,Dynamics.Kernel.one_sub_hitProb_le_of_drift}\leanok
\uses{def:hitprob,thm:multdrift}
Let $V\ge0$ vanish on $B$ and be at least $V_{\min}>0$ off $B$, with
$\E[V(X_{t+1})\mid X_t=x]\le\rho V(x)$ for $x\notin B$. Then
$\Pr_a(T_B>t)\le\rho^tV(a)/V_{\min}$.
\end{lemma}
\begin{proof}\leanok
In the chain stopped on entering $B$, being outside $B$ at time $t$ is the
event $T_B>t$; apply Theorem~\ref{thm:multdrift} to the stopped chain with
$\delta=1-\rho$.
\end{proof}
\begin{theorem}[Hitting-time bound; Doerr, Goldberg, Minder, Sauerwald, Scheideler, SPAA 2011, Claim 2.9]\label{thm:drifthitting}
\lean{Dynamics.Kernel.drift_hitting,Dynamics.Kernel.drift_hitting_log}\leanok
\uses{lem:hitdrift}
Let $c_1>1$ and $c_2,c_3,c_4,c_6>0$. There is $c_5>0$ such that for every
finite chain and observable $X\in\{0,\dots,q\}$ with $c_4\log q\le q$: if
from every state below the target $\Pr[X_{t+1}\ge\min\{c_1X_t,q\}]\ge
1-e^{-c_2X_t}$, and $\Pr[X_{t+1}\ge1]\ge c_3$ when $X_t=0$, then
$X$ reaches $c_4\log q$ within $c_5\log q+\log_{c_1}(c_4\log q)$ steps with
probability at least $1-q^{-c_6}$ (equivalently within $C\log q$ steps).
\end{theorem}
\begin{proof}\leanok
The paper omits the proof. Let $p=\min\{c_3,1-e^{-c_2}\}$, $\theta=1/p$ and
$\beta=c_2/2$. For a large threshold $x_0$ take
$g(x)=1-\eta(\theta^x-1)$ for $x\le x_0$, with $\eta$ chosen so that
$g(x_0)=1/2$, and $g(x)=\tfrac12e^{-\beta(x-x_0)}$ for $x\ge x_0$; let
$V=g(X)$ below the target and $V=0$ on it. Below $x_0$ a step goes up by at
least one with probability at least $p$, and
$pg(x+1)+1-p\le\rho g(x)$ with $\rho=1-(1-p)\eta/2$; above $x_0$ growth by
$c_1$ fails with probability at most $e^{-c_2x}$, and the expected potential is
at most $\tfrac38e^{-\beta(x-x_0)}\le\rho g(x)$. Since $V\le1$ and
$V\ge\tfrac12q^{-\beta c_4}$ below the target, Lemma~\ref{lem:hitdrift} gives
$\Pr(T>t)\le2\rho^tq^{\beta c_4}\le q^{-c_6}$ for
$t\ge(1+\beta c_4+c_6)\log q/\log(1/\rho)$.
\end{proof}


%%FILE epidemics/blueprint/src/content.tex
\begin{abstract}
We formalize the pathwise correspondence between the Reed--Frost (Independent Cascade) epidemic
and bond percolation, after Kempe, Kleinberg and Tardos: with one coin per edge, the nodes
infected in round $t$ are exactly those at distance $t$ from the initial set in the graph of
open edges. Consequently the final outbreak is the set of nodes connected to the initial set by
open edges, and with i.i.d.\ Bernoulli$(p)$ coins the probability of eventual infection is a
bond-percolation connection probability. Below the threshold, on graphs of maximum degree $d$ with
$p(d-1)\le1-\varepsilon$, every cluster has at most $(10/\varepsilon^2)\log n$ vertices with
probability at least $1-1/n$ (Becchetti et al.), so a subcritical Reed--Frost epidemic stays
small and dies out within $O(\log n)$ rounds. We then formalize the supercritical phase of the
Erd\H{o}s--R\'enyi random graph $G(n,(1+\varepsilon)/n)$ (bond percolation on $K_n$) by the
depth-first-search argument of Krivelevich and Sudakov: with probability at least $1-C/n$ there is
a path of length $\varepsilon^2 n/5$ and a component of $\varepsilon n/2$ vertices; read through
the coupling, a Reed--Frost epidemic on $K_n$ with $R_0=pn>1$ infects $\Omega(n)$ nodes with
probability $\Omega(1)$. The development is in the finite layer of the shared \texttt{dynamics}
package.

We also formalize the deterministic Kermack--McKendrick SIR model: for every solution of
$s'=-\beta si$, $i'=\beta si-\gamma i$, $r'=\gamma i$ on $[0,\infty)$ with $r(0)=0$, the threshold
theorem ($i$ initially increases iff $R_0s(0)>1$, with $R_0=\beta/\gamma$), the epidemic peak, and
the final-size equation $s_\infty=s(0)e^{-R_0(1-s_\infty)}$, following Hethcote (SIAM Review
2000, Theorem~2.1). This part uses Mathlib's real analysis.

We also formalize the upper bounds of Doerr and Kostrygin's analysis of general rumor-spreading
processes (\emph{Randomized rumor spreading revisited}, ICALP 2017): the exponential growth and
shrinking regimes and the total spreading time $\log_{1+\gamma}n+\frac1\rho\ln n+O(1)$, with
exponential tails. Lemma~20 of that paper needs a major correction; we prove corrected versions.

We also formalize the duality between the coalescing-branching random walk (COBRA) and the BIPS
epidemic of Cooper, Radzik and Rivera (Theorem~4), pathwise and by time reversal of i.i.d.\ rounds,
and their bounds on the COBRA cover time on regular expanders: $O(\log n/(1-\lambda)^3)$ when
$1-\lambda\ge C_0\sqrt{\log n/n}$ (Theorems~1 and~2), and $O(\log n)$ with branching factor
$1+\rho$ for constant $\lambda<1$ and $\rho>0$ (Theorem~3), with high probability and in
expectation.
\end{abstract}

\section{The model}\label{sec:kkt-live-edge}
This section and the three following ones (distances, pathwise layers, random coins) formalize
the live-edge coupling of Kempe, Kleinberg and Tardos (KDD 2003) for the Independent Cascade
model, also Theorem~A.3 of Becchetti, Clementi, Denni, Pasquale, Trevisan and Ziccardi
(arXiv:2103.16398).

\begin{definition}[Percolated graph and epidemic]\label{def:reedfrost}
\lean{Epidemics.perc,Epidemics.SIR,Epidemics.step,Epidemics.run,Epidemics.setDist}\leanok
Every edge $e$ of a finite graph $G$ carries a coin $\omega(e)\in\{\mathrm{true},\mathrm{false}\}$;
the percolated graph $G_\omega$ keeps the edges whose coin is true. In one round, a susceptible node
becomes infected if it has an infected neighbour across an open edge, and every infected node
recovers. Initially the set $I_0$ is infected and nobody is recovered. We write $d(v)\in\mathbb N\cup\{\infty\}$
for the distance from $I_0$ to $v$ in $G_\omega$.
\end{definition}

\section{Distances from a set}
\begin{lemma}[Distance facts]\label{lem:dist}
\lean{Epidemics.setDist_eq_subtype,Epidemics.setDist_eq_zero_iff,Epidemics.setDist_le_add_one,Epidemics.exists_edist_eq_setDist,Epidemics.exists_adj_of_setDist_succ,Epidemics.setDist_ne_top_iff,Epidemics.setDist_lt_card}\leanok\uses{def:reedfrost}
$d(v)=0$ iff $v\in I_0$; if $u$ and $v$ are adjacent in $G_\omega$ then $d(v)\le d(u)+1$; if
$d(v)=t+1$ then $v$ has a neighbour $u$ in $G_\omega$ with $d(u)=t$; $d(v)<\infty$ iff $v$ is
connected to $I_0$ in $G_\omega$, and then $d(v)<|V|$.
\end{lemma}
\begin{proof}\leanok
The infimum over $I_0$ is attained; take the second-to-last vertex of a shortest walk; a shortest
walk is a path.
\end{proof}

\section{Pathwise layers}
\begin{theorem}[Infected and recovered sets]\label{thm:layers}
\lean{Epidemics.infected_recovered_iff,Epidemics.infected_iff,Epidemics.recovered_iff}\leanok\uses{lem:dist}
For every coin assignment and every $t$, the nodes infected in round $t$ are those with $d(v)=t$,
and the nodes recovered by round $t$ are those with $d(v)<t$.
\end{theorem}
\begin{proof}\leanok
Both statements together, by induction on $t$, using Lemma~\ref{lem:dist}.
\end{proof}
\begin{corollary}[Extinction and final size]\label{cor:final}
\lean{Epidemics.no_infected_of_dist_lt,Epidemics.extinct_of_dist_lt,Epidemics.extinct,Epidemics.final_recovered_iff}\leanok\uses{thm:layers}
Nobody is infected in round $t$ as soon as every finite $d(v)$ is $<t$, in particular for
$t\ge|V|$; the nodes recovered after $|V|$ rounds are those connected to $I_0$ in $G_\omega$.
\end{corollary}
\begin{proof}\leanok
Immediate from Theorem~\ref{thm:layers} and $d(v)<|V|$ when finite.
\end{proof}

\section{Random coins}
\begin{theorem}[Reed--Frost $\Leftrightarrow$ bond percolation]\label{thm:perc}
\lean{Epidemics.bernoulli,Epidemics.coins,Epidemics.coins_prob_open,Epidemics.prob_infected_eq_prob_connected}\leanok\uses{cor:final}
With independent Bernoulli$(p)$ coins, each edge is open with probability $p$, and the
probability that $v$ is eventually infected equals the probability that $v$ is connected to $I_0$
by open edges.
\end{theorem}
\begin{proof}\leanok
The two events coincide for every coin assignment (Corollary~\ref{cor:final}).
\end{proof}
The coupling is pathwise because every edge is used at most once: once its coin has been seen
from an infected endpoint, that endpoint is recovered for good. This is why a single coin per edge
gives the same process, in distribution, as fresh coins in every round (that distributional
statement is not formalized here).

\section{Subcritical percolation and small outbreaks (EPI-2)}\label{sec:bcdptz-small-world}
After Becchetti, Clementi, Denni, Pasquale, Trevisan, Ziccardi (arXiv:2103.16398), Theorem 2.3
and its proof, Theorem E.1. Throughout, the degrees of $G$ are at most $d$, $n = |V|$, and
$p(d-1) \le 1 - \varepsilon$ with $0 < \varepsilon < 1$. The small-world part of the paper is
formalized in the next section (EPI-6).

A few points of the paper need a minor correction, and the corrected versions are what is
formalized: the Chernoff step in the proof of Theorem~E.1, which does not account for the extra
trial at the source (the formal bound is $e^{\varepsilon}\exp(-\varepsilon^2 t/2)$, still of the
form $\exp(-\Omega(\varepsilon^2 t))$ that the theorem states), and a typo in its union bound
(\href{https://github.com/formal-dynamics/leanamics/blob/main/epidemics/FORMALIZATION_DIFFERENCES.md#subcritical-percolation-and-small-outbreaks-subcritical-epi-2}{details}),
and, in the small-world part, the hypothesis of claim~2 of Theorem~2.5 and the starting set of
the exploration in Algorithm~3
(\href{https://github.com/formal-dynamics/leanamics/blob/main/epidemics/FORMALIZATION_DIFFERENCES.md#small-world-networks-below-the-percolation-threshold-smallworld-epi-6}{details}).

\begin{lemma}[Binomial tail and Chernoff bound]\label{lem:chernoff}
\lean{Epidemics.binTail,Epidemics.binTail_succ_succ,Epidemics.one_sub_mul_exp_le,Epidemics.binomial_tail_le}\leanok\uses{thm:perc}
At least $t$ successes among $t(d-1)+1$ independent Bernoulli$(p)$ trials occur with probability
at most $\exp(\varepsilon - \varepsilon^2 t/2)$.
\end{lemma}
\begin{proof}\leanok
Markov's inequality on $e^{\varepsilon X}$ and the bound $\mathbb E[e^{\varepsilon X}] \le
\exp(\mu(e^{\varepsilon}-1))$ for a sum of independent trials with mean $\mu \le
(1-\varepsilon)(t+1)$ (the Chernoff layer of the shared \texttt{dynamics} package, before
optimizing the tilt: \texttt{Distribution.prob\_ge\_le\_exp}), then $(1-\varepsilon)e^{\varepsilon} \le 1 - \varepsilon^2/2$.
\end{proof}

\begin{theorem}[Deferred decisions]\label{thm:deferred}
\lean{Epidemics.reachSet,Epidemics.closeOff,Epidemics.frontier,Epidemics.prob_reachSet_le_binTail,Epidemics.prob_cluster_gt_le_binTail,Epidemics.prob_cluster_gt_le_binomial}\leanok\uses{lem:chernoff}
The component of $s$ in $G_p$ has more than $t$ vertices with probability at most that of at
least $t$ successes in $t(d-1)+1$ independent Bernoulli$(p)$ trials.
\end{theorem}
\begin{proof}\leanok
For an exploration state (discovered set $D$, examined pairs $X$, forced closed), induction on a
budget $m \ge \mathrm{frontier}(D, X) + (k-1)(d-1)$: conditioning on the coin of one frontier edge,
a closed coin removes that edge from the frontier, and an open one adds a vertex and at most $d-1$
new frontier edges at it while removing the examined one. This is the recursion of the binomial
tail. The initial frontier of $\{s\}$ has at most $d$ edges.
\end{proof}

\begin{theorem}[Theorems E.1 and 2.3]\label{thm:subcritical}
\lean{Epidemics.prob_cluster_gt_le,Epidemics.prob_components_le_of_tail,Epidemics.card_mul_exp_le,Epidemics.prob_components_small}\leanok\uses{thm:deferred}
The component of any vertex has more than $t$ vertices with probability at most
$\exp(\varepsilon - \varepsilon^2 t/2)$, and with probability at least $1 - 1/n$ every component of
$G_p$ has at most $(10/\varepsilon^2)\log n$ vertices.
\end{theorem}
\begin{proof}\leanok
Theorem~\ref{thm:deferred} and Lemma~\ref{lem:chernoff}, then a union bound over the vertices.
\end{proof}

\begin{corollary}[Subcritical Reed--Frost and Erd\H{o}s--R\'enyi]\label{cor:subcritical}
\lean{Epidemics.card_recovered_le,Epidemics.infected_eq_empty_of_ncard_le,Epidemics.reedFrost_subcritical,Epidemics.erdosRenyi_subcritical}\leanok\uses{thm:subcritical,cor:final}
With probability at least $1 - 1/n$, the Reed--Frost epidemic with $R_0 = p(d-1) \le 1 -
\varepsilon$ infects at most $|I_0| (10/\varepsilon^2)\log n$ nodes and nobody is infected in round
$\lfloor (10/\varepsilon^2)\log n \rfloor$; and every component of $G(n, c/n)$ with
$c \le 1 - \varepsilon$ has at most $(10/\varepsilon^2)\log n$ vertices.
\end{corollary}
\begin{proof}\leanok
The final outbreak is the union of the components of $I_0$, and distances inside a component are
smaller than its size (Corollary~\ref{cor:final}). $G(n, c/n)$ is bond percolation on $K_n$, with
$d = n - 1$ and $p(d-1) \le c$.
\end{proof}

\section{Small-world networks below the threshold (EPI-6)}
This section continues Section~\ref{sec:bcdptz-small-world}.
After Becchetti, Clementi, Denni, Pasquale, Trevisan, Ziccardi (arXiv:2103.16398), Definitions
1.1 and 1.2, claim 2 of Theorems 2.1, 2.2, 2.4, 2.5, and Lemma C.1. $\mathrm{SWG}(n, q)$ is the
cycle $C_n$ together with the edges of $G(n, q)$; $3\text{-}\mathrm{SWG}(n)$ is the cycle together
with a uniformly random perfect matching. Probabilities are over the graph and the percolation.

\begin{definition}[Small-world graphs and the threshold]\label{def:swg}
\lean{Epidemics.swg,Epidemics.PerfectMatching,Epidemics.uniformMatching,Epidemics.swg3,Epidemics.swgThreshold,Epidemics.swgThreshold_spec,Epidemics.lt_swgThreshold_iff,Epidemics.swgThreshold_one,Epidemics.swg3_degree_le}\leanok\uses{thm:perc}
$\mathrm{SWG}(n, q)$ is $C_n \cup \mathrm{perc}(K_n, b)$ with $b$ i.i.d.\ Bernoulli$(q)$, and
$3\text{-}\mathrm{SWG}(n)$ is $C_n \cup M$ with $M$ uniform among the perfect matchings of $K_n$ ($n$
even). The critical value $p^* = (\sqrt{c^2+6c+1} - c - 1)/(2c)$ is the root of $c p (1 + p) =
1 - p$, so that $p < p^*$ if and only if $c p(1+p) < 1-p$; for $c = 1$, $p^* = \sqrt 2 - 1$. The
graph $3\text{-}\mathrm{SWG}(n)$ has maximum degree $3$.
\end{definition}

\begin{lemma}[Collapsing the two coin families]\label{lem:swg-collapse}
\lean{Epidemics.swgRate,Epidemics.perc_swg_eq,Epidemics.expect_expect_coord,Epidemics.expect_prob_swg_eq}\leanok\uses{def:swg}
The percolation graph of $\mathrm{SWG}(n, q)$ with parameter $p$ has the law of the percolation of
$K_n$ with independent coins, of probability $p$ on the cycle edges and $pq$ on the other pairs.
\end{lemma}
\begin{proof}\leanok
A pair is an open edge if and only if its percolation coin is open and it is a cycle edge or an
open bridge; these combined coins are independent across pairs, with the stated probabilities.
\end{proof}

\begin{lemma}[Supermartingale for adaptive observations]\label{lem:swg-mart}
\lean{Epidemics.hist,Epidemics.histProd,Epidemics.expect_prod_hist_le_one,Epidemics.expect_prod_hist_le_one_of_fresh,Epidemics.prob_le_expect_div}\leanok
If an adaptive process reads, at each step, coordinates of an independent family not read before,
and every factor $f_h \ge 0$ satisfies $\mathbb E[f_h(\mathrm{obs}_h)] \le 1$ for every history $h$,
then the product of the factors along the process has expectation at most $1$.
\end{lemma}
\begin{proof}\leanok
Condition on the history: the event that the history is $h$ depends only on the coordinates read
along $h$, which are independent of the fresh ones (principle of deferred decisions); induction on
the number of steps.
\end{proof}

\begin{theorem}[Weighted exploration]\label{thm:swg-explore}
\lean{Epidemics.Explore.state,Epidemics.Explore.Inv,Epidemics.Explore.pot,Epidemics.Explore.prob_cluster_gt_le}\leanok\uses{lem:swg-mart}
Explore the cluster of $s$ one node at a time, with independent Bernoulli$(r_e)$ coins, giving
each discovered node the weight $h_e \in [w_{\min}, w_{\max}]$ of the pair through which it was
discovered. If the expected weight discovered from every node of weight $w$ is at most $\lambda w$,
then for $\theta \ge 0$ and $\kappa \le 2$ with $\theta w_{\max} \le \kappa - 1$ and $\kappa\lambda
\le 1$, the cluster of $s$ has more than $t$ nodes with probability at most
$\exp(-\theta((1-\kappa\lambda) t w_{\min} - w_{\max}))$.
\end{theorem}
\begin{proof}\leanok
The exponential of $\theta$ times (discovered weight $-\ \kappa\lambda\,\cdot$ processed weight) is
a supermartingale, by $1 - r + r e^x \le e^{\kappa r x}$ for $0 \le x \le \kappa - 1 \le 1$ and
Lemma~\ref{lem:swg-mart}. A cluster with more than $t$ nodes leaves $t$ processed and $t + 1$
discovered nodes after $t$ steps; Markov's inequality concludes.
\end{proof}

\begin{theorem}[Lemma C.1 and Theorem 2.1, claim 2]\label{thm:swg-sub}
\lean{Epidemics.swgWeight,Epidemics.swg_step,Epidemics.prob_components_le_of_tail',Epidemics.swg_components_small,Epidemics.swg_subcritical}\leanok\uses{lem:swg-collapse,thm:swg-explore}
Let $c > 0$, $p_0 = p^* - \varepsilon > 0$ and $\delta = 1 - p_0 - c p_0 (1 + p_0)$. For $n \ge 2$, $q =
c/n$ and $p \le p_0$, with probability at least $1 - 2/n$ every component of the percolated
$\mathrm{SWG}(n, q)$ has at most $(64/\delta^2) \log n$ nodes.
\end{theorem}
\begin{proof}\leanok
Weights $1$ (discovered through a cycle edge: at most one undiscovered cycle neighbour) and $1 +
p_0$ (otherwise: at most two), and at most $n$ further pairs each open with probability $pc/n$:
the expected discovered weight is at most $(1 - \delta/2)$ times the weight. Apply
Theorem~\ref{thm:swg-explore} with $\theta = \delta/8$, $\kappa = 1 + \delta/4$ and a union bound
over the nodes.
\end{proof}

\begin{corollary}[Theorem 2.2, claim 2]\label{cor:swg3-sub}
\lean{Epidemics.swg3_subcritical}\leanok\uses{def:swg,thm:subcritical}
For $p < 1/2 - \varepsilon$, with probability at least $1 - 1/n$ every component of the percolated
$3\text{-}\mathrm{SWG}(n)$ has at most $(10/(2\varepsilon)^2) \log n$ nodes.
\end{corollary}
\begin{proof}\leanok
Theorem~\ref{thm:subcritical} with $d = 3$ and $2\varepsilon$ in place of $\varepsilon$.
\end{proof}

\begin{corollary}[Theorems 2.4 and 2.5, claim 2]\label{cor:swg-rf}
\lean{Epidemics.reedFrost_of_components_le,Epidemics.swg_reedFrost_subcritical,Epidemics.swg3_reedFrost_subcritical}\leanok\uses{thm:swg-sub,cor:swg3-sub,cor:subcritical}
Below the threshold ($p < p^* - \varepsilon$ on $\mathrm{SWG}(n, c/n)$, $p < 1/2 - \varepsilon$ on
$3\text{-}\mathrm{SWG}(n)$), with probability at least $1 - C/n$ the Reed--Frost epidemic from $I_0$
is over within $\beta \log n$ rounds and at most $\beta |I_0| \log n$ nodes are recovered at the end.
\end{corollary}
\begin{proof}\leanok
Deterministic consequences of the component bounds (Corollary~\ref{cor:subcritical}).
\end{proof}

\section{The supercritical giant component}\label{sec:krivelevich-sudakov}
Source: M.~Krivelevich and B.~Sudakov, \emph{The phase transition in random graphs: a simple
proof}, Random Structures \& Algorithms 43 (2013), arXiv:1201.6529. Throughout, $G(n,p)$ is the
percolated complete graph $(K_V)_\omega$ with i.i.d.\ Bernoulli$(p)$ coins, $n=|V|$, and
$p=(1+\varepsilon)/n$ is written $pn=1+\varepsilon$.

\begin{lemma}[Principle of deferred decisions]\label{lem:deferred}
\lean{Epidemics.queryAnswers,Epidemics.FreshUpTo,Epidemics.prob_forall_eval_eq,Epidemics.queryAnswers_eq_iff,Epidemics.prob_queryAnswers_eq,Epidemics.expect_queryAnswers,Epidemics.prob_queryAnswers}\leanok
If an adaptive strategy queries the coins one at a time, the next coordinate being a function of
the answers so far, and never queries a coordinate twice among its first $m$ queries, then its
first $m$ answers are i.i.d.\ Bernoulli$(p)$.
\end{lemma}
\begin{proof}\leanok
The answers equal a list $L$ iff the coordinates queried along $L$, which are determined by $L$
and distinct, take the values of $L$: a cylinder event of probability $\prod_k w(L_k)$.
\end{proof}

\begin{definition}[Depth-first search fed with answers]\label{def:dfs}
\lean{Epidemics.DFS.State,Epidemics.DFS.State.pending,Epidemics.DFS.State.move,Epidemics.DFS.State.answer,Epidemics.DFS.State.settle,Epidemics.DFS.ofAnswers,Epidemics.DFS.nextQuery}\leanok
The search of Krivelevich--Sudakov (Section~2) with the sets $S$ (done), $U$ (stack) and $T$
(unvisited), stopped before each query; after $U=T=\emptyset$ the remaining pairs are queried.
\end{definition}

\begin{lemma}[Properties of the search]\label{lem:dfs}
\lean{Epidemics.DFS.State.Inv,Epidemics.DFS.State.move_settle,Epidemics.DFS.fresh_nextQuery,Epidemics.DFS.card_queried_ofAnswers,Epidemics.DFS.mem_found_iff_of_queryAnswers,Epidemics.DFS.queried_of_mem_done_of_mem_unvisited,Epidemics.DFS.stack_chain_ofAnswers,Epidemics.DFS.length_stack_le,Epidemics.DFS.count_true_le_card,Epidemics.DFS.card_union_append_le,Epidemics.DFS.reachable_of_mem_comp,Epidemics.DFS.count_drop_le_card_comp,Epidemics.DFS.exists_of_epochs_lt}\leanok\uses{def:dfs,lem:deferred}
Every pair is queried at most once; all pairs between $S$ and $T$ have been queried negatively;
$U$ spans a path; $|U|\le 1+\sum X_i$ and, while $T\ne\emptyset$, $|S\cup U|\ge\sum X_i$; the
vertices of one epoch are connected; when an epoch starts, the explored set $D$ has all its pairs
with $V\setminus D$ queried.
\end{lemma}
\begin{proof}\leanok
An invariant preserved by every move and answer; moves without query terminate because
$2|S|+|U|\le 2n$ increases.
\end{proof}

\begin{lemma}[Deterministic analysis]\label{lem:analysis}
\lean{Epidemics.DFS.three_mul_explored_lt,Epidemics.DFS.le_length_stack,Epidemics.DFS.exists_epoch_start,Epidemics.DFS.le_card_comp,Epidemics.DFS.exists_path_of_stack,Epidemics.DFS.card_comp_le_ncard}\leanok\uses{lem:dfs}
The deterministic steps of the proofs of Theorems~1 and~2 of the paper: $|S\cup U|<n/3$ at time
$N_0$, $|U|\ge \ell n+1$, and the current epoch started before $t_1$.
\end{lemma}
\begin{proof}\leanok
In each case $|S||T|$ (or $|D|(n-|D|)$) would exceed the number of queries.
\end{proof}

\begin{lemma}[Krivelevich--Sudakov, Lemma~1(2), with Chernoff]\label{lem:ks1}
\lean{Epidemics.prob_count_take_far,Epidemics.Good,Epidemics.prob_not_good_le}\leanok
Among $N_0$ i.i.d.\ Bernoulli$(p)$ trials, the number of successes deviates from $N_0p$ by
$\delta N_0p$ with probability at most $2e^{-\delta^2N_0p/3}$ (Chernoff bounds of FND-3 instead of
Chebyshev, as in the paper's Discussion, item~1); the same holds at every time of a window.
\end{lemma}
\begin{proof}\leanok
Chernoff bounds and the union bound.
\end{proof}

\begin{theorem}[Long path; Theorem~1(2)]\label{thm:path}
\lean{Epidemics.core_path,Epidemics.exists_long_path}\leanok\uses{lem:analysis,lem:ks1}
For every small enough $\varepsilon>0$ there is $C$ such that $G(n,(1+\varepsilon)/n)$ contains a
path with at least $\varepsilon^2n/5$ edges with probability at least $1-C/n$.
\end{theorem}
\begin{proof}\leanok
Run the search for $N_0=\lfloor\varepsilon n^2/2\rfloor$ queries; on the typical event the stack is
a path of the required length.
\end{proof}

\begin{theorem}[Giant component; Theorem~2 and every $\varepsilon$]\label{thm:giant}
\lean{Epidemics.core_component,Epidemics.exists_giant_component,Epidemics.exists_linear_component}\leanok\uses{lem:analysis,lem:ks1}
For every small enough $\varepsilon>0$ there is $C$ such that $G(n,(1+\varepsilon)/n)$ has a
component with at least $\varepsilon n/2$ vertices with probability at least $1-C/n$; for every
$\varepsilon>0$ there are $c>0$ and $C$ such that it has a component with at least $cn$ vertices
with probability at least $1-C/n$.
\end{theorem}
\begin{proof}\leanok
No epoch starts between the times $t_1=\lfloor\eta n^2\rfloor$ and $N_0=\lfloor\theta n^2\rfloor$,
so the positive answers in between lie in one component; $\theta=\varepsilon/2$ for the first
statement, $\theta$ small in terms of $\varepsilon$ for the second.
\end{proof}

\begin{corollary}[Supercritical Reed--Frost epidemics]\label{cor:outbreak}
\lean{Epidemics.card_final_recovered,Epidemics.coins_prob_perm,Epidemics.prob_component_ge,Epidemics.reedFrost_large_outbreak_explicit,Epidemics.reedFrost_large_outbreak}\leanok\uses{thm:giant,thm:perc}
On $K_n$ with $pn=R_0>1$, the Reed--Frost epidemic started from any single node infects at least
$cn$ nodes with probability at least $q>0$ for $n\ge n_0$; for $R_0=1+\varepsilon$ with
$\varepsilon$ small, at least $\varepsilon n/2$ nodes with probability at least
$\varepsilon/2-C/n$.
\end{corollary}
\begin{proof}\leanok
The final outbreak is the component of the initial node (Corollary~\ref{cor:final}); by symmetry
of $K_n$ this node lies in a component of size $\ge k$ with probability at least $k/n$ times the
probability that such a component exists.
\end{proof}

\section{The Kermack--McKendrick SIR model}\label{sec:kermack-mckendrick}
The SIR model of Kermack and McKendrick (1927), following Hethcote, \emph{The mathematics of
infectious diseases} (SIAM Review 2000), system (2.2) and Theorem~2.1.
\begin{definition}[Solutions]\label{def:km}
\lean{Epidemics.KermackMcKendrick.R₀,Epidemics.KermackMcKendrick.sirField,Epidemics.KermackMcKendrick.IsSolution}\leanok
Let $\beta,\gamma>0$ and $R_0=\beta/\gamma$. A solution is a triple $(s,i,r)$ of real functions
such that $t\mapsto(s,i,r)(t)$ is an integral curve on $[0,\infty)$ (one-sided derivative at $0$)
of the field $(s,i,r)\mapsto(-\beta si,\ \beta si-\gamma i,\ \gamma i)$, with $s(0),i(0)>0$,
$r(0)=0$ and $s(0)+i(0)+r(0)=1$. Solutions are assumed, not constructed.
\end{definition}

\begin{theorem}[Invariants]\label{thm:km-basic}
\lean{Epidemics.KermackMcKendrick.IsSolution.sum_eq_one,Epidemics.KermackMcKendrick.IsSolution.s_pos,Epidemics.KermackMcKendrick.IsSolution.i_pos,Epidemics.KermackMcKendrick.IsSolution.r_nonneg,Epidemics.KermackMcKendrick.IsSolution.strictAntiOn_s,Epidemics.KermackMcKendrick.IsSolution.strictMonoOn_r,Epidemics.KermackMcKendrick.IsSolution.s_eq_mul_exp}\leanok\uses{def:km}
On $[0,\infty)$: $s+i+r=1$; $s,i>0$ and $r\ge0$; $s$ is strictly decreasing and $r$ strictly
increasing; and $s(t)=s(0)e^{-R_0r(t)}$.
\end{theorem}
\begin{proof}\leanok
The derivatives of $s+i+r$ and of $se^{R_0r}$ vanish ($R_0\gamma=\beta$), and a function with
zero right derivative is constant. For $i>0$, compare $i$ with the barrier
$b(t)=\tfrac{i(0)}2e^{-\gamma t}$: where they touch, $b'=-\gamma i<\beta si-\gamma i=i'$ since
$s>0$, so $i\ge b$. Monotonicity follows from the signs of $s'=-\beta si$ and $r'=\gamma i$.
\end{proof}

\begin{theorem}[Threshold]\label{thm:km-threshold}
\lean{Epidemics.KermackMcKendrick.IsSolution.strictAntiOn_i,Epidemics.KermackMcKendrick.IsSolution.initially_increasing_iff}\leanok\uses{thm:km-basic}
If $R_0s(0)\le1$, then $i$ is strictly decreasing on $[0,\infty)$. The function $i$ is strictly
increasing on some $[0,\varepsilon]$, $\varepsilon>0$, iff $R_0s(0)>1$.
\end{theorem}
\begin{proof}\leanok
$i'=i(\beta s-\gamma)$ and $s(t)<s(0)$ for $t>0$. If $R_0s(0)>1$, then $\beta s-\gamma>0$ near $0$
by continuity.
\end{proof}

\begin{theorem}[Limits and final size]\label{thm:km-final}
\lean{Epidemics.KermackMcKendrick.IsSolution.tendsto_i,Epidemics.KermackMcKendrick.IsSolution.exists_tendsto_s,Epidemics.KermackMcKendrick.IsSolution.tendsto_r,Epidemics.KermackMcKendrick.IsSolution.final_size,Epidemics.KermackMcKendrick.IsSolution.limit_pos,Epidemics.KermackMcKendrick.IsSolution.R₀_mul_limit_lt_one,Epidemics.KermackMcKendrick.IsSolution.final_size_unique,Epidemics.KermackMcKendrick.IsSolution.final_size_unique_of_le_one}\leanok\uses{thm:km-threshold}
$i(t)\to0$, $s(t)\to s_\infty$ and $r(t)\to1-s_\infty$; moreover
$s_\infty=s(0)e^{-R_0(1-s_\infty)}$, $0<s_\infty<1/R_0$, and $s_\infty$ is the unique root of this
equation in $(0,1/R_0]$ and in $(0,1]$.
\end{theorem}
\begin{proof}\leanok
$s$ and $r$ are monotone and bounded, so $i=1-s-r$ converges; if $i\to L>0$, the mean value
theorem makes $r$ grow by more than $1$. Pass to the limit in $s=s(0)e^{-R_0r}$. If
$R_0s_\infty\ge1$ then $i'>0$ forever, contradicting $i\to0$. Uniqueness: $\log x-R_0x$ is strictly
increasing on $(0,1/R_0]$ and strictly decreasing on $[1/R_0,\infty)$, and at $x=1$ it exceeds its
value $\log s(0)-R_0$ at the roots because $s(0)<1$.
\end{proof}

\begin{theorem}[Peak]\label{thm:km-peak}
\lean{Epidemics.KermackMcKendrick.IsSolution.exists_peak}\leanok\uses{thm:km-final}
If $R_0s(0)>1$, there is $t_{\max}>0$ with $s(t_{\max})=1/R_0$ such that $i$ is strictly
increasing on $[0,t_{\max}]$ and strictly decreasing on $[t_{\max},\infty)$, and
$i(t_{\max})=i(0)+s(0)-1/R_0-\log(R_0s(0))/R_0$.
\end{theorem}
\begin{proof}\leanok
Intermediate value theorem between $s(0)>1/R_0$ and $s_\infty<1/R_0$; the sign of $\beta s-\gamma$
changes at $t_{\max}$. The value follows from $\log s=\log s(0)-R_0r$ and $r=1-s-i$.
\end{proof}

\section{Kurtz's law of large numbers for SIR, in discrete time}\label{sec:kurtz}
After Kurtz (J. Appl. Probab.~7, 1970) and Wormald (1999, Theorem~5.1), for the SIR model of
Section~\ref{sec:kermack-mckendrick}.
\begin{definition}[The uniformized SIR chain]\label{def:kurtz}
\lean{Epidemics.Kurtz.Compartment,Epidemics.Kurtz.Config,Epidemics.Kurtz.count,Epidemics.Kurtz.scaled,Epidemics.Kurtz.Round,Epidemics.Kurtz.step,Epidemics.Kurtz.chain,Epidemics.Kurtz.deviationProb}\leanok\uses{def:km}
Let $\beta,\gamma\in\mathbb N$. A configuration assigns to each of $N$ agents a compartment
$S$, $I$ or $R$; $X=(S,I,R)/N$ are the scaled counts. One step draws $(u,v)$ uniformly in
$[N]^2$ and one of $\beta+\gamma$ equally likely clocks: on one of the $\beta$ infection clocks an
infected $u$ infects a susceptible $v$, on one of the $\gamma$ recovery clocks an infected $u$
recovers. This is the continuous-time chain ($S+I\to2I$ at rate $\beta SI/N$, $I\to R$ at rate
$\gamma I$) uniformized at rate $(\beta+\gamma)N$. The deviation probability is the probability
that $\|X_k-x(k/((\beta+\gamma)N))\|_\infty>\theta$ for some $k\le n$.
\end{definition}

\begin{lemma}[Drift and bounded increments]\label{lem:kurtz-drift}
\lean{Epidemics.Kurtz.drift_susceptible,Epidemics.Kurtz.drift_infected,Epidemics.Kurtz.drift_recovered,Epidemics.Kurtz.dist_scaled_step_le,Epidemics.Kurtz.iterate_chain}\leanok\uses{def:kurtz}
$\mathbb E[X_{k+1}-X_k\mid X_k=x]=F(x)/((\beta+\gamma)N)$ with $F$ the Kermack--McKendrick field,
and $\|X_{k+1}-X_k\|_\infty\le1/N$.
\end{lemma}
\begin{proof}\leanok
Among the $N^2(\beta+\gamma)$ rounds, $\beta IS$ move an agent from $S$ to $I$ and $\gamma NI$ move
one from $I$ to $R$; a round changes the compartment of at most one agent.
\end{proof}

\begin{lemma}[Maximal Azuma--Hoeffding]\label{lem:azuma}
\lean{Epidemics.Kurtz.incrementSum,Epidemics.Kurtz.expList_azuma}\leanok
If increments $D(X_j,\rho_j)$ along i.i.d.\ uniform rounds have zero conditional mean and
$|D|\le c$, their partial sums $M_k$ satisfy
$\mathbb P(\exists k\le n,\ M_k\ge\lambda)\le e^{-\lambda^2/(2nc^2)}$.
\end{lemma}
\begin{proof}\leanok
Hoeffding's lemma ($\mathbb E e^{\theta D}\le\cosh(\theta c)\le e^{\theta^2c^2/2}$ by convexity)
makes $e^{\theta M_k-k\theta^2c^2/2}$ a nonnegative supermartingale; Ville's maximal inequality,
proved by induction on $n$ peeling off the first round, bounds the probability that it reaches
$\mu$ by $1/\mu$; take $\theta=\lambda/(nc^2)$.
\end{proof}

\begin{theorem}[Law of large numbers]\label{thm:kurtz}
\lean{Epidemics.Kurtz.law_of_large_numbers,Epidemics.Kurtz.tendsto_deviationProb}\leanok\uses{lem:kurtz-drift,lem:azuma,thm:km-basic}
Let $\beta,\gamma,T>0$. There are $C,c>0$ and $L$ such that for every $N$, every initial
configuration, every solution $x$ of the Kermack--McKendrick system started in the simplex and
every $\varepsilon>0$, with probability at least $1-Ce^{-c\varepsilon^2N}$,
$\|X_k-x(k/((\beta+\gamma)N))\|_\infty\le L\|X_0-x(0)\|_\infty+\varepsilon$ for all
$k\le T(\beta+\gamma)N$. Hence, if $X_0\to x(0)$, the scaled chain converges to $x$ in probability,
uniformly on $[0,T]$.
\end{theorem}
\begin{proof}\leanok
The simplex is invariant under the ODE (first integral for $s$, integrating factor for $i$); there
$F$ is $(2\beta+\gamma)$-Lipschitz and bounded by $\beta+\gamma$, so one Euler step errs by
$O(h^2)$, $h=1/((\beta+\gamma)N)$. By Lemma~\ref{lem:kurtz-drift} each coordinate of
$X_k-X_0-h\sum_{j<k}F(X_j)$ is a martingale with increments at most $2/N$; on the event where the
six one-sided maximal deviations stay below $\delta=\varepsilon/(2L)$, the discrete Gr\"onwall
inequality gives the bound with $L=e^{(2\beta+\gamma)T}$, once $\varepsilon N\ge2LT(2\beta+\gamma)$;
Lemma~\ref{lem:azuma} bounds the complement by $6e^{-c\varepsilon^2N}$. Small $N$ and
$\varepsilon\ge2$ are absorbed in $C$.
\end{proof}

\section{Randomized rumor spreading revisited}\label{sec:doerr-kostrygin}
We follow Doerr and Kostrygin, \emph{Randomized rumor spreading revisited} (ICALP 2017; long
version arXiv:2303.11150), with the numbering of the long version.

Lemma~20 of the paper needs a major correction; corrected versions are formalized
(Lemma~\ref{lem:rumor-jump},
\href{https://github.com/formal-dynamics/leanamics/blob/main/epidemics/FORMALIZATION_DIFFERENCES.md#lemma-20-of-the-paper-needs-a-major-correction}{details}).
The proof of Theorem~43 and the parameters of the push, pull and push--pull instances need a
minor correction, and the corrected versions are formalized
(\href{https://github.com/formal-dynamics/leanamics/blob/main/epidemics/FORMALIZATION_DIFFERENCES.md#theorem-43-and-the-push-pull-and-pushpull-protocols}{details}).

\begin{definition}[Rumor-spreading processes]\label{def:rumor-process}
\lean{Epidemics.Revisited.RumorProcess,Epidemics.Revisited.RumorProcess.informProb,Epidemics.Revisited.RumorProcess.cov,Epidemics.Revisited.RumorProcess.Homogeneous,Epidemics.Revisited.RumorProcess.notYet}\leanok
A rumor-spreading process on $n$ nodes is a Markov kernel on the set $S$ of informed nodes under
which informed nodes stay informed. For a round started from $S$, $\mathrm{informProb}(S,x)$ is
the probability that $x$ is informed after the round and $\mathrm{cov}(S,x,y)$ the covariance of
the events that $x$ and $y$ are. The tail of the spreading time is
$\mathrm{notYet}(m,t,S)=\Pr[\text{fewer than }m\text{ nodes informed after }t\text{ rounds from
}S]=\Pr[T(|S|,m)>t]$; expected times are bounded through all partial sums of
$\sum_t\Pr[T>t]$. Homogeneity (Definition~6) is defined but not assumed.
\end{definition}

\begin{definition}[Upper growth and shrinking conditions]\label{def:rumor-conditions}
\lean{Epidemics.Revisited.RumorProcess.UpperGrowth,Epidemics.Revisited.RumorProcess.UpperShrinking}\leanok\uses{def:rumor-process}
Upper exponential growth (Definition~9): for $1\le|S|=k<fn$, every uninformed node is informed
with probability at least $\gamma\frac kn(1-a\frac kn-\frac b{\ln n})$ and covariances are at
most $ck/n^2$. Upper exponential shrinking (Definition~11): when $u=n-|S|\le gn$, every
uninformed node stays uninformed with probability at most $e^{-\rho}+au/n$ and covariances are
at most $c/u$.
\end{definition}

\begin{lemma}[Lemma 9]\label{lem:rumor-variance}
\lean{Epidemics.Revisited.variance_card_le}\leanok\uses{def:rumor-process}
If distinct uninformed nodes have covariance at most $c$, the number of informed nodes after one
round has variance at most its expected increase plus $c(n-|S|)^2$.
\end{lemma}
\begin{proof}\leanok
Expand the variance of a sum of indicators.
\end{proof}

\begin{lemma}[Lemma 19]\label{lem:rumor-connect}
\lean{Epidemics.Revisited.connect_tail,Epidemics.Revisited.connect_expect}\leanok\uses{def:rumor-process}
If every uninformed node is informed with probability at least $p$ in every round started with
$\ell\le|S|<m$ informed nodes, then $\Pr[T(\ell,m)>r]\le\frac{n-\ell}{n-m}(1-p)^r$ and
$E[T(\ell,m)]\le\frac{n-\ell}{n-m}\cdot\frac1p$.
\end{lemma}
\begin{proof}\leanok
The potential $n-|S|$ (set to $0$ once $|S|\ge m$) contracts by $1-p$ in expectation.
\end{proof}

\begin{theorem}[Theorem 21: exponential growth]\label{thm:rumor-growth}
\lean{Epidemics.Revisited.growth_upper_tail,Epidemics.Revisited.growth_upper_expect}\leanok\uses{def:rumor-conditions,lem:rumor-variance,lem:rumor-connect}
Under the upper exponential growth conditions, with $\gamma$ between two positive constants and
$af<1$, from any nonempty $S$: $\Pr[T(|S|,fn)>\lceil\log_{1+\gamma}n\rceil+r]\le Ae^{-\alpha r}$ and
$E[T(|S|,fn)]\le\log_{1+\gamma}n+B$, with constants depending only on the parameters.
\end{theorem}
\begin{proof}\leanok
Phases $k_{j+1}=k_j+E_0(k_j)$ with a round target $E_0(k)=E(k)-Ak^{3/4}$; Cantelli's inequality
(Lemma~\ref{lem:rumor-variance}) bounds the probability of missing a target, and a phase potential
on the iterated kernel replaces the paper's domination by geometric variables. Lemma~\ref{lem:rumor-connect}
crosses the last stretch up to $fn$.
\end{proof}

\begin{lemma}[Lemma 20, corrected]\label{lem:rumor-jump}
\lean{Epidemics.Revisited.JumpsOver,Epidemics.Revisited.RumorProcess.jumpProb,Epidemics.Revisited.overshoot_round,Epidemics.Revisited.jumpProb_le}\leanok\uses{lem:rumor-variance}
If every round started with $1\le|S|<fn$ has $p_k\le p$ and covariances at most $c/n$, then
(one round) for every $f'\in\,]f+p(1-f),1[$, a round from $|S|<fn$ ends with at least $f'n$
informed nodes with probability at most $C/n$; and (path form) there is $f'\in\,]f,1[$ such that
the probability of jumping over $[fn,f'n[$ within $t$ rounds is at most $\frac Cn$ times the
expected number of these rounds started below $fn$, i.e.\ $O(E[T(|S|,fn)]/n)$.
\end{lemma}
\begin{proof}\leanok
Chebyshev's inequality with Lemma~\ref{lem:rumor-variance} for one round, then a union bound over
the rounds along the trajectory.
\end{proof}
The paper's Lemma~20 claims probability $1-O(1/n)$ of landing in $[fn,f'n]$; this needs a
major correction, as the process that from every $|S|<fn$ informs all nodes with probability $c/n$ (and otherwise none)
satisfies the hypotheses and jumps from $1$ to $n$ with probability one
(\href{https://github.com/formal-dynamics/leanamics/blob/main/epidemics/FORMALIZATION_DIFFERENCES.md#lemma-20-of-the-paper-needs-a-major-correction}{details}).

\begin{theorem}[Theorem 31: exponential shrinking]\label{thm:rumor-shrinking}
\lean{Epidemics.Revisited.shrinking_upper_tail,Epidemics.Revisited.shrinking_upper_expect}\leanok\uses{def:rumor-conditions,lem:rumor-variance,lem:rumor-connect}
Under the upper exponential shrinking conditions, with $\rho$ between two positive constants and
$e^{-\rho_{lo}}+ag<1$, from any $S$ with $n-|S|\le gn$:
$\Pr[T(|S|,n)>\lceil\ln n/\rho\rceil+r]\le Ae^{-\alpha r}$ and $E[T(|S|,n)]\le\ln n/\rho+B$.
\end{theorem}
\begin{proof}\leanok
Lemma~\ref{lem:rumor-connect} first brings the number $u$ of uninformed nodes below $g_0n$; then
the potential $\Phi=u+\frac\beta nu^2$ satisfies $E[\Phi']\le e^{-\rho}(1+K/n)\Phi$
(Lemma~\ref{lem:rumor-variance} for the second moment), which replaces the paper's phase calculus.
\end{proof}

\begin{theorem}[Total spreading time]\label{thm:rumor-total}
\lean{Epidemics.Revisited.spreading_upper_tail,Epidemics.Revisited.spreading_upper_expect}\leanok\uses{thm:rumor-growth,thm:rumor-shrinking,lem:rumor-connect}
Under the growth conditions on $[1,fn[$, the shrinking conditions below $gn$ uninformed nodes and
a probability at least $p>0$ of being informed in between, from any nonempty $S$:
$\Pr[T>\lceil\log_{1+\gamma}n\rceil+\lceil\ln n/\rho\rceil+r]\le Ae^{-\alpha r}$ and
$E[T]\le\log_{1+\gamma}n+\ln n/\rho+B$.
\end{theorem}
\begin{proof}\leanok
Compose the three exponential tails (Theorem~\ref{thm:rumor-growth}, Lemma~\ref{lem:rumor-connect},
Theorem~\ref{thm:rumor-shrinking}) by the Markov property at fixed times, and sum the tail.
\end{proof}

\begin{definition}[Upper double exponential shrinking conditions]\label{def:rumor-double}
\lean{Epidemics.Revisited.RumorProcess.UpperDoubleShrinking,Epidemics.Revisited.RumorProcess.FastFinishing}\leanok\uses{def:rumor-process}
Definition~13: when $n^{1-\alpha}\le u=n-|S|\le gn$, every uninformed node stays uninformed with
probability at most $a(u/n)^{\ell-1}$ and covariances are at most $cn/u^2$. Fast finishing (the
second hypothesis of Theorem~43): when $u\le n^{1-\alpha}$, every uninformed node stays
uninformed with probability at most $n^{-\tau}$.
\end{definition}

\begin{theorem}[Theorem 43: double exponential shrinking]\label{thm:rumor-double}
\lean{Epidemics.Revisited.double_shrinking_upper_tail,Epidemics.Revisited.double_shrinking_upper_expect}\leanok\uses{def:rumor-double,lem:rumor-variance}
Let $\ell>1$, $a,c\ge0$, $g,\alpha\in[0,1]$, $ag^{\ell-1}<1$ and $\tau>0$. Under the upper double
exponential shrinking conditions and fast finishing, from any $S$ with $n-|S|\le gn$:
$\Pr[T(|S|,n)>\lceil\log_\ell\ln n\rceil+r]\le Cn^{A'-\alpha'r}$ and
$E[T(|S|,n)]\le\log_\ell\ln n+B$.
\end{theorem}
\begin{proof}\leanok
Three stages. Geometric targets $g\mu^jn$, $\mu=(1+ag^{\ell-1})/2$, bring $u$ below $g_0n$ in a
constant number of phases; double exponential targets $\varepsilon_{j+1}=2a_1\varepsilon_j^\ell$,
i.e.\ $\varepsilon_j=e^{-(\kappa+\ell^jD_0)}$, bring it below $\varepsilon_Jn\le n^{1-\beta/(2\ell)}$
in $J\le\log_\ell\ln n$ phases; below that, every uninformed node stays uninformed with
probability at most $n^{-\tau_3}$ and Markov's inequality finishes. A phase fails with
probability at most $n^{-2\delta}$ (Chebyshev with Lemma~\ref{lem:rumor-variance} above
$n^{1-\alpha}$ uninformed nodes, Markov with fast finishing below), and the phase potential
$\lambda^{J-j}$, $\lambda=n^\delta$, contracts by $n^{-2\delta}+n^{-\delta}$ per round. The
paper's first stage uses Lemma~19, whose tail rate does not depend on $n$; the per-round
Chebyshev bound gives the polynomial tail.
\end{proof}

\begin{theorem}[Total spreading time, double exponential shrinking]\label{thm:rumor-total-double}
\lean{Epidemics.Revisited.spreading_upper_tail_double,Epidemics.Revisited.spreading_upper_expect_double}\leanok\uses{thm:rumor-growth,thm:rumor-double,lem:rumor-connect}
Under the growth conditions on $[1,fn[$, the double exponential shrinking conditions and fast
finishing below $gn$ uninformed nodes ($g>0$), and a probability at least $p>0$ of being informed
in between, from any nonempty $S$:
$\Pr[T>\lceil\log_{1+\gamma}n\rceil+\lceil\log_\ell\ln n\rceil+r]\le Ae^{-\kappa r}$ and
$E[T]\le\log_{1+\gamma}n+\log_\ell\ln n+B$.
\end{theorem}
\begin{proof}\leanok
As Theorem~\ref{thm:rumor-total}, with Theorem~\ref{thm:rumor-double} in the last stage (its
polynomial tail is an exponential tail once $\ln n\ge1$).
\end{proof}

\begin{definition}[Push, pull and push--pull on $K_n$]\label{def:rumor-protocols}
\lean{Epidemics.Revisited.calls,Epidemics.Revisited.pushRound,Epidemics.Revisited.pullRound,Epidemics.Revisited.pushPullRound,Epidemics.Revisited.push,Epidemics.Revisited.pull,Epidemics.Revisited.pushPull}\leanok\uses{def:rumor-process}
In a round every node calls a uniformly random node (itself included), independently. Push:
informed nodes inform the nodes they call. Pull: uninformed nodes calling an informed node become
informed. Push--pull: both.
\end{definition}

\begin{lemma}[One-round quantities and conditions]\label{lem:rumor-protocol-conditions}
\lean{Epidemics.Revisited.push_informProb,Epidemics.Revisited.pull_informProb,Epidemics.Revisited.pushPull_informProb,Epidemics.Revisited.push_homogeneous,Epidemics.Revisited.pull_homogeneous,Epidemics.Revisited.pushPull_homogeneous,Epidemics.Revisited.push_cov_nonpos,Epidemics.Revisited.pull_cov_eq_zero,Epidemics.Revisited.pushPull_cov_nonpos,Epidemics.Revisited.push_upperGrowth,Epidemics.Revisited.push_upperShrinking,Epidemics.Revisited.pull_upperGrowth,Epidemics.Revisited.pull_upperDoubleShrinking,Epidemics.Revisited.pull_fastFinishing,Epidemics.Revisited.pushPull_upperGrowth,Epidemics.Revisited.pushPull_upperDoubleShrinking,Epidemics.Revisited.pushPull_fastFinishing}\leanok\uses{def:rumor-protocols,def:rumor-conditions,def:rumor-double}
With $k=|S|$ informed nodes, an uninformed node is informed with probability
$1-(1-\frac1n)^k$ (push), $\frac kn$ (pull), $1-(1-\frac1n)^k(1-\frac kn)$ (push--pull);
distinct uninformed nodes are nonpositively correlated (independent for pull). Hence push
satisfies the upper growth conditions with $\gamma=1$, $a=\frac12$ and the upper shrinking
conditions with $\rho=1$, $a=\frac2e$, $g=\frac12$; pull and push--pull satisfy the growth
conditions with $\gamma=1$, $a=0$ and $\gamma=2$, $a=\frac34$, and the double exponential
shrinking conditions with $\ell=2$, $a=1$, $g=\alpha=\frac12$ and fast finishing with
$\tau=\frac12$ (all with $b=c=0$, $f=\frac12$).
\end{lemma}
\begin{proof}\leanok
The calls are independent uniform coordinates, so the events ``$x$ stays uninformed'' and
``$x$ and $y$ stay uninformed'' have product probabilities. The bounds use
$(1-x)^k\le1-kx+(kx)^2/2$ and $e^t\le1+2t$ on $[0,1]$.
\end{proof}

\begin{theorem}[Theorems 51, 52, 53: push, pull, push--pull, upper bounds]\label{thm:rumor-protocols}
\lean{Epidemics.Revisited.push_spreading_tail,Epidemics.Revisited.push_spreading_expect,Epidemics.Revisited.pull_spreading_tail,Epidemics.Revisited.pull_spreading_expect,Epidemics.Revisited.pushPull_spreading_tail,Epidemics.Revisited.pushPull_spreading_expect}\leanok\uses{lem:rumor-protocol-conditions,thm:rumor-total,thm:rumor-total-double}
From any nonempty set of informed nodes on $K_n$, all nodes are informed within
$\lceil\log_2n\rceil+\lceil\ln n\rceil+r$ (push), $\lceil\log_2n\rceil+\lceil\log_2\ln n\rceil+r$
(pull) and $\lceil\log_3n\rceil+\lceil\log_2\ln n\rceil+r$ (push--pull) rounds, except with
probability at most $Ae^{-\kappa r}$; the expected times are at most $\log_2n+\ln n+B$,
$\log_2n+\log_2\ln n+B$ and $\log_3n+\log_2\ln n+B$.
\end{theorem}
\begin{proof}\leanok
Instantiate Theorem~\ref{thm:rumor-total} (push) and Theorem~\ref{thm:rumor-total-double}
(pull, push--pull) with Lemma~\ref{lem:rumor-protocol-conditions}; the middle range is empty since
$f=g=\frac12$.
\end{proof}

\section{Rumor spreading revisited: lower bounds for push, pull and push--pull}\label{sec:doerr-kostrygin-lower}
The lower halves of Theorems~51 to~53 of Doerr and Kostrygin (same paper and numbering as
Section~\ref{sec:doerr-kostrygin}): started from one informed node, push, pull and push--pull
need $\log_2n+\ln n$, $\log_2n+\log_2\ln n$ and $\log_3n+\log_2\ln n$ rounds minus $O(1)$, in
expectation and with exponential tails. With Theorem~\ref{thm:rumor-protocols} this determines the
spreading times up to additive constants. The general lower bounds (Theorems~27, 38 and~48) are
not formalized: the paper joins them through Lemma~20, which needs a major correction, and the
formal proofs take a direct route for the instances. The growth phase is bounded by the first
moment, the final phase by a round-by-round lower envelope, and the two are joined at a fixed time.

The paper's joining of the phases also needs a minor correction (Theorems~38 and~48 start from an
exact number of uninformed nodes, which the process started from one node need not visit), and so
do Corollary~17 (not for $u=0$), the statement of Lemma~42 ($q(u_{j+1})$ for $q(u_j)$), some typos
in Lemmas~39 and~40, and a constant in the proof of Theorem~27; none affects the statements
(\href{https://github.com/formal-dynamics/leanamics/blob/main/epidemics/FORMALIZATION_DIFFERENCES.md#lower-bounds-of-theorems-51-to-53-lowerlean}{details}).

\begin{lemma}[Lower envelopes and joining at a fixed time]\label{lem:rumor-lower-generic}
\lean{Epidemics.Revisited.envelope_le,Epidemics.Revisited.envelope_seq_le,Epidemics.Revisited.reach_le_of_expect_card_le,Epidemics.Revisited.reach_add_le}\leanok\uses{def:rumor-process}
Let $V$ be an observable of a Markov kernel and $(g_i)$, $(\delta_i)$ sequences with
$\delta_i\ge0$. If from every state with $V\ge g_i$ the observable falls below $g_{i+1}$ in one
round with probability at most $\delta_i$, then from a state with $V\ge g_0$ it is below $g_t$
after $t$ rounds with probability at most $\sum_{i<t}\delta_i$ (the round-by-round form of
Lemmas~42 and~50). For a rumor-spreading process: if one round multiplies the expected number of
informed nodes by at most $\mu$, then $\Pr[T(|S|,m)\le t]\le\mu^t|S|/m$; and if every state with
fewer than $m'$ informed nodes reaches $m$ within $t$ rounds with probability at most $B$, then
$\Pr[T(|S|,m)\le s+t]\le\Pr[T(|S|,m')\le s]+B$.
\end{lemma}
\begin{proof}\leanok
Induction on $t$ for the envelope; Markov's inequality on $E|S_t|\le\mu^t|S|$; the Markov
property at time $s$.
\end{proof}

\begin{lemma}[Expectations and growth]\label{lem:rumor-lower-growth}
\lean{Epidemics.Revisited.push_expect_card_le,Epidemics.Revisited.push_expect_uninformed,Epidemics.Revisited.push_growth_lower,Epidemics.Revisited.pull_expect_card_le,Epidemics.Revisited.pull_expect_uninformed,Epidemics.Revisited.pull_growth_lower,Epidemics.Revisited.pushPull_expect_card_le,Epidemics.Revisited.pushPull_expect_uninformed,Epidemics.Revisited.pushPull_growth_lower}\leanok\uses{lem:rumor-protocol-conditions,lem:rumor-lower-generic}
After one round from $S$, $E|S'|\le2|S|$ (push, pull) and $E|S'|\le3|S|$ (push--pull), and
$E[n-|S'|]=(n-|S|)(1-\frac1n)^{|S|}$, $(n-|S|)(1-\frac{|S|}n)$ and
$(n-|S|)(1-\frac1n)^{|S|}(1-\frac{|S|}n)$. Hence $\Pr[T(|S|,m)\le t]\le2^t|S|/m$ (push, pull) and
$3^t|S|/m$ (push--pull).
\end{lemma}
\begin{proof}\leanok
Sum the exact probabilities of Lemma~\ref{lem:rumor-protocol-conditions} over the uninformed
nodes, with $1-(1-\frac1n)^k\le\frac kn$; then Lemma~\ref{lem:rumor-lower-generic}.
\end{proof}

\begin{lemma}[One round of the final phase; Lemmas 39 and 49]\label{lem:rumor-lower-round}
\lean{Epidemics.Revisited.push_round_lower,Epidemics.Revisited.pull_round_lower,Epidemics.Revisited.pushPull_round_lower}\leanok\uses{lem:rumor-protocol-conditions,lem:rumor-variance,lem:rumor-lower-growth}
Let $u=n-|S|$ and $u'$ the number of uninformed nodes after one round. Push, at any level
$4\le v\le u$: $\Pr[u'<v/e-v^{3/4}]\le v^{-1/2}$. Pull: $\Pr[u'<u^2/(2n)]\le4n^2/u^3$.
Push--pull: $\Pr[u'<u^2/(2en)]\le4e^2n^2/u^3$.
\end{lemma}
\begin{proof}\leanok
Chebyshev's inequality on the lower tail of $u'$, with $\mathrm{Var}[u']\le u$
(Lemma~\ref{lem:rumor-variance} with nonpositive covariances) and $E[u']\ge u/e$ (push, from
$(1-\frac1n)^k\ge e^{-1}$ for $k<n$), $E[u']=u^2/n$ (pull), $E[u']\ge u^2/(en)$ (push--pull).
\end{proof}

\begin{theorem}[Final phase of push; Theorem 38 for push]\label{thm:rumor-lower-final-push}
\lean{Epidemics.Revisited.push_final_lower}\leanok\uses{lem:rumor-lower-round,lem:rumor-lower-generic}
For every $n$, every $S$ with $u=n-|S|$ uninformed nodes and every $t$:
$\Pr[T(|S|,n)\le t]\le1600\,e^{(t-\ln u)/2}$.
\end{theorem}
\begin{proof}\leanok
Lower envelope (Lemma~\ref{lem:rumor-lower-generic}) along the explicit targets
$g_i=q_i^3(q_i-20)$, $q_i=(ue^{-i})^{1/4}$, which replace the paper's recursive $u_{j+1}=E_0(u_j)$
(Lemma~40). The step $g(qe^{-1/4})\le g(q)/e-g(q)^{3/4}$ reduces to $c^4-20c+20\le0$ for
$c=e^{1/4}$; while $q_i\ge40$, round $i$ fails with probability at most $2/q_i^2$
(Lemma~\ref{lem:rumor-lower-round}) and these failure probabilities sum to at most
$4e^{(t-\ln u)/2}$; when $q_t<40$ the bound exceeds $1$.
\end{proof}

\begin{theorem}[Final phase of pull and push--pull; Theorem 48 with $\ell=2$]\label{thm:rumor-lower-final-double}
\lean{Epidemics.Revisited.pull_final_lower,Epidemics.Revisited.pushPull_final_lower}\leanok\uses{lem:rumor-lower-round,lem:rumor-lower-generic}
For $n\ge3$, every $S$ with $2|S|\le n$ and every $t\le\log_2\ln n-5$:
$\Pr[T(|S|,n)\le t]\le128\,n^{-1/2}$ (pull) and $128e^2n^{-1/2}$ (push--pull).
\end{theorem}
\begin{proof}\leanok
Lower envelope along the double exponential targets $g_i=2n(1/4)^{2^i}$ (pull) and
$2en(1/(4e))^{2^i}$ (push--pull), the paper's $\varepsilon_{j+1}=E(\varepsilon_j)/2$, with
Lemma~\ref{lem:rumor-lower-round}. The paper's exponent $-1+2\alpha\ell$ is replaced by the weaker
$-1/2$, which comes from the cruder variance bound and suffices below.
\end{proof}

\begin{theorem}[Theorems 51, 52, 53: push, pull, push--pull, lower bounds]\label{thm:rumor-protocols-lower}
\lean{Epidemics.Revisited.push_spreading_lower_tail,Epidemics.Revisited.push_spreading_lower_expect,Epidemics.Revisited.pull_spreading_lower_tail,Epidemics.Revisited.pull_spreading_lower_expect,Epidemics.Revisited.pushPull_spreading_lower_tail,Epidemics.Revisited.pushPull_spreading_lower_expect}\leanok\uses{lem:rumor-lower-growth,thm:rumor-lower-final-push,thm:rumor-lower-final-double,lem:rumor-lower-generic}
There are $A$, $B$, $\kappa>0$ and $N$ such that, for $n\ge N$ and one informed node on $K_n$, all
nodes are informed within $\lfloor\log_2n\rfloor+\lfloor\ln n\rfloor-r$ (push),
$\lfloor\log_2n\rfloor+\lfloor\log_2\ln n\rfloor-r$ (pull) and
$\lfloor\log_3n\rfloor+\lfloor\log_2\ln n\rfloor-r$ (push--pull) rounds with probability at most
$Ae^{-\kappa r}$; the expected times are at least $\log_2n+\ln n-B$, $\log_2n+\log_2\ln n-B$ and
$\log_3n+\log_2\ln n-B$ (stated for every partial sum $\sum_{t<R}\Pr[T>t]$ with $R$ at least the
bound).
\end{theorem}
\begin{proof}\leanok
Join the two phases at a fixed time (Lemma~\ref{lem:rumor-lower-generic}) with $m'=n/2$. Push:
after $\lfloor\log_2n\rfloor-\lceil r/2\rceil$ rounds at most $n/2$ nodes are informed except with
probability $O(2^{-r/2})$ (Lemma~\ref{lem:rumor-lower-growth}), and from there
Theorem~\ref{thm:rumor-lower-final-push} bounds the rest. Pull and push--pull: the same with
$\lfloor\log_2\ln n\rfloor-5$ rounds for the final phase (Theorem~\ref{thm:rumor-lower-final-double}),
whose $n^{-1/2}$ becomes $e^{-(\ln2/4)r}$ since a nontrivial time forces $r\le2\log_2n$. The tails
have rate $\kappa=\ln2/4$; the expectations follow by summing the tail series.
\end{proof}

\section{COBRA and BIPS: the model}\label{sec:crr-cobra-bips}
We follow Cooper, Radzik and Rivera, \emph{The coalescing-branching random walk on expanders and
the dual epidemic process}, PODC 2016 (arXiv:1602.05768). The duality itself (their Theorem~4)
is in the next section.
\begin{definition}[Rounds of neighbour choices]\label{def:choices}
\lean{Epidemics.Choices,Epidemics.instFintypeChoices,Epidemics.choices_nonempty}\leanok
In a round $r$, every vertex $x$ of a finite graph $G$ samples $k$ neighbours
$r(x,0),\dots,r(x,k-1)$. A uniform round is exactly the paper's sampling: $k$ uniform neighbours
with replacement, independently for every vertex; $t$ i.i.d.\ uniform rounds are averaged with
\texttt{Dynamics.expList}. If $G$ is connected and has at least two vertices, rounds exist.
\end{definition}
\begin{proof}\leanok
Every vertex of a connected graph with at least two vertices has a neighbour.
\end{proof}
\begin{definition}[COBRA and BIPS]\label{def:cobra}
\lean{Epidemics.cobraStep,Epidemics.cobraRun,Epidemics.bipsStep,Epidemics.bipsRun}\leanok\uses{def:choices}
COBRA from $C_0=C$: $C_{s+1}=\{r_{s+1}(x,i) : x\in C_s,\ i<k\}$. BIPS with persistent source $v$
from $A_0$: $A_{s+1}=\{v\}\cup\{u : \exists i,\ r_{s+1}(u,i)\in A_s\}$.
\end{definition}

\section{COBRA--BIPS duality}
\begin{lemma}[Time reversal of i.i.d.\ rounds]\label{lem:reverse}
\lean{Dynamics.expList_comp_reverse,Epidemics.expList_congr_length}\leanok
For every functional $F$ of $T$ i.i.d.\ uniform rounds,
$\mathbb E[F(r_T,\dots,r_1)]=\mathbb E[F(r_1,\dots,r_T)]$.
\end{lemma}
\begin{proof}\leanok
The uniform average over $\mathrm{Fin}\,T\to\alpha$ is invariant under $\omega\mapsto\omega\circ\mathrm{rev}$
(the time-reversal lemma of the shared \texttt{dynamics} package).
\end{proof}
\begin{lemma}[Pathwise duality]\label{lem:pathwise}
\lean{Epidemics.cobra_hit_iff_bips_reverse}\leanok\uses{def:cobra}
For all rounds $r_1,\dots,r_t$: COBRA from $C$ along $r_1,\dots,r_t$ visits $v$ at some time
$s\le t$ (time $0$ included) iff BIPS from $\{v\}$ along $r_t,\dots,r_1$ infects a vertex of $C$.
\end{lemma}
\begin{proof}\leanok
Induction on $t$, generalizing $C$. Both sides say that a chain $x_0\in C$,
$x_j=r_j(x_{j-1},i_j)$, reaches $v$ within $t$ steps: COBRA unfolds the first round, BIPS the
last round of the reversed sequence.
\end{proof}
\begin{theorem}[Theorem 4 of Cooper--Radzik--Rivera]\label{thm:duality}
\lean{Epidemics.cobra_bips_duality,Epidemics.cobra_bips_duality_singleton}\leanok\uses{lem:reverse,lem:pathwise}
For every finite graph, every $k$, every vertex $v$, set $C$ and time $t$:
$\hat{\mathbb P}(\mathrm{Hit}_C(v)>t\mid C_0=C)=\mathbb P(C\cap A_t=\emptyset\mid A_0=\{v\})$. In
particular (equation (2)), $\hat{\mathbb P}(\mathrm{Hit}_u(v)>t)=\mathbb P(u\notin A_t\mid A_0=\{v\})$.
The paper assumes $G$ connected and regular and $k\ge1$; none of this is needed.
\end{theorem}
\begin{proof}\leanok
Lemma~\ref{lem:pathwise} identifies the two events after reversing the rounds, and
Lemma~\ref{lem:reverse} says that reversing the rounds preserves their law. The paper proves the
theorem instead by induction on $t$ with one-step conditioning.
\end{proof}

\section{COBRA cover time on expanders}\label{sec:crr-cover}
Here $G$ is $r$-regular with $r>0$ on $n$ vertices, $P=A(G)/r$, $\lambda=\max(\lambda_2,|\lambda_n|)$,
$k\ge2$, and $d_A(x)=|N(x)\cap A|$. Expectations of unbounded times $\tau$ are written as tail
sums: every partial sum $\sum_{s<H}\Pr(\tau>s)$ is bounded uniformly in $H$.

\begin{definition}[Second eigenvalue]\label{def:lambda}
\lean{Epidemics.transitionMatrix,Epidemics.walkEigenvalues,Epidemics.lambdaG}\leanok
$\lambda=\max(\lambda_2,|\lambda_n|)$ for the sorted eigenvalues of $P$ ($0$ if $n<2$).
\end{definition}

\begin{lemma}[Spectral core of Lemma~1]\label{lem:crr-spectral}
\lean{Epidemics.sum_sq_neighbor_le}\leanok\uses{def:lambda}
For every set $A$: $\sum_x (d_A(x)/r)^2\le\lambda^2|A|+(1-\lambda^2)|A|^2/n$.
\end{lemma}
\begin{proof}\leanok
$\sum_x(d_A(x)/r)^2=\|P1_A\|^2$; expand $1_A$ in an orthonormal eigenbasis of $P$: the top
eigenvector is constant with coefficient $|A|/\sqrt n$, and every other eigenvalue has square
at most $\lambda^2$. If $\lambda\ge1$, use $d_A(x)\le r$ and $\sum_x d_A(x)=r|A|$.
\end{proof}

\begin{lemma}[Lemma~1, MGF and Chernoff bound]\label{lem:crr-growth}
\lean{Epidemics.bips_expected_growth,Epidemics.bips_mgf_le,Epidemics.bips_chernoff_lower}\leanok\uses{lem:crr-spectral,def:cobra}
One BIPS round from $A$ gives $\E|A'|\ge|A|(1+(1-\lambda^2)(1-|A|/n))$,
$\E e^{-\varphi|A'|}\le\exp(-(1-e^{-\varphi})\E|A'|)$ for every real $\varphi$, and
$\Pr(|A'|\le(1-\delta)\mu)\le e^{-\delta^2\mu/2}$ for $\mu\le\E|A'|$, $0<\delta<1$.
\end{lemma}
\begin{proof}\leanok
Vertices are infected independently; $u\ne v$ with probability $1-(1-d_A(u)/r)^k\ge
2d_A(u)/r-(d_A(u)/r)^2$. Sum and apply Lemma~\ref{lem:crr-spectral}. The MGF factors over the
vertices and $1-x\le e^{-x}$; the Chernoff bound is Markov's inequality with
$\varphi=-\log(1-\delta)$.
\end{proof}

\begin{theorem}[Lemmas~2 to~4, for an abstract growth process]\label{thm:crr-phases}
\lean{Epidemics.GrowthProcess,Epidemics.round_small_phase,Epidemics.round_large_phase,Epidemics.round_end_phase,Epidemics.bips_small_phase,Epidemics.bips_large_phase,Epidemics.bips_end_phase}\leanok\uses{lem:crr-growth}
Let a round-driven set process have a persistent source, the MGF bound of
Lemma~\ref{lem:crr-growth}, growth $\E|A'|\ge|A|(1+c(1-|A|/n))$ with $0<c\le1$, be monotone, and
have $V$ absorbing. From $\{v\}$ the size exceeds $m\le n/2$ within $13m/c+24C\log n/c^2$ rounds
except with probability $n^{-C}$; from $|A_0|\ge4000\log n/c^2$ it reaches $9n/10$ within
$T\ge24\log n/c$ rounds except with probability $T n^{-5}$; from $|A_0|\ge9n/10$ (if
$4000\log n/c^2\le 9n/10$) it covers $V$ after $T\ge8\log n/c$ rounds except with probability
$n^{-5}$. BIPS satisfies the hypotheses with $c=1-\lambda$.
\end{theorem}
\begin{proof}\leanok
Lemma~2: the exponential supermartingale $e^{-\varphi|A_t|}$ on the event of staying small,
$\varphi=\log(1+c/2)$. Lemma~3: a round with $|A'|<|A|(1+c/23)$ has probability $\le n^{-5}$;
union bound over rounds with moving targets. Lemma~4: staying above $9n/10$ and the recursion
$\E B_{t+1}\le\theta\E B_t+tn^{-7}$ for the number $B_t$ of healthy vertices. The paper's
Lemmas~3 and~4 need minor corrections ($\log_2$ versus $\ln$ doublings; a missing factor $B_0$).
\end{proof}

\begin{theorem}[Theorem~2]\label{thm:crr-bips}
\lean{Epidemics.round_first_hit,Epidemics.GrowthProcess.fail_le,Epidemics.GrowthProcess.tail_sum_le_log,Epidemics.bips_growthProcess,Epidemics.bips_infection_time,Epidemics.bips_infection_time_expectation}\leanok\uses{thm:crr-phases}
If $1-\lambda\ge128\sqrt{\log n/n}$, BIPS from $\{v\}$ has $A_T\ne V$ with probability at most
$3/n^3$ for $T\ge60000\log n/(1-\lambda)^3$, and $\sum_{s<H}\Pr(A_s\ne V)\le
130000\log n/(1-\lambda)^3$.
\end{theorem}
\begin{proof}\leanok
Chain the three phases at their first hitting times with $m=\ceil{4000\log n/(1-\lambda)^2}$ and
$C=3$. (The paper's proof takes $m=4000\log n/(1-\lambda^2)$, which needs a minor correction:
Lemma~3 needs $|A|\ge4000\log n/(1-\lambda)^2$.) For the expectation, restart from the current
set: monotonicity and the persistent source give $\Pr(A_{s+T_0}\ne V)\le\Pr(A_s\ne V)/2$.
\end{proof}

\begin{theorem}[Theorem~1]\label{thm:crr-cobra}
\lean{Epidemics.cover_fail_le_log,Epidemics.cover_tail_sum_le_log,Epidemics.cobra_cover_time,Epidemics.cobra_cover_time_expectation}\leanok\uses{thm:crr-bips,thm:duality}
If $1-\lambda\ge128\sqrt{\log n/n}$, COBRA from $\{u\}$ has $\bigcup_{t=1}^T C_t\ne V$ with
probability at most $3/n^2$ for $T\ge60002\log n/(1-\lambda)^3$, and every partial tail sum of
the cover time is at most $130000\log n/(1-\lambda)^3$.
\end{theorem}
\begin{proof}\leanok
Union bound over the missed vertex $w$. Splitting off the first round, missing $w$ at the times
$1,\dots,T$ is not hitting $w$ within $T-1$ rounds from $C_1\ne\emptyset$, which by
Theorem~\ref{thm:duality} has probability $\Pr(C_1\cap A_{T-1}=\emptyset)\le\Pr(A_{T-1}\ne V)$.
For the expectation, restart COBRA from the nonempty set $C_s$ (equation~(1)).
\end{proof}

\begin{definition}[Branching factor $1+p/q$]\label{def:coin}
\lean{Epidemics.CoinChoices,Epidemics.coinTargets,Epidemics.cobraCoinStep,Epidemics.cobraCoinRun,Epidemics.bipsCoinStep,Epidemics.bipsCoinRun}\leanok\uses{def:choices}
A round gives every vertex two neighbour samples and a uniform coin in $\{0,\dots,q-1\}$; a vertex
uses its second sample iff its coin is $<p$. COBRA pushes to the used samples, BIPS contacts them.
\end{definition}

\begin{theorem}[Theorem~4, Corollary~1 and Theorem~3]\label{thm:crr-coin}
\lean{Epidemics.tgt_duality,Epidemics.cobraCoin_bips_duality,Epidemics.coin_hit_prob,Epidemics.bipsCoin_expected_growth,Epidemics.bipsCoin_growthProcess,Epidemics.bipsCoin_infection_time,Epidemics.cobra_cover_time_branching,Epidemics.cobra_cover_time_branching_expectation}\leanok\uses{def:coin,thm:crr-bips,thm:crr-cobra}
The duality holds for the coin processes; with $\rho=p/q$,
$\E|A'|\ge|A|(1+\rho(1-\lambda^2)(1-|A|/n))$; and for constants $\lambda_0<1$, $\rho_0>0$ there
are $C,N$ such that for $n\ge N$, $\lambda\le\lambda_0$ and $\rho\in[\rho_0,1]$ the COBRA cover time
with branching factor $1+\rho$ is at most $C\log n$ with probability $1-C/n^2$ and in expectation.
\end{theorem}
\begin{proof}\leanok
Both coin processes are target-set processes, whose duality is proved as Theorem~4. A vertex
$u\ne v$ is infected with probability $(1+\rho)P_u-\rho P_u^2$, $P_u=d_A(u)/r$; sum and apply
Lemma~\ref{lem:crr-spectral}. Then Theorems~\ref{thm:crr-phases} to~\ref{thm:crr-cobra} apply
with $c=\min(1,\rho_0(1-\max(\lambda_0,0)))$, and $128\sqrt{\log n/n}\le c$ for $n$ large.
\end{proof}



%%FILE median/blueprint/src/content.tex
\begin{abstract}
We formalize the median dynamics of Doerr, Goldberg, Minder, Sauerwald and Scheideler (SPAA
2011): every node adopts the median of its own value and the values of two random nodes. The
median commutes with monotone maps, so thresholding the process at any value gives the binary
median process, i.e.\ 2-Choices, driven by the same samples. For two values we prove consensus
from a gap of order $\sqrt{n\log n}$ within $O(\log n)$ rounds with high probability, then
consensus from any binary configuration and, by a union bound over the thresholds, from any
configuration (Theorem~1 of the paper, without adversary). With an odd number of equally
supported values, the middle one wins within $O(\log m+\log\log n)$ rounds (the odd case of
Theorem~21). On $d$-regular graphs with $\lambda_G\le 3/5-\varepsilon$, two-sample voting
eliminates a minority of size at most $(\varepsilon/5)n$ within $O(\log n)$ rounds with high
probability (Theorem~4 of Cooper, Elsässer and Radzik, ICALP 2014), and on every $d$-regular graph
an initial imbalance $\nu_0\ge K\lambda_G$ suffices, with a failure probability that tends to $0$
once $\nu_0^2n$ is large compared with $\log\log n$ (their Theorem~2, which needs a major correction
for very small $\lambda_G$). For 2-Choices with any number
$k$ of colours, the plurality colour wins within $O((n/c_1)\log n)$ rounds from a gap of order
$\sqrt{n\log n}$ to every other colour (Theorem~1.2 of Elsässer, Friedetzky, Kaaser,
Mallmann-Trenn and Trinker, arXiv:1602.04667, in arXiv v5 numbering, without their hypothesis
$k=O(n^\varepsilon)$), and from colours of $O(\log n)$ nodes each consensus needs
$\Omega(n/\log n)$ rounds (Theorem~5 of arXiv:1702.04921 (v1), Berenbrink, Clementi, Elsässer,
Kling, Mallmann-Trenn and Natale, PODC 2017).
\end{abstract}

\section{The model}\label{sec:dgmss-median}
This section, the four following ones and the one against an adaptive adversary formalize
Doerr, Goldberg, Minder, Sauerwald and Scheideler, \emph{Stabilizing consensus with the power of
two choices} (SPAA 2011). Claim~2.9 of the paper, a hitting-time lemma, is formalized in the
shared \texttt{dynamics} package, with two points made explicit there: the target $c_4\log q$ is
at most $q$, and $c_5$ also depends on $c_1,c_2,c_3$
(\href{https://github.com/formal-dynamics/leanamics/blob/main/dynamics/FORMALIZATION_DIFFERENCES.md#hitting-time-lemma-drifthitting-maj-8-doerr-et-al-spaa-2011-claim-29}{details}).
\begin{definition}[Median dynamics]\label{def:median}
\lean{Median.med3,Median.Config,Median.Round,Median.step,Median.run,Median.ones,Median.Consensus,Median.notConsensus,Median.kernel}\leanok
Values lie in a linear order. In a round every node samples two nodes uniformly at random
(independently, with replacement) and adopts the median of its own value and the two sampled
values. With two values this is 2-Choices: a node keeps its value unless both samples hold the
other one.
\end{definition}

\section{Structure}
\begin{theorem}[Threshold reduction]\label{thm:threshold}
\lean{Median.med3_monotone,Median.threshold_step,Median.threshold_run}\leanok\uses{def:median}
The median commutes with monotone maps; hence for every threshold $\theta$ the indicator of
$\{x_v\ge\theta\}$ evolves, round by round and for the same samples, as the binary median process.
\end{theorem}
\begin{proof}\leanok
$\mathrm{med}(a,b,c)=\max(\min(a,b),\min(\max(a,b),c))$, and monotone maps commute with $\min$ and $\max$.
\end{proof}
\begin{theorem}[Validity, range, consensus]\label{thm:basic}
\lean{Median.step_mem,Median.step_mem_Icc,Median.step_of_consensus,Median.absorbed}\leanok\uses{def:median}
Every new value is a current value, values stay within any interval containing all current
values, consensus configurations are fixed, and consensus is reached almost surely.
\end{theorem}
\begin{proof}\leanok
The median is one of its arguments; the round in which every node samples node $0$ twice gives
consensus and has positive probability; conclude with the shared absorption theorem.
\end{proof}
\begin{theorem}[Binary expectation]\label{thm:expect}
\lean{Median.expected_ones}\leanok\uses{def:median}
With two values and a fraction $p$ of \texttt{true}, the expected number of \texttt{true} after one
round is $n(3p^2-2p^3)$, the same cubic as for 3-Majority.
\end{theorem}
\begin{proof}\leanok
A \texttt{true} node stays unless both samples are \texttt{false}, a \texttt{false} node switches iff
both samples are \texttt{true}.
\end{proof}

\section{Consensus from a vanishing bias}
\begin{theorem}[2-Choices, w.h.p.]\label{thm:whp}
\lean{Median.binary_consensus,Median.consensus_whp}\leanok\uses{thm:expect}
If $\log n\ge 128$ and the \texttt{true} nodes outnumber the \texttt{false} ones by at least
$128\sqrt{n\log n}$, all nodes hold \texttt{true} after $\lceil 128\log n\rceil$ rounds with
probability at least $1-128/n$.
\end{theorem}
\begin{proof}\leanok
Phases as in the plurality package: while it is small, the gap grows by a factor $5/4$ per phase
except with probability $O(n^{-2})$ (Bernstein); then the minority shrinks by $7/8$ per phase down
to $O(\log n)$; a final round reaches consensus except with probability $O(n^{-1/2})$ (Markov). The
phases are combined with the nested-phases lemma of the shared layer. The threshold
$\log n\ge128$, i.e.\ $n\ge e^{128}$, comes from these crude constants.
\end{proof}

\section{Consensus from any configuration}
\begin{theorem}[2-Choices from any start]\label{thm:binary-any}
\lean{Median.binary_any_start}\leanok\uses{thm:whp}
There is a constant $C>0$ such that, whenever $\log n\ge C$, from every binary configuration
(possibly perfectly balanced) all nodes agree after $\lceil C\log n\rceil$ rounds except with
probability at most $C/n$.
\end{theorem}
\begin{proof}\leanok
Take $C=2^{18}$ and write $g$ for the number of \texttt{true} nodes minus the number of
\texttt{false} nodes. \emph{Escape:} the potential $e^{-|g|/(256\sqrt n)}$ contracts in
expectation by $e^{-1/65536}$ per round, up to an additive error $e^{1/131072-\sqrt n/512}$:
away from balance by Hoeffding's lemma applied at every node, near balance because the next gap
has variance at least $n/10$, so by Paley--Zygmund it leaves $[-\sqrt n/4,\sqrt n/4]$ with
probability at least $9/64$. Iterating over $\lceil 2^{17}\log n\rceil$ rounds and applying
Markov's inequality, $|g|\ge 128\sqrt{n\log n}$ except with probability $2/n$.
\emph{Consensus:} Theorem~\ref{thm:whp}, applied to the configuration or to its flip, finishes
within $\lceil 128\log n\rceil$ rounds except with probability $128/n$.
\end{proof}
\begin{theorem}[Reduction to two values]\label{thm:of-binary}
\lean{Median.consensus_of_binary}\leanok\uses{thm:threshold}
If every binary configuration fails to reach consensus within $T$ rounds with probability at
most $\varepsilon$, then a configuration with $m$ distinct values fails with probability at most
$(m-1)\varepsilon$.
\end{theorem}
\begin{proof}\leanok
If a run has not reached consensus, two nodes hold values $a<b$, and the threshold configuration
$\{x_v\ge b\}$, which follows the same rounds by Theorem~\ref{thm:threshold}, has not reached
consensus either. Take a union bound over the $m-1$ thresholds above the minimum value.
\end{proof}
\begin{theorem}[Consensus from any configuration]\label{thm:any}
\lean{Median.median_consensus_any}\leanok\uses{thm:binary-any,thm:of-binary}
There is a constant $C>0$ such that, whenever $\log n\ge C$, from every configuration with values
in a linear order (any number of distinct values) all nodes agree after $\lceil C\log n\rceil$
rounds except with probability at most $C/n$. This is Theorem~1 of Doerr et al.\ without the
adversary.
\end{theorem}
\begin{proof}\leanok
Consensus is absorbing, so three blocks of $\lceil C_0\log n\rceil$ rounds turn the binary failure
bound $C_0/n$ of Theorem~\ref{thm:binary-any} into $(C_0/n)^3$. The union bound of
Theorem~\ref{thm:of-binary} over at most $n-1$ thresholds gives $n(C_0/n)^3\le (3C_0+3)/n$, since
$C_0^3\le n$ when $\log n\ge 3C_0+3$.
\end{proof}

\section{Many values: the middle one wins fast}
\begin{theorem}[Fast binary consensus from a large gap]\label{thm:binary-fast}
\lean{Median.binary_consensus_fast}\leanok\uses{def:median}
There is a constant $C>0$ such that, whenever $\log n\ge C$, if the \texttt{true} nodes outnumber
the \texttt{false} ones by at least $\Delta\ge C\sqrt{n\log n}$, all nodes hold \texttt{true} after
$\lceil C(\log(n/\Delta)+\log\log n)\rceil$ rounds with probability at least $1-C/n$.
\end{theorem}
\begin{proof}\leanok
Take $C=128$ and chain one-round moves, each failing with probability at most $n^{-2}$, through
four segments: the gap grows by a factor $5/4$ per round until the minority is below $n/4$
($O(\log(n/\Delta))$ rounds); six rounds bring the minority below $n/8$; it then shrinks
quadratically, $m\mapsto\approx 3m^2/n$ (Bernstein), down to $512\log n$ ($O(\log\log n)$
rounds); eight more rounds give consensus. Consensus is absorbing, so the remaining rounds are
harmless, and a union bound over the rounds gives the failure probability.
\end{proof}
\begin{theorem}[Margins around a value]\label{thm:median-fast}
\lean{Median.median_consensus_fast}\leanok\uses{thm:threshold,thm:binary-fast}
There is a constant $C>0$ such that, whenever $\log n\ge C$: if, for a value $v$, the nodes holding a value
$\ge v$ outnumber those holding a value $<v$ by at least $\Delta\ge C\sqrt{n\log n}$, and
symmetrically the nodes holding a value $\le v$ outnumber those holding a value $>v$ by at least
$\Delta$, then all nodes hold $v$ after $\lceil C(\log(n/\Delta)+\log\log n)\rceil$ rounds with
probability at least $1-2C/n$.
\end{theorem}
\begin{proof}\leanok
The indicators of $\{x_u\ge v\}$ and of $\{x_u\le v\}$ both evolve as 2-Choices with the same
samples (the second because the median is self-dual), with gaps the two margins. When both reach
consensus on \texttt{true}, all nodes hold $v$; conclude with Theorem~\ref{thm:binary-fast}
(with the same $C=128$) and a union bound.
\end{proof}
\begin{theorem}[An odd number of equally supported values]\label{thm:odd-split}
\lean{Median.odd_split_consensus}\leanok\uses{thm:median-fast}
There is a constant $C>0$ such that, whenever $\log n\ge C$: if the nodes hold $2k+1$ distinct
values, each held by $n/(2k+1)$ nodes, and $C(2k+1)\sqrt{n\log n}\le n$, then the middle value is
held by all nodes after $\lceil C(\log(2k+1)+\log\log n)\rceil$ rounds with probability at least
$1-C/n$. This is the deterministic core of the odd case of Theorem~21 of Doerr et al.
\end{theorem}
\begin{proof}\leanok
Both margins around the middle value equal $n/(2k+1)\ge C_0\sqrt{n\log n}$, and
$\log(n/\Delta)=\log(2k+1)$; apply Theorem~\ref{thm:median-fast} with $C=2C_0$.
\end{proof}

\section{Two-sample voting on expanders}\label{sec:cer-two-choices}
This section formalizes Theorem~4 of Cooper, Elsässer and Radzik, \emph{The power of two choices
in distributed voting} (ICALP 2014, arXiv:1404.7479); lemma numbers follow the arXiv version.

\begin{definition}[Two-sample voting on a graph]\label{def:graph-voting}
\lean{Median.GraphRound,Median.graphStep,Median.graphRun,Median.graphStep_bool,Median.minority,Median.edgeCount}\leanok\uses{def:median}
On a simple graph $G$, in a round every vertex samples two neighbours independently and uniformly
at random (with replacement) and adopts the median of its own opinion and the two sampled ones;
with two opinions it adopts the sampled opinion when the two samples agree. For a majority
opinion $a$, the minority $B$ is the set of vertices holding another opinion. $E(S,T)$ is the
number of ordered pairs $(u,v)\in S\times T$ of adjacent vertices, so $E(S,S)=2|E(S)|$.
\end{definition}

\begin{definition}[Second eigenvalue]\label{def:lambdaG}
\lean{Median.transitionMatrix,Median.transitionMatrix_isHermitian,Median.walkEigenvalues,Median.lambdaG}\leanok
For a $d$-regular graph, $\lambda_1\ge\lambda_2\ge\dots\ge\lambda_n$ are the eigenvalues of the
transition matrix $P=A/d$ and $\lambda_G=\max\{\lambda_2,|\lambda_n|\}$.
\end{definition}

\begin{lemma}[Expander mixing lemma, the paper's Lemma~3]\label{lem:eml}
\lean{Median.expander_mixing}\leanok\uses{def:graph-voting,def:lambdaG}
On every $d$-regular $n$-vertex graph and for all vertex sets $S,T$,
$|E(S,T)-d|S||T|/n|\le\lambda_G\,d\sqrt{|S||T|}$.
\end{lemma}
\begin{proof}\leanok
Expand $1_S\cdot P1_T=\sum_k\mu_k s_kt_k$ in an orthonormal eigenbasis. If $\lambda_G<1$, every
eigenvector other than the top one has eigenvalue at most $\lambda_2<1$, hence is orthogonal to
the all-ones vector, which is fixed by $P$; so the top term is $|S||T|/n$, and Cauchy--Schwarz
bounds the others by $\lambda_G\sqrt{|S||T|}$. If $\lambda_G\ge1$, both $E(S,T)$ and $d|S||T|/n$
lie in $[0,d\sqrt{|S||T|}]$.
\end{proof}

\begin{lemma}[Small sets are sparse]\label{lem:sparse}
\lean{Median.sparse_of_lambdaG}\leanok\uses{lem:eml}
If $\lambda_G\le 3/5-\varepsilon$ and $|S|\le\varepsilon n$, then $E(S,S)\le (3/5)\,d|S|$.
\end{lemma}
\begin{proof}\leanok
By Lemma~\ref{lem:eml} with $T=S$,
$E(S,S)\le d|S|^2/n+\lambda_Gd|S|\le(\varepsilon+3/5-\varepsilon)d|S|$.
\end{proof}

\begin{lemma}[One round, the paper's Lemma~5 at $\alpha=3/10$]\label{lem:phase2}
\lean{Median.expected_minority_step,Median.phaseII_step}\leanok\uses{def:graph-voting}
If every superset $S\supseteq B$ with $|S|\le(1+1/\alpha)|B|$ spans at most $\alpha d|S|$ edges,
then after one round $\mathbb E|B'|\le(1-(1-2\alpha)(1-3\alpha))|B|=(24/25)|B|$, and
$|B'|\le(49/50)|B|$ with probability at least $1-e^{-|B|/4850}$.
\end{lemma}
\begin{proof}\leanok
Let $p_v$ be the fraction of neighbours of $v$ in $B$ and $y=\sum_{v\in B}(1-p_v)$, which equals
$\sum_{v\notin B}p_v$ by double counting. Then
$\mathbb E|B'|=\sum_{v\notin B}p_v^2+\sum_{v\in B}(1-(1-p_v)^2)$. Cauchy--Schwarz gives
$\sum_{v\in B}(1-p_v)^2\ge y^2/|B|$. With $C=\{v\notin B:p_v>3/10\}$, the chord bound
$p^2\le(3/10)p+(p-3/10)^+$ and the sparsity of $B\cup C$ (which has at most $(13/3)|B|$
vertices) give $\sum_{v\notin B}p_v^2\le(3/10)y+(y-(2/5)|B|)/2$. Altogether
$\mathbb E|B'|\le(24/25)|B|-(y-(2/5)|B|)^2/|B|$. The new minority is a sum of independent
indicators; the multiplicative Chernoff bound with mean bound $(24/25)m$ and $\delta=1/48$ gives
$\Pr(|B'|\ge(49/50)m)\le e^{-m/4850}$ for every $m\ge|B|$.
\end{proof}

\begin{theorem}[Theorem~4]\label{thm:expander}
\lean{Median.two_choices_expander_explicit,Median.two_choices_expander,Median.two_choices_failure_tendsto}\leanok\uses{lem:sparse,lem:phase2}
Let $G$ be a $d$-regular $n$-vertex graph with $d>0$ and $\lambda_G\le3/5-\varepsilon$, $\varepsilon>0$,
and let the minority have size at most $(\varepsilon/5)n$. After $T$ rounds every vertex holds
the majority opinion except with probability at most
$(24/25)^T|B|+Te^{-\varepsilon n/24250}$. For $T=\lceil C\log n\rceil$ with $C=25000$ the failure
probability is at most $1/n+(C\log n+1)e^{-\varepsilon n/C}$, which tends to $0$ for fixed
$\varepsilon$.
\end{theorem}
\begin{proof}\leanok
While $|B|\le(\varepsilon/5)n$, every superset of $B$ of size at most $(13/3)|B|\le\varepsilon n$
is sparse (Lemma~\ref{lem:sparse}), so by Lemma~\ref{lem:phase2} the expected minority contracts by
$24/25$ in every round and one round leaves this region with probability at most
$e^{-\varepsilon n/24250}$ (take $m=(\varepsilon/5)n$). By induction on $T$, a $[0,1]$-valued
observable bounded by the potential $|B|$ in the region has expectation at most
$(24/25)^T|B_0|+Te^{-\varepsilon n/24250}$; apply it to the indicator of not being in consensus.
For the $O(\log n)$ form, $(24/25)^T\le e^{-T/25}\le n^{-2}$ and $|B_0|\le n$.
\end{proof}

\section{Two-sample voting from a small imbalance}\label{sec:cer-small-imbalance}
This section formalizes Theorem~2 of Cooper, Elsässer and Radzik (ICALP 2014, arXiv:1404.7479),
with its Phase~I (Lemma~2 and Corollary~2), on top of
Section~\ref{sec:cer-two-choices}; numbering follows the arXiv version
(\href{https://github.com/formal-dynamics/leanamics/blob/main/median/FORMALIZATION_DIFFERENCES.md#two-sample-voting-on-expanders-from-a-small-imbalance-expandergeneral-roadmap-maj-5-b}{differences
from the paper}). Theorem~2 of the paper needs a major correction for very small $\lambda_G$, and
the corrected statement is formalized (Remark~\ref{rem:cer-correction}). Throughout, $a$ is the
majority opinion, $A$ and $B$ are the numbers of vertices holding $a$ and not holding $a$,
$n=A+B$, and primes denote the values after one round.

\begin{definition}[Imbalance and flows]\label{def:cer-imbalance}
\lean{Median.ExpanderGeneral.majority,Median.ExpanderGeneral.majority_add_minority,Median.ExpanderGeneral.imbalance,Median.ExpanderGeneral.phaseEta,Median.ExpanderGeneral.gainCount,Median.ExpanderGeneral.lossCount,Median.ExpanderGeneral.minority_add_gainCount,Median.ExpanderGeneral.phaseIRounds}\leanok\uses{def:graph-voting}
The imbalance is $\nu=(A-B)/n$, and $\eta=\alpha n/\sqrt{AB}$ for a parameter $\alpha$. The flows
of a round are $\Delta_{BA}$, the number of vertices of $B$ that hold $a$ after the round, and
$\Delta_{AB}$, the number of vertices of $A$ that no longer hold $a$; so
$B'+\Delta_{BA}=B+\Delta_{AB}$. The length of Phase~I is
$T_1(\nu,c)=\ceil{\log_{5/4}(1/(2\nu))}+\ceil{\log_{4/3}(1/(4c))}$.
\end{definition}

\begin{definition}[The mixing hypothesis of Lemma~2]\label{def:cer-mixing-prop}
\lean{Median.ExpanderGeneral.MixingProp}\leanok\uses{def:graph-voting}
A $d$-regular graph satisfies the mixing hypothesis with parameters $\alpha$ and $c$ if
$|E(X,Y)-d|X||Y|/n|\le\alpha d\sqrt{|X||Y|}$ for all disjoint $X,Y$ with $|Y|\ge cn$ and
$|X|\ge(2/3)\alpha c^{3/2}n$.
\end{definition}

\begin{lemma}[Lemma~3]\label{lem:cer-mixing-prop}
\lean{Median.ExpanderGeneral.mixingProp_of_lambdaG}\leanok\uses{def:cer-mixing-prop,lem:eml}
If $\lambda_G\le\alpha$, the mixing hypothesis holds with $\alpha$ and every $c$.
\end{lemma}
\begin{proof}\leanok
The expander mixing lemma (Lemma~\ref{lem:eml}) with $\lambda_G\le\alpha$.
\end{proof}

In the next three lemmas, $G$ is $d$-regular with $d>0$, $0<c\le1/2$,
$0<\alpha\le c^{3/2}/36$, the mixing hypothesis holds with $\alpha$ and $c$, and $cn\le B\le A$.

\begin{lemma}[Expected flows]\label{lem:cer-flows-expect}
\lean{Median.ExpanderGeneral.avg_gainCount,Median.ExpanderGeneral.avg_lossCount,Median.ExpanderGeneral.phaseEta_bounds,Median.ExpanderGeneral.sum_gain_ge,Median.ExpanderGeneral.sum_sq_nbCount_le,Median.ExpanderGeneral.expected_gain_ge,Median.ExpanderGeneral.expected_loss_le}\leanok\uses{def:cer-imbalance,def:cer-mixing-prop}
$\E\Delta_{BA}\ge(A^2B/n^2)(1-2\eta)$ and $\E\Delta_{AB}\le(AB^2/n^2)(1+15\eta)$.
\end{lemma}
\begin{proof}\leanok
Let $p_v$ be the fraction of neighbours of $v$ in $B$. A vertex of $B$ adopts $a$ with probability
$(1-p_v)^2$ and a vertex of $A$ leaves $a$ with probability $p_v^2$, so
$\E\Delta_{BA}=\sum_{v\in B}(1-p_v)^2$ and $\E\Delta_{AB}=\sum_{v\in A}p_v^2$. For the first,
Cauchy--Schwarz gives $\sum_{v\in B}(1-p_v)^2\ge(E(B,A)/d)^2/B$, and the mixing hypothesis gives
$E(B,A)\ge(dAB/n)(1-\eta)$. For the second, the hypothesis bounds the size of
$\{v\in A:p_v\ge(1+2^j\eta)B/n\}$ by $A/4^j$; summing $p_v^2$ as a layer cake over these sets,
up to the $m$ with $4^m\le n^2/(\eta B^2)<4^{m+1}$, gives the bound, since $(1+3m)\eta\le5$.
\end{proof}

\begin{lemma}[Tails of the flows]\label{lem:cer-flows-tail}
\lean{Median.ExpanderGeneral.gainCount_eq_sum,Median.ExpanderGeneral.lossCount_eq_sum,Median.ExpanderGeneral.gain_tail,Median.ExpanderGeneral.loss_tail}\leanok\uses{lem:cer-flows-expect}
$\Pr(\Delta_{BA}\le(A^2B/n^2)(1-3\eta))\le e^{-\alpha^2cn/6}$ and
$\Pr(\Delta_{AB}\ge(AB^2/n^2)(1+17\eta))\le e^{-\alpha^2c^2n/2}$.
\end{lemma}
\begin{proof}\leanok
Both flows are sums of independent indicators over the vertices of $B$ and of $A$. Apply the
multiplicative Chernoff bounds with the lower and upper bounds on the means of
Lemma~\ref{lem:cer-flows-expect}, and $\eta\ge2\alpha$.
\end{proof}

\begin{lemma}[One round of Phase~I]\label{lem:cer-phaseI-step}
\lean{Median.ExpanderGeneral.phaseI_step,Median.ExpanderGeneral.growth_small,Median.ExpanderGeneral.growth_large}\leanok\uses{lem:cer-flows-tail}
With probability at least $1-e^{-\alpha^2cn/6}-e^{-\alpha^2c^2n/2}$,
$\nu'\ge\nu+\nu(1-\nu^2)/2-12\alpha/\sqrt{1-\nu^2}$. Any such $\nu'$ satisfies
$\nu'\ge(5/4)\nu$ if $120\alpha\le\nu\le1/2$, and $1-\nu'\le(3/4)(1-\nu)$ if
$1/2\le\nu\le1-2c$.
\end{lemma}
\begin{proof}\leanok
$\nu'-\nu=2(\Delta_{BA}-\Delta_{AB})/n$. Outside the two events of Lemma~\ref{lem:cer-flows-tail},
$\Delta_{BA}-\Delta_{AB}\ge AB(A-B)/n^2-\eta(3A^2B+17AB^2)/n^2$; with $A=n(1+\nu)/2$ and
$B=n(1-\nu)/2$ this is the bound (the paper's (ncnwd-Appx)). The two growth statements are
elementary inequalities.
\end{proof}

\begin{theorem}[Phase~I: Lemma~2 and Corollary~2]\label{thm:cer-phaseI}
\lean{Median.ExpanderGeneral.phaseI_of_step,Median.ExpanderGeneral.phaseI,Median.ExpanderGeneral.phaseI_expander}\leanok\uses{lem:cer-phaseI-step,lem:cer-mixing-prop}
Let $G$ be $d$-regular with $d>0$, $0<c\le1/2$ and $0<\alpha\le c^{3/2}/36$, and let the mixing
hypothesis hold with $\alpha$ and $c$ (for instance $\lambda_G\le\alpha$). If $\nu_0\ge120\alpha$,
then at some time $t\le T_1=T_1(\nu_0,c)$ the minority has at most $cn$ vertices, except with
probability at most $T_1(e^{-\alpha^2cn/6}+e^{-\alpha^2c^2n/2})$.
\end{theorem}
\begin{proof}\leanok
Freeze the process at the first time the minority has at most $cn$ vertices. While it is not
frozen, $B\ge cn$ and $\nu\le1-2c$, and by Lemma~\ref{lem:cer-phaseI-step} the imbalance stays
above the deterministic envelope $(5/4)^t\nu_0$ up to $1/2$, then $1-(3/4)^{t-t_1}/2$, except
with probability $e^{-\alpha^2cn/6}+e^{-\alpha^2c^2n/2}$ per round. After $T_1$ rounds the envelope
exceeds $1-2c$, so the process is frozen; a union bound over the rounds gives the claim. With
$\lambda_G\le\alpha$, Lemma~\ref{lem:cer-mixing-prop} gives the mixing hypothesis.
\end{proof}

\begin{theorem}[Theorem~2, corrected]\label{thm:cer-theorem2}
\lean{Median.ExpanderGeneral.not_consensus_le,Median.ExpanderGeneral.two_choices_expander_general_explicit,Median.ExpanderGeneral.phaseIRounds_le,Median.ExpanderGeneral.general_of_explicit,Median.ExpanderGeneral.two_choices_expander_general}\leanok\uses{thm:cer-phaseI,thm:expander}
Let $G$ be $d$-regular with $d>0$ and $\nu_0>0$, $4000\lambda_G\le\nu_0$, and let
$\alpha=\nu_0/4000$ and $T_1=T_1(\nu_0,1/20)$. For every $T_2$, after $T_1+T_2$ rounds every
vertex holds $a$, except with probability at most
$T_1(e^{-\alpha^2n/120}+e^{-\alpha^2n/800})+(24/25)^{T_2}n+(T_1+T_2)e^{-n/97000}$. Hence there are
absolute constants $K,C>0$ ($K=4000$, $C=2\cdot10^{10}$) such that, if $K\lambda_G\le\nu_0$,
after $\ceil{C\log n}$ rounds every vertex holds $a$ except with probability at most
$1/n+(2C\log n+C)e^{-\nu_0^2n/C}$.
\end{theorem}
\begin{proof}\leanok
Phase~I (Theorem~\ref{thm:cer-phaseI} with $c=1/20$ and $\alpha=\nu_0/4000\ge\lambda_G$; note
$\alpha\le c^{3/2}/36$ and $120\alpha\le\nu_0$) brings the minority to at most $n/20$ at some time
$t\le T_1$. From then on, Theorem~4 (Theorem~\ref{thm:expander} with $\varepsilon=1/4$, which
applies since $\lambda_G\le1/4000\le3/5-1/4$) finishes; the composition at the random time $t$ is
an induction on the length of Phase~I. For the $O(\log n)$ form, $T_1\le5\log n+7$ (as
$\nu_0\ge1/n$), and the remaining rounds make $(24/25)^{T_2}n\le1/n$.
\end{proof}

\begin{remark}[Theorem~2 of the paper needs a major correction]\label{rem:cer-correction}
The paper states Theorem~2 with success probability $1-o(1)$ whenever $\nu_0\ge K\lambda_G$, on
every $d$-regular graph. Its proof applies Lemma~2 with $\alpha=\lambda_G$, whose failure
probability is $e^{-\Theta(\lambda_G^2n)}$; this is not $o(1)$ when $\lambda_G$ is of order
$1/\sqrt n$ or smaller. On the complete graph, $\lambda_G=1/(n-1)$, so a constant difference
$A-B\in\{K+2,K+3\}$ satisfies the hypothesis for $n\ge K+3$; the first round produces
fluctuations of order $\sqrt n$ in $A-B$, and by symmetry the initial majority then wins with
probability tending to $1/2$. Theorem~\ref{thm:cer-theorem2} is the corrected statement: it keeps
the hypothesis $\nu_0\ge K\lambda_G$ and bounds the failure probability in terms of $\nu_0^2n$,
which tends to $0$ once $\nu_0^2n$ is large compared with $\log\log n$
(\href{https://github.com/formal-dynamics/leanamics/blob/main/median/FORMALIZATION_DIFFERENCES.md#theorem-2-of-the-paper-needs-a-major-correction}{details}).
\end{remark}

\section{Almost stable consensus against an adaptive adversary}
\begin{definition}[Adaptive adversary]\label{def:adversary}
\lean{Median.Adversary,Median.Adversary.Bounded,Median.Adversary.UsesValues,Median.runAdv,Median.IsAdvRun,Median.AlmostConsensus,Median.notAlmostStable}\leanok\uses{def:median}
An adversary maps the list of rounds played so far (the current one last) and the output of the
median rule in the current round to a new configuration; it is $F$-bounded if it recolours at
most $F$ nodes, and it uses values of $S$ if every value it writes lies in $S$. The run against
the adversary applies, in every round, the median rule and then the adversary. More generally, a
perturbed run is any function of the rounds played so far that starts at $x$ and, after every
round, differs from the output of the median rule in at most $F$ nodes. A configuration is in
almost consensus on $b$ with $K$ exceptions if at most $K$ nodes hold a value other than $b$;
almost stable consensus during the times $T_0,\dots,T_0+H$ means almost consensus on the same
value $b$ at all these times.
\end{definition}
\begin{theorem}[2-Choices against an adaptive adversary]\label{thm:binary-adv}
\lean{Median.runAdv_isAdvRun,Median.binary_almost_stable_of_isAdvRun,Median.binary_almost_stable}\leanok\uses{def:adversary,thm:binary-any}
There is a constant $C>0$ such that, whenever $\log n\ge C$ and $CF\le\sqrt n$, for every binary
configuration, every perturbed run with budget $F$ (in particular, the run against every
$F$-bounded adversary) and every $H\ge0$, almost stable consensus with $C(F+\log n)$ exceptions
holds during the times $\lceil C\log n\rceil,\dots,\lceil C\log n\rceil+H$, except with
probability at most $(C\log n+H)/n^2$. This is the two-value case (Theorem~10 of the 2009
version) of the main theorem of Doerr et al.
\end{theorem}
\begin{proof}\leanok
Take $C=2^{20}$, so $F\le\sqrt n/1024$, and let $g$ be the gap and $K=\max(16F,1024\log n)$.
Recolouring $F$ nodes moves $g$ by at most $2F$ and the minority by at most $F$.
\emph{Escape:} the potential $e^{-|g|/(256\sqrt n)}$ of Theorem~\ref{thm:binary-any} grows by at
most $e^{1/131072}$ under the recolourings, so it still contracts by $e^{-1/131072}$ per round up
to an additive $e^{1/65536-\sqrt n/512}$; after $\lceil 2^{19}\log n\rceil$ rounds Markov's
inequality gives $|g|\ge 128\sqrt{n\log n}$ except with probability $2/n^2$. \emph{Growth and
saturation} (for $g>0$; otherwise flip): a chain of one-round moves, each failing with
probability at most $n^{-2}$ for every recolouring, along sets indexed by time: the gap grows from
$G$ to $\min(9G/8,17n/32)$ (Hoeffding: $5G/4$ before the recolouring), then the minority shrinks
from $t$ to $\max(15t/16,K)$ (Bernstein: $\max(7t/8,512\log n)$ before the recolouring) and stays
below $K$. A union bound over the rounds of the chain and of the window gives the bound.
\end{proof}
\begin{theorem}[Almost stable consensus against an adaptive adversary]\label{thm:median-adv}
\lean{Median.median_almost_stable}\leanok\uses{thm:binary-adv,thm:threshold}
There is a constant $C>0$ such that, whenever $\log n\ge C$ and $CF\le\sqrt n$: let $S$ be a set
of $m$ legal values containing the values of a configuration $x$. Against every $F$-bounded
adversary that only writes values of $S$, and for every $H\ge0$, almost stable consensus with
$C(F+\log n)$ exceptions holds during the times $\lceil C\log n\rceil,\dots,\lceil C\log n\rceil+H$,
except with probability at most $(m-1)(C\log n+H)/n^2$.
\end{theorem}
\begin{proof}\leanok
Take $C=2^{20}$. All values of the run lie in $S$. For every legal value $b$, the threshold
$u\mapsto[b\le x_u]$ of the run is a perturbed run of 2-Choices with the same budget, because the
median commutes with monotone maps and thresholding does not increase the Hamming distance. If,
for every legal $b$ above the minimum $m_0$, the threshold run stays in almost consensus on some
$c_b$ with $K$ exceptions, let $w$ be the largest legal value with $w=m_0$ or $c_w=\texttt{true}$:
the nodes below $w$ are exceptions of the threshold at $w$, and the nodes above $w$ are exceptions
of the threshold at the next legal value, so the run is in almost consensus on $w$ with $2K$
exceptions. Conclude with the internal bound of Theorem~\ref{thm:binary-adv} and a union bound
over the $m-1$ thresholds.
\end{proof}

\section{Plurality consensus with $k$ colours}\label{sec:efkmt-plurality}
This section formalizes the synchronous upper bound of Elsässer, Friedetzky, Kaaser,
Mallmann-Trenn and Trinker, arXiv:1602.04667 (arXiv v5 numbering, v5 being titled \emph{Rapid
asynchronous plurality consensus}): Theorem~1.2 and Section~2.1, without the adversary and without the
hypothesis $k=O(n^\varepsilon)$
(\href{https://github.com/formal-dynamics/leanamics/blob/main/median/FORMALIZATION_DIFFERENCES.md#plurality-consensus-with-k-colours-twochoices-roadmap-maj-4}{differences
from the paper}). Throughout, $c_i$ is the number of nodes of colour $i$, primes denote the values
after one round, and $S=\sum_j c_j^2$.

\begin{definition}[$k$-party 2-Choices]\label{def:two-choices}
\lean{Median.TwoChoices.rule,Median.TwoChoices.step,Median.TwoChoices.run,Median.TwoChoices.count,Median.TwoChoices.twice,Median.TwoChoices.kernel}\leanok\uses{def:median}
Colours lie in any type with decidable equality. In a round every node samples two nodes uniformly
at random (independently, with replacement) and adopts their colour if the two samples hold the
same colour; otherwise it keeps its own. $c_i$ is the number of nodes of colour $i$, and the number
of nodes whose two samples both hold $i$ is denoted $\mathrm{twice}_i$.
\end{definition}

\begin{lemma}[Basic facts and two colours]\label{lem:tc-basic}
\lean{Median.TwoChoices.step_of_ne,Median.TwoChoices.step_of_eq,Median.TwoChoices.step_of_consensus,Median.TwoChoices.count_step_le,Median.TwoChoices.count_step_eq_zero,Median.TwoChoices.rule_bool,Median.TwoChoices.step_bool,Median.TwoChoices.run_bool,Median.TwoChoices.count_true,Median.TwoChoices.run_dominates}\leanok\uses{def:two-choices}
Consensus configurations are fixed; $c_i'\le c_i+\mathrm{twice}_i$, so an extinct colour stays
extinct. On two colours the rule is the median of the own colour and the two samples. For the
same rounds, every node that holds \texttt{true} in the binary median process started from the
indicator of colour $i$ holds colour $i$.
\end{lemma}
\begin{proof}\leanok
A node that does not hold $i$ adopts $i$ only if both samples hold $i$. With two colours, the new
colour is the majority among the own colour and the two samples, that is, their median. For the
domination, the indicator of $\{\text{new colour}=i\}$ is at least the median of the three
indicators, and the median is monotone; induct on the rounds.
\end{proof}

\begin{lemma}[One-round expectations]\label{lem:tc-expect}
\lean{Median.TwoChoices.avg_pair,Median.TwoChoices.avg_step_node,Median.TwoChoices.expected_count,Median.TwoChoices.expected_count_mul,Median.TwoChoices.sum_count,Median.TwoChoices.expected_count_mono,Median.TwoChoices.expected_gap,Median.TwoChoices.sum_sq_le_aggregate,Median.TwoChoices.sum_sq_le_max,Median.TwoChoices.expected_gap_ge}\leanok\uses{def:two-choices}
$\E c_i'=c_i(1+c_i/n-S/n^2)$; hence $c_r\le c_s$ implies $\E c_r'\le\E c_s'$
(Observation~2.1), and $\E(c_i'-c_j')=(c_i-c_j)(1+(c_i+c_j)/n-S/n^2)$. If every colour other
than $i$ has at most $b$ nodes, then $S\le c_i^2+(n-c_i)b$. If $a$ and $b$ are the sizes of the
largest and the second largest colour, then $\E(a'-b')\ge(a-b)(1+(a/n)(1-a/n))$.
\end{lemma}
\begin{proof}\leanok
A node ends in colour $i$ with probability $(c_i/n)^2+[x_v=i](1-S/n^2)$; sum over the nodes. The
map $c\mapsto c(1+c/n-S/n^2)$ is increasing on $c\ge0$ since $S\le n^2$. The aggregation bound is
$\sum_{j\ne i}c_j^2\le b\sum_{j\ne i}c_j=b(n-c_i)$, and with $S\le a^2+nb$ the gap formula gives
the last bound.
\end{proof}

\begin{lemma}[One-round concentration]\label{lem:tc-conc}
\lean{Median.TwoChoices.gap_dev_le,Median.TwoChoices.count_dev_le,Median.TwoChoices.gap_tail_twelve,Median.TwoChoices.count_tail_twelve}\leanok\uses{lem:tc-expect}
Let colour $i$ be a largest colour, with $a=c_i>0$ nodes, and $j\ne i$. For every $\lambda\ge0$,
$\Pr(a'-c_j'\le\E(a'-c_j')-\lambda)\le\exp(-\lambda^2/(20a^2/n+4\lambda/3))$ and
$\Pr(a'\le\E a'-\lambda)\le\exp(-\lambda^2/(4a^2/n+2\lambda/3))$. In particular, if
$128\sqrt{n\log n}\le g\le a$, both probabilities are at most $n^{-12}$ for $\lambda=ga/(8n)$ and
$\lambda=ag/(4n)$ respectively.
\end{lemma}
\begin{proof}\leanok
Both quantities are sums of independent per-node contributions bounded by $1$ in absolute value. A
node of colour $i$ or $j$ contributes only when it leaves its colour, which has probability at most
$S/n^2\le a/n$, and any other node joins $i$ or $j$ with probability at most $2a^2/n^2$. Bounding
each variance by the second moment about the node's current contribution gives
$\mathrm{Var}(a'-c_j')\le10a^2/n$ and $\mathrm{Var}(a')\le2a^2/n$; apply Bernstein's inequality.
The exponents are then at least $(g^2/(1344n))\ge12\log n$.
\end{proof}

\begin{lemma}[The distance increases, Lemma~2.2]\label{lem:tc-distance}
\lean{Median.TwoChoices.distance_increases}\leanok\uses{lem:tc-expect,lem:tc-conc}
There are constants $z,C>0$ such that, whenever $\log n\ge C$: if colour $i$ is the largest, with
$a=c_i\le n/2$ nodes, $j$ is the second largest, with $b=c_j$ nodes, and $a-b\ge z\sqrt{n\log n}$,
then $a'-b'>(a-b)(1+a/(4n))$ with probability at least $1-C/n^2$.
\end{lemma}
\begin{proof}\leanok
Take $z=128$ and $C=2$. By Lemma~\ref{lem:tc-expect} and $a\le n/2$,
$\E(a'-b')\ge(a-b)(1+a/(2n))$; by Lemma~\ref{lem:tc-conc} with $\lambda=(a-b)a/(8n)$, the gap falls
below $\E(a'-b')-\lambda$ with probability at most $n^{-12}$. No bound on the number of colours is
needed.
\end{proof}

\begin{lemma}[The coupling, Lemma~2.3, as a stochastic domination]\label{lem:tc-stochdom}
\lean{Median.TwoChoices.count_stochDom}\leanok\uses{def:two-choices}
If $c_c\le c_b$, then $\Pr(c_c'\ge t)\le\Pr(c_b'\ge t)$ for every $t$.
\end{lemma}
\begin{proof}\leanok
Choose a permutation $\tau$ of the nodes sending every node of colour $c$ to a node of colour $b$.
For every node $v$, the sample pairs that make $v$ end in colour $c$ are at most as many as those
that make $\tau(v)$ end in colour $b$: a node of another colour adopts a colour of size $m$ on the
$m^2$ pairs with both samples on it, a node of that colour keeps it on the $n^2-S+m^2$ pairs that
do not agree on another colour, and both counts increase with $m$. Choose a permutation of the
sample pairs embedding the first set into the second.
The induced bijection of rounds preserves the uniform distribution and maps the rounds with
$c_c'\ge t$ into the rounds with $c_b'\ge t$. The paper's coupling, read literally, needs a minor
correction
(\href{https://github.com/formal-dynamics/leanamics/blob/main/median/FORMALIZATION_DIFFERENCES.md#remarks-on-the-sources}{details}).
\end{proof}

\begin{lemma}[One round of the growth phase]\label{lem:tc-growth-round}
\lean{Median.TwoChoices.growth_round}\leanok\uses{lem:tc-expect,lem:tc-conc}
There are constants $z,C>0$ such that, whenever $\log n\ge C$: if colour $i$ has $a=c_i\le3n/4$
nodes and leads every other colour by at least $g\ge z\sqrt{n\log n}$, then with probability at
least $1-C/n^2$ after one round $a'\ge a$ and colour $i$ leads every other colour by at least
$g(1+a/(8n))$.
\end{lemma}
\begin{proof}\leanok
Take $z=128$ and $C=2$, and let $b$ be the second largest colour. For every other colour $j$,
Observation~2.1 gives $\E c_j'\le\E b'$, so by Lemma~\ref{lem:tc-expect}
$\E(a'-c_j')\ge(a-b)(1+(a/n)(1-a/n))\ge g(1+a/(4n))$ since $a\le3n/4$. Lemma~\ref{lem:tc-conc}
with $\lambda=ga/(8n)$ bounds the failure for $j$ by $n^{-12}$. The aggregation bound gives
$\E a'-a\ge ag/(4n)$, so $a'<a$ has probability at most $n^{-12}$. A union bound over the at most
$n$ colours present gives the failure probability $2/n^2$. This replaces the coupling of
Lemma~2.3, and the range $a\le3n/4$ (instead of $n/2$) closes the gap before the two-colour
phase.
\end{proof}

\begin{lemma}[The growth phase]\label{lem:tc-growth-phase}
\lean{Median.TwoChoices.growth_phase}\leanok\uses{lem:tc-growth-round}
There are constants $z,C>0$ such that, whenever $\log n\ge C$: if colour $i$ leads every other
colour by at least $z\sqrt{n\log n}$, then after $\lceil C(n/c_i)\log n\rceil$ rounds it holds at
least $3n/4$ nodes with probability at least $1-C/n$.
\end{lemma}
\begin{proof}\leanok
Take $z=C=128$, $a_0=c_i$, $g_0$ the initial lead and $q=1+a_0/(8n)$. Consider the moving targets
``colour $i$ holds $3n/4$ nodes, or it holds at least $a_0$ nodes and leads every other colour by
$g_0q^t$''. Each round leaves them with probability at most $2/n^2$, by
Lemma~\ref{lem:tc-growth-round} while $a\le3n/4$, and because a colour holding $3n/4$ nodes keeps
them except with probability $n^{-2}$ (Bernstein, drift from the aggregation bound). Since
$q^{\lceil128(n/a_0)\log n\rceil}>n$, a lead of $g_0q^t$ is impossible at the end, so the colour
holds $3n/4$ nodes.
\end{proof}

\begin{lemma}[Finishing phase]\label{lem:tc-finish}
\lean{Median.TwoChoices.finish_phase,Median.TwoChoices.expList_consensus_mono}\leanok\uses{lem:tc-basic,thm:whp}
If $\log n\ge128$ and colour $i$ holds at least $3n/4$ nodes, all nodes hold $i$ after
$\lceil128\log n\rceil$ rounds with probability at least $1-128/n$.
\end{lemma}
\begin{proof}\leanok
By Lemma~\ref{lem:tc-basic}, colour $i$ dominates the binary median process started from its
indicator, whose gap is at least $n/2\ge128\sqrt{n\log n}$; apply Theorem~\ref{thm:whp}.
\end{proof}

\begin{theorem}[Plurality consensus, Theorem~1.2 without the adversary]\label{thm:tc-plurality}
\lean{Median.TwoChoices.plurality_whp,Median.TwoChoices.plurality_whp_k}\leanok\uses{lem:tc-growth-phase,lem:tc-finish}
There are constants $z,C>0$ such that, whenever $\log n\ge C$, for every number $k$ of colours: if
colour $i$ leads every other colour by at least $z\sqrt{n\log n}$, all nodes hold $i$ after
$\lceil C(n/c_i)\log n\rceil$ rounds, and hence after $\lceil Ck\log n\rceil$ rounds, with
probability at least $1-C/n$.
\end{theorem}
\begin{proof}\leanok
Take $z=128$ and $C=384$: Lemma~\ref{lem:tc-growth-phase}, then Lemma~\ref{lem:tc-finish} from the
configuration reached, with failure at most $256/n$; consensus is absorbing, so more rounds do not
hurt. For the second form, $c_i\ge n/k$. The paper's proof of Theorem~1.2 needs a minor
correction (the range $n/2<a<(1/2+\varepsilon_1)n$ and the time bound), which
Lemma~\ref{lem:tc-growth-round} supplies.
\end{proof}

\begin{theorem}[Two colours and almost-sure consensus]\label{thm:tc-two}
\lean{Median.TwoChoices.two_colours_consensus_whp,Median.TwoChoices.absorbed}\leanok\uses{lem:tc-basic,thm:whp}
With two colours, a gap of $C\sqrt{n\log n}$ gives consensus on the majority after
$\lceil C\log n\rceil$ rounds with probability at least $1-C/n$, whenever $\log n\ge C$. With
finitely many colours, consensus is reached almost surely.
\end{theorem}
\begin{proof}\leanok
On two colours the process is the binary median process (Lemma~\ref{lem:tc-basic}); apply
Theorem~\ref{thm:whp}. The round in which every node samples node $0$ twice gives consensus and
has positive probability; conclude with the shared absorption theorem.
\end{proof}

\section{A lower bound for 2-Choices}\label{sec:bcekmn-lower}
This section formalizes the 2-Choices lower bound of Berenbrink, Clementi, Elsässer, Kling,
Mallmann-Trenn and Natale, \emph{Ignore or comply? On breaking symmetry in consensus} (PODC 2017):
Theorem~5 of arXiv:1702.04921 (v1) and the 2-Choices half of its Theorem~1
(Simplified)
(\href{https://github.com/formal-dynamics/leanamics/blob/main/median/FORMALIZATION_DIFFERENCES.md#the-lower-bound-for-2-choices-twochoiceslower-roadmap-maj-6-a}{differences
from the paper}).

\begin{lemma}[Few nodes change colour]\label{lem:tc-few-changes}
\lean{Median.TwoChoices.avg_exp_twice,Median.TwoChoices.expected_see_distinct,Median.TwoChoices.expected_changed_le}\leanok\uses{def:two-choices}
The number of nodes that see colour $i$ twice is a sum of $n$ independent Bernoulli variables of
parameter $(c_i/n)^2$, so $\E e^{\theta\,\mathrm{twice}_i}=(1+(c_i/n)^2(e^\theta-1))^n$. The expected
number of nodes whose two samples hold different colours is $n-S/n$, and if every colour has at
most $\ell$ nodes, the expected number of nodes that change colour is at most $\ell$.
\end{lemma}
\begin{proof}\leanok
The samples of different nodes are independent. Two samples hold the same colour with probability
$S/n^2$, and only such nodes can change colour; $S\le\ell n$.
\end{proof}

\begin{lemma}[One colour does not grow fast]\label{lem:tc-escape}
\lean{Median.TwoChoices.colour_escape_le}\leanok\uses{lem:tc-basic,lem:tc-few-changes}
If $c_i\le L$, then for every $T$, the probability that $c_i(t)>L$ for some $t\le T$ is at most
$\exp(-(L-c_i)+(e-1)TL^2/n)$.
\end{lemma}
\begin{proof}\leanok
Induction on the horizon. While $c_i\le L$, $c_i'\le c_i+\mathrm{twice}_i$ (Lemma~\ref{lem:tc-basic})
and, with $\theta=1$, $\E e^{\mathrm{twice}_i}\le\exp((e-1)n(c_i/n)^2)\le\exp((e-1)L^2/n)$, so
$e^{c_i(t)-(e-1)tL^2/n}$, stopped when $c_i$ exceeds $L$, is a supermartingale; conclude with
Markov's inequality. This replaces the paper's domination by a process with binomial increments.
\end{proof}

\begin{theorem}[Theorem~5]\label{thm:tc-lower-strong}
\lean{Median.TwoChoices.lower_bound_strong}\leanok\uses{lem:tc-escape}
There is a constant $\gamma_0$ such that for every $\gamma\ge\gamma_0$: if every colour has at most
$\ell$ nodes, $\ell'=\max\{2\ell,\gamma\log n\}$ and $T<n/(\gamma\ell')$, the probability that some
colour has more than $\ell'$ nodes at some time $t\le T$ is at most $1/n$.
\end{theorem}
\begin{proof}\leanok
Take $\gamma_0=8$. By Lemma~\ref{lem:tc-escape} with $L=\ell'$, since $L-c_i\ge\ell'/2$ and
$T\ell'^2/n<\ell'/\gamma$, one colour exceeds $\ell'$ with probability at most
$\exp(-\ell'(1/2-(e-1)/\gamma))\le n^{-2}$; take a union bound over the at most $n$ colours
present. A constant in the paper's proof needs a minor correction ($\gamma\ge36$ instead of
$\gamma\ge18$ in its Chernoff step); the theorem is unaffected.
\end{proof}

\begin{theorem}[2-Choices needs $\Omega(n/\log n)$ rounds, Theorem~1 (Simplified)]\label{thm:tc-lower}
\lean{Median.TwoChoices.consensus_time_lower}\leanok\uses{thm:tc-lower-strong}
For every $\beta>0$ there is $C>0$ such that, from every configuration in which every colour has at
most $\beta\log n$ nodes, 2-Choices reaches consensus at some time $t\le T$ with probability at
most $1/n$, for every $T$ with $(T+1)C\log n<n$.
\end{theorem}
\begin{proof}\leanok
Take $\gamma=\max\{8,2\beta\}$ and $C=\gamma^2$. Then $\ell'=\gamma\log n$, the window of
Theorem~\ref{thm:tc-lower-strong} contains $T$, and consensus would put $n>\ell'$ nodes on one
colour.
\end{proof}


%%FILE moran/blueprint/src/content.tex
\begin{abstract}
We formalize the Birth--death Moran process on a finite graph and the ``if'' direction of the
isothermal theorem of Lieberman, Hauert and Nowak (2005): on a connected regular graph, a
mutant population of size $k$ with fitness $r\neq 1$ takes over the whole population with
probability $(1-r^{-k})/(1-r^{-n})$, and with probability $k/n$ when $r=1$. The complete graph
gives Moran's formula (1958). As in the voter project, eventual probabilities are suprema of
increasing finite-time probabilities; no path-space measure is used.
\end{abstract}

\section{The model}\label{sec:lhn-isothermal}
This section and the three following ones formalize the ``if'' direction of the isothermal
theorem of Lieberman, Hauert and Nowak, \emph{Evolutionary dynamics on graphs} (Nature 2005),
with Moran's formula (1958) on the complete graph; the main theorem is in the section on fixation
probabilities.
\begin{definition}[Birth--death Moran kernel]\label{def:moran}
\lean{Moran.Config,Moran.mutants,Moran.fitness,Moran.totalFitness,Moran.target,Moran.pairDist,Moran.moranKernel}\leanok
Mutants have fitness $r>0$ and residents fitness $1$. A step chooses a parent $u$ with
probability proportional to its fitness and an offspring position $w$ uniformly among the
neighbours of $u$ (or $w=u$ if $u$ is isolated); the individual at $w$ is replaced by a copy
of the parent.
\end{definition}
\begin{lemma}[The kernel is well defined]\label{lem:well-defined}
\lean{Moran.target_nonneg,Moran.target_sum_one,Moran.totalFitness_pos,Moran.pair_weight_nonneg,Moran.pair_weight_sum_one}\leanok\uses{def:moran}
The weights of (parent, offspring position) pairs are nonnegative and sum to one.
\end{lemma}
\begin{proof}\leanok
Sum over the offspring position first: each parent's placement weights sum to one.
The remaining sum of fitnesses over the total fitness is one.
\end{proof}

\section{Invariance on regular graphs}
\begin{lemma}[Birth mass is $r$ times death mass]\label{lem:birth-death}
\lean{Moran.orientedCut,Moran.orientedCut_symm,Moran.birthMass,Moran.deathMass,Moran.typed_mass,Moran.birth_eq_r_death}\leanok\uses{lem:well-defined}
On a $d$-regular graph, the probability that a step adds a mutant equals $r$ times the
probability that it removes one: both are proportional to the number of edges between the two
types, counted from either side.
\end{lemma}
\begin{proof}\leanok
For $d>0$ both masses are the oriented cut divided by $Fd$, weighted by the parent's fitness;
the cut is symmetric under swapping endpoints. For $d=0$ both masses vanish.
\end{proof}
\begin{theorem}[Invariant potentials]\label{thm:invariance}
\lean{Moran.apply_affine,Moran.moran_potential_invariant,Moran.moran_mutants_invariant}\leanok\uses{lem:birth-death}
On a regular graph, $(1/r)^{\#\text{mutants}}$ is preserved in expectation by one step; when
$r=1$, so is the number of mutants.
\end{theorem}
\begin{proof}\leanok
A step changes the number of mutants by $\pm1$ or not at all, so its expectation is affine in
the birth and death masses, and these cancel by the previous lemma.
\end{proof}

\section{Absorption on connected graphs}
\begin{theorem}[Fixation or extinction]\label{thm:absorption}
\lean{Moran.unfixed,Moran.moran_constant,Moran.unfixed_step,Moran.unfixed_access,Moran.moran_unfixed_tendsto}\leanok\uses{lem:well-defined}
On a connected graph, the probability that neither type has taken over tends to zero.
\end{theorem}
\begin{proof}\leanok
Constant configurations are absorbing. From a mixed configuration, a mutant adjacent to a
resident (which exists by connectivity) reproduces onto it with positive probability; by
induction on the number of residents, fixation is reachable. Apply the shared finite-chain
absorption theorem.
\end{proof}

\section{Fixation probabilities}
\begin{lemma}[Fixation from an invariant]\label{lem:fixation-limit}
\lean{Moran.allMutant,Moran.fixation,Moran.iterate_allMutant,Moran.unfixed_eq,Moran.iterate_unfixed,Moran.fixation_eq_of_invariant}\leanok
Let an observable be invariant in expectation, equal to $1$ on the all-mutant and to $0$ on the
all-resident configuration. If the all-mutant configuration is absorbing and the probability
that neither type has taken over tends to zero, the observable is the fixation probability.
\end{lemma}
\begin{proof}\leanok
This is the finite-horizon optional stopping theorem of the shared \texttt{dynamics} package
(\texttt{Dynamics.Kernel.iSup\_event\_of\_invariant}, roadmap FND-4) with the two constant
configurations as targets: the fixation and unfixed indicators are the corresponding events.
\end{proof}
\begin{theorem}[Isothermal theorem]\label{thm:isothermal}
\lean{Moran.moranPsi,Moran.moranPsi_invariant,Moran.moranPsi_all,Moran.moranPsi_none,Moran.isothermal}\leanok\uses{thm:invariance,thm:absorption,lem:fixation-limit}
On a connected regular graph with $r\neq1$, the fixation probability from $k$ mutants is
$(1-r^{-k})/(1-r^{-n})$.
\end{theorem}
\begin{proof}\leanok
Apply the previous lemma to $\psi=(1-(1/r)^{\#\text{mutants}})/(1-(1/r)^n)$, invariant by the
invariance theorem and equal to $1$ and $0$ on the two constant configurations; the
absorption theorem gives the vanishing survival probability.
\end{proof}
\begin{theorem}[Neutral case]\label{thm:neutral}
\lean{Moran.neutralPsi,Moran.neutralPsi_invariant,Moran.neutralPsi_all,Moran.neutralPsi_none,Moran.isothermal_neutral}\leanok\uses{thm:invariance,thm:absorption,lem:fixation-limit}
On a connected regular graph with $r=1$, the fixation probability is $k/n$.
\end{theorem}
\begin{proof}\leanok
The same argument with $\psi=\#\text{mutants}/n$.
\end{proof}
\begin{theorem}[Moran's formula]\label{thm:moran}
\lean{Moran.moran_formula}\leanok\uses{thm:isothermal}
On the complete graph, the fixation probability from $k$ mutants is $(1-r^{-k})/(1-r^{-n})$.
\end{theorem}
\begin{proof}\leanok
The complete graph is connected and $(n-1)$-regular.
\end{proof}
The converse of the isothermal theorem does not hold in general (Galanis, Göbel, Goldberg, Lapinskas and
Richerby, Proposition~12), so only this direction is formalized.

\section{Push versus pull on arbitrary graphs}\label{sec:ars-push-pull}
On the complete graph the neutral push and pull voter processes coincide; on an irregular graph
they fix a mutant set with different probabilities, weighted by $1/\deg$ and by $\deg$
respectively. In population genetics this is the difference between Birth--death and
death--Birth updating. Antal, Redner and Sood, \emph{Evolutionary dynamics on
degree-heterogeneous graphs} (Phys.\ Rev.\ Lett.\ 96, 2006), derive these weights by a
mean-field approximation; here they are proved exactly, on every connected graph.
\begin{definition}[Death--Birth kernel]\label{def:pull}
\lean{Moran.pullWeight,Moran.pullDist,Moran.pullKernel,Moran.pull_weight_nonneg,Moran.pull_weight_sum_one}\leanok\uses{def:moran}
In one pull (death--Birth) step, a uniformly random vertex $u$ adopts the type of a uniformly
random neighbour $w$ (an isolated vertex keeps its type), so the pair $(u,w)$ has weight
$\frac1n\cdot\frac1{\deg u}$ for $w\sim u$. The neutral push (Birth--death) process is the
Moran kernel with $r=1$.
\end{definition}
\begin{definition}[Reproductive values]\label{def:repro}
\lean{Moran.pushValue,Moran.pullValue}\leanok
For a configuration with mutant set $S$, $\Phi_{\mathrm{push}}(S)=\sum_{v\in S}1/\deg v$ and
$\Phi_{\mathrm{pull}}(S)=\sum_{v\in S}\deg v$.
\end{definition}
\begin{theorem}[Reproductive values are invariant]\label{thm:pushpull-invariance}
\lean{Moran.pushValue_update,Moran.pullValue_update,Moran.push_increment,Moran.pull_increment,Moran.sum_swap_neg,Moran.push_value_invariant,Moran.pull_value_invariant}\leanok\uses{def:pull,def:repro}
On every finite graph, $\Phi_{\mathrm{push}}$ is invariant in expectation under neutral push and
$\Phi_{\mathrm{pull}}$ is invariant in expectation under pull.
\end{theorem}
\begin{proof}\leanok
Write $x_v\in\{0,1\}$ for the type of $v$. Under push, the pair (parent $u$, target $w\sim u$)
has weight $\frac1n\frac1{\deg u}$ and changes $\Phi_{\mathrm{push}}$ by $(x_u-x_w)/\deg w$, so
the expected change is $\frac1n\sum_{u\sim w}(x_u-x_w)/(\deg u\deg w)$. Under pull, $u$ copies
$w\sim u$ with weight $\frac1n\frac1{\deg u}$ and changes $\Phi_{\mathrm{pull}}$ by
$(x_w-x_u)\deg u$, so the expected change is $\frac1n\sum_{u\sim w}(x_w-x_u)$. Both sums are
over ordered adjacent pairs of an antisymmetric function, hence vanish.
\end{proof}
\begin{theorem}[Fixation or extinction under pull]\label{thm:pull-absorption}
\lean{Moran.pull_constant,Moran.pull_unfixed_step,Moran.pull_move_weight_pos,Moran.pull_unfixed_access,Moran.pull_allMutant_step,Moran.pull_unfixed_tendsto}\leanok\uses{def:pull}
On a connected graph, the probability that the pull process has not reached a constant
configuration after $t$ steps tends to $0$.
\end{theorem}
\begin{proof}\leanok
As for push: constant configurations are absorbing, and from any other configuration a resident
adjacent to a mutant copies it with positive probability, so the number of mutants reaches $n$
with positive probability; the shared absorption theorem concludes.
\end{proof}
\begin{theorem}[Push and pull fixation]\label{thm:pushpull}
\lean{Moran.pushRatio_invariant,Moran.pullRatio_invariant,Moran.pushValue_false,Moran.pullValue_false,Moran.pushValue_total_pos,Moran.pullValue_total_pos,Moran.push_fixation,Moran.pull_fixation}\leanok\uses{thm:pushpull-invariance,thm:absorption,thm:pull-absorption,lem:fixation-limit}
On a connected graph with at least two vertices, the fixation probability of a mutant set $S$
is $\sum_{v\in S}\frac1{\deg v}\big/\sum_{v\in V}\frac1{\deg v}$ under neutral push and
$\sum_{v\in S}\deg v/(2|E|)$ under pull.
\end{theorem}
\begin{proof}\leanok
Apply the fixation-limit lemma to the normalized reproductive values, which are invariant and
equal to $1$ and $0$ on the two constant configurations; the denominators are positive because
every degree is at least $1$.
\end{proof}

\section{Exact fixation on the star}\label{sec:broom-rychtar}
The star with $n$ leaves is Mathlib's \texttt{starGraph} $c$ on $n+1$ vertices: the centre $c$
is adjacent to every leaf and leaves only to $c$. It is not regular, and it amplifies selection
(Lieberman, Hauert and Nowak 2005); its exact fixation probability is due to Broom and
Rycht\'a\v{r} (2008, \S5). Write
\[ q=\frac{n+r}{r(nr+1)},\qquad \kappa=\frac{nr+1}{r(n+r)}=\frac1{r^2q}. \]
\begin{definition}[Star potential]\label{def:star-potential}
\lean{Moran.leafMutants,Moran.starRatio,Moran.starCentreWeight,Moran.starPotential}\leanok
For a configuration with $i$ mutant leaves, $\Phi=q^i$ if the centre is a resident and
$\Phi=\kappa q^i$ if it is a mutant.
\end{definition}
\begin{theorem}[The star potential is invariant]\label{thm:star-invariance}
\lean{Moran.star_sum_pairs,Moran.star_edge_cancel,Moran.starPotential_apply,Moran.star_potential_invariant}\leanok\uses{def:moran,def:star-potential}
For every $r>0$, $\Phi$ is invariant in expectation under one Birth--death step on the star.
\end{theorem}
\begin{proof}\leanok
Only the ordered pairs $(c,\ell)$ and $(\ell,c)$ for a leaf $\ell$ carry weight, namely
$\mathrm{fit}(c)/(nF)$ and $\mathrm{fit}(\ell)/F$. Along each edge the two contributions cancel:
if $\ell$ is a mutant and $c$ a resident they give $\frac{q^{i-1}}F\big(\frac{1-q}n+rq(\kappa-1)\big)=0$,
and if $c$ is a mutant and $\ell$ a resident they give
$\frac{q^i}F\big(1-\kappa+\frac{r\kappa(q-1)}n\big)=0$; equal types change nothing.
\end{proof}
\begin{theorem}[Fixation on the star]\label{thm:star}
\lean{Moran.star_fixation_eq,Moran.star_fixation,Moran.star_fixation_leaf,Moran.star_fixation_centre,Moran.star_fixation_uniform}\leanok\uses{thm:star-invariance,thm:absorption,lem:fixation-limit}
For $r\ne1$, the fixation probability from $s$ is $(1-\Phi(s))/(1-\kappa q^n)$. A single mutant
fixes with probability $(1-q)/(1-\kappa q^n)$ from a leaf, $(1-\kappa)/(1-\kappa q^n)$ from the
centre, and $(n(1-q)+(1-\kappa))/((n+1)(1-\kappa q^n))$ from a uniformly random vertex.
\end{theorem}
\begin{proof}\leanok
For $r>1$ we have $q\le1$ and $\kappa<1$, so $\kappa q^n\le\Phi\le1$; for $r<1$ all inequalities
reverse. Hence the normalized potential is $1$ at fixation, $0$ at extinction and in $[0,1]$
otherwise, and the fixation-limit lemma applies.
\end{proof}
\begin{theorem}[The formula of Broom and Rycht\'a\v{r}]\label{thm:star-sum}
\lean{Moran.star_closed_eq_sum,Moran.star_sum_neutral,Moran.neutral_uniform_fixation,Moran.star_fixation_uniform_sum}\leanok\uses{thm:star,thm:pushpull}
For every $r>0$, a single mutant at a uniformly random vertex fixes with probability
\[ \frac{n\cdot\frac{nr}{nr+1}+\frac r{r+n}}{(n+1)\big(1+\frac n{n+r}\sum_{j=1}^{n-1}q^j\big)}. \]
\end{theorem}
\begin{proof}\leanok
For $r\ne1$ and $n\ge1$, sum the geometric series: $1+\frac n{n+r}\sum_{j=1}^{n-1}q^j=
r^2(1-\kappa q^n)/(r^2-1)$. For $r=1$ both sides are $1/(n+1)$, by push--pull fixation; for
$n=0$ both are $1$.
\end{proof}
\begin{theorem}[The star is an amplifier]\label{thm:star-amplifier}
\lean{Moran.geom_gap_pos,Moran.star_delta_pos,Moran.star_amplifier_sign,Moran.star_amplifier,Moran.star_amplifier_deleterious}\leanok\uses{thm:star}
For every $n\ge2$, the fixation probability from a uniformly random vertex is larger than
Moran's $(1-1/r)/(1-r^{-(n+1)})$ for $r>1$ and smaller for $0<r<1$.
\end{theorem}
\begin{proof}\leanok
Put $s=1/r$ and $t=(1+ns)/(n+s)$, so that $q=st$ and $\kappa=s/t$. The difference has the sign
of $(1-s)\big(E(1-s^{n+1})-s^{n+1}(1-t^{n-1})\big)$ with $E=s^2(n-1)^2/((n+s)(1+ns))$. By
Bernoulli, $1-t^{n-1}\le(n-1)(1-t)=(n-1)^2(1-s)/(n+s)$, which reduces the claim to
$s^{n-1}(1-s)(1+ns)<1-s^{n+1}$. Its slack is $\sum_{k=0}^{n-2}(1-s)(s^k-s^n)$, a nonempty sum of
positive terms on both sides of $s=1$.
\end{proof}
\begin{theorem}[Large stars amplify $r$ to $r^2$]\label{thm:star-limit}
\lean{Moran.starRatio_le_inv,Moran.star_closed_tendsto,Moran.star_fixation_uniform_tendsto}\leanok\uses{thm:star}
For $r>1$, the fixation probability from a uniformly random vertex of the star with $n$ leaves
tends to $1-1/r^2$ as $n\to\infty$.
\end{theorem}
\begin{proof}\leanok
$n(1-q)/(n+1)=\frac n{n+1}\cdot\frac n{n+1/r}\cdot(1-r^{-2})\to1-r^{-2}$,
$(1-\kappa)/(n+1)\to0$, and $0\le\kappa q^n\le r^{-n}\to0$ because $q\le1/r$ for $n\ge1$.
\end{proof}


%%FILE plurality/blueprint/src/content.tex
% Blueprint for the formalization of
%   L. Becchetti, A. Clementi, E. Natale, F. Pasquale, R. Silvestri, L. Trevisan,
%   "Simple Dynamics for Plurality Consensus", SPAA 2014 (arXiv:1310.2858).
% Numbering follows the SPAA 2014 proceedings version. Every statement carries
% \lean{...} / \leanok / \uses{...} tags matched to the Lean declarations in
% Plurality/*.lean (namespace `Plurality`). Statements without \leanok are
% not formalized yet. Where the Lean proof departs from the paper, this file
% follows the Lean proof; the departures are listed in
% FORMALIZATION_DIFFERENCES.md.

\begin{abstract}
\noindent
$n$ anonymous nodes each support one of $k$ colors. In every round each node
samples three nodes uniformly at random (with repetition, possibly itself)
and adopts the majority color among them, or the first one if all three
differ. If some color $m$ leads every other color by a sufficiently large
\emph{bias} $s$, the process reaches consensus on $m$. The paper proves that
this takes $O(\lambda\log n)$ rounds w.h.p.\ whenever $c_m \ge n/\lambda$
and $s \ge c\sqrt{\lambda n\log n}$ (Theorem~\ref{thm:3_8}), that
$\Omega(k\log n)$ rounds are needed from balanced configurations, that
3-majority is essentially the only 3-input rule that solves plurality
consensus, and that sampling $h > 3$ nodes speeds the process up by at most
a factor $O(h^2)$.

The upper bound (Section~\ref{sec:upper}), including the headline bound
$O(\min\{k,(n/\log n)^{1/3}\}\log n)$ of Corollary~\ref{cor:3_10}, is fully
formalized, with no
\texttt{sorry}, on the finite probability layer of the shared
\texttt{dynamics} package: all randomness is uniform over finite types,
multi-round probabilities are iterates of a finite Markov kernel, and the
concentration inequalities (Chernoff, Hoeffding, Bernstein) are proved from
scratch. The lower bounds of Section~\ref{sec:lower} are formalized as well,
with two exceptions where the paper's argument does not go through: the
$\Omega(k\log n)$ bound is proved for $k \le n^{1/4-\delta}$ rather than up to
$(n/\log n)^{1/4}$, and part~(a) of Theorem~4.8 (Theorem~\ref{thm:4_8_a}) leaves open rules
with $\Delta_r, \Delta_b \le 1$. Observation~3.9 (an adversary) is not
formalized. Beyond the paper, Section~\ref{sec:any_start} proves that with two opinions
3-majority reaches consensus from any configuration, including a perfectly balanced one, within
$O(\log n)$ rounds w.h.p.
\end{abstract}

\section{The model}\label{sec:model}

\begin{definition}[Colorings and color counts]\label{def:config}
  \lean{Plurality.Config, Plurality.count, Plurality.sum_count}
  \leanok
  A \emph{coloring} is a map $x \colon [n] \to [k]$. The number of nodes of
  color $j$ is $c_j = |\{v : x(v) = j\}|$, and $\sum_j c_j = n$. The vector
  $c = (c_1,\dots,c_k)$ is the paper's $k$-color distribution ($k$-cd).
\end{definition}

\begin{definition}[3-input rules and the 3-majority rule]\label{def:rule}
  \lean{Plurality.maj3, Plurality.maj3_mem}
  \leanok
  A \emph{3-input rule} is a map $f \colon [k]^3 \to [k]$ with
  $f(a,b,c) \in \{a,b,c\}$. The \emph{3-majority rule} is
  $\mathrm{maj}_3(a,b,c) = b$ if $b = c$ and $a$ otherwise; it returns the
  majority color, and the first color if all three differ.
\end{definition}

\begin{definition}[One round and the process]\label{def:step}
  \lean{Plurality.stepWith, Plurality.step, Plurality.runWith, Plurality.run}
  \leanok
  \uses{def:config, def:rule}
  A round is a map $r \colon [n] \to [n]^3$ (the three samples of every node),
  drawn uniformly. The next coloring under rule $f$ is
  $v \mapsto f(x(r_v^1), x(r_v^2), x(r_v^3))$. The process after $T$ rounds
  folds this step over $T$ independent uniform rounds.
\end{definition}

\begin{definition}[Consensus]\label{def:mono}
  \lean{Plurality.Mono, Plurality.mono_iff_count, Plurality.stepWith_mono, Plurality.run_mono}
  \leanok
  \uses{def:step}
  $x$ is in \emph{consensus on $m$} if every node has color $m$, i.e.\
  $c_m = n$. For every 3-input rule, consensus is absorbing.
\end{definition}

\begin{definition}[The Markov kernel]\label{def:kernel}
  \lean{Plurality.kernel, Plurality.kernel_event, Dynamics.Kernel.ofStep, Dynamics.Kernel.iterate_ofStep}
  \leanok
  \uses{def:step}
  The process with rule $f$ is the finite Markov kernel on colorings whose
  row at $x$ is the law of one round of $f$ from $x$. Its $T$-step event
  probabilities are the expectations over $T$ independent uniform rounds.
\end{definition}

\section{The next expected coloring}\label{sec:expected}

\begin{lemma}[Adoption probability]\label{lem:adopt}
  \lean{Plurality.count_stepWith_eq_sum, Plurality.avg_adopt, Plurality.sum_triple_maj3, Plurality.avg_adopt_maj3}
  \leanok
  \uses{def:step}
  The new count of color $j$ is a sum of $n$ independent indicators, one per
  node. Under 3-majority, a node adopts $j$ with probability
  $x_j^2 + x_j(1 - \sum_h x_h^2)$, where $x_h = c_h/n$.
\end{lemma}

\begin{proof}
  The three samples are independent with law $(x_h)_h$. Summing the
  indicator of $\mathrm{maj}_3(a,b,c) = j$ first over $c$, then $b$, then $a$
  gives the polynomial.
\end{proof}

\begin{lemma}[Lemma 2.1, next expected coloring]\label{lem:2_1}
  \lean{Plurality.expected_count}
  \leanok
  \uses{lem:adopt}
  For every coloring and every color $j$,
  \[
    \mu_j(c) = \E[C_{j,t+1} \mid C_t = c]
      = c_j\Bigl(1 + \frac{1}{n^2}\Bigl(n c_j - \sum_h c_h^2\Bigr)\Bigr).
  \]
\end{lemma}

\begin{proof}
  Multiply the adoption probability of Lemma~\ref{lem:adopt} by $n$.
\end{proof}

\section{The upper bound}\label{sec:upper}

This section formalizes Theorem~3.8 of Becchetti, Clementi, Natale, Pasquale, Silvestri and
Trevisan, \emph{Simple Dynamics for Plurality Consensus} (SPAA 2014), and its Corollaries 3.10
to~3.12. The SPAA proof of Theorem~3.8 needs a minor correction (small repairs: the saturation
phases as defined are not shown to be preserved, and a union bound over the competing colors is
missing), and the repaired proof is what is formalized
(\href{https://github.com/formal-dynamics/leanamics/blob/main/plurality/FORMALIZATION_DIFFERENCES.md#1-theorem-38-how-the-lean-proof-differs-from-spaa}{details}).

\subsection{Plurality, bias, \texorpdfstring{$\alpha$}{alpha} and \texorpdfstring{$\gamma$}{gamma}}

\begin{definition}[Section 3 quantities]\label{def:quantities}
  \lean{Plurality.maxc, Plurality.argmaxSet, Plurality.secondc, Plurality.bias, Plurality.alpha, Plurality.gamma, Plurality.mu}
  \leanok
  For a $k$-cd $c$ with $\sum_h c_h = n$: $m(c) = \max_h c_h$,
  $M(c) = \{j : c_j = m(c)\}$, the \emph{bias}
  $s(c) = m(c) - \max_{h \notin M(c)} c_h$ if $|M(c)| = 1$ and $0$ otherwise,
  $\alpha(c) = (n - m(c))\,s(c)/n^2$ and
  $\gamma(c) = (n\,m(c) - \sum_h c_h^2)/n^2 - \alpha(c)$.
\end{definition}

\begin{lemma}[Every other color trails by the bias]\label{lem:bias_gap}
  \lean{Plurality.add_bias_le, Plurality.exists_add_bias_eq, Plurality.gamma_mul_sq}
  \leanok
  \uses{def:quantities}
  If $c_m = m(c)$ then $c_h + s(c) \le m(c)$ for every $h \ne m$, with
  equality for some $h$ when $k \ge 2$. Moreover
  $\gamma(c)\,n^2 = \sum_{h \ne m} c_h\,(m(c) - s(c) - c_h)$.
\end{lemma}

\begin{lemma}[Lemma 3.1]\label{lem:3_1}
  \lean{Plurality.lemma_3_1_a, Plurality.lemma_3_1_b, Plurality.lemma_3_1_c}
  \leanok
  \uses{lem:bias_gap}
  For every $k$-cd $c$:
  (a) $0 \le s(c) \le m(c) - (n - m(c))/(k-1)$;
  (b) $0 \le \alpha(c) \le \min\{s(c)/n,\ 1/4\}$;
  (c) $0 \le \gamma(c) \le 1/8$.
\end{lemma}

\begin{proof}
  (a) The $k-1$ other colors each have at most $m(c) - s(c)$ nodes.
  (b) $(n - m)m \le n^2/4$ and $s \le m$.
  (c) Use the identity of Lemma~\ref{lem:bias_gap}: every summand is
  nonnegative, the color attaining the bias contributes $0$, and the rest is
  at most $d(n - m - d) \le n^2/8$ with $d = m(c) - s(c)$.
\end{proof}

\begin{lemma}[Lemma 3.2]\label{lem:3_2}
  \lean{Plurality.lemma_3_2_a, Plurality.lemma_3_2_b, Plurality.lemma_3_2_c}
  \leanok
  \uses{def:quantities, lem:bias_gap}
  Let $m \in M(c)$ and let $\ell \ne m$ have the most nodes among the other
  colors. Then (a) $\mu_m(c) = c_m(1 + \gamma + \alpha)$;
  (b) $\mu_j(c) = c_j(1 + \gamma + \alpha - (m(c) - c_j)/n)$ for every $j$;
  (c) $\mu_\ell(c) = c_\ell(1 + \gamma + \alpha - s(c)/n)$.
\end{lemma}

\subsection{Concentration}

\begin{lemma}[Concentration for sums of independent coordinates]\label{lem:tails}
  \lean{Dynamics.avg_bernstein, Dynamics.avg_hoeffding, Dynamics.avg_chernoff_lower_mul,
    Dynamics.avg_markov_one, Dynamics.avg_chernoff_upper_log}
  \leanok
  For $X = \sum_i Y_i(\omega_i)$ with independent uniform coordinates:
  Bernstein's inequality
  $\Pr(X \ge \E X + \lambda) \le \exp(-\lambda^2/(2\sigma^2(1 + b\lambda/(3\sigma^2))))$
  when $Y_i - \E Y_i \le b$ and $\sum_i \mathrm{Var}\,Y_i \le \sigma^2$;
  Hoeffding's inequality for $\{0,1\}$ coordinates; the multiplicative
  Chernoff lower tail $\Pr(X \le (1-\delta)\mu) \le \exp(-\delta^2\mu/2)$;
  Markov's inequality $\Pr(X \ge 1) \le \E X$ for nonnegative coordinates;
  and the closed-form Chernoff upper tail
  $\Pr(X \ge k) \le \exp(k - \mu - k\log(k/\mu))$ for $\E X \le \mu \le k$.
  All are proved once in the shared \texttt{dynamics} library
  (\texttt{Dynamics/Concentration.lean}, \texttt{Dynamics/Chernoff.lean},
  \texttt{Dynamics/Tail.lean}).
\end{lemma}

\subsection{Growth of the bias}

\begin{lemma}[Lemma 3.3, increasing rate of the bias]\label{lem:3_3}
  \lean{Plurality.lemma_3_3}
  \leanok
  \uses{lem:2_1, lem:3_1, lem:3_2, lem:tails}
  If $M(c) = \{m\}$ then for every $j \ne m$,
  \[
    \Pr\Bigl(C_{m,t+1} - C_{j,t+1} \le s(c)\Bigl(1 + \gamma(c)
      + \frac{c_m\alpha(c)}{2s(c)}\Bigr) \Bigm| C_t = c\Bigr)
      \le \exp\Bigl(-\frac{c_m\alpha(c)^2}{25}\Bigr).
  \]
\end{lemma}

\begin{proof}
  By Lemmas~\ref{lem:3_1} and~\ref{lem:3_2},
  $\E[C_{m,t+1} - C_{j,t+1}] \ge s(1+\gamma) + c_m\alpha$. The difference is
  a sum of per-node variables in $\{-1,0,1\}$ with variance at most $3c_m$;
  apply Bernstein's inequality with $b = 2$ and $\lambda = c_m\alpha/2$.
\end{proof}

\begin{lemma}[Lemma 3.4]\label{lem:3_4}
  \lean{Plurality.lemma_3_4}
  \leanok
  \uses{lem:3_1, lem:3_2, lem:tails}
  If $M(c) = \{m\}$ then
  $\Pr(C_{m,t+1} \le c_m(1 + \gamma + \alpha/2) \mid C_t = c)
  \le \exp(-c_m\alpha^2/11)$.
\end{lemma}

\begin{proof}
  The multiplicative Chernoff lower tail with
  $\delta = \alpha/(2(1 + \gamma + \alpha))$.
\end{proof}

\begin{lemma}[Lemma 3.5, large plurality and large bias]\label{lem:3_5}
  \lean{Plurality.lemma_3_5}
  \leanok
  \uses{lem:3_3, lem:3_4}
  If $M(c) = \{m\}$, $0 < \lambda \le 2/3$, $\lambda n \le c_m \le (2/3)n$
  and $s(c) \ge 22\sqrt{(1/\lambda)\,n\log n}$, then for every $j \ne m$
  \[
    \Pr\bigl(C_{m,t+1} - C_{j,t+1} \le s(c)(1 + \lambda/6)\bigr) \le 1/n^2
    \quad\text{and}\quad
    \Pr\bigl(C_{m,t+1} \le c_m\bigr) \le 1/n^2.
  \]
\end{lemma}

\begin{proof}
  Under the hypotheses $c_m\alpha/(2s) \ge \lambda/6$ and
  $c_m\alpha^2/25 \ge 2\log n$; conclude with Lemmas~\ref{lem:3_3}
  and~\ref{lem:3_4}.
\end{proof}

\subsection{Saturation}

\begin{lemma}[Lemma 3.6]\label{lem:3_6}
  \lean{Plurality.lemma_3_6}
  \leanok
  \uses{lem:2_1, lem:3_1}
  For every $m \in M(c)$, writing $\bar C_{m,t+1} = n - C_{m,t+1}$,
  \[
    (n - c_m)\Bigl(1 - \frac{c_m^2}{n^2}\Bigr)
      \le \E[\bar C_{m,t+1} \mid C_t = c]
      \le (n - c_m)\Bigl(1 - \frac{s(c)\,c_m}{n^2}\Bigr).
  \]
\end{lemma}

\begin{lemma}[Lemma 3.7, very large plurality]\label{lem:3_7}
  \lean{Plurality.lemma_3_7_i, Plurality.lemma_3_7_ii, Plurality.bias_ge_of_singleton}
  \leanok
  \uses{lem:3_6, lem:tails}
  Let $\log n \ge 40$, $M(c) = \{m\}$.
  (i) If $s(c) \ge n/3$ and $n - c_m \ge n^{1/4}\log n$ then
  $\Pr(\bar C_{m,t+1} \ge \frac{17}{18}(n - c_m)) \le 1/n^2$.
  (ii) If $n - c_m < n^{1/4}\log n$ then $\Pr(\bar C_{m,t+1} > 0) \le n^{-1/5}$
  and $\Pr(\bar C_{m,t+1} \ge n^{1/4}\log n) \le 1/n^2$.
\end{lemma}

\begin{proof}
  The paper omits this proof. (i) By Lemma~\ref{lem:3_6} and
  $s\,c_m \ge n^2/9$, $\E\bar C_{m,t+1} \le \frac89(n - c_m)$; apply the
  closed-form Chernoff upper tail at ratio $17/16$. (ii) Here
  $s(c) \ge n - 2(n - c_m)$, so $\E\bar C_{m,t+1} \le 3(n-c_m)^2/n
  \le 3\log^2 n\, e^{-(\log n)/2}$; apply Markov's inequality for the first
  bound and the Chernoff upper tail for the second.
\end{proof}

\subsection{Nested phases and the main theorem}

\begin{lemma}[Lemma A.4, nested phases]\label{lem:A_4}
  \lean{Dynamics.Kernel.nested_phases}
  \leanok
  Let $A_1 \supseteq \cdots \supseteq A_T$ be sets of states of a finite
  Markov chain such that from every state of $A_i$ the chain stays in $A_i$
  with probability $\ge 1-\varepsilon$ and, for $i < T$, moves into
  $A_{i+1}$ with probability $\ge 1 - \nu$. Then from any state of $A_1$,
  after $\ell T$ steps the chain is in $A_T$ with probability at least
  $1 - T(\ell\varepsilon + \nu^\ell)$.
\end{lemma}

\begin{proof}
  The paper omits this proof. Phase by phase: within $\ell$ steps the chain
  leaves $A_i$ with probability at most $\ell\varepsilon$, and fails to jump
  into $A_{i+1}$ in each of the $\ell$ steps with probability at most
  $\nu^\ell$. The hypothesis $\varepsilon \le \nu$ of the paper is not needed.
\end{proof}

\begin{definition}[The phases]\label{def:phases}
  \lean{Plurality.growthSet, Plurality.satSet, Plurality.monoSet}
  \leanok
  \uses{def:quantities, def:mono}
  Fix $m$, $\lambda$, $\Lambda$ and let $q = 1 + 1/(6\lambda)$,
  $B = n^{1/4}\log n$.
  \emph{Growth phase $i$}: $c_m > \frac23 n$, or $M(c) = \{m\}$,
  $c_m \ge n/\lambda$ and $s(c) \ge \Lambda q^i$.
  \emph{Saturation phase $j$}:
  $n - c_m < \max\{\frac n3 (\frac{17}{18})^j,\ B\}$.
  \emph{Consensus}: $c_m = n$.
\end{definition}

\begin{remark}
  The paper's saturation phases also require $s(c) \ge n/3$, which
  Lemma~\ref{lem:3_7} does not show to be preserved. Here the saturation
  phases only bound $n - c_m < n/3$, which already implies $M(c) = \{m\}$
  and $s(c) \ge n/3$ (Lemma~\ref{lem:dissent}).
\end{remark}

\begin{lemma}[A clear majority]\label{lem:dissent}
  \lean{Plurality.of_dissent_lt, Plurality.argmaxSet_eq_of_lt, Plurality.le_bias_of_gap, Plurality.count_step_eq_zero}
  \leanok
  \uses{def:quantities}
  If $n - c_m < n/3$ then $M(c) = \{m\}$ and $s(c) \ge n/3$. If $m$ leads
  every other color by at least $g > 0$ and $c_m \ge g$, then $M(c) = \{m\}$
  and $s(c) \ge g$. A color without nodes stays without nodes.
\end{lemma}

\begin{lemma}[One round of each phase]\label{lem:phase_steps}
  \lean{Plurality.growth_step, Plurality.sat_step, Plurality.mono_step, Plurality.mono_stay, Plurality.growthSet_succ_sub, Plurality.satSet_succ_sub, Plurality.growthSet_sub_of_lt, Plurality.satSet_zero_sub}
  \leanok
  \uses{def:phases, def:kernel, lem:3_5, lem:3_7, lem:dissent}
  Let $\log n \ge 40$, $k \ge 2$, $\lambda \ge 3$ and
  $\Lambda \ge 22\sqrt{\lambda n \log n}$.
  From growth phase $i$ the chain enters growth phase $i+1$ with probability
  $\ge 1 - 1/n$. From saturation phase $j$ it enters saturation phase $j+1$
  with probability $\ge 1 - 1/n^2$. From fewer than $B$ dissenting nodes it
  reaches consensus with probability $\ge 1 - n^{-1/5}$, and consensus is
  kept with probability $1$. The phase sets are nested, and growth phase $i$
  lies in saturation phase $0$ once $\Lambda q^i > n$.
\end{lemma}

\begin{proof}
  \emph{Growth, case $c_m > \frac23 n$}: saturation phase $0$ moves into
  saturation phase $1$, which lies in $\{c_m > \frac23 n\}$.
  \emph{Growth, otherwise}: apply Lemma~\ref{lem:3_5} with $1/\lambda$ to each
  other color $j$ that currently has a node, and to the growth of $c_m$. There
  are fewer than $n - c_m$ such colors, so a union bound gives failure
  probability at most $n/n^2$. Colors without nodes stay empty, and the color
  attaining the bias has a node (otherwise $c_m = n$), so after a successful
  round $m$ leads every other color by $s(c)\,q$ and
  Lemma~\ref{lem:dissent} applies.
  \emph{Saturation}: Lemma~\ref{lem:3_7}(i) above $B$ dissenters and
  Lemma~\ref{lem:3_7}(ii) below.
\end{proof}

\begin{lemma}[Number of phases]\label{lem:phases}
  \lean{Plurality.growthPhases, Plurality.satPhases, Plurality.phases, Plurality.phases_le}
  \leanok
  \uses{def:phases}
  With $T_1 = \lceil \log n/\log q\rceil$ growth phases,
  $T_2 = \lceil \log n/\log(18/17)\rceil + 1$ saturation phases and one
  consensus phase, $T = T_1 + T_2 + 1 \le 13\lambda\log n$ for $\lambda \ge 3$
  and $\log n \ge 40$.
\end{lemma}

\begin{proof}
  $\log(1+x) \ge x/(1+x)$ gives $\log q \ge 1/(6\lambda+1)$ and
  $\log(18/17) \ge 1/18$.
\end{proof}

\begin{theorem}[Theorem 3.8, the general upper bound]\label{thm:3_8}
  \lean{Plurality.theorem_3_8, Plurality.theorem_3_8_bigO}
  \leanok
  \uses{lem:A_4, lem:phase_steps, lem:phases}
  Let $\log n \ge 40$, $k \ge 2$ and $\lambda \ge 3$. If $M(c) = \{m\}$,
  $c_m \ge n/\lambda$ and $s(c) \ge 22\sqrt{\lambda n\log n}$, then after
  $10T \le 130\lambda\log n$ rounds all nodes support $m$ with probability
  at least $1 - 11T/n \ge 1 - 143\lambda\log n/n$.
\end{theorem}

\begin{proof}
  Apply Lemma~\ref{lem:A_4} to the kernel of 3-majority with the $T_1$
  growth phases, the $T_2$ saturation phases and consensus,
  $\varepsilon = 1/n$, $\nu = n^{-1/5}$ and $\ell = 10$, so that
  $\nu^\ell = n^{-2}$. The initial coloring is in growth phase $0$; growth
  phase $T_1$ lies in saturation phase $0$ because $\Lambda q^{T_1} \ge 22n$;
  saturation phase $T_2 - 1$ lies below $B$ dissenters because
  $\frac n3(\frac{17}{18})^{T_2-1} \le \frac13$.
\end{proof}

\begin{remark}
  The paper states the failure probability as $11T/n^2$; the extra factor
  $n$ comes from the union bound over competing colors in
  Lemma~\ref{lem:phase_steps}, which the paper omits. The paper's
  hypothesis $\lambda < \sqrt n$ is not used; it only makes the hypothesis
  on $s(c)$ satisfiable.
\end{remark}

\begin{corollary}[Corollary 3.10]\label{cor:3_10}
  \lean{Plurality.corollary_3_10, Plurality.lamK, Plurality.le_mul_maxc}
  \leanok
  \uses{thm:3_8}
  If $M(c) = \{m\}$ and
  $s(c) \ge 22\sqrt{\min\{2k, (n/\log n)^{1/3}\}\,n\log n}$, then 3-majority
  reaches consensus on $m$ in $O(\min\{2k, (n/\log n)^{1/3}\}\log n)$
  rounds w.h.p. Precisely, for $\log n \ge 40$ and $k \ge 2$: at most
  $130\lambda\log n$ rounds with probability at least
  $1 - 143\lambda\log n/n$.
\end{corollary}

\begin{proof}
  \leanok
  Take $\lambda = \min\{2k, (n/\log n)^{1/3}\}$. If $\lambda = 2k$ then
  $c_m \ge n/k$ by pigeonhole; otherwise $\lambda^3 = n/\log n$, so
  $22\sqrt{\lambda n\log n} = 22n/\lambda$ and $c_m \ge s(c) \ge n/\lambda$.
  Moreover $\lambda \ge 3$ because $k \ge 2$ and $n = e^{\log n} \ge 27\log n$.
\end{proof}

\begin{corollary}[Corollaries 3.11 and 3.12]\label{cor:3_11}
  \lean{Plurality.corollary_3_11, Plurality.corollary_3_12}
  \leanok
  \uses{thm:3_8}
  If $c_m \ge n/\log^\ell n$ and $s(c) \ge 22\sqrt{n\log^{\ell+1}n}$,
  consensus takes $O(\log^{\ell+1} n)$ rounds w.h.p. If $c_m \ge n/\beta$
  and $s(c) \ge 22\sqrt{\beta n\log n}$ for a constant $\beta \ge 3$, it
  takes $O(\log n)$ rounds w.h.p.
\end{corollary}

\begin{proof}
  \leanok
  Take $\lambda = \log^\ell n$ (with $\ell \ge 1$, so that $\lambda \ge 3$)
  and $\lambda = \beta$ in Theorem~\ref{thm:3_8}.
\end{proof}

\subsection{Two colors}

\begin{theorem}[Binary case]\label{thm:binary}
  \lean{Plurality.colorSet_step, Plurality.colorSet_run, Plurality.binary_consensus_whp}
  \leanok
  \uses{def:step}
  With $k = 2$ the set of nodes of color $1$ evolves exactly as the binary
  3-majority process of the \texttt{three\_majority} package. Hence, if
  $c_1 \ge \frac35 n$ and $\log n \ge 30$, all nodes support color $1$ after
  $O(\log n)$ rounds with probability at least $1 - 500/n$.
\end{theorem}

\begin{corollary}[Two opinions from a vanishing bias]\label{cor:binary_vanishing}
  \lean{Plurality.ofSet, Plurality.colorSet_ofSet, Plurality.count_ofSet_one, Plurality.count_ofSet_zero, Plurality.majority3_vanishing_bias}
  \leanok
  \uses{thm:3_8, thm:binary}
  In the binary 3-majority process, if the initial majority leads by at least
  $22\sqrt{3n\log n}$ nodes (a fraction $\frac12 + O(\sqrt{\log n/n})$) and $\log n \ge 40$,
  all nodes adopt it within at most $390\log n$ rounds with probability at least
  $1 - 429\log n/n$.
\end{corollary}
\begin{proof}
  \leanok
  Theorem~\ref{thm:3_8} with $k = 2$ and $\lambda = 3$: the majority color is the unique
  plurality, holds more than $n/2 \ge n/3$ nodes, and its bias is the gap; the conclusion
  transfers to the binary process through Theorem~\ref{thm:binary}.
\end{proof}

\subsection{Two colors from any configuration}\label{sec:any_start}

This part goes beyond the paper: it follows the symmetry-breaking argument for the binary
median dynamics of Doerr, Goldberg, Minder, Sauerwald and Scheideler (SPAA 2011), as described in
Becchetti, Clementi, Natale, \emph{Consensus dynamics: an overview} (SIGACT News 2020, \S4,
Case~3), here for binary 3-majority. Let $I$ be the set of nodes with opinion $1$ and
$s = |I| - (n - |I|)$ the gap. After one round, $s' = 2\sum_v Y_v - n$ for independent
$\{0,1\}$-valued $Y_v$ with mean $p = 3x^2 - 2x^3$, $x = |I|/n$, so
$\mathbb{E}s' = s\,(3/2 - s^2/(2n^2))$.

\begin{lemma}[Jump near balance]\label{lem:jump_near_balance}
  \lean{Plurality.gap, Plurality.binKernel, Plurality.avg_pz_zero_one, Plurality.jump_near_balance}
  \leanok
  \uses{thm:binary}
  If $n \ge 5$ and $|s| \le 4\sqrt n/25$, then $|s'| \ge \sqrt n/5$ with probability at
  least $9/64$.
\end{lemma}
\begin{proof}
  \leanok
  The centered sum $S = \sum_v (Y_v - p)$ has variance $\sigma^2 = np(1-p) \ge n/5$ and fourth
  moment at most $\sigma^2 + 3\sigma^4$, so by Paley--Zygmund $4S^2 \ge \sigma^2$ with
  probability at least $9/64$. Then $|s'| = |2S + \mathbb{E}s'| \ge \sqrt{n/5} - \frac32|s|
  \ge \sqrt n/5$.
\end{proof}

\begin{lemma}[Growth away from balance]\label{lem:growth_far}
  \lean{Plurality.mean_ge, Plurality.growth_far}
  \leanok
  \uses{thm:binary}
  If $0 \le s \le n/2$ and $\lambda \ge 0$, then $s' > \frac{11}{8}s - 2\lambda$ with
  probability at least $1 - e^{-2\lambda^2/n}$.
\end{lemma}
\begin{proof}
  \leanok
  $\mathbb{E}s' \ge \frac{11}{8}s$ for $0 \le s \le n/2$, and Hoeffding's inequality bounds the
  lower tail of $\sum_v Y_v$.
\end{proof}

\begin{theorem}[Symmetry breaking]\label{thm:symmetry_breaking}
  \lean{Plurality.gapUnits, Plurality.units_near, Plurality.units_far, Plurality.units_grow, Plurality.hitProb_mono, Plurality.majority3_symmetry_breaking}
  \leanok
  \uses{lem:jump_near_balance, lem:growth_far}
  There is $C > 0$ such that, for $\log n \ge 40$, from any configuration the gap reaches
  $22\sqrt{3n\log n}$ in absolute value within any $t \ge C\log n$ rounds with probability at
  least $1 - 1/n$.
\end{theorem}
\begin{proof}
  \leanok
  Apply the hitting-time bound of Doerr et al.\ (their Claim~2.9,
  \texttt{Dynamics.Kernel.drift\_hitting\_log}) to $X = \lfloor |s|/(\sqrt n/100)\rfloor$ with
  $q = n$, $c_1 = 5/4$, $c_2 = 1/320000$, $c_3 = 9/64$, $c_4 = 1000$, $c_6 = 1$; by the
  symmetry between the opinions assume $s \ge 0$. If $X = 0$, Lemma~\ref{lem:jump_near_balance}
  gives $X' \ge 20 \ge 1$ with probability $9/64$. If $1 \le X < 16$, the same lemma gives
  $X' \ge 20 \ge \frac54 X$ with probability $9/64 \ge 1 - e^{-X/320000}$. If
  $16 \le X < 1000\log n$, then $s \le n/2$, and Lemma~\ref{lem:growth_far} with
  $\lambda = X\sqrt n/3200$ gives $s' > \frac{21}{16}X\sqrt n/100$, hence $X' \ge \frac54 X$,
  except with probability $e^{-X^2/5120000} \le e^{-X/320000}$. Reaching $X \ge 1000\log n$
  means $|s| \ge 10\sqrt n\log n \ge 22\sqrt{3n\log n}$.
\end{proof}

\begin{lemma}[Hit, then absorb]\label{lem:hit_absorb}
  \lean{Plurality.hitProb_le_one, Plurality.hitProb_nonneg, Plurality.hitProb_sub_le_event, Plurality.event_amplify}
  \leanok
  Let $P$ be an absorbing event of a finite Markov chain. If from every state of $B$ the chain is
  in $P$ after $T$ steps with probability at least $1 - \varepsilon$, then from any state it is in
  $P$ at time $t + T$ with probability at least $\Pr(T_B \le t) - \varepsilon$. If moreover from
  every state the chain is in $P$ after $T_1$ steps with probability at least $1 - \varepsilon$,
  and from $a$ after $T_0$ steps with probability at least $1 - \delta$, then from $a$ it is in
  $P$ after $T_0 + T_1$ steps with probability at least $1 - \delta\varepsilon$.
\end{lemma}
\begin{proof}
  \leanok
  The first claim by induction on $t$, through the recursion of $\Pr(T_B \le t)$ and the
  monotonicity in time of an absorbing event; the second by the Markov property at time $T_0$.
\end{proof}

\begin{theorem}[Two opinions from any configuration]\label{thm:any_start}
  \lean{Plurality.step_compl, Plurality.consensus_absorbing, Plurality.event_consensus_of_gap, Plurality.majority3_any_start, Plurality.majority3_any_start_whp}
  \leanok
  \uses{thm:symmetry_breaking, cor:binary_vanishing, lem:hit_absorb}
  There is $C > 0$ such that, for $\log n \ge 40$, from any configuration (in particular from a
  perfectly balanced one) all nodes hold the same opinion after any $T \ge C\log n$ rounds with
  probability at least $1 - C\log n/n$; and, for $\log n \ge C$, with probability at least
  $1 - 1/n$.
\end{theorem}
\begin{proof}
  \leanok
  Theorem~\ref{thm:symmetry_breaking}, then Corollary~\ref{cor:binary_vanishing} applied to the
  set of nodes of opinion $1$ or to its complement, composed by Lemma~\ref{lem:hit_absorb};
  consensus on either opinion is absorbing, so further rounds keep it. For the second bound, two
  consecutive blocks each fail with probability at most $\varepsilon = C_0\log n/n$, and
  $\varepsilon^2 \le 1/n$.
\end{proof}

\section{Lower bounds}\label{sec:lower}

\subsection{Lower bound for 3-majority}

\begin{lemma}[Escaping a moving target]\label{lem:escape}
  \lean{Dynamics.expList_escape, Dynamics.expList_sum}
  \leanok
  If a round-based process starts in $G_0$ and, for every $t < T$, one round
  from any state of $G_t$ misses $G_{t+1}$ with probability at most $p$, then
  after $T$ rounds it lies outside $G_T$ with probability at most $Tp$.
\end{lemma}

\begin{proof}
  \leanok
  Induction on $T$, conditioning on the first round.
\end{proof}

\begin{lemma}[Lemma 4.1]\label{lem:4_1}
  \lean{Plurality.lemma_4_1, Plurality.mu_le_of_count_le, Plurality.sq_div_le_sum_sq}
  \leanok
  \uses{lem:2_1, lem:tails}
  Let $k\sqrt{n\log n} \le b \le n/k$. If $c_j \le n/k + b$, then
  $\Pr(C_{j,t+1} \ge n/k + (1 + 3/k)b \mid C_t = c) \le 1/n^2$.
\end{lemma}

\begin{proof}
  \leanok
  By Cauchy--Schwarz $\sum_h c_h^2 \ge n^2/k$, so
  $\mu_j \le c_j(1 + c_j/n - 1/k)$. Writing $c_j = n/k + a$ with $a \le b$,
  this is at most $n/k + (1+1/k)b + b^2/n \le n/k + (1+2/k)b$. Hoeffding's
  inequality with deviation $b/k \ge \sqrt{n\log n}$ gives $1/n^2$.
\end{proof}

\begin{theorem}[Theorem 4.2, lower bound for 3-majority]\label{thm:4_2}
  \lean{Plurality.theorem_4_2, Plurality.theorem_4_2_log, Plurality.color_stays_small, Plurality.balancedSet}
  \leanok
  \uses{lem:4_1, lem:escape}
  Let $k \ge 3$ and suppose every color starts with at most $n/k + b$ nodes,
  where $b \ge k\sqrt{n\log n}$. If $b(1 + 3/k)^T \le n/k$, in particular if
  $T \le \frac k3 \log\frac{n}{kb}$, then the coloring after $T$ rounds is
  monochromatic with probability at most $kT/n^2$.
\end{theorem}

\begin{proof}
  \leanok
  For each color apply Lemma~\ref{lem:escape} to the thresholds
  $n/k + b(1+3/k)^t$, with Lemma~\ref{lem:4_1} for each round. A
  monochromatic coloring has a color with $n > 2n/k$ nodes, above its final
  threshold; sum over the $k$ colors.
\end{proof}

\begin{remark}
  With $b = \max\{(n/k)^{1-\varepsilon}, k\sqrt{n\log n}\}$ this gives
  $\Omega(k\log(n/(kb)))$ rounds, which is $\Omega(k\log n)$ for
  $k \le n^{1/4-\delta}$. The paper states $\Omega(k\log n)$ for all
  $k \le (n/\log n)^{1/4}$, but at that end $n/k = k\sqrt{n\log n}$ and the
  argument gives nothing. The paper's proof also chooses the thresholds after
  seeing the trajectory; the union bound needs the fixed schedule used here.
\end{remark}

\subsection{Which 3-input rules solve plurality consensus}

\begin{definition}[Rules, $\Delta$, $\delta$ and solvers]\label{def:good_rules}
  \lean{Plurality.Rule, Plurality.Conservative, Plurality.ClearMajority, Plurality.majCount, Plurality.deltaCount, Plurality.UniformRule, Plurality.Solver}
  \leanok
  \uses{def:rule, def:kernel}
  A rule is \emph{conservative} if it returns one of its inputs. It has the
  \emph{clear-majority} property if it returns the majority color whenever two
  inputs agree. For colors $r \ne b$, $\Delta_r$ counts the orderings of
  $(r,r,b)$ on which the rule returns $r$. For distinct $r,g,b$, $\delta_c$
  counts the orderings of $(r,g,b)$ on which it returns $c$; the rule is
  \emph{uniform} if $\delta_r = 2$ for all distinct triples. It is an
  \emph{$(s,\varepsilon)$-solver} if from every coloring whose unique plurality
  color $m$ has bias at least $s$, consensus on $m$ is eventually reached with
  probability at least $1-\varepsilon$.
\end{definition}

\begin{lemma}[Consensus needs survival]\label{lem:alive}
  \lean{Plurality.mono_le_alive, Plurality.count_runWith_eq_zero}
  \leanok
  \uses{def:good_rules}
  Under a conservative rule a color without nodes never returns. Hence for all
  $T, T_0$, the probability of consensus on $m$ at time $T$ is at most the
  probability that $m$ has a node at time $T_0$.
\end{lemma}

\begin{theorem}[Theorem 4.8(a), clear majority]\label{thm:4_8_a}
  \lean{Plurality.theorem_4_8_a, Plurality.not_solver_of_drift, Plurality.expected_red_le, Plurality.expected_red_le_iter}
  \leanok
  \uses{def:good_rules, lem:adopt}
  Let $8 \mid n$. With only colors $r, b$ present and $x$ the fraction of $r$,
  a node adopts $r$ with probability $p(x)$, where
  $p(x) - x = x(1-x)\bigl(x(\Delta_r - 2) + (1-x)(2 - \Delta_b)\bigr)$. If
  $\Delta_r \le 2 \le \Delta_b$, the number of $r$ nodes is a supermartingale,
  so from $5n/8$ nodes of $r$ (bias $n/4$) consensus on $r$ has probability at
  most $5/8$ at every time: the rule is not an $(n/4, 1/4)$-solver.
  Consequently, for an $(n/4, 1/4)$-solver and every pair $r \ne b$, either
  $\Delta_r = \Delta_b = 3$ or $\Delta_r, \Delta_b \le 1$.
\end{theorem}

\begin{proof}
  \leanok
  $\E[X_T] \le X_0$ by induction on $T$, and $\Pr(X_T = n) \le \E[X_T]/n$.
\end{proof}

\begin{remark}
  The paper claims that an $(n/4,1/4)$-solver has the clear-majority
  property, i.e.\ that $\Delta_r, \Delta_b \le 1$ is also impossible. Its
  proof uses the supermartingale property, but checks it only when red is the
  majority; for $\Delta_r, \Delta_b \le 1$ it fails when red is the minority,
  because the process is pulled to the interior point
  $x^* = (2 - \Delta_b)/(4 - \Delta_r - \Delta_b)$. That case is not
  formalized.
\end{remark}

\begin{lemma}[Adopting $r$ among three colors]\label{lem:adopt3}
  \lean{Plurality.adopt_r_le, Plurality.avg_adopt_r_le, Plurality.expected_count_le, Plurality.jump_prob_le}
  \leanok
  \uses{def:good_rules, lem:adopt, lem:tails}
  Let the rule be conservative with the clear-majority property, let only
  $r, g, b$ be present and $\delta_r \le 1$. Then
  $p(r) \le 3x^2 - 2x^3 + \delta_r x_r x_g x_b \le x(3x - 2x^2 + (1-x)^2/4)$,
  which is at most $0.97x$ for $x \le 2/5$; and from $x \le 2/5$ one round
  ends above $2n/5$ with probability at most $e^{-9n/31250}$.
\end{lemma}

\begin{lemma}[$r$ dies out]\label{lem:extinct}
  \lean{Plurality.frozenStep, Plurality.potential, Plurality.frozen_potential_le, Plurality.frozen_stop_le, Plurality.foldl_frozen_eq, Plurality.alive_le}
  \leanok
  \uses{lem:adopt3, lem:escape}
  Under the hypotheses of Lemma~\ref{lem:adopt3}, starting with at most $2n/5$
  nodes of $r$, the probability that $r$ still has a node after $T$ rounds is
  at most $0.97^T c_r + T e^{-9n/31250}$.
\end{lemma}

\begin{proof}
  \leanok
  Consider the chain frozen once $r$ exceeds $2n/5$. On it,
  $\Phi = C_r\,\mathbb 1[C_r \le 2n/5]$ contracts by $0.97$ in every state,
  it is frozen within $T$ rounds with probability at most $Te^{-9n/31250}$,
  and it agrees with the real chain until it is frozen.
\end{proof}

\begin{theorem}[Theorem 4.8(b), uniform property]\label{thm:4_8_b}
  \lean{Plurality.theorem_4_8_b, Plurality.not_solver_of_nonuniform, Plurality.extinctRounds, Plurality.deltaCount_sum}
  \leanok
  \uses{lem:extinct, lem:alive}
  Let $60 \mid n$ and $\log n \ge 20$. A rule with the clear-majority property
  that is an $(n/20, 1/4)$-solver has the uniform property.
\end{theorem}

\begin{proof}
  \leanok
  If a triple is not uniform, some color $r$ of it has $\delta_r \le 1$,
  since the three counts sum to $6$. Start from $23n/60$ nodes of $r$, $n/3$
  of $g$ and $17n/60$ of $b$: $r$ is the unique plurality color with bias
  $n/20$. By Lemmas~\ref{lem:extinct} and~\ref{lem:alive}, with
  $T_0 = \lceil\log(8n)/\log(100/97)\rceil$, consensus on $r$ has probability
  at most $1/8 + 1/8$ at every time.
\end{proof}

\begin{remark}
  The paper's proof (Lemma 4.10) is a sketch that treats $\delta_r = 1$ with
  $(\delta_g,\delta_b) \in \{(3,2),(4,1)\}$ and asserts that $r$ ``will not
  converge'' by iterating a one-round estimate. The proof here covers every
  non-uniform triple and makes the iteration rigorous.
\end{remark}

\subsection{A lower bound for \texorpdfstring{$h$}{h}-plurality}

\begin{definition}[$h$-plurality]\label{def:hplur}
  \lean{Plurality.HSample, Plurality.HRound, Plurality.mult, Plurality.hplur, Plurality.hstep, Plurality.hrun}
  \leanok
  \uses{def:config}
  Every node samples $h$ nodes and a uniformly random ordering $\sigma$ of the
  $h$ sample positions, and adopts the color of the position of smallest
  $\sigma$-rank among those whose color is most frequent. All tied colors
  appear equally often, so each is chosen with the same probability.
\end{definition}

\begin{lemma}[Adoption under $h$-plurality]\label{lem:hadopt}
  \lean{Plurality.hplur_first, Plurality.avg_hadopt_le, Plurality.avg_eval_two, Plurality.avg_prod_iter}
  \leanok
  \uses{def:hplur}
  A node adopts $j$ with probability at most $x_j + \frac{h^2}{2}x_j^2$.
\end{lemma}

\begin{proof}
  \leanok
  If $j$ wins and is seen exactly once, all $h$ sampled colors are distinct
  and $j$ is at the position ranked first, which has probability $x_j$.
  Otherwise $j$ is seen at least twice, and
  $\Pr(N_j \ge 2) \le \E[N_j(N_j-1)]/2 = \binom h2 x_j^2$.
\end{proof}

\begin{lemma}[Lemma 4.11]\label{lem:4_11}
  \lean{Plurality.lemma_4_11, Plurality.lemma_4_11_paper, Plurality.exp_bound_le}
  \leanok
  \uses{lem:hadopt, lem:tails}
  If $c_j \le a \le 2n/k$, then
  $\Pr(C_{j,t+1} \ge a(1 + 2h^2/k)) \le \exp(-2(ah^2/k)^2/n)$, which is at
  most $e^{-2nh^4/k^4}$ when $a \ge n/k$.
\end{lemma}

\begin{remark}
  The paper states the factor $1 + h^2/k$ but its proof gives
  $1 + 2h^2/k$, which is also what Theorem~\ref{thm:4_12} uses. Allowing
  $c_j \le a$ instead of $c_j = a$ covers colors below $n/k$, which the
  paper's proof of Theorem~\ref{thm:4_12} needs but the lemma excludes.
\end{remark}

\begin{theorem}[Theorem 4.12, lower bound for $h$-plurality]\label{thm:4_12}
  \lean{Plurality.theorem_4_12, Plurality.theorem_4_12_log, Plurality.hBalancedSet}
  \leanok
  \uses{lem:4_11, lem:escape}
  Let $k \ge 3$. If every color starts with at most $\frac32\cdot\frac nk$
  nodes and $T \le \frac{k}{2h^2}\log\frac43$, then after $T$ rounds of
  $h$-plurality the coloring is monochromatic with probability at most
  $kT e^{-2nh^4/k^4}$.
\end{theorem}

\begin{proof}
  \leanok
  As for Theorem~\ref{thm:4_2}, with thresholds
  $\frac32\cdot\frac nk(1 + 2h^2/k)^t \le 2n/k$.
\end{proof}


%%FILE rumor_spread/blueprint/src/content.tex
% Content of the blueprint for the rumor-spreading (uniform push) formalization.
% Mirrors latex/rumor_push.tex lemma-for-lemma, but every statement here also
% carries \lean{...} / \leanok / \uses{...} tags so that leanblueprint can
% build the dependency graph and check coverage against RumorSpread/*.lean.

\begin{abstract}
\noindent
We give a complete and elementary proof that the uniform \emph{push} protocol
on the complete graph $K_n$, started from a single informed node, informs all
$n$ nodes within $O(\log n)$ rounds with high probability. The probability
space is finite and uniform throughout; the only analytic ingredients are
Bernoulli's inequality and Markov's inequality, and the only concentration
tool is a one-line exponential-moment induction for adaptive Bernoulli
trials. No measure theory, no PMF/\texttt{ENNReal}, no Chernoff bounds, and
no martingales are used. This argument has been \emph{fully formalized} (no
\texttt{sorry}) in Lean~4 + Mathlib; the main theorem is
\texttt{RumorPush.push\_informs\_all\_whp}.
With the same random calls, the \emph{pull} and \emph{push--pull} protocols
of Karp, Schindelhauer, Shenker and V\"ocking also inform all nodes within
$O(\log n)$ rounds with high probability (Section~\ref{sec:pull}); for pull,
a second-moment bound replaces the reverse Markov inequality.
\end{abstract}

\section{The model}\label{sec:model}

Fix $n \ge 2$ and let $V = \{0,1,\dots,n-1\}$ be the vertex set of the
complete graph $K_n$.

\begin{definition}[Round randomness]\label{def:round}
  \lean{RumorPush.Tgt}
  \leanok
  A \emph{round configuration} is a function
  $\omega \colon V \to V$ with $\omega(v) \ne v$ for all $v$.
  Let $R_n$ denote the (finite) set of round configurations; it carries the
  uniform probability measure, i.e.\ every node picks a uniformly random
  \emph{target} among the other $n-1$ nodes, independently of all other
  nodes.
\end{definition}

\begin{definition}[Push step]\label{def:step}
  \lean{RumorPush.step}
  \leanok
  \uses{def:round}
  For an \emph{informed set} $I \subseteq V$ and a round configuration
  $\omega \in R_n$, the next informed set is
  \[
    \mathrm{step}(I,\omega) \;=\; I \,\cup\, \omega(I)
    \;=\; I \cup \{\omega(v) : v \in I\}.
  \]
  (Only informed nodes actually transmit; letting \emph{all} nodes draw a
  target and ignoring the choices of uninformed nodes is an equivalent and
  formalization-friendlier convention.)
\end{definition}

\begin{definition}[Process]\label{def:process}
  \lean{RumorPush.run}
  \leanok
  \uses{def:step}
  Fix a horizon $T \in \N$ and a start vertex $v_0 \in V$. The sample space
  is $\Omega = R_n^{\,T}$ (uniform measure, i.e.\ i.i.d.\ uniform rounds).
  For $\omega = (\omega_0,\dots,\omega_{T-1}) \in \Omega$ define
  $I_0(\omega) = \{v_0\}$ and
  $I_{t+1}(\omega) = \mathrm{step}(I_t(\omega), \omega_t)$ for $0 \le t < T$.
  We write $m_t = |I_t|$ and $U_t = n - m_t$ for the numbers of informed and
  uninformed nodes. Since $\Omega$ is finite, all probabilities are finite
  ratios and all expectations are finite sums; ``conditioning on $I_t$''
  means partitioning $\Omega$ according to the first $t$ coordinates. In the
  formalization, a trajectory is simply a \texttt{List} of round
  configurations and $\mathrm{run}\ I\ l$ folds \texttt{step} over the list
  $l$.
\end{definition}

Two trivial but load-bearing facts, both immediate from
Definition~\ref{def:step}:

\begin{lemma}[Monotonicity]\label{lem:mono}
  \lean{RumorPush.subset_step, RumorPush.card_le_card_step,
    RumorPush.subset_run}
  \leanok
  \uses{def:process}
  $I_t(\omega) \subseteq I_{t+1}(\omega)$ for all $t$ and all $\omega$;
  hence $m_t$ is non-decreasing and $U_t$ is non-increasing in $t$,
  pointwise on $\Omega$.
\end{lemma}

\begin{lemma}[At most doubling]\label{lem:atmost}
  \lean{RumorPush.card_step_le}
  \leanok
  \uses{def:step}
  $|I_{t+1}(\omega) \smallsetminus I_t(\omega)| \le |I_t(\omega)|$; equivalently
  $(\mathrm{step}(I,\omega)).\mathrm{card} \le 2\,I.\mathrm{card}$.
\end{lemma}

\section{Main theorem}\label{sec:main-statement}

Throughout, $\log$ denotes the natural logarithm.

\begin{theorem}[Push protocol informs $K_n$ in $O(\log n)$ rounds w.h.p.]
  \label{thm:main}
  \lean{RumorPush.push_informs_all_whp, RumorPush.push_informs_all_whp'}
  \leanok
  \uses{lem:gluing, lem:constants, lem:phase2cor}
  Let $n \ge 2$, let $v_0 \in V$, and set
  \[
    T_1 \;=\; \ceil{117 \log n} + 23, \qquad
    T_2 \;=\; \ceil{6 \log n}, \qquad
    T \;=\; T_1 + T_2 .
  \]
  Then $\Pr\bigl[\, I_T = V \,\bigr] \;\ge\; 1 - \frac{2}{n}$.
  In particular all nodes are informed after $O(\log n)$ rounds with high
  probability.
\end{theorem}

\begin{remark}
  The constants are deliberately crude: the proof is organized to minimize
  the analytic toolkit, not the round count. The sharp answer is
  $\log_2 n + \log n + o(\log n)$ (Frieze--Grimmett \cite{FG85}, Pittel
  \cite{Pittel87}). Replacing $2/n$ by $n^{-\alpha}$ for any fixed
  $\alpha \ge 1$ only changes the constants in $T_1, T_2$.
\end{remark}

The proof has two phases, glued by Lemma~\ref{lem:mono}:
\begin{itemize}
\item \textbf{Growth phase} (Section~\ref{sec:growth}): while $m_t \le n/2$,
  each round multiplies $m_t$ by $\ge 9/8$ with probability $\ge 1/8$;
  an exponential-moment bound for adaptive trials shows that after $T_1$
  rounds, $m_{T_1} > n/2$ except with probability $\le 1/n$.
\item \textbf{Saturation phase} (Section~\ref{sec:saturation}): once
  $m_t \ge n/2$, the expected number of uninformed nodes contracts by a
  factor $2/3$ per round; after $T_2$ further rounds Markov's inequality
  gives $U_T = 0$ except with probability $\le 1/n$.
\end{itemize}

\section{Finite uniform probability}\label{sec:prelim}

All randomness in this development is uniform over finite types, so instead
of measure theory or \texttt{PMF} we work with plain finite sums. The
definitions and their basic laws are those of the shared \texttt{dynamics}
library (\texttt{Dynamics.Uniform}), which also identifies them with Mathlib's
finite expectation \texttt{Finset.expect} and with the independent product of
uniform distributions (\texttt{Dynamics.Bridge}).

\begin{definition}[Average]\label{def:avg}
  \lean{Dynamics.avg, Dynamics.avg_eq_expect}
  \leanok
  For a finite type $\alpha$ and $f \colon \alpha \to \R$,
  $\mathrm{avg}(f) := \bigl(\sum_{a} f(a)\bigr) / |\alpha|$ is the
  expectation of $f$ under the uniform distribution on $\alpha$.
\end{definition}

\begin{definition}[Trajectory expectation]\label{def:expList}
  \lean{Dynamics.expList}
  \leanok
  \uses{def:avg}
  For a finite type $\alpha$, $\mathrm{expList}\ \alpha\ T\ F$ is the
  expectation of a functional $F$ of a trajectory of $T$ i.i.d.\ uniform
  draws from $\alpha$, defined \emph{by recursion on $T$}: the head draw is
  averaged out first, i.e.\ $\mathrm{expList}\ \alpha\ 0\ F = F(\,[]\,)$ and
  $\mathrm{expList}\ \alpha\ (T{+}1)\ F
    = \mathrm{avg}\bigl(a \mapsto \mathrm{expList}\ \alpha\ T\ (l \mapsto
    F(a :: l))\bigr)$.
  This makes ``conditioning on the first round'' a definitional unfolding
  rather than a measure-theoretic construction.
\end{definition}

\begin{lemma}[Basic properties]\label{lem:avg-basic}
  \lean{Dynamics.avg_le_avg, Dynamics.avg_add, Dynamics.avg_sub,
    Dynamics.avg_const_mul, Dynamics.avg_const,
    Dynamics.expList_le_expList, Dynamics.expList_add,
    Dynamics.expList_const_mul, Dynamics.expList_const}
  \leanok
  \uses{def:avg, def:expList}
  $\mathrm{avg}$ and $\mathrm{expList}$ are monotone, additive, and
  homogeneous, and constants pass through unchanged. These pointwise
  inequalities pushed through $\mathrm{avg\_le\_avg}$ /
  $\mathrm{expList\_le\_expList}$ are the only form in which Markov's
  inequality and union bounds are used below (Lemmas~\ref{lem:markov}
  and~\ref{lem:revmarkov}); they are not separately formalized lemmas.
\end{lemma}

\begin{theorem}[Faithfulness of the model]\label{thm:faithful}
  \lean{RumorPush.expList_eq_avg_ofFn}
  \leanok
  \uses{def:expList}
  $\mathrm{expList}\ \alpha\ T\ F$ equals the uniform average of $F$ over
  \emph{all} length-$T$ sequences of draws, i.e.\ over the product
  probability space $\mathrm{Fin}\ T \to \alpha$ of $T$ i.i.d.\ uniform
  rounds. This confirms that the recursive definition of
  Definition~\ref{def:expList} is the textbook object.
\end{theorem}

\section{Elementary inequalities}\label{sec:ineq}

We write $x := \tfrac{1}{n-1}$, so $0 < x \le 1$ for $n \ge 2$.

\begin{lemma}[Bernoulli]\label{lem:bernoulli}
  \lean{RumorPush.one_sub_mul_le_pow}
  \leanok
  For all real $y \ge -1$ (we only use $y\in[-1,1]$) and $m \in \N$:
  $(1+y)^m \ge 1 + my$. Equivalently, with $y = -x$:
  $1 - mx \le (1-x)^m$.
\end{lemma}

\begin{lemma}[Second-order lower bound, ``Bonferroni'']\label{lem:bonferroni}
  \lean{RumorPush.mul_sub_sq_le_one_sub_pow, RumorPush.half_mul_le_one_sub_pow}
  \leanok
  \uses{lem:bernoulli}
  For $0 \le x \le 1$ and $m \in \N$:
  \[
    1-(1-x)^m \;\ge\; mx - \binom{m}{2}x^2 \;\ge\; mx\Bigl(1-\frac{mx}{2}\Bigr),
  \]
  and consequently $mx/2 \le 1-(1-x)^m$ whenever $mx \le 1$.
\end{lemma}

\begin{lemma}[Upper bound without $e$]\label{lem:upper}
  \lean{RumorPush.pow_one_sub_le_one_div}
  \leanok
  \uses{lem:bernoulli}
  For $0 \le x \le 1$ and $m \in \N$: $(1-x)^m \le \dfrac{1}{1+mx}$.
\end{lemma}

\begin{lemma}[Logarithm bounds]\label{lem:logbounds}
  \lean{RumorPush.one_sub_inv_le_log}
  \leanok
  For $x > 0$: $1 - \tfrac1x \le \log x$. (The formalization only needs this
  left inequality, applied at $x = 9/8,\ 16/15,\ 3/2$, together with
  Mathlib's numeric bound $\log 2 < 0.6931471808$.)
\end{lemma}

\begin{lemma}[Markov's inequality, finite uniform version]\label{lem:markov}
  \lean{Dynamics.avg_le_avg}
  \leanok
  \uses{def:avg}
  Let $\Omega$ be a finite nonempty set with the uniform measure,
  $f \colon \Omega \to \R$ with $f \ge 0$, and $a > 0$. Then
  $\Pr[f \ge a] \le \E[f]/a$. In particular, if $f$ takes values in $\N$,
  then $\Pr[f \ne 0] \le \E[f]$.
\end{lemma}

\begin{lemma}[Reverse Markov for bounded variables]\label{lem:revmarkov}
  \lean{Dynamics.avg_le_avg}
  \leanok
  \uses{def:avg}
  Let $f \colon \Omega \to \R$ with $0 \le f \le b$ pointwise, $b > 0$, and
  let $0 < a < b$. Then $\Pr[f \ge a] \ge \dfrac{\E[f] - a}{b}$. (As with
  Lemma~\ref{lem:markov}, the formalization does not isolate this as a
  named lemma: each use site supplies its own pointwise bound and pushes it
  through $\mathrm{avg\_le\_avg}$, e.g.\ in the proof of
  Lemma~\ref{lem:goodround}.)
\end{lemma}

\section{One-round estimates}\label{sec:oneround}

Fix an informed set $I$ with $|I| = m \ge 1$ and let $\omega \in R_n$ be a
uniform round configuration. All statements in this section concern this
single round; by independence of the rounds they apply verbatim to the
conditional law of round $t$ given $I_t = I$.

\begin{lemma}[Contact probability]\label{lem:contact}
  \lean{RumorPush.avg_not_contacted}
  \leanok
  \uses{def:step, def:avg}
  For any fixed $u \notin I$,
  \[
    \Pr_\omega\bigl[u \in \mathrm{step}(I,\omega)\bigr]
    \;=\; 1 - \Bigl(1-\frac{1}{n-1}\Bigr)^{m}.
  \]
  The Lean proof is the only place where a probability is \emph{computed}
  rather than estimated: it counts round configurations avoiding $u$ on $I$
  via an explicit equivalence to a product of subtypes and
  \texttt{Fintype.card\_pi}.
\end{lemma}

\begin{lemma}[Expected growth]\label{lem:expgrowth}
  \lean{RumorPush.avg_card_step}
  \leanok
  \uses{lem:contact, lem:bonferroni}
  Let $N(\omega) = |\mathrm{step}(I,\omega)| - m$ be the number of newly
  informed nodes. The formalization gives the \emph{exact} formula
  \[
    \E_\omega[|\mathrm{step}(I,\omega)|] \;=\;
      m + (n-m)\bigl(1-(1-x)^m\bigr),
  \]
  from which, if $1 \le m \le n/2$: $\E_\omega[N] \ge m/4$ (using
  Lemma~\ref{lem:bonferroni} and $n - m \ge n/2$, $mx \le 1$).
\end{lemma}

\begin{lemma}[A good round has constant probability]\label{lem:goodround}
  \lean{RumorPush.prob_goodRound}
  \leanok
  \uses{lem:expgrowth, lem:atmost, lem:revmarkov}
  If $1 \le m \le n/2$, then
  $\Pr_\omega\bigl[\,|\mathrm{step}(I,\omega)| \ge \tfrac{9}{8}\,m\,\bigr]
    \ge \tfrac18$.
\end{lemma}

\begin{lemma}[Contraction above half]\label{lem:contract}
  \lean{RumorPush.avg_uninformed_le}
  \leanok
  \uses{lem:contact, lem:upper}
  If $m \ge n/2$, then for every fixed $u \notin I$,
  $\Pr_\omega[u \notin \mathrm{step}(I,\omega)] = (1-x)^m \le \tfrac23$, and
  consequently, with $U' := n - |\mathrm{step}(I,\omega)|$,
  $\E_\omega[U'] \le \tfrac23(n-m)$.
\end{lemma}

\section{Adaptive Bernoulli trials}\label{sec:adaptive}

\begin{definition}[Good rounds]\label{def:good}
  \lean{RumorPush.goodRound, RumorPush.goodCount}
  \leanok
  \uses{def:step}
  A round is \emph{good} if it multiplies the informed set by $\ge 9/8$, or
  if the informed set already exceeds $n/2$:
  \[
    X_t \;=\; \mathbf{1}\Bigl[\, m_{t+1} \ge \tfrac98\, m_t
    \;\;\text{or}\;\; m_t > \tfrac n2 \,\Bigr] \in \{0,1\}.
  \]
  $\mathrm{goodCount}$ counts the good rounds along a trajectory. By
  Lemma~\ref{lem:goodround} (applied when $m_t \le n/2$; trivial otherwise),
  $\Pr[X_t = 1 \mid \omega_0,\dots,\omega_{t-1}] \ge \tfrac18$ for every
  prefix, since $\omega_t$ is independent of the prefix and
  $\mathrm{expList}$ conditions on it definitionally.
\end{definition}

\begin{lemma}[Exponential moment for adaptive trials]\label{lem:adaptive}
  \lean{RumorPush.expList_half_pow_goodCount}
  \leanok
  \uses{def:expList, def:good, lem:goodround}
  For every $T$ and every nonempty $I$,
  \[
    \E\Bigl[\bigl(\tfrac12\bigr)^{\mathrm{goodCount}(I,\,\cdot\,)}\Bigr]
      \;\le\; \Bigl(\frac{15}{16}\Bigr)^{T}
  \]
  over $T$ rounds started at $I$. (This specializes the textbook
  ``$\E[s^G] \le (1-p(1-s))^T$ for adaptive $\{0,1\}$-trials with
  per-step success probability $\ge p$'' at $p = 1/8$, $s = 1/2$; the
  formalization proves the specialized statement directly by induction on
  $T$ rather than the general one, since only this instance is needed. The
  ensuing tail bound $\Pr[G < L] \le 2^{L}(15/16)^{T}$ — via
  $s^{L}\Pr[G<L] \le \E[s^G]$ — is likewise not isolated as a separate
  lemma; it is proved inline inside Lemma~\ref{lem:phase1}
  (\texttt{phase1}).
\end{lemma}

\section{Growth phase}\label{sec:growth}

\begin{lemma}[Enough good rounds force half-saturation]\label{lem:force}
  \lean{RumorPush.pow_goodCount_mul_card_le_card_run}
  \leanok
  \uses{def:good, lem:mono}
  Let $L \in \N$ satisfy $(9/8)^L > n/2$. If the informed set never exceeds
  $n/2$ up to time $T_1$, then along that trajectory
  $m_{T_1} \ge (9/8)^{\mathrm{goodCount}} \cdot m_0$: every good round
  multiplies $m_t$ by $\ge 9/8$ and every round is non-decreasing
  (Lemma~\ref{lem:mono}), so $\mathrm{goodCount} \ge L$ would force
  $m_{T_1} > n/2$, a contradiction — hence
  $\mathrm{goodCount} < L \Rightarrow m_{T_1} \le n/2$ is impossible to rule
  out only via the deterministic growth bound stated here.
\end{lemma}

\begin{lemma}[Phase 1]\label{lem:phase1}
  \lean{RumorPush.phase1}
  \leanok
  \uses{lem:force, lem:adaptive}
  If $(9/8)^L \ge n$, then after $T_1$ rounds the informed set is still
  $\le n/2$ with probability at most $2^{L}\,(15/16)^{T_1}$.
\end{lemma}

\begin{lemma}[Admissible constants]\label{lem:constants}
  \lean{RumorPush.numeric_A, RumorPush.numeric_B}
  \leanok
  \uses{lem:phase1, lem:logbounds}
  With $L = \ceil{9 \log n} + 1$ we have $(9/8)^L \ge n$
  (\texttt{numeric\_A}), and with $T_1 = \ceil{117 \log n} + 23$ we have
  $2^{L}(15/16)^{T_1} \le 1/n$ (\texttt{numeric\_B}), using
  $\log\tfrac98 \ge \tfrac19$, $\log\tfrac{16}{15} \ge \tfrac1{16}$ and
  $\log 2 < 0.6932$. Combined with Lemma~\ref{lem:phase1}:
  $\Pr[m_{T_1} \le n/2] \le 1/n$. The clean statement to formalize is that
  \emph{any} $L, T_1$ satisfying these two numeric inequalities work; the
  values above merely exhibit admissible ones.
\end{lemma}

\section{Saturation phase}\label{sec:saturation}

\begin{lemma}[Phase 2]\label{lem:phase2}
  \lean{RumorPush.saturation}
  \leanok
  \uses{lem:contract}
  Let $A = \{m_{T_1} > n/2\}$ (an event determined by the first $T_1$
  coordinates). Then for every $s \ge 0$,
  $\E\bigl[\, U_{T_1 + s}\, \mathbf{1}_A \,\bigr] \le (2/3)^{s}\, n$, by
  induction on $s$ using Lemma~\ref{lem:contract} at each step (monotonicity
  keeps $m_{T_1+s} > n/2$ on $A$).
\end{lemma}

\begin{corollary}\label{lem:phase2cor}
  \lean{RumorPush.numeric_C}
  \leanok
  \uses{lem:phase2, lem:markov, lem:logbounds}
  With $T_2 = \ceil{6 \log n}$: $\Pr[\, U_T \ge 1 \text{ and } A \,]
    \le n\,(2/3)^{T_2} \le 1/n$. ($U_T$ is $\N$-valued, so Markov applied to
  $U_T \mathbf{1}_A$ and Lemma~\ref{lem:phase2} give the first inequality;
  $(2/3)^{T_2} n \le 1/n$ — \texttt{numeric\_C} — follows from
  $\log\tfrac32 \ge \tfrac13$ and $T_2 \ge 6\log n$.)
\end{corollary}

\section{Proof of the main theorem}\label{sec:proof-main}

\begin{lemma}[Gluing the two phases]\label{lem:gluing}
  \lean{RumorPush.prNotAllInformed, RumorPush.prNotAllInformed_le}
  \leanok
  \uses{lem:phase1, lem:phase2, lem:mono}
  Let $\mathrm{prNotAllInformed}(n,v_0,T)$ be the probability that some node
  is still uninformed after $T$ rounds started at $v_0$. For all
  $L, T_1, T_2$ with $(9/8)^L \le n$:
  \[
    \mathrm{prNotAllInformed}(n,v_0,T_1+T_2) \;\le\;
      2^{L}\Bigl(\frac{15}{16}\Bigr)^{T_1} + \Bigl(\frac23\Bigr)^{T_2} n.
  \]
  Failure requires either phase 1 to fail (informed set still $\le n/2$
  after $T_1$ rounds, Lemma~\ref{lem:phase1}) or phase 2 to leave an
  uninformed node on the complementary event, controlled by
  Lemma~\ref{lem:phase2} via Markov.
\end{lemma}

With $T_1 = \ceil{117\log n} + 23$, $T_2 = \ceil{6\log n}$, $T = T_1+T_2$,
Theorem~\ref{thm:main} follows from Lemma~\ref{lem:gluing},
Lemma~\ref{lem:constants} and Corollary~\ref{lem:phase2cor}:
\[
  \Pr[I_T \ne V] \;=\; \Pr[U_T \ge 1]
  \;\le\; \Pr[\,\overline{A}\,] + \Pr[\,U_T \ge 1 \text{ and } A\,]
  \;\le\; \frac1n + \frac1n \;=\; \frac2n .
\]

\section{PULL and PUSH--PULL}\label{sec:pull}

We now treat the two other protocols of Karp, Schindelhauer, Shenker and
V\"ocking \cite{KSSV00}, driven by the \emph{same} random calls: in a round
configuration $\omega \in R_n$ every node $v$ calls $\omega(v)$. In PULL the
caller learns the rumor from its callee; in PUSH--PULL the rumor travels in
both directions along every call. As for PUSH, partners are chosen among the
other $n-1$ nodes (KSSV allow self-calls; only constants change), and the
round count is a crude $O(\log n)$ instead of the sharp
$\log_3 n + O(\log\log n)$ of \cite[Thm.~2.1]{KSSV00}.

\begin{definition}[PULL and PUSH--PULL rounds]\label{def:pull}
  \lean{RumorPush.pullStep, RumorPush.pullRun, RumorPush.pushPullStep,
    RumorPush.pushPullRun, RumorPush.prPullNotAllInformed,
    RumorPush.prPushPullNotAllInformed}
  \leanok
  \uses{def:process}
  $\mathrm{pullStep}(I,\omega) = I \cup \{v : \omega(v) \in I\}$ and
  $\mathrm{pushPullStep}(I,\omega) = \mathrm{step}(I,\omega) \cup
  \mathrm{pullStep}(I,\omega)$; trajectories fold these steps over a list of
  i.i.d.\ uniform rounds, and the failure probabilities are the
  probabilities that some node is still uninformed after $T$ rounds started
  from $\{v_0\}$.
\end{definition}

\begin{lemma}[Pathwise domination]\label{lem:pull-domination}
  \lean{RumorPush.run_subset_pushPullRun, RumorPush.pullRun_subset_pushPullRun}
  \leanok
  \uses{def:pull, lem:mono}
  With the same calls, the PUSH--PULL informed set contains the PUSH and the
  PULL informed sets after every number of rounds.
\end{lemma}

\begin{lemma}[Independent calls]\label{lem:pull-indep}
  \lean{RumorPush.avg_pi_prod, RumorPush.avg_tgt_mul_prod,
    RumorPush.avg_call_mem}
  \leanok
  \uses{def:round, def:avg}
  The calls $\omega(v)$, $v \in V$, of a uniform round are independent, each
  uniform among the other $n-1$ nodes; in particular a caller $v \notin I$
  calls into $I$ with probability $|I|/(n-1)$.
\end{lemma}

\begin{lemma}[One PULL or PUSH--PULL round, exactly]\label{lem:pull-oneround}
  \lean{RumorPush.pull_prob_stuck_singleton, RumorPush.pull_avg_card_step,
    RumorPush.pull_avg_uninformed, RumorPush.pull_avg_uninformed_le,
    RumorPush.pushPull_avg_uninformed}
  \leanok
  \uses{lem:pull-indep, lem:contact}
  Let $m = |I|$ and $u = n - m$. From a single informed node a PULL round
  informs nobody with probability $(1 - \frac1{n-1})^{n-1}$. After one PULL
  round, $\E|I'| = m + u\,\frac{m}{n-1}$ and
  $\E[n - |I'|] = \frac{u(u-1)}{n-1} \le \frac{u^2}{n}$ (quadratic
  shrinking). After one PUSH--PULL round,
  $\E[n - |I'|] = \frac{u(u-1)}{n-1}\bigl(1 - \frac1{n-1}\bigr)^{m}$.
\end{lemma}

\begin{lemma}[The two-phase argument for any inflationary step]
  \label{lem:pull-phases}
  \lean{RumorPush.goodRoundOf, RumorPush.goodCountOf,
    RumorPush.expList_half_pow_goodCountOf, RumorPush.phase1Of,
    RumorPush.saturationOf, RumorPush.notAllOf_le,
    RumorPush.notAllOf_le_two_div}
  \leanok
  \uses{lem:adaptive, lem:phase1, lem:phase2, lem:gluing, lem:constants,
    lem:phase2cor}
  Let $f(I,a) \supseteq I$ be a spreading step driven by i.i.d.\ uniform
  rounds $a$. Suppose that from every nonempty $I$ a round is good
  ($|f(I,a)| \ge \frac98 |I|$, or $|I| > n/2$) with probability at least
  $1/8$, and that for $|I| \ge n/2$ the expected number of uninformed nodes
  contracts by $2/3$. Then after $\ceil{160 \log n}$ rounds from $\{v_0\}$,
  some node is uninformed with probability at most $2/n$. The proof is the
  PUSH proof of Sections~\ref{sec:adaptive}--\ref{sec:proof-main}, which uses
  no other property of $\mathrm{step}$.
\end{lemma}

\begin{lemma}[The two inputs for PULL]\label{lem:pull-good}
  \lean{RumorPush.pull_avg_new_sq_le, RumorPush.pull_prob_goodRound,
    RumorPush.pull_avg_uninformed_contract}
  \leanok
  \uses{lem:pull-indep, lem:pull-oneround}
  A PULL round is good with probability at least $1/8$, and for
  $|I| \ge n/2$ it contracts the expected uninformed count by $2/3$.
\end{lemma}

\begin{proof}
  The contraction follows from $\frac{u(u-1)}{n-1} \le \frac23 u$ for
  $u \le n/2$. For good rounds, let $X$ be the number of uninformed callers
  that call into $I$, so $|I'| = |I| + X$. By Lemma~\ref{lem:pull-indep},
  $\mu = \E X = \frac{(n-m)m}{n-1} \ge \max(m/2, 1/2)$ for $m \le n/2$ and,
  the calls being pairwise independent, $\E X^2 \le \mu + \mu^2$. A round
  that is not good has $X < m/8 \le \mu/4$, so pointwise
  $X \le \mu/4 + X^2/(2t) + (t/2)\,\mathbf 1[\text{good}]$ with
  $t = 4(1+\mu)/3$; taking expectations gives
  $\Pr[\text{good}] \ge \frac{3\mu}{4t} \ge 3/16$ (a Paley--Zygmund bound).
  Unlike PUSH, the reverse Markov argument of Lemma~\ref{lem:goodround} does
  not apply, since a PULL round can inform many more than $|I|$ new nodes.
\end{proof}

\begin{theorem}[PULL and PUSH--PULL inform $K_n$ in $O(\log n)$ rounds w.h.p.]
  \label{thm:pull}
  \lean{RumorPush.pull_informs_all_whp, RumorPush.pull_informs_all_whp',
    RumorPush.pushPull_informs_all_whp, RumorPush.pushPull_informs_all_whp'}
  \leanok
  \uses{lem:pull-phases, lem:pull-good, lem:pull-domination}
  Let $n \ge 2$ and $v_0 \in V$. After $\ceil{160 \log n}$ rounds of PULL,
  or of PUSH--PULL, started from $\{v_0\}$, all nodes are informed with
  probability at least $1 - 2/n$.
\end{theorem}

\begin{proof}
  For PULL, apply Lemma~\ref{lem:pull-phases} with the inputs of
  Lemma~\ref{lem:pull-good}. For PUSH--PULL, the failure event is contained
  in the PULL failure event by Lemma~\ref{lem:pull-domination}.
\end{proof}

\section{The Lean formalization}\label{sec:lean}

The argument above is fully formalized in the accompanying repository
(library \texttt{RumorSpread}, namespace \texttt{RumorPush}); the build
contains no \texttt{sorry} and the main theorem depends only on the three
standard axioms (\texttt{propext}, \texttt{Classical.choice},
\texttt{Quot.sound}). File structure (dependency order):
\texttt{Bounds.lean} (Section~\ref{sec:ineq}), \texttt{Dynamics.Uniform}
of the shared \texttt{dynamics} library (Section~\ref{sec:prelim}), \texttt{Model.lean}
(Section~\ref{sec:model}), \texttt{OneRound.lean}
(Section~\ref{sec:oneround}), \texttt{Growth.lean}
(Sections~\ref{sec:adaptive}--\ref{sec:growth}), \texttt{Saturation.lean}
(Section~\ref{sec:saturation}), \texttt{Main.lean} (main theorem and
numeric lemmas), \texttt{Equivalence.lean} (Theorem~\ref{thm:faithful}).
PULL and PUSH--PULL (Section~\ref{sec:pull}): \texttt{PullFoldl.lean} and
\texttt{PullModel.lean} (Definition~\ref{def:pull},
Lemma~\ref{lem:pull-domination}), \texttt{PullIndep.lean}
(Lemma~\ref{lem:pull-indep}), \texttt{PullOneRound.lean}
(Lemma~\ref{lem:pull-oneround}), \texttt{PullPhases.lean}
(Lemma~\ref{lem:pull-phases}), \texttt{PullGood.lean}
(Lemma~\ref{lem:pull-good}), \texttt{PullMain.lean} and
\texttt{PullPushPull.lean} (Theorem~\ref{thm:pull}).

\begin{thebibliography}{9}
\bibitem{FG85} A.~Frieze, G.~Grimmett.
\emph{The shortest-path problem for graphs with random arc-lengths}.
Discrete Applied Mathematics 10 (1985) 57--77.
\bibitem{Pittel87} B.~Pittel.
\emph{On spreading a rumor}. SIAM J.\ Appl.\ Math.\ 47 (1987) 213--223.
\bibitem{KSSV00} R.~Karp, C.~Schindelhauer, S.~Shenker, B.~V\"ocking.
\emph{Randomized rumor spreading}. FOCS 2000.
\end{thebibliography}


%%FILE undecided/blueprint/src/content.tex
\begin{abstract}
We formalize the basic properties of the synchronous undecided-state dynamics with two opinions
on the complete graph: the exact one-round expectations of the three counts, the resulting
growth of the bias by the factor $1+q/n$ (the undecided nodes amplify the current majority), and
almost-sure absorption in a monochromatic configuration.
\end{abstract}

\section{The model}
\begin{definition}[Undecided-state dynamics]\label{def:usd}
\lean{Undecided.Op,Undecided.update,Undecided.Config,Undecided.step,Undecided.count,Undecided.kernel}\leanok
Each of $n$ nodes holds opinion $a$, opinion $b$, or is undecided. In a round every node samples a
node uniformly at random (with replacement, possibly itself): an undecided node adopts the sampled
opinion, and a decided node that samples the other opinion becomes undecided. We write $a$, $b$,
$q$ for the numbers of $a$-, $b$- and undecided nodes.
\end{definition}

\section{One round in expectation}
\begin{lemma}[Per-node probabilities]\label{lem:node}
\lean{Undecided.count_eq_sum,Undecided.count_step_eq_sum,Undecided.avg_indicator_eq,Undecided.node_prob_a,Undecided.node_prob_b,Undecided.node_prob_u}\leanok\uses{def:usd}
The state of a node after a round depends only on the node it samples, which is uniform; so an
$a$-node stays $a$ with probability $1-b/n$ and an undecided node becomes $a$ with probability $a/n$.
\end{lemma}
\begin{proof}\leanok
The average over all sampling functions of a function of one coordinate is the uniform average.
\end{proof}
\begin{theorem}[Expected counts and bias]\label{thm:expect}
\lean{Undecided.expected_count_a,Undecided.expected_count_b,Undecided.expected_count_u,Undecided.expected_bias}\leanok\uses{lem:node}
After one round, the expected numbers of $a$-, $b$- and undecided nodes are $a(n-b+q)/n$,
$b(n-a+q)/n$ and $(q^2+2ab)/n$; hence the expected bias is $(a-b)(1+q/n)$.
\end{theorem}
\begin{proof}\leanok
Sum Lemma~\ref{lem:node} over the nodes; $a(n-b+q)-b(n-a+q)=(a-b)(n+q)$.
\end{proof}

\section{Absorption}
\begin{theorem}[Almost-sure absorption]\label{thm:absorbed}
\lean{Undecided.Mono,Undecided.notMono,Undecided.step_of_mono,Undecided.notMono_step,Undecided.access_from_decided,Undecided.absorbed}\leanok\uses{def:usd}
Monochromatic configurations are fixed points, and from every configuration the probability of
not being monochromatic after $t$ rounds tends to $0$.
\end{theorem}
\begin{proof}\leanok
If some node $w$ is decided, the round in which everybody samples $w$, repeated twice, makes the
configuration monochromatic; it has positive probability. Conclude with the shared absorption
theorem.
\end{proof}

\section{The majority phase (UND-1)}\label{sec:cggnps-majority}
Becchetti, Clementi, Natale, Pasquale and Silvestri (\emph{Plurality consensus in the gossip
model}, SODA 2015, Theorem~11 with $k=2$) prove convergence to the majority in $O(\log n)$ rounds
w.h.p.\ from configurations without undecided nodes and with a bias of order $n$. The stronger
form below, from any configuration with bias $\Omega(\sqrt{n\log n})$, is Theorem~3.2 of Clementi,
Ghaffari, Gual\`a, Natale, Pasquale and Scornavacca (MFCS 2018). Write $s=a-b$ for the bias,
$\Lambda=\sqrt{n\log n}$, and call a round \emph{good} if the new counts of $a$-, $b$- and
undecided nodes are within $\Lambda$ of their expectations (from below for $a$, from above for
$b$, both ways for the undecided nodes).

\begin{lemma}[Bad rounds]\label{lem:badround}
\lean{Undecided.count_add,Undecided.tail_upper,Undecided.tail_lower,Undecided.bad_round,Undecided.four_exp_sqrt}\leanok\uses{thm:expect}
A round is bad with probability at most $4e^{-2\Lambda^2/n}=4/n^2$.
\end{lemma}
\begin{proof}\leanok
Each count after the round is a sum of independent indicators, one per node, depending only on
the node it samples; apply Hoeffding's inequality to the four tails.
\end{proof}

\begin{lemma}[Good rounds]\label{lem:goodround}
\lean{Undecided.bias_next,Undecided.undec_next,Undecided.contract,Undecided.first_det,Undecided.growth_det,Undecided.bridge_det,Undecided.fin_det}\leanok\uses{thm:expect}
After a good round, $s'>s(1+q/n)-2\Lambda$ and $q'>n/3-s^2/(2n)-\Lambda$. If $b+2q\le 2n/3$ and
$q\le n/3$, the potential $\Psi=12b+q$ satisfies $\mathbb E[\Psi']\le\frac56\Psi$, hence
$\Psi'<\frac56\Psi+13\Lambda$ after a good round.
\end{lemma}
\begin{proof}\leanok
$a(n-b+q)-b(n-a+q)=(a-b)(n+q)$; $(q^2+2ab)/n-n/3+s^2/(2n)=3(q-n/3)^2/(2n)$; and
$12b(b+2q)+q^2+2ab\le 8bn+qn/3+2bn$.
\end{proof}

\begin{theorem}[Phases]\label{thm:phases}
\lean{Undecided.missP,Undecided.missP_comp,Undecided.growth_stage,Undecided.bridge_stage,Undecided.fin_stage}\leanok\uses{lem:badround,lem:goodround}
Let $2000\Lambda\le n$ and $p=4/n^2$.
(Growth) From $s\ge 402\Lambda$, after $T+1$ rounds $s\ge\min(400\Lambda(201/200)^T,7n/10)$
except with probability $(T+1)p$.
(Bridge) From $s\ge 7n/10$, after $18$ rounds $\Psi\le n/3$ except with probability $18p$.
(Final) From $\Psi\le n/3$, not all nodes hold $a$ after $T$ rounds with probability at most
$(5/6)^T\Psi+Tp$.
\end{theorem}
\begin{proof}\leanok
Growth: the sets $\{s\ge g,\ (q\ge n/100\text{ or }s\ge g+n/20)\}$ are entered with $g=400\Lambda$
after one good round and $g$ grows by $201/200$ per good round (if $s\le 4n/5$ then $q'\ge n/100$);
union bound over rounds with moving targets. Bridge: while $s\ge 2n/3$, $\Psi$ shrinks by $9/10$
per good round and $s$ loses at most $2\Lambda$. Final: induction on $T$, splitting on whether
the first round stays in $\{\Psi\le n/3\}$; and $\Psi\ge 1$ unless all nodes hold $a$.
\end{proof}

\begin{theorem}[Majority phase]\label{thm:majority}
\lean{Undecided.majority_explicit,Undecided.majority_whp,Undecided.swapOp,Undecided.foldl_swap,Undecided.majority_whp_abs,Undecided.majority_whp_of_ratio}\leanok\uses{thm:phases}
Let $C=10^4$. For every $n$ with $\log n\ge C$ and every configuration (any number of undecided
nodes) with $|a-b|\ge C\sqrt{n\log n}$, all nodes hold the initial majority opinion after
$\lceil C\log n\rceil$ rounds with probability at least $1-2/n$. Without undecided nodes and with
$a\ge(1+\alpha)b$, the same holds with a constant depending on $\alpha$ (SODA 2015, Theorem~11
for two colours).
\end{theorem}
\begin{proof}\leanok
Run $T_1+1$ growth rounds with $T_1=\lceil\log n/\log(201/200)\rceil$, $18$ bridge rounds and the
remaining $T_3\ge 12\log n$ final rounds, so that $(5/6)^{T_3}\le n^{-2}$; the failure
probability is at most $\lceil C\log n\rceil p+n^{-2}\cdot n/3\le 2/n$. A $b$-majority follows by
exchanging the opinions. With $a\ge(1+\alpha)b$ and $q=0$ the bias is at least $\alpha n/(2+\alpha)$,
which exceeds $C\sqrt{n\log n}$ once $\log n\ge C(1+3(2+\alpha)/\alpha)$; extra rounds do not hurt
since the all-$a$ configuration is absorbing.
\end{proof}

\section{Many colours: plurality consensus (UND-3)}\label{sec:bcnps-plurality}
Becchetti, Clementi, Natale, Pasquale and Silvestri (\emph{Plurality consensus in the gossip
model}, SODA 2015) analyse the undecided-state dynamics with $k$ colours: a state is a colour or
the undecided state (Table~1), and with $c_i$ nodes of colour $i$ and $q$ undecided nodes one
round gives in expectation $\mu_i=c_i(c_i+2q)/n$ nodes of colour $i$ and
$\mu_q=(q^2+(n-q)^2-\sum_i c_i^2)/n$ undecided nodes (their (3)--(4)). The monochromatic distance
is $\mathrm{md}(c)=\sum_i(c_i/c_1)^2\in[1,k]$, $c_1$ the plurality. For $k=2$ the dynamics is the
binary one.

The range of $k$ in Theorem~11 of the paper needs a minor correction: its proof invokes Lemma~10,
which assumes $k=O((n/\log n)^{1/4})$, and the round-by-round argument closes for
$k\le\varepsilon(\alpha)(n/\log n)^{1/3}$. The formalized statement takes
$Ck\le(n/\log n)^{1/3}$ with $C=C(\alpha)$
(\href{https://github.com/formal-dynamics/leanamics/blob/main/undecided/FORMALIZATION_DIFFERENCES.md#the-dynamics-with-k-colours-plurality-roadmap-und-3}{details}).

\begin{theorem}[Model]\label{thm:pmodel}
\lean{Undecided.Plurality.update,Undecided.Plurality.step,Undecided.Plurality.md,Undecided.Plurality.expected_count_some,Undecided.Plurality.expected_count_none,Undecided.Plurality.md_le_card,Undecided.Plurality.one_le_md,Undecided.Plurality.foldl_two}\leanok
Equations (3)--(4) hold, $1\le\mathrm{md}(c)\le k$, and the two-colour dynamics is the binary
dynamics under undecided $\mapsto u$, colour $0\mapsto a$, colour $1\mapsto b$.
\end{theorem}
\begin{proof}\leanok
Linearity over the nodes, as in the binary case; the binary update is checked on the nine pairs
of states.
\end{proof}

\begin{lemma}[Good rounds]\label{lem:pgood}
\lean{Undecided.Plurality.tail_up,Undecided.Plurality.tail_down,Undecided.Plurality.bad_le,Undecided.Plurality.thr_step,Undecided.Plurality.inv_step,Undecided.Plurality.phi_step,Undecided.Plurality.fpsi_step,Undecided.Plurality.first_step}\leanok\uses{thm:pmodel}
Let $\ell=3\log n$ and call a round good if every count is within
$\sqrt{3\ell(\mu+2\ell)}$ of its mean $\mu$ (Bernstein); a round is bad with probability at most
$(k+4)/n^3$. Let $\theta=1/(1+\alpha)$, $\kappa=1-\theta$, $B=2\,\mathrm{md}(x_0)$, $S$ the nodes
of colours other than $m$, $c=c_m$, $\Phi=\mu_m$. If $n\ge 2\cdot10^{12}K^2k^3\log n$ with
$K=(1+\alpha)^2/\alpha$, a good round keeps the invariant
$c_i+12\sqrt{\ell c}\le\theta c$ ($i\ne m$), $S+12B\sqrt{\ell c}\le Bc$, multiplies $\Phi$ by
$1+\kappa/(8000B)$ unless $4S+q\le n/10$, and keeps $\{4S+q\le n/10\}$; the first round from
$q=0$, $(1+\alpha)c_i\le c_m$ enters the invariant with $\Phi\ge n/(2k^2)$.
\end{lemma}
\begin{proof}\leanok
For the ratios: if $X\le tc/2$ one round cannot double it; if fewer than $n/20$ nodes are
decided, $c$ grows by $16/9$, which absorbs the slack; otherwise $c\ge n/(20(1+B))$ and the
deterministic drift $X\mu_m-\mu_Xc\ge\frac t2\kappa c^3/n$ beats the deviations, which needs
$c^3\gtrsim n^2\log n$, i.e.\ $k\lesssim(n/\log n)^{1/3}$. For $\Phi$:
$\mu_m+2\mu_q=n+(n-c-2q)^2/n+2W/n$ with $W=\sum_{j\ne m}c_j(c-c_j)\ge\kappa cS$ (the plurality
drift, Lemma~1 of the paper). In the first round $q=0$, so the ratios square and
$\mu_S=(\mathrm{md}-1)\mu_m$.
\end{proof}

\begin{theorem}[Plurality consensus]\label{thm:plurality}
\lean{Undecided.Plurality.main_stage,Undecided.Plurality.fin_stage,Undecided.Plurality.plurality_explicit,Undecided.Plurality.plurality_whp}\leanok\uses{lem:pgood}
Let $\alpha>0$ and $C=10^5((1+\alpha)^2/\alpha)^2$. For every $n$ with $\log n\ge C$, every $k$
with $Ck\le(n/\log n)^{1/3}$ and every configuration without undecided nodes in which colour $m$
has at least $1+\alpha$ times as many nodes as every other colour, all nodes hold $m$ after
$\lceil C\,\mathrm{md}(x)\log n\rceil$ rounds with probability at least $1-6/n$ (SODA 2015,
Theorem~11; the range of $k$ is discussed in \texttt{FORMALIZATION\_DIFFERENCES.md}).
\end{theorem}
\begin{proof}\leanok
One round into the invariant, then $T_2=\lceil\log(4k^2)/\log(1+\kappa/(8000B))\rceil\le
96000K\,\mathrm{md}\log n+1$ rounds with moving targets $\Phi\ge\frac n{2k^2}(1+\kappa/(8000B))^t$
or $4S+q\le n/10$, which end in $\{4S+q\le n/10\}$ since $\Phi\le n$; then
$T_3=\lceil 8\log n\rceil$ rounds in which $4S+q$ contracts by $3/4$ in expectation. The failure
probability is at most $T(k+4)/n^3+(3/4)^{T_3}n/10\le 6/n$.
\end{proof}

\section{Many colours: the \texorpdfstring{$\Omega(\mathrm{md}(c))$}{Omega(md(c))} lower bound (UND-3)}\label{sec:bcnps-lower-bound}
The same paper (SODA 2015, Theorem~8) shows that, without any bias, the dynamics needs
$\Omega(\mathrm{md}(c))$ rounds. With $R(c)=\sum_i c_i/c_1$ and $\Lambda(c)=R(c)^2/\mathrm{md}(c)$:
after the first round most nodes are undecided (Lemma~3); their number $q$ then descends to about
$n/2$ within $O(\log n)$ rounds while no colour exceeds $\gamma n/\mathrm{md}$ (Lemma~6); from
there no colour grows faster than by a factor $1+O(1/\mathrm{md})$ per round (Lemma~7). All
constants are existential, ``w.h.p.'' is $1-C/n$, and $k$ ranges over $Ck\le(n/\log n)^{a}$ with
the paper's exponent $a$.

Lemmas~6 and~7 need minor corrections: the window $|Q-n/2|\le 2\gamma^2/\mathrm{md}(\bar c)$ should
read $2\gamma^2n/\mathrm{md}(\bar c)$; the proof of (17) uses the initial instead of the current
monochromatic distance; in Lemma~7 the growth factor is $1+(4\gamma^2+2\gamma+1)/\mathrm{md}$, the
lower bound (23) on $E[\Delta']$ reverses the Cauchy--Schwarz inequality, and the proof needs
$\gamma\ge1$ (any $\gamma>1/2$ would do) rather than an arbitrary $\gamma>0$; and both lemmas must
bound the largest colour, not only the initial plurality
(\href{https://github.com/formal-dynamics/leanamics/blob/main/undecided/FORMALIZATION_DIFFERENCES.md#the-lower-bound-with-k-colours-lowerbound-roadmap-und-3}{details}).

\begin{lemma}[Ratios and expectations]\label{lem:lb-basic}
\lean{Undecided.Plurality.ratioR,Undecided.Plurality.ratioLam,Undecided.Plurality.md_le_ratioR,Undecided.Plurality.ratioR_le_card,Undecided.Plurality.md_le_ratioLam,Undecided.Plurality.ratioLam_le_card,Undecided.Plurality.ratioR_of_und_zero,Undecided.Plurality.ratioLam_of_und_zero,Undecided.Plurality.mu_eq_half,Undecided.Plurality.muU_sub_half}\leanok\uses{thm:pmodel}
$\mathrm{md}\le R\le k$ and $\mathrm{md}\le\Lambda\le k$; without undecided nodes $R=n/c_1$ and
$\Lambda=n^2/\sum_ic_i^2$. With $\delta=q-n/2$: $\mu_i=(1+(2\delta+c_i)/n)\,c_i$ and
$E[Q'-n/2]=(2\delta^2-\sum_jc_j^2)/n$ (the paper's (20) and (19)).
\end{lemma}
\begin{proof}\leanok
$c_i/c_1\in[0,1]$ gives $(c_i/c_1)^2\le c_i/c_1\le1$; Cauchy--Schwarz gives $R^2\le k\,\mathrm{md}$.
The expectations are (3)--(4) rewritten with $q=n/2+\delta$.
\end{proof}

\begin{lemma}[First round; Lemma~3]\label{lem:lb-first}
\lean{Undecided.Plurality.first_round}\leanok\uses{lem:lb-basic,lem:pgood}
Let $Ck\le(n/\log n)^{1/2}$. From a configuration $x$ without undecided nodes, with plurality
colour $m$, after one round, with probability at least $1-C/n$: $c_m\ge n/(2R^2)$, every colour has
at most $2n/R^2$ nodes, and $n(1-2/\Lambda)\le q\le n(1-1/(2\Lambda))$ ($R$, $\Lambda$ of $x$).
\end{lemma}
\begin{proof}\leanok
With $q=0$, $\mu_i=c_i^2/n\le\mu_m=n/R^2$ and $\mu_q=n-\sum_ic_i^2/n=n(1-1/\Lambda)$; in a good
round (Lemma~\ref{lem:pgood}, deviation level $3\log n$) the deviations are small compared with
$\mu_m\ge n/k^2\ge C^2\log n$ and with $n/\Lambda\ge n/k$. Proved with $C=100$.
\end{proof}

\begin{lemma}[Descent of the undecided nodes; Lemma~6]\label{lem:lb-descent}
\lean{Undecided.Plurality.undecided_square,Undecided.Plurality.undecided_not_below,Undecided.Plurality.descent}\leanok\uses{lem:lb-basic,lem:pgood}
Let $Ck\le(n/\log n)^{1/6}$. If $q=(1+\delta)n/2$ with $1-\delta\ge1/(2k)$, then after one round
$Q'\le(1+\delta^2)n/2$ with probability at least $1-C/n^2$ (16). If every colour is at most
$\gamma n/D$, with $\gamma\ge1$ and $0<D\le k$, then $Q'\ge n/2-2\gamma^2n/D$ with probability at
least $1-C/n^2$ (17, corrected). Hence there are $\gamma\ge1$ and $C$ such that, for $x$ without
undecided nodes with $\mathrm{md}(x)\ge C$ and $y$ satisfying the conclusion of
Lemma~\ref{lem:lb-first}, there is a round $t\le C\log n$ such that, with probability at least
$1-C/n$, every colour is at most $\gamma n/\mathrm{md}(x)$ at each round $s\le t$, and at round $t$
also $|q-n/2|\le2\gamma^2n/\mathrm{md}(x)$.
\end{lemma}
\begin{proof}\leanok
(16): $\mu_q=(1+\delta^2)n/2-\sum_jc_j^2/n$ and $\sum_jc_j^2\ge(n-q)^2/k$. (17):
$\mu_q\ge n/2-\sum_jc_j^2/n$ and $\sum_jc_j^2\le(\max_jc_j)(n-q)\le\gamma n^2/D$. The round $t$ is
deterministic: the first $s$ with $(1-1/\Lambda)^{2^s}\le4\gamma^2/\mathrm{md}$. The colours follow
an explicit envelope (at most $8n/D$, then at most $\gamma n/D$ with $\gamma=24$), iterated with
moving targets. Proved with $\gamma=24$ and $C=10^6$.
\end{proof}

\begin{lemma}[Plateau; Lemma~7]\label{lem:lb-plateau}
\lean{Undecided.Plurality.plateau_step,Undecided.Plurality.plateau}\leanok\uses{lem:lb-basic,lem:pgood}
Let $\gamma\ge1$ and $Ck\le(n/\log n)^{1/4}$. If $C\le D\le k$, $|q-n/2|\le2\gamma^2n/D$ and every
colour is at most $\gamma n/D$, then for every $T\le D/C$, after $T$ rounds, with probability at
least $1-C/n$, every colour is at most $2\gamma n/D$ and $|q-n/2|\le2\gamma^2n/D$.
\end{lemma}
\begin{proof}\leanok
One round multiplies the bound $B\in[\gamma n/D,2\gamma n/D]$ on the colours by at most
$1+(4\gamma^2+2\gamma+1)/D$ and keeps the window, with probability at least $1-C/n^2$: by (20),
$2\delta+c_i\le(4\gamma^2+2\gamma)n/D$, and by (19) $E[\Delta']\le 8\gamma^4n/D^2$ and
$E[\Delta']\ge-(\max_jc_j)(n-q)/n\ge-(4\gamma/3)n/D$. Iterate for $T\le D/C$ rounds, during which
the factor stays at most $2$. Proved with $C=30\gamma^2$.
\end{proof}

\begin{theorem}[Lower bound; Theorem~8]\label{thm:lb}
\lean{Undecided.Plurality.lower_bound_whp,Undecided.Plurality.lower_bound_consensus}\leanok\uses{lem:lb-first,lem:lb-descent,lem:lb-plateau}
There is $C>0$ such that, for every $n$ with $\log n\ge C$, every $k$ with $Ck\le(n/\log n)^{1/6}$
and every configuration $x$ without undecided nodes: for every $T$ with $CT\le\mathrm{md}(x)$, after
$T$ rounds every colour has at most $Cn/\mathrm{md}(x)$ nodes with probability at least $1-C/n$; and
for every $T$ with $C(T+1)\le\mathrm{md}(x)$, all nodes hold the same colour after $T$ rounds with
probability at most $C/n$. Since monochromatic configurations are absorbing, the convergence time
is $\Omega(\mathrm{md}(x))$ with probability at least $1-C/n$.
\end{theorem}
\begin{proof}\leanok
Compose Lemma~\ref{lem:lb-first}, Lemma~\ref{lem:lb-descent} and Lemma~\ref{lem:lb-plateau}
(with the $\gamma$ of Lemma~\ref{lem:lb-descent}, $D=\mathrm{md}(x)$) by the Markov property; a
round $T$ up to the descent round is covered by the bound at every round of
Lemma~\ref{lem:lb-descent}. A monochromatic configuration has $n>Cn/\mathrm{md}(x)$ nodes of one
colour once $\mathrm{md}(x)>C$. Proved with $C=1017429$ for the first form and $2C$ for the second.
\end{proof}

\section{The sequential dynamics: approximate majority}\label{sec:aae-approx-majority}
\begin{definition}[Sequential undecided-state dynamics]\label{def:seq}
\lean{Undecided.Sequential.Interaction,Undecided.Sequential.step,Undecided.Sequential.run}\leanok
\uses{def:usd}
At each step an ordered pair (initiator, responder) of distinct agents is drawn uniformly among
the $n(n-1)$ such pairs, independently of the past; only the responder updates, by the rule of
Definition~\ref{def:usd} applied to the initiator's state. This is the approximate-majority
protocol of Angluin, Aspnes and Eisenstat (Distributed Computing 21, 2008), with $x$, $y$ and
blank written $a$, $b$ and $u$.
\end{definition}
\begin{lemma}[One-step formula]\label{lem:seq-step}
\lean{Undecided.Sequential.avg_transition,Undecided.Sequential.avg_step_counts}\leanok
\uses{def:seq}
The expectation of a function of the next configuration that depends only on the states of the
initiator and the responder is its value on idle interactions plus the four state-changing
transitions $xb$, $yb$, $xy$, $yx$, weighted by $aq$, $bq$, $ab$, $ab$ pairs out of $n(n-1)$
(with $q$ the number of blank agents).
\end{lemma}
\begin{lemma}[Weighted supermartingales]\label{lem:seq-mart}
\lean{Undecided.Sequential.expList_exp_pathSum_le,Undecided.Sequential.expList_ind_le_of_weight}
\leanok
If one step never increases $e^{\varphi}F$ in expectation, then $e^{\sum_{\text{path}}\varphi}F$
has expectation at most $F$ at every finite time, and Markov's inequality bounds the probability
that the path sum of $\varphi$ is large.
\end{lemma}
\begin{lemma}[Six potentials]\label{lem:seq-pot}
\lean{Undecided.Sequential.stepSC,Undecided.Sequential.stepC,Undecided.Sequential.stepB,
Undecided.Sequential.stepX,Undecided.Sequential.stepY,Undecided.Sequential.stepM}\leanok
\uses{lem:seq-step}
The potentials $1/((a-b)^2+4n)$ (state-changing interactions), $1$ (central region), $1/(a+b)$
(blank corner), $3b+q+1$ and $3a+q+1$ (the two opinion corners), and $e^{-\lambda\max(a-b,0)}$
(the gap stays positive) satisfy the hypothesis of Lemma~\ref{lem:seq-mart} with suitable
weights.
\end{lemma}
\begin{theorem}[Convergence; AAE Theorem 1]\label{thm:seq-conv}
\lean{Undecided.Sequential.prob_notCons_le,Undecided.Sequential.consensus_whp}\leanok
\uses{lem:seq-mart,lem:seq-pot}
For every $c$ there is $C$ such that from every non-blank configuration of $n\ge2$ agents, after
$T\ge Cn\log n$ interactions all agents agree, with probability at least $1-C/n^c$.
\end{theorem}
\begin{theorem}[Correctness; AAE Theorem 2]\label{thm:seq-maj}
\lean{Undecided.Sequential.prob_notAllX_le,Undecided.Sequential.majority_whp,
Undecided.Sequential.approximate_majority}\leanok
\uses{thm:seq-conv}
For every $c$ there is $C$ such that if initially $a-b\ge C\sqrt n\log n$, then after
$T\ge Cn\log n$ interactions all agents hold $a$, with probability at least $1-C/n^c$.
\end{theorem}


%%FILE voter/blueprint/src/content.tex
\begin{abstract}
We formalize Sections 2.1--2.3 of Hassin and Peleg,
\emph{Distributed Probabilistic Polling and Applications to Proportionate Agreement},
Information and Computation 171 (2001), 248--268.
The synchronous process uses independent neighbor samples with directional weights.
On a finite connected nonbipartite undirected graph it reaches consensus with
probability one, and each color wins with its initial stationary vertex weight.
No path-space measure is needed: eventual probabilities are suprema of increasing
finite-time consensus probabilities.
\end{abstract}

\section{Shared foundation (generic library results)}
\begin{definition}[Finite weighted expectations]\label{def:weights}
\lean{Dynamics.Distribution,Dynamics.Kernel}\leanok
A distribution has nonnegative real weights summing to one. A kernel is a
family of such distributions, one for each state. Expectations are finite sums.
These are shared library definitions, not numbered results of the paper.
\end{definition}
\begin{theorem}[Stationary distribution exists]\label{thm:stationary}
\lean{Dynamics.Kernel.exists_stationary}\leanok\uses{def:weights}
Every finite nonempty stochastic matrix has a stationary distribution.
\end{theorem}
\begin{proof}\leanok
The shared library uses Ces\`aro averages of evolved distributions and
compactness of the finite probability simplex. This supplies the existence
fact recalled before Lemma 2.2; uniqueness is not required.
\end{proof}
\begin{theorem}[Finite-chain absorption]\label{thm:generic-absorption}
\lean{Dynamics.Kernel.exists_uniform_block,Dynamics.Kernel.geometric_blocks,Dynamics.Kernel.finite_absorption}\leanok\uses{def:weights}
An absorbing target accessible from every state of a finite chain admits a
uniform positive success probability in a fixed number of rounds. The survival
probability decays geometrically in blocks and tends to zero.
\end{theorem}
\begin{proof}\leanok
Take a common block length exceeding all finitely many access times, then a
maximum survival probability strictly below one. Iterate the resulting contraction.
\end{proof}
\begin{theorem}[Finite-horizon optional stopping]\label{thm:optional-stopping}
\lean{Dynamics.Kernel.event_error_of_invariant}\leanok\uses{def:weights}
If an observable equals $\varphi_A$ on $A$, $\varphi_B$ on $B$, lies in
$[\varphi_B+m,\varphi_B+M]$ elsewhere, and keeps its expectation up to time $T$, then
$\varphi(x_0)-\varphi_B-(\varphi_A-\varphi_B)\Pr(X_T\in A)$ lies between $m$ and $M$
times the probability of lying in neither event at time $T$.
\end{theorem}
\begin{proof}\leanok
A shared library result (roadmap FND-4): decompose the observable pointwise into its values
on $A$ and off $A\cup B$, take expectations, and bound the second part.
\end{proof}

\section{Section 2.1: model}\label{sec:hassin-peleg}
This section and the following ones formalize Sections 2.1 to 2.4 of Hassin and Peleg,
\emph{Distributed Probabilistic Polling and Applications to Proportionate Agreement}
(Information and Computation 171, 2001): absorption and the consensus probability (Section 2.2),
multiple colors (Section 2.3), the self-loop Remark of p.~254 (lazy voter and Wright--Fisher),
and the consensus time on the complete graph and on connected graphs (Section 2.4).
\begin{definition}[Synchronous copying]\label{def:step}
\lean{Voter.step,Voter.round,Voter.transition}\leanok\uses{def:weights}
A configuration is $s:V\to C$. Independently sample $r(i)$ according to row
$H_i$, then set the next configuration to $i\mapsto s(r(i))$.
\end{definition}
\begin{lemma}[Configuration transition formula]\label{lem:product}
\lean{Voter.transition_product,Voter.transition_product_bool}\leanok\uses{def:step}
For configurations $s,t$, the transition probability is
\[P_{st}=\prod_i\sum_{j:s(j)=t(i)}H_{ij}.\]
For Boolean colors a white target vertex contributes $\sum_j H_{ij}s(j)$;
a black target contributes its complement, giving the paper's product formula.
\end{lemma}
\begin{proof}\leanok
The configuration distribution is the pushforward of the independently sampled
neighbors. Factor its equality indicator into a product over vertices.
\end{proof}

\section{Section 2.2: absorption and consensus probability}
\begin{lemma}[Monochromatic-edge propagation]\label{lem:propagation}
\lean{Voter.monochromatic_edge,Voter.propagate_region,Voter.possible_consensus}\leanok\uses{def:step}
On a connected nonbipartite graph every Boolean configuration has a possible
finite path to a constant configuration.
\end{lemma}
\begin{proof}\leanok
A coloring without a monochromatic edge would be a proper two-coloring.
Starting with a monochromatic edge, each vertex of its region can keep that
color by sampling an internal neighbor. Connectedness supplies an edge leaving
every proper region. One more vertex can copy across that edge. Induction on
the number of vertices outside the region gives consensus.
\end{proof}
\begin{lemma}[Lemma 2.1: absorption]\label{lem:21}
\lean{Voter.transition_constant,Voter.possible_positive,Voter.consensus_tendsto}\leanok\uses{lem:propagation,thm:generic-absorption,lem:product}
Assume $H_{ij}>0$ whenever $i,j$ are adjacent; extra support, such as
self-loops, is allowed. Both constant configurations
are absorbing, and the probability of remaining nonconstant tends to zero.
This is the finite-time survival formulation of the paper's transience claim.
\end{lemma}
\begin{proof}\leanok
Every possible copying round has positive probability. Thus graph propagation
gives accessibility. Apply the shared finite-chain absorption theorem.
\end{proof}
\begin{lemma}[Lemma 2.3: invariant stationary weight]\label{lem:23}
\lean{Voter.round_mass,Voter.iterate_mass}\leanok\uses{def:step,thm:stationary}
For any stationary distribution $\pi$, the expected white mass at every finite
time is the initial mass $\sum_i\pi_i s(i)$. No graph connectivity assumption
is imposed on this invariant.
\end{lemma}
\begin{proof}\leanok
A vertex's expected next color indicator is $\sum_jH_{ij}s(j)$.
Interchanging finite sums and using $\pi H=\pi$ proves one-step preservation;
induction proves the iterated statement.
\end{proof}
\begin{lemma}[Lemma 2.2: limiting identification]\label{lem:22}
\lean{Voter.whiteMass_bounds,Voter.whiteMass_const,Voter.survival_eq,Voter.iterate_survival,Voter.colorProbability_eq_event,Voter.whiteProbability_error,Voter.whiteProbability_tendsto}\leanok\uses{lem:21,lem:23,thm:optional-stopping}
The discrepancy between expected stationary white mass and all-white
probability lies between zero and nonconsensus probability. Consequently
all-white probability converges to initial stationary white mass.
The implementation combines the paper's limiting identification with the
invariant from Lemma 2.3 when stating the limit.
\end{lemma}
\begin{proof}\leanok
Pointwise, the white mass agrees with the all-white indicator on constant
configurations and lies in $[0,1]$ elsewhere. Finite-horizon optional stopping,
with the two consensus configurations as targets and the invariant of Lemma 2.3,
bounds the discrepancy; squeeze it using Lemma 2.1.
\end{proof}
\begin{theorem}[Theorem 2.1]\label{thm:21}
\lean{Voter.eventualColor,Voter.colorProbability_mono,Voter.consensus_probability}\leanok\uses{lem:22}
Define eventual all-white probability as the supremum of the increasing
finite-time all-white probabilities. It equals the initial stationary white mass.
\end{theorem}
\begin{proof}\leanok
Consensus is absorbing, so the probabilities increase and are bounded by one.
Their supremum is their limit. Identify this limit using Lemma 2.2.
\end{proof}
\begin{theorem}[Corollary 2.2 and regular graphs]\label{cor:22}
\lean{Voter.degree_stationary,Voter.uniform_consensus_probability,Voter.regular_consensus_probability}\leanok\uses{thm:21}
For uniform neighbor sampling, $\pi_i=d_i/(2|E|)$, so all-white consensus has
probability $\sum_{i:s(i)=1}d_i/(2|E|)$. On a regular graph this is the initial
fraction of white vertices.
\end{theorem}
\begin{proof}\leanok
The degree-sum formula normalizes these weights. In the stationarity equation,
the degree cancels the reciprocal sampling degree; undirected adjacency gives
the remaining column sum. Equal degrees then give uniform vertex weights.
\end{proof}

\section{Section 2.3: multiple colors}
\begin{theorem}[Consensus in a specified color]\label{thm:colors}
\lean{Voter.step_project,Voter.iterate_project,Voter.colorProbability_project,Voter.color_consensus_probability,Voter.sum_color_consensus_probability}\leanok\uses{thm:21}
For a finite color set and any specified color $c$, its eventual consensus
probability is $\sum_{i:s(i)=c}\pi_i$. These probabilities sum to one.
\end{theorem}
\begin{proof}\leanok
The color-indicator projection commutes with every copying step and with
iterated expectations. Its all-white event is exactly all-$c$ consensus.
Apply Theorem 2.1 to this Boolean projection.
\end{proof}

\section{Section 2.4 on the complete graph: coalescing walks and consensus time}
\begin{definition}[Rounds, backward map, Wright--Fisher sampling]\label{def:backward}
\lean{Voter.runRounds,Voter.backward,Voter.disagreement,Voter.wfKernel}\leanok\uses{def:step}
A run applies rounds $r_1,\dots,r_T\colon V\to V$ in order. Its backward map is
$r_1\circ r_2\circ\cdots\circ r_T$. The disagreement indicator is $0$ on constant
configurations and $1$ otherwise. Wright--Fisher sampling picks every row uniformly
on $V$, self included, i.e.\ the complete graph with self-loops.
\end{definition}
\begin{lemma}[Duality with coalescing walks]\label{lem:duality}
\lean{Voter.runRounds_eq_comp,Voter.disagreement_runRounds_eq_zero,Voter.disagreement_eq_survival}\leanok\uses{def:backward}
After rounds $r_1,\dots,r_T$ the configuration is $s\circ r_1\circ\cdots\circ r_T$:
running the voter forward reads the initial colours through the backward map, whose
point images are coalescing random walks. If the backward map is constant, the
configuration is a consensus.
\end{lemma}
\begin{proof}\leanok
Induction on the list of rounds, since one round maps $s$ to $s\circ r$.
\end{proof}
\begin{lemma}[Rounds are independent and uniform]\label{lem:wf-rounds}
\lean{Voter.independent_wfKernel_expect,Voter.iterate_wfKernel_eq_expList}\leanok\uses{def:backward,lem:product}
Under Wright--Fisher sampling a round is uniform on $V\to V$, so the $T$-step
configuration kernel is the average over $T$ independent uniform rounds.
\end{lemma}
\begin{proof}\leanok
The product weights are all $n^{-n}$; induct on $T$, peeling off the first round.
\end{proof}
\begin{lemma}[Meeting probability]\label{lem:meeting}
\lean{Voter.avg_maps_agree,Voter.avg_coord_disagree,Voter.expList_backward_indicator,Voter.expList_backward_ne}\leanok\uses{lem:wf-rounds}
Two coordinates of a uniform map $V\to V$ agree with probability $1/n$. Hence the
backward walks from two distinct vertices have not met after $T$ rounds with
probability exactly $(1-1/n)^T$.
\end{lemma}
\begin{proof}\leanok
Peel off the first round, swap it with the remaining average, and multiply by the
one-round probability of not meeting; induct on $T$.
\end{proof}
\begin{theorem}[Consensus time]\label{thm:consensus-time}
\lean{Voter.disagreement_runRounds_le,Voter.iterate_disagreement_le,Voter.voter_consensus_whp}\leanok\uses{lem:duality,lem:wf-rounds,lem:meeting}
The probability that no consensus is reached after $T$ rounds is at most
$(n-1)(1-1/n)^T$. Consequently consensus is reached within $2n\log n$ rounds with
probability at least $1-1/n$.
\end{theorem}
\begin{proof}\leanok
Consensus holds once every backward walk has met the walk of a fixed vertex; apply the
union bound over the other $n-1$ vertices and the meeting probability, then
$(1-1/n)^T\le e^{-T/n}\le n^{-2}$ for $T\ge 2n\log n$.
\end{proof}

\section{Self-loops, the lazy voter and Wright--Fisher (VOT-2)}
\begin{lemma}[Propagation through the kernel support]\label{lem:kernel-propagation}
\lean{Voter.KernelPossible,Voter.exists_weight_pos,Voter.kernelPossible_positive,Voter.reachable_success_of_kernelPossible,Voter.propagate_kernel_region,Voter.kernelPossible_consensus_of_selfLoop,Voter.consensus_tendsto_of_selfLoop}\leanok\uses{lem:propagation,lem:21,thm:generic-absorption}
Call a round possible for $H$ when every vertex samples a vertex of positive weight in its
row. On a connected graph whose edges are charged by $H$, a monochromatic region in which
every vertex charges a vertex of the region grows to consensus through possible rounds.
If $H_{vv}>0$ for one vertex $v$, the region $\{v\}$ qualifies, so nonconsensus probability
tends to zero, bipartite graphs included.
\end{lemma}
\begin{proof}\leanok
As in Lemma~2.1: a boundary vertex copies its neighbour in the region, vertices of the region
copy inside it, and all other vertices copy any vertex of positive weight. Possible rounds
have positive probability, and the finite-chain absorption theorem applies.
\end{proof}
\begin{theorem}[Remark, p.~254]\label{thm:selfloop}
\lean{Voter.eventualColor_eq_whiteMass_of_tendsto,Voter.consensus_probability_of_selfLoop,Voter.color_consensus_probability_of_selfLoop}\leanok\uses{lem:kernel-propagation,lem:22,thm:colors}
On a connected graph, bipartite or not, if $H$ charges every edge and $H_{vv}>0$ for some
vertex $v$, then consensus in colour $c$ has probability $\sum_{i:s(i)=c}\pi_i$ for every
stationary $\pi$.
\end{theorem}
\begin{proof}\leanok
Lemma~2.2 only needs nonconsensus probability to vanish; then project colours as in
Section~2.3.
\end{proof}
\begin{theorem}[Lazy voter]\label{thm:lazy}
\lean{Voter.lazy,Voter.lazyNeighbor,Voter.lazy_weight,Voter.lazyNeighbor_weight,Voter.lazy_stationary_iff,Voter.lazy_selfLoop_pos,Voter.lazy_weight_pos,Voter.lazy_consensus_probability,Voter.lazyNeighbor_consensus_probability}\leanok\uses{thm:selfloop,cor:22}
The lazy kernel $(I+H)/2$ has the stationary distributions of $H$ and positive self-loops.
On any connected graph the lazy voter $(I+D^{-1}A)/2$ reaches consensus in colour $c$ with
probability $\sum_{i:s(i)=c} d_i/2m$.
\end{theorem}
\begin{proof}\leanok
Apply the Remark to $(I+H)/2$; for uniform neighbour sampling use the degree distribution.
\end{proof}
\begin{theorem}[Neutral Wright--Fisher fixation]\label{thm:wright-fisher}
\lean{Voter.wfKernel_pos,Voter.wfKernel_stationary,Voter.uniform_prob_eq_card,Voter.top_not_colorable_two,Voter.wrightFisher_fixation_of_three_le,Voter.wrightFisher_fixation}\leanok\uses{thm:colors,thm:selfloop,def:backward}
If every offspring picks a uniform parent, itself included, allele $a$ fixes with probability
$c_a/n$, where $c_a$ is its initial number of copies.
\end{theorem}
\begin{proof}\leanok
The uniform distribution is stationary and every entry is positive. For $n\ge 3$ the
complete graph is nonbipartite and Section~2.3 applies; for every $n$ the Remark applies.
\end{proof}

\section{Section 2.4 on connected graphs: meeting time and $O(n^3\log n)$ consensus}\label{sec:kms-meeting-time}
Theorem~2.5 of Hassin and Peleg (Theorem~8 of the survey of Becchetti, Clementi and Natale,
\emph{Consensus dynamics: an overview}, SIGACT News 2020) is formalized for the lazy voter, with
the meeting-time argument of Kanade, Mallmann-Trenn and Sauerwald, \emph{On coalescence time in
graphs} (SODA 2019), Proposition~B.9 ($t_{\mathrm{meet}}\le 4t_{\mathrm{hit}}$). For the
synchronous voter with plain uniform-neighbour sampling, Theorem~8 of the survey needs a minor
correction (a nonbipartiteness hypothesis, or a lazy walk); this section formalizes the lazy
version, and Section~\ref{sec:hp-plain-walk} the plain walk on nonbipartite graphs
(\href{https://github.com/formal-dynamics/leanamics/blob/main/voter/FORMALIZATION_DIFFERENCES.md#41-survey-theorem-8-needs-a-minor-correction-the-lazy-walk-is-formalized}{details}).
\begin{definition}[Two-token walk, lazy sampling]\label{def:pairwalk}
\lean{Voter.pairWalk,Voter.leftRounds}\leanok\uses{def:backward,thm:lazy}
For a sampling kernel $H$, two tokens move with the rounds of the voter dynamics: the
token at $x$ moves to $r(x)$, where $r$ is drawn from the independent product of the rows
of $H$. Off the diagonal the tokens make independent $H$-steps; on the diagonal they move
together. The lazy kernel of a graph samples the vertex itself with probability $1/2$ and
otherwise a uniformly random neighbour. Rounds can also be composed on the left,
$W\mapsto r\circ W$.
\end{definition}
\begin{lemma}[Duality for every kernel]\label{lem:duality-general}
\lean{Voter.independent_expect_pair,Voter.iterate_transition_eq_leftRounds,Voter.iterate_leftRounds_pair,Voter.iterate_disagreement_le_pairWalk}\leanok\uses{def:pairwalk,lem:duality}
For every sampling kernel, the probability that no consensus is reached after $T$ rounds
is at most the sum over $v\neq u_0$ of the probability that the tokens started at $v$ and
$u_0$ are apart after $T$ steps.
\end{lemma}
\begin{proof}\leanok
The voter dynamics on maps started from the identity composes the rounds on the right;
composing them on the left gives the same finite-time expectations, since the rounds are
independent and identically distributed (induction on $T$, exchanging two finite
expectations). The left composition read at two vertices is the two-token walk. Conclude
with the pointwise union bound of Theorem~\ref{thm:consensus-time}.
\end{proof}
\begin{theorem}[Hassin--Peleg Theorem 2.4, tail form]\label{thm:hp24}
\lean{Voter.pairWalk_iterate_outside_mul_le,Voter.iterate_disagreement_le_of_meeting}\leanok\uses{lem:duality-general}
If from every pair of vertices the tokens are apart after $T_0$ steps with probability at
most $1/2$, then consensus fails after $kT_0$ rounds with probability at most
$(n-1)2^{-k}$.
\end{theorem}
\begin{proof}\leanok
The diagonal is absorbing, so the probabilities of being apart are submultiplicative in
blocks of $T_0$ steps; apply Lemma~\ref{lem:duality-general}.
\end{proof}
\begin{lemma}[Hitting times of the lazy walk; Fact 2.3]\label{lem:hitting}
\lean{Voter.le_zero_of_lapMatrix_le,Voter.exists_hitting,Voter.hitting,Voter.hitting_nonneg,Voter.hitting_symm,Voter.hitting_add_hitting_le}\leanok
On a connected graph with Laplacian $L$ and $\mathrm{vol}=\sum_z d_z$, the system
$h(y)=0$, $(Lh)(x)=2d_x$ for $x\neq y$ has a nonnegative solution $h_y$, the expected
hitting time of $y$ for the lazy walk. It satisfies
$\sum_z d_z h_y(z)-\mathrm{vol}\,h_y(x)=\sum_z d_z h_x(z)-\mathrm{vol}\,h_x(y)$ and
$h_y(x)+h_x(y)\le 4\,\mathrm{vol}\,(n-1)\le 4n^3$.
\end{lemma}
\begin{proof}\leanok
The maximum principle makes the linear system injective, hence surjective. The identity is
the symmetry of $L$. For the bound, $\psi=h_y-h_x$ has energy
$\psi^\top L\psi=2\,\mathrm{vol}\,(h_y(x)+h_x(y))$; apply Cauchy--Schwarz along a path from
$x$ to $y$.
\end{proof}
\begin{lemma}[Sequential and synchronous walks]\label{lem:sequential}
\lean{Voter.seqSurvival,Voter.seqSurvival_mul_le,Voter.iterate_outside_succ_le,Voter.iterate_outside_le_three_quarters}\leanok\uses{def:pairwalk}
Let the two tokens move one at a time, meeting being checked after each move. If a
nonnegative potential $F$ drops by one in expectation under a step of either token, the
sequential walks are apart after $T$ steps with probability at most $F/T$. For a lazy
kernel, the synchronous tokens are apart after $T+1$ steps with probability at most the
average of the indicator of starting apart and the sequential survival; hence at most $3/4$
once $T\ge 2\sup F$.
\end{lemma}
\begin{proof}\leanok
Two inductions on $T$. When the first token lands on the old position of the second, the
second stays put with probability at least $1/2$.
\end{proof}
\begin{theorem}[Meeting time of lazy walks; Lemma 2.4]\label{thm:meeting-lazy}
\lean{Voter.meetingPotential,Voter.lazy_meeting_core,Voter.lazy_meeting_le_half,Voter.lazy_apart_le_pow,Voter.lazy_meeting_le_pow}\leanok\uses{lem:hitting,lem:sequential,thm:hp24}
On a connected graph with $n$ vertices, two lazy walks are apart after
$3(16n^3+1)\le 51n^3$ steps with probability at most $1/2$, and after $51kn^3$ steps with
probability at most $2^{-k}$, from every pair of start vertices.
\end{theorem}
\begin{proof}\leanok
Apply Lemma~\ref{lem:sequential} to the potential $h_y(x)-\sum_z d_z h_y(z)/\mathrm{vol}+4n^3$
of Lemma~\ref{lem:hitting}, bounded by $8n^3$, then submultiplicativity.
\end{proof}
\begin{theorem}[Hassin--Peleg Theorem 2.5; Survey Theorem 8]\label{thm:hp25}
\lean{Voter.card_sub_one_mul_half_pow_le,Voter.lazy_voter_consensus_whp}\leanok\uses{thm:meeting-lazy,lem:duality-general}
On every connected graph with $n$ vertices and from every initial colouring, the lazy voter
dynamics reaches consensus within $255\,n^3\log n$ rounds with probability at least $1-1/n$.
\end{theorem}
\begin{proof}\leanok
With $k=\lfloor T/(51n^3)\rfloor$, Lemma~\ref{lem:duality-general} and
Theorem~\ref{thm:meeting-lazy} bound the failure probability by $(n-1)2^{-k}\le 1/n$.
\end{proof}

\section{Section 2.4 for the plain walk: the bipartite double cover}\label{sec:hp-plain-walk}
Hassin and Peleg's own setting is the plain synchronous voter, in which every vertex copies a
uniformly random neighbour, on a connected nonbipartite graph. Their Lemma~2.4 bounds the
meeting time through the bipartite double cover $\tilde G$, quoting Tetali and Winkler
(\emph{Simultaneous reversible Markov chains}, 1993). We follow the reduction to $\tilde G$, but
prove the meeting bound with the potential of Coppersmith, Tetali and Winkler
(\emph{Collisions among random walks on a graph}, SIAM J.\ Discrete Math.\ 1993) on $\tilde G$.
The proof of Lemma~2.4 needs a minor correction: it uses $\mathrm{Diam}(\tilde G)\le n$, but the
diameter of $\tilde G$ can be $2n-3$ (a path ending in a triangle); $\mathrm{Diam}(\tilde G)\le
2n-1$ only changes the constant, and our route does not use the diameter
(\href{https://github.com/formal-dynamics/leanamics/blob/main/voter/FORMALIZATION_DIFFERENCES.md#43-the-plain-walk-on-connected-nonbipartite-graphs-lemma-24-and-theorem-25}{details}).
Hitting times $h$ are those of Lemma~\ref{lem:hitting}, twice the hitting times of the plain walk.
\begin{definition}[Bipartite double cover]\label{def:double-cover}
\lean{Voter.doubleCover,Voter.doubleCover_adj,Voter.doubleCoverFlip,Voter.doubleCover_degree,Voter.doubleCover_degree_pos,Voter.volume_doubleCover}\leanok\uses{def:step}
The graph $\tilde G=G\times K_2$ has vertices $(a,i)$, $a\in V$, $i\in\{0,1\}$, and
$(a,i)\sim(b,j)$ iff $a\sim b$ in $G$ and $i\neq j$. Swapping the layers is an automorphism,
degrees are those of $G$, and $\mathrm{vol}\,\tilde G=2\,\mathrm{vol}\,G$.
\end{definition}
\begin{lemma}[Connectivity of the double cover]\label{lem:double-cover-connected}
\lean{Voter.doubleCover_reachable_of_walk,Voter.doubleCover_reachable_flip,Voter.doubleCover_connected}\leanok\uses{def:double-cover}
If $G$ is connected and not $2$-colourable, then $\tilde G$ is connected.
\end{lemma}
\begin{proof}\leanok
Walks of $G$ lift to $\tilde G$, flipping the layer at every step. If $\tilde G$ were
disconnected, ``$(w,0)$ is reachable from $(x_0,0)$'' would be a proper $2$-colouring of $G$.
\end{proof}
\begin{lemma}[Hitting times on the double cover]\label{lem:hitting-flip}
\lean{Voter.hitting_unique,Voter.hitting_doubleCover_flip,Voter.uniformNeighbor_expect,Voter.uniformNeighbor_expect_hitting,Voter.uniformNeighbor_doubleCover_expect}\leanok\uses{lem:hitting,def:double-cover,lem:double-cover-connected,cor:22}
The hitting-time system has a unique solution, so the hitting times of $\tilde G$ satisfy
$\tilde h_{(y,1-j)}(x,1-i)=\tilde h_{(y,j)}(x,i)$. A plain step away from $y$ lowers $h_y$ by $2$
in expectation, and a plain step of $\tilde G$ from $(x,i)$ is a plain step of $G$ from $x$
followed by a flip of the layer.
\end{lemma}
\begin{proof}\leanok
Swapping the layers maps solutions to solutions; the maximum principle gives uniqueness. The
first-step equation is the row $x$ of the Laplacian system.
\end{proof}
\begin{lemma}[Additive drift for two tokens]\label{lem:additive-drift}
\lean{Voter.pairWalk_apart_mul_le}\leanok\uses{def:pairwalk}
If a nonnegative potential $F$ drops by at least $c$ in expectation under one step of the
two-token walk from every pair of distinct vertices, then
$cT\cdot\Pr[\text{apart after }T\text{ steps}]\le F$.
\end{lemma}
\begin{proof}\leanok
Truncate $F$ to $0$ on the diagonal, show
$\mathbb E[F'_T]+c\sum_{t<T}\Pr[\text{apart at }t]\le F'$ by induction on $T$, and use that
being apart is antitone in $t$.
\end{proof}
\begin{lemma}[Exact drift of the potential]\label{lem:plain-drift}
\lean{Voter.plainPotential,Voter.doubleCover_hitting_le,Voter.plainPotential_nonneg,Voter.plainPotential_le,Voter.doubleCover_hitting_sum_flip,Voter.meetingPotential_doubleCover_flip,Voter.expect_meetingPotential_fst,Voter.expect_meetingPotential_snd,Voter.plainPotential_drift}\leanok\uses{lem:hitting-flip,thm:meeting-lazy}
Let $\Phi$ be the Coppersmith--Tetali--Winkler potential of $\tilde G$ (as in
Theorem~\ref{thm:meeting-lazy}) and $\Psi(x,y)=\Phi((x,0),(y,0))$. Every hitting time of
$\tilde G$ is at most $16n^3$, $0\le\Psi\le 32n^3$, and for $x\neq y$ one synchronous plain step
of the two tokens lowers $\Psi$ by exactly $4$ in expectation.
\end{lemma}
\begin{proof}\leanok
The bound is the commute bound of Lemma~\ref{lem:hitting} on $\tilde G$ ($2n$ vertices, volume
at most $2n^2$). A synchronous step at $(x,y)$ is, on $\tilde G$, a move of the first token from
$(x,0)$ to layer $1$ followed by a move of the second token from $(y,0)$ to layer $1$. Between
the two moves the tokens lie in different layers, hence are distinct, and each move lowers
$\Phi$ by $2$. The layer symmetry of Lemma~\ref{lem:hitting-flip} returns both tokens to
layer $0$.
\end{proof}
\begin{theorem}[Meeting time of plain walks; Lemma 2.4]\label{thm:meeting-plain}
\lean{Voter.plain_apart_mul_le,Voter.plain_meeting_core,Voter.plain_apart_le_pow,Voter.plain_meeting_le_half}\leanok\uses{lem:additive-drift,lem:plain-drift,lem:double-cover-connected,thm:hp24}
On a connected nonbipartite graph with $n$ vertices, two synchronous plain walks are apart after
$T$ steps with probability at most $8n^3/T$; hence after $16n^3$ steps with probability at most
$1/2$, and after $16kn^3$ steps with probability at most $2^{-k}$, from every pair of start
vertices.
\end{theorem}
\begin{proof}\leanok
Lemma~\ref{lem:additive-drift} with $c=4$ and Lemma~\ref{lem:plain-drift} give
$4T\cdot\Pr[\text{apart}]\le 32n^3$; then submultiplicativity in blocks.
\end{proof}
\begin{theorem}[Hassin--Peleg Theorem 2.5, plain walk; Survey Theorem 8 with the nonbipartiteness hypothesis]\label{thm:hp25-plain}
\lean{Voter.plain_disagreement_le,Voter.plain_voter_consensus_whp}\leanok\uses{thm:meeting-plain,lem:duality-general,thm:hp24,thm:hp25}
On every connected nonbipartite graph with $n$ vertices and from every initial colouring, the
plain voter dynamics fails to reach consensus after $16kn^3$ rounds with probability at most
$(n-1)2^{-k}$, and reaches consensus within $80\,n^3\log n$ rounds with probability at least
$1-1/n$.
\end{theorem}
\begin{proof}\leanok
As for Theorem~\ref{thm:hp25}, with $k=\lfloor T/(16n^3)\rfloor$ and
Theorem~\ref{thm:meeting-plain}.
\end{proof}

\section{Consensus time via conductance (VOT-5)}\label{sec:bgkm-conductance}
Source: Berenbrink, Giakkoupis, Kermarrec, Mallmann-Trenn, \emph{Bounds on the voter model in
dynamic networks}, ICALP 2016, arXiv:1603.01895 (BGKM16). The process is the lazy voter
$(I+D^{-1}A)/2$ of Theorem~\ref{thm:lazy}.

Lemma~2.1 of BGKM16 needs a minor correction (its sum should run over the minority side, which is
what its proof gives), and so do the statement of Lemma~2.4 (it prints $1/n^2$ for $1-1/n^2$)
and the sign of an exponent in its proof; the corrected versions are formalized
(\href{https://github.com/formal-dynamics/leanamics/blob/main/voter/FORMALIZATION_DIFFERENCES.md#51-lemma-21-of-bgkm16-needs-a-minor-correction}{details
on Lemma~2.1},
\href{https://github.com/formal-dynamics/leanamics/blob/main/voter/FORMALIZATION_DIFFERENCES.md#5-consensus-time-via-conductance-berenbrinkgiakkoupiskermarrecmallmann-trenn}{details
on Lemma~2.4} in Section~5.5 there).
\begin{definition}[Volume, conductance, minority side, potential]\label{def:conductance}
\lean{Voter.vol,Voter.conductance,Voter.discordant,Voter.minority,Voter.potential,Voter.dynamicLazy}\leanok
$\mathrm{vol}(S)=\sum_{u\in S}d_u$; the conductance $\varphi(G)$ is the minimum of
$|\mathrm{cut}(U,V\setminus U)|/\mathrm{vol}(U)$ over $0<\mathrm{vol}(U)\le m$; $\lambda_u$ is the
number of neighbours of $u$ with another opinion; for two opinions $s_t$ is the side of smaller
volume and $\Psi=\sqrt{\mathrm{vol}(s_t)}$.
\end{definition}
\begin{lemma}[Independent coordinates]\label{lem:indep}
\lean{Voter.independent_expect_update,Voter.independent_expect_mul_of_update,Voter.independent_expect_sum_sq,Voter.independent_expect_sum_cube,Voter.independent_expect_comp_sum_le}\leanok
Resampling one coordinate of an independent point does not change its law. Consequently
centred coordinate functions have additive second and third moments, and if replacing each
coordinate function improves $\mathbb{E}f(z+\cdot)$ for all admissible $z$, then the replacement of
all of them improves $\mathbb{E}f(c+\sum_i g_i(x_i))$ (BGKM16 Lemma~A.1).
\end{lemma}
\begin{proof}\leanok
Reindex the pairs $(x,a)$ by $(x[i\mapsto a],x_i)$; induct over the set of replaced coordinates.
\end{proof}
\begin{lemma}[BGKM16 Lemma~2.1, corrected]\label{lem:bgkm-21}
\lean{Voter.potential_drift,Voter.potential_drift_conductance,Voter.expect_sqrt_le_comp,Voter.expect_compSum_sq_ge,Voter.expect_compSum_cube_le,Voter.sqrt_add_le_taylor,Voter.sqrt_chord}\leanok\uses{def:conductance,lem:indep,thm:lazy}
If $s_t\neq\emptyset$, then
$\mathbb{E}[\Psi_{t+1}\mid s_t]\le\Psi-\sum_{u\in s_t}\lambda_u d_u/(32\Psi^3)
\le\Psi-d_{\min}\varphi/(32\Psi)$. The paper prints the sum over all of $V$; this needs a minor correction (the
inequality does not hold on the star $K_{1,15}$), and its proof gives the sum over $s_t$
(\href{https://github.com/formal-dynamics/leanamics/blob/main/voter/FORMALIZATION_DIFFERENCES.md#51-lemma-21-of-bgkm16-needs-a-minor-correction}{details}).
\end{lemma}
\begin{proof}\leanok
$\Psi_{t+1}$ is at most the root of the new volume $P+\sum_u X_u$ of the opinion of $s_t$.
Replace $X_u$ for $u\notin s_t$ by $Y_u=\lambda_u\mathbf{1}[\text{move}]$ (concavity of
$\sqrt{\cdot}$), expand $\sqrt{P+x}$ to third order, and use $\mathbb{E}\Delta'=0$,
$\mathbb{E}\Delta'^2\ge\sum_{s_t}\lambda_ud_u/4$, $\mathbb{E}\Delta'^3\le0$. The conductance gives
$\sum_{s_t}\lambda_ud_u\ge d_{\min}\varphi\,\mathrm{vol}(s_t)$.
\end{proof}
\begin{theorem}[BGKM16 Lemma~2.2 and Theorem~1.1~(i), two opinions]\label{thm:bgkm-two}
\lean{Voter.lazy_consensus_of_minority,Voter.lazy_consensus_conductance,Voter.lazy_expected_consensus_time,Voter.dynamic_consensus_conductance,Voter.sum_iterate_le_of_block}\leanok\uses{lem:bgkm-21}
If $128\,\mathrm{vol}(s_0)\le d_{\min}\varphi T$, the opinions still disagree at time $T$ with
probability at most $1/2$; in particular after $128m/(d_{\min}\varphi)$ rounds, and the expected
consensus time is at most twice that. On dynamic graphs with fixed degrees the same holds once
$128\,\mathrm{vol}(s_0)\le d_{\min}\sum_{t<T}\varphi_t$.
\end{theorem}
\begin{proof}\leanok
The drift lemma of the shared library (FND-5) with drift $d_{\min}\varphi_t/32$; restarting
bounds the expected time.
\end{proof}
\begin{theorem}[BGKM16 Theorem~1.1~(i), any number of opinions]\label{thm:bgkm-many}
\lean{Voter.phase_step,Voter.level_time_le,Voter.expected_time_many_bound,Voter.lazy_expected_consensus_time_many,Voter.lazy_consensus_conductance_many}\leanok\uses{thm:bgkm-two,thm:colors}
There is a constant $b$ such that the expected consensus time is at most $b\,m/(d_{\min}\varphi)$
and consensus holds within $b\,m/(d_{\min}\varphi)$ rounds with probability at least $1/2$, for
every number of opinions.
\end{theorem}
\begin{proof}\leanok
Project each opinion to the two-opinion process ``$i$ against the rest''. From $\ell\le\theta$
opinions, after $384\,\mathrm{vol}(V)/(\theta d_{\min}\varphi)$ rounds at most $5\theta/6$
remain with probability at least $1/3$ (BGKM16 Lemma~2.3). Summing over the levels
$\theta_j=n(5/6)^j$ gives a geometric series; Markov's inequality gives the probability form.
\end{proof}
\begin{lemma}[Multiplicative drift of the potential]\label{lem:bgkm-mul}
\lean{Voter.potential_drift_mul,Voter.sq_conductance_mul_vol_le,Voter.one_le_potential,Voter.potential_le_card}\leanok\uses{lem:bgkm-21}
On a graph with $n$ vertices and conductance $\varphi$, from every two-opinion configuration,
$\mathbb{E}[\Psi_{t+1}\mid s_t]\le(1-\varphi^2/(32n))\,\Psi(s_t)$. Moreover $1\le\Psi\le n$
before consensus.
\end{lemma}
\begin{proof}\leanok
$\lambda_ud_u\ge\lambda_u^2$ and Cauchy--Schwarz over the at most $n$ vertices of $s_t$ give
$\sum_{s_t}\lambda_ud_u\ge(\sum_{s_t}\lambda_u)^2/n\ge\varphi^2\Psi^4/n$; insert this into
Lemma~\ref{lem:bgkm-21}. Volumes are integers and $\mathrm{vol}(s_t)\le m\le n^2$.
\end{proof}
\begin{theorem}[BGKM16 Lemma~2.4 and Theorem~1.1~(ii)]\label{thm:bgkm-sq}
\lean{Voter.lazy_consensus_conductance_sq,Voter.dynamic_consensus_conductance_sq,Voter.lazy_consensus_conductance_sq_many,Voter.dynamic_consensus_conductance_sq_many,Voter.lazy_expected_consensus_time_sq,Voter.iterateSeq_pos_le_of_mul_drift,Voter.prod_one_sub_le}\leanok\uses{lem:bgkm-mul,thm:colors}
If $96\,n\ln n\le\sum_{t<T}\varphi_t^2$ (static graph: $96\,n\ln n\le\varphi^2T$), the opinions
still disagree at time $T$ with probability at most $1/n^2$ for two opinions and at most $1/n$
for any number of opinions, on static graphs and on dynamic graphs with fixed degrees. On a static
graph the expected consensus time is at most $2T$.
\end{theorem}
\begin{proof}\leanok
The multiplicative drift lemma of the shared library (FND-5) gives
$\Pr(\Psi_T>0)\le\Psi(s_0)\prod_{t<T}(1-\varphi_t^2/(32n))\le n\cdot e^{-3\ln n}$. For many
opinions apply this to each projection ``$i$ against the rest'' (the projection commutes with a
round, whatever graph the adversary chooses) and take a union bound over the at most $n$ opinions
present initially; restarting bounds the expected time.
\end{proof}
\begin{theorem}[BGKM16 Theorem~1.1]\label{thm:bgkm-min}
\lean{Voter.lazy_consensus_conductance_min,Voter.dynamic_consensus_conductance_min}\leanok\uses{thm:bgkm-two,thm:bgkm-many,thm:bgkm-sq}
The consensus time is at most $\min\{\tau,\tau'\}$ with probability at least $1/2$: on a static
graph for every number of opinions, with $\tau=b\,m/(d_{\min}\varphi)$ and
$\tau'=b\,n\ln n/\varphi^2$ for a constant $b$; on a dynamic graph with fixed degrees for two
opinions, with $\tau$ the first time with $d_{\min}\sum_{t<\tau}\varphi_t\ge128\,\mathrm{vol}(s_0)$
and $\tau'$ the first time with $\sum_{t<\tau'}\varphi_t^2\ge96\,n\ln n$.
\end{theorem}
\begin{proof}\leanok
Combine Theorems~\ref{thm:bgkm-two}, \ref{thm:bgkm-many} and \ref{thm:bgkm-sq}, using
$1/n\le1/2$ for $n\ge2$; with one vertex there is nothing to agree on.
\end{proof}

\section{Examples and scope}
\begin{lemma}[Triangle and a nonuniform matrix]\label{lem:examples}
\lean{Voter.Examples.triangle_one_white,Voter.Examples.biased_stationary}\leanok\uses{cor:22,lem:23}
One white vertex on a triangle wins with probability $1/3$.
The stochastic matrix whose two rows are $(1/3,2/3)$ has stationary
weights $(1/3,2/3)$, providing a nonuniform invariant example.
\end{lemma}
\begin{lemma}[Wright--Fisher sampling]\label{lem:wf}
\lean{Voter.Examples.wrightFisher,Voter.Examples.wrightFisher_stationary,Voter.Examples.wrightFisher_one_third}\leanok\uses{thm:21}
On three individuals, let every offspring pick a uniform parent, itself included.
This kernel has self-loops, so it is supported beyond the edges of the triangle.
Theorem~2.1 still applies, and one allele copy fixes with probability $1/3$.
\end{lemma}
\begin{lemma}[Bipartite two-cycle]\label{lem:cycle}
\lean{Voter.Examples.twoVertex_unique_neighbor,Voter.Examples.twoVertex_flip,Voter.Examples.twoVertex_cycle,Voter.Examples.twoVertex_never_consensus}\leanok\uses{def:step}
On the two-vertex graph, every vertex must copy the other vertex.
The alternating coloring flips each round and returns after two rounds;
neither coloring is constant. Thus connectedness alone is insufficient.
\end{lemma}
Hassin--Peleg Theorem~2.5 is formalized for the lazy walk on every connected graph
(Theorem~\ref{thm:hp25}) and for the plain walk on connected nonbipartite graphs
(Theorem~\ref{thm:hp25-plain}). Dynamic networks and extremal coalition results are future work
beyond this milestone. Chernoff bounds remain in the
3-majority package until another theorem requires their extraction.

